% Research contribution of agent /root/monodromy.  Input into a parent document.
% Status is local to the assertions below; this is not a claim to solve all of A.3--A.5.
\section{A.3--A.5: proved statements and remaining gates}
\label{sec:monodromy-contribution}
\status{Incomplete universal determination; the individually labelled results below are proved.  Lead A's full-linear assertion is refuted, its positivity assertion is proved with an explicit threshold, and its plane-curve point-count consequence is proved.}

\subsection{Conventions and the generic rank for projective hypersurfaces}

Let $S_d=\mathbf P(H^0(\mathbf P^r,\mathcal O(d)))$, and let
$\pi:\mathcal Y_d\to S_d$ be the universal degree-$d$ hypersurface.
For $r\geq2$, the sheaf $\mathcal H=R^{r-1}\pi_*\mathcal O_{\mathcal Y_d}$
is locally free of rank
\[
 h_d=\begin{cases}\binom{d-1}{r}&d\geq r+1,\\0&d\leq r.\end{cases}
\]
This follows from the universal divisor exact sequence and the cohomology
of line bundles on projective space; formation of $\mathcal H$ commutes
with all base changes.  Absolute Frobenius induces a map
$F_S^*\mathcal H\to\mathcal H$.  On its invertibility locus $W_d$, the
kernel of $F-1$ is a locally constant $\mathbf F_p$-sheaf of rank $h_d$.
On the smooth locus it is $R^{r-1}\pi_*\mathbf F_p$, by the
Artin--Schreier sequence.  For $X=\mathbf P^r$, this is the vanishing
local system in the question, since $H^{r-1}(\mathbf P^r,\mathcal O)=0$.

\begin{proposition}[All degrees, all characteristics]\label{prop:projective-generic-hw-rank}
\status{Proved under stated hypotheses}
The generic degree-$d$ hypersurface of $\mathbf P^r$ has bijective
Frobenius on $H^{r-1}(\mathcal O)$, and its generic stable rank is $h_d$.
In particular, for plane curves the generic rank is
$g_d=(d-1)(d-2)/2$, and the smallest positivity threshold is $d_0=3$.
\end{proposition}

\begin{proof}
Work over an algebraic closure and take $d$ hyperplanes in general
position.  Their union $D$ is a simple normal crossings divisor.
The alternating resolution of $\mathcal O_D$ by the structure sheaves
of intersections has only $H^0$ in its first hypercohomology page,
because every nonempty intersection is a projective space.  Frobenius
is bijective on each of these $H^0$'s and hence on the cohomology of
their complex.  It is therefore bijective on
$H^{r-1}(D,\mathcal O_D)$.  Thus $W_d$ is nonempty.  The smooth locus
is a nonempty dense open of the irreducible scheme $S_d$, so it meets
$W_d$.  Invertibility is open, which proves the assertion.  This proof
establishes only bijectivity on the specified coherent cohomology;
it does not conflate this condition with every formulation of
ordinarity for higher dimensional varieties.
\end{proof}

\par\smallskip\noindent\textbf{Comparison with Illusie.}
Illusie (1990), Theorem~0.1, p.~375, proves the stronger statement that
the general smooth complete intersection of every specified dimension
and multidegree is ordinary: its Newton and Hodge polygons coincide.
The theorem and the explanation of the earlier Koblitz argument on
pp.~375--376 were inspected in the publisher's two-page preview of
the reprint.  In particular, generic Hasse--Witt bijectivity follows;
the proposition above gives a separate complete proof of exactly the
part needed here.  The parameter space in Illusie's theorem varies all
defining equations.  It does not, by restriction of an open set alone,
prove generic ordinarity for hypersurface sections of every fixed $X$.

\subsection{The determinant character and the Hasse invariant}
\label{subsec:hasse-determinant}

Let $S$ be connected and let $(\mathcal E,F)$ be a vector bundle of
rank $h$ in characteristic $p$, with $F:F_S^*\mathcal E\simeq\mathcal E$.
Write $V=\ker(F-1)$ and $\mathcal L=\det\mathcal E$.  Taking
determinants gives an invertible section
\[
 \operatorname{Ha}_F\in
 H^0\bigl(S,\operatorname{Hom}(F_S^*\mathcal L,\mathcal L)\bigr)
 =H^0(S,\mathcal L^{\otimes(1-p)}).
\]
This is the Hasse invariant of the specified Frobenius module.  It
determines the determinant character exactly, including its arithmetic
part: $\det V$ is the rank-one Lang local system
$\ker(F_{\mathcal L}-1)$.  If $e$ trivializes $\mathcal L$ and
$F_{\mathcal L}(e)=h e$, a fixed nonzero vector $ze$ satisfies
\begin{equation}\label{eq:determinant-kummer}
 z^{p-1}=h^{-1}.
\end{equation}
Thus the $\mathbf F_p^\times$-torsor of frames of $\det V$ is the
Kummer torsor \eqref{eq:determinant-kummer}; a deck transformation
$z\mapsto az$ acts on the determinant fibre by $a$.  Changing the
trivialization of $\mathcal L$ changes this equation by a
$(p-1)$st power, as required.  Over $\mathbf F_q$, geometric Frobenius
acts on $\det V$ by the determinant of the $q$-power coherent
Frobenius; arithmetic Frobenius is its inverse.  This description
holds for the stable bundle of the relative vanishing system on its
constant-rank locus as well as for the full ordinary coherent module.
It does not assert that the image of the Kummer character is always
all of $\mathbf F_p^\times$: that depends on its Kummer class.

\subsection{The exact plane-curve monodromy in degrees at most four}
\label{subsec:small-plane-monodromy}

For $d=1,2$, the rank is zero and both monodromy groups are trivial.
For $d=3$, both images are $\mathrm{GL}_1(\mathbf F_p)$: for odd $p$
this follows from the Legendre subfamily and its simple Hasse roots,
proved below; for $p=2$ the assertion means the trivial group on a
one-dimensional space.  A nonzero rank-one representation is
irreducible.

For $d=4$, both images are $\mathrm{GL}_3(\mathbf F_p)$, and the
representation is irreducible.  Here is why the moduli-space theorem
actually applies.  Every nonhyperelliptic smooth genus-three curve is
canonically a smooth plane quartic, and conversely.  The stack quotient
of the parameter space of smooth plane quartics by $\mathrm{PGL}_3$
is the nonhyperelliptic open substack of $\mathcal M_3$.
The map from the parameter space to this quotient is a
$\mathrm{PGL}_3$-torsor in the stack sense, with geometrically connected
fibres, and is surjective on etale fundamental groups: a connected
finite etale cover of the quotient remains connected after pullback
along a torsor under a connected algebraic group.  Restriction to
the ordinary locus preserves this assertion.  That locus is a
nonempty open of the ordinary locus of $\mathcal M_3$, hence its
fundamental group has the same finite monodromy image.
Achter--Pries (2008), Proposition~4.8, says that this image on the
etale $p$-torsion of the Jacobian is $\mathrm{GL}_3(\mathbf F_p)$.
Dualizing gives the same full group on $V$.
The arithmetic image contains the geometric image and is already
a subgroup of the same finite general linear group, so it too is full.
No analogous moduli-space surjectivity assertion is used for degrees
$d\geq5$.  In characteristic two the separate contact-curve argument
and group induction instead determine every degree, in
Theorem~\ref{thm:plane-char2-all-groups}.

\subsection{A universal obstruction in characteristic two}
\label{subsec:odd-plane-obstruction}

Let $p=2$ and let $d\geq3$ be odd.  Write
\[
 f=\sum_{|a|=d}c_aX^{a_1}Y^{a_2}Z^{a_3},\qquad
 I_d=\{u\in\mathbf Z_{>0}^3:|u|=d\}.
\]
Put $c_a=0$ if a coordinate of $a$ is negative.  On a coefficient
chart of $S_d$, use the Cech basis of
$H^2(\mathbf P^2,\mathcal O(-d))$ indexed by $I_d$.
The matrix of the semilinear Frobenius on $H^1(C_f,\mathcal O)$ is
\begin{equation}\label{eq:plane-hw-char2}
 H_{u,v}=c_{2v-u}\quad(u,v\in I_d),\qquad F(x)=Hx^{[2]}.
\end{equation}
Indeed, the connecting map from
$0\to\mathcal O(-d)\xrightarrow{f}\mathcal O\to\mathcal O_{C_f}\to0$
identifies Frobenius with multiplication by $f$ after squaring the
Cech representative $X^{-v_1}Y^{-v_2}Z^{-v_3}$.

Define a column $w=(w_u)$ by
\[
 w_u=\begin{cases}c_u&u_1,u_2,u_3\text{ all odd},\\0&\text{otherwise}.
 \end{cases}
\]
Then the following is an identity over the polynomial ring in the
coefficients, not merely a numerical identity at selected curves:
\begin{equation}\label{eq:universal-covector}
 H^t w=w^{[2]}.
\end{equation}
For its $v$-coordinate, pair the term $c_{2v-u}c_u$ indexed by $u$
with the term indexed by $2v-u$.  A nonzero term has both indices
positive and all odd, so the involution preserves the indexing set.
All two-element orbits cancel in characteristic two.  The only possible
fixed index is $u=v$, contributing $c_v^2$ exactly when $v$ is all odd.

If $H$ is invertible, then $w\ne0$.  Otherwise all coefficients with
three odd exponents vanish, and every row with all-odd index $u$
in \eqref{eq:plane-hw-char2} vanishes.  Such an index exists because
$(1,1,d-2)\in I_d$.  This contradicts invertibility.

For a Lang fixed vector $x$, namely $Hx^{[2]}=x$, equation
\eqref{eq:universal-covector} gives
\[
 (w^tx)^2=(w^{[2]})^tx^{[2]}=w^tHx^{[2]}=w^tx.
\]
Thus $x\mapsto w^tx$ is an $\mathbf F_2$-valued linear functional on
the local system $V=\ker(F-1)$.  It is nonzero, since the Lang fixed
vectors form a basis after an etale cover.  It is a section of
$V^\vee$, so it is fixed by both geometric and arithmetic monodromy.
Working on a dense coefficient chart suffices: its intersection with
the ordinary smooth locus is dense, and the map on fundamental groups
from a nonempty open of a normal integral scheme is surjective.

\begin{theorem}[Refutation of full linear monodromy]\label{thm:odd-plane-char2-obstruction}
\status{Proved under stated hypotheses}
For every finite field of characteristic two and every odd $d\geq5$,
let $\bar\eta$ be a geometric generic point of the ordinary smooth
degree-$d$ plane-curve locus $U_d$, and put
$V=H^1_{\mathrm{et}}(C_{\bar\eta},\mathbf F_2)$.
Both images
\[
 G_{\mathrm{geom}}=\mathrm{im}\bigl(\pi_1(U_d\otimes\bar k,\bar\eta)
 \to\mathrm{GL}(V)\bigr),\qquad
 G_{\mathrm{arith}}=\mathrm{im}\bigl(\pi_1(U_d,\bar\eta)
 \to\mathrm{GL}(V)\bigr)
\]
fix a nonzero covector.  Consequently both are contained in the
proper subgroup
\[
 \mathrm{Stab}(\ell)\simeq
 \mathbf F_2^{\,g_d-1}\rtimes\mathrm{GL}_{g_d-1}(\mathbf F_2),
\]
and $V$ is reducible, with invariant hyperplane $\ker\ell$.
There is no integer $d_1(\mathbf P^2)$ for which the proposed full
linear geometric or arithmetic monodromy holds for every $d\geq d_1$
in characteristic two.
\end{theorem}

The theorem proves containment, not equality with the displayed
stabilizer.  Determining the exact image remains a separate gate.
The determinant character in this characteristic is automatically
trivial, because $\mathbf F_2^\times=\{1\}$.

\par\smallskip\noindent\textbf{Geometric interpretation and literature check.}
The obstruction agrees with the difference between the effective
geometric theta characteristic
$\mathcal O_C((d-3)/2)$ and the canonical characteristic-two theta
characteristic, which is non-effective for an ordinary curve.
Cais, Ellenberg, and Zureick-Brown (2013, Theorem~4.2) prove the resulting
rational two-torsion statement for odd degree plane curves.
Castryck, Streng, and Testa (2014, Corollary~13, Remark~14, and
Theorem~16) give the ordinary and Hasse--Witt formulations.
The explicit identity above avoids any descent ambiguity for theta
characteristics over imperfect coefficient fields.

\par\smallskip\noindent\textbf{Explicit checked fibre.}
Over $\mathbf F_2$, take
\begin{align*}
f={}&x^5+x^4y+x^3y^2+x^2y^3+x^2y^2z+xy^4+xyz^3+xz^4\\
 &+y^4z+y^3z^2+yz^4.
\end{align*}
In the ordered basis
$(1,1,3),(1,2,2),(1,3,1),(2,1,2),(2,2,1),(3,1,1)$, its matrix is
\[
H=\begin{pmatrix}
1&0&0&0&0&0\\1&0&1&0&1&1\\0&1&0&0&0&0\\
1&1&0&0&1&1\\0&0&1&0&1&0\\0&0&0&1&0&0
\end{pmatrix},\qquad w=(1,0,0,0,0,0)^t.
\]
Its rank is six.  The ideal $(f,f_x,f_y,f_z)$ has Groebner basis
$\{1\}$ on each of the three standard projective charts, proving
geometric smoothness.  Direct enumeration over
$\mathbf F_{2^m}$ for $m=1,\ldots,6$ gives respectively
$3,3,9,23,38,33$ points; the traces of $H^m$ modulo two are
$0,0,0,0,1,0$.  The program \texttt{code/plane\_char2.py} and its JSON
output record the input, field arithmetic, certificates, and timings.
The fibre is a check of the formulas; the universal identity, rather
than this one fibre, proves the generic monodromy obstruction.

\subsection{Plane curves: every point-count residue in every open}
\label{subsec:plane-residues}

\begin{theorem}\label{thm:plane-all-residues}
\status{Proved under stated hypotheses}
Let $k$ be a finite field of characteristic $p$.  For every $d\geq4$,
every nonempty Zariski open $U\subset S_d$, and every $a\in\mathbf F_p$,
there exist a finite extension $k'/k$ and a smooth plane curve
$C\in U(k')$ of degree $d$ with $\#C(k')\equiv a\pmod p$.
For $d=1,2$ the only residue for smooth curves is $1$.
For $d=3$, the residues on the ordinary smooth locus are exactly
$\mathbf F_p\setminus\{1\}$, and every nonempty open of that locus
attains every such residue after a finite extension.  Supersingular
smooth cubics give residue $1$, so all residues occur among all smooth
cubics; the every-open assertion for all residues fails in degree three.
\end{theorem}
\begin{proof}

\par\smallskip\noindent\textbf{Elliptic input.}
For every $\lambda\in\mathbf F_p^\times$, there is an ordinary elliptic
curve $E/\mathbf F_Q$, with $\mathbf F_Q/k$ finite, for which geometric
Frobenius on $H^1_{\mathrm{et}}(E_{\bar k},\mathbf F_p)$ is multiplication
by $\lambda$.  For $p=2$ take any ordinary elliptic curve.  For odd $p$,
write $m=(p-1)/2$.  The Legendre family
$y^2=x(x-1)(x-t)$ has Hasse polynomial
\[
 A(t)=(-1)^m\sum_{i=0}^{m}\binom mi^2t^i.
\]
Its roots are simple and avoid $0,1$: its values there are nonzero,
and it satisfies
$t(1-t)A''+(1-2t)A'-m^2A=0$ modulo $p$.
A repeated root away from $0,1$ would, by successive differentiation,
force all derivatives through degree $m<p$ to vanish, a contradiction.
At a simple root, the Kummer cover $z^{p-1}=A(t)^{-1}$ is totally
ramified of degree $p-1$.  Its geometric monodromy is therefore the
whole $\mathbf F_p^\times$.  Finite-field Chebotarev on the ordinary
Legendre locus gives every scalar $\lambda$.
Moreover we may replace $Q$ by $Q^e$ with arbitrarily large
$e\equiv1\pmod{p-1}$ without changing $\lambda$.

\par\smallskip\noindent\textbf{A nodal curve with prescribed residue.}
Put $r=d-3\geq1$ and embed $E$ as a plane cubic by $|3O|$.
Choose rational lines $L_1,\ldots,L_r$ in general position such that
all intersections are transverse, no three components meet, the
intersection $E\cap L_i$ is split for $i\geq2$, and
$E\cap L_1$ has either cycle type $3$ or cycle type $1+2$ under
Frobenius.  Write $t=0$ or $1$ for the number of rational points in
this last intersection.  Such lines exist after an extension of the
above kind, as follows.

Split lines arise from distinct $P,Q,R\in E(\mathbf F_Q)$ with
$P+Q+R=0$, giving order $Q^2$ choices.  Lines of type $1+2$ arise from
$P\in E(\mathbf F_{Q^2})\setminus E(\mathbf F_Q)$ and
$R=-(P+\pi P)\in E(\mathbf F_Q)$; the three points are distinct,
collinear, and Frobenius invariant.  There are order $Q^2$ choices.
For type $3$, use nonrational points in
$\ker(1+\pi+\pi^2)$.  This is the kernel of a separable isogeny,
is contained in $E(\mathbf F_{Q^3})$, and has cardinality
$\#E(\mathbf F_{Q^3})/\#E(\mathbf F_Q)$, of order $Q^2$.
At most nine of these points are rational, since rational points in
the kernel lie in $E[3]$.  Thus there are order $Q^2$ suitable lines
of each required type.  The finitely many incidence and tangency
conditions excluded at each choice remove only order $Q$ rational
lines.  This proves existence for sufficiently large $Q$.

The connected nodal curve
$C_0=E\cup L_1\cup\cdots\cup L_r$ has degree $d$.
Its normalization sequence
\[
0\longrightarrow\mathcal O_{C_0}\longrightarrow
\mathcal O_E\oplus\bigoplus_i\mathcal O_{L_i}
\longrightarrow\bigoplus_{P\text{ node}}k(P)\longrightarrow0
\]
shows that Frobenius on $H^1(C_0,\mathcal O_{C_0})$ is bijective:
it is bijective on the finite reduced node scheme, on constants of
the normalization, and on $H^1(E,\mathcal O_E)$.
Inclusion--exclusion gives
\begin{align*}
 \#C_0(\mathbf F_Q)
 &\equiv (1-\lambda)+r-\{t+3(r-1)\}-\binom r2\\
 &\equiv 4-2r-\binom r2-t-\lambda\pmod p.
\end{align*}
As $t\in\{0,1\}$ and $\lambda\in\mathbf F_p^\times$ vary, these
residues fill $\mathbf F_p$, including when $p=2$.

\par\smallskip\noindent\textbf{Passage to an arbitrary smooth open.}
The ordinary coherent-cohomology locus $W_d\subset S_d$ introduced
above is geometrically integral and contains $C_0$.
The finite etale local system $\ker(F-1)$ is defined over all of $W_d$,
including its singular fibres.  After base changing to $\mathbf F_Q$,
the Frobenius class of $C_0$ belongs to the finite arithmetic
monodromy group.  Chebotarev gives that same class at a closed point
of any nonempty open of $W_d$.  Apply this to
$U\cap W_d\cap S_d^{\mathrm{smooth}}$.
The coherent Frobenius trace formula
\[
 \#C(\mathbf F_{Q^e})\equiv
 1-\operatorname{Tr}\bigl(\operatorname{Frob}_{Q^e}^{\mathrm{geom}},
 H^1_{\mathrm{et}}(C_{\bar k},\mathbf F_p)\bigr)\pmod p
\]
also applies to the proper nodal curve $C_0$, so equality of the
classes gives the desired residue.  More precisely, the class at
$C_0$ maps to degree one in the finite constant-field quotient, and
Chebotarev supplies it at degrees congruent to one in that quotient;
the permitted residue fields are finite extensions of $k$.
This argument never assumes full general linear monodromy for
degree-$d$ plane curves.
\end{proof}

\par\smallskip\noindent\textbf{Primary-source verification of the two inputs.}
Katz--Sarnak (1999), Theorem~9.7.13, p.~290, is the finite-monodromy
equidistribution theorem, with the constant-field coset condition made
explicit in 9.7.1--9.7.9.  To apply its characteristic-zero formulation
to our finite $\mathbf F_p$-representation, take a faithful finite
dimensional characteristic-zero representation of its finite image
(for example the regular representation), choose $\ell\ne p$, and
regard it as a lisse $\overline{\mathbf Q}_\ell$-sheaf.  It is pure of
weight zero, and its finite geometric and arithmetic groups have the
required normality.  The theorem gives a positive limiting proportion
for every nonempty conjugacy class in the required coset.  Restriction
to a nonempty open preserves the finite monodromy image because the
base is normal.  Katz (1973), SGA~7~II, Expose~XXII, Corollary~3.2,
formulas (3.2.1)--(3.2.2), pp.~410--411, gives precisely the trace
congruence for every proper scheme over a finite field, including
our nodal curve.  Proposition~2.0 and formula (2.2.4) of the same
expose identify its coherent and etale formulations.

\subsection{An explicit positivity threshold for every smooth ambient variety}
\label{subsec:explicit-positivity}

Here $X\subset\mathbf P^N$ is smooth projective of pure dimension
$n\geq2$ over a perfect field of characteristic $p$, and $L=\mathcal O_X(1)$.
Geometric connectedness is not needed for this proposition.  All
dimensions in the following definitions may be computed after extending
the ground field to an algebraic closure.  Put
\[
 D=\deg_L X,\qquad
 a=\dim H^{n-1}(X,\mathcal O_X)_{\mathrm{ss}},
\]
and, for $1\leq m\leq n-2$, define the explicit cohomological bounds
\[
 B_m=\sum_{j=0}^{m}\binom mj
       h^{n-m-1+j}(X,L^{-j}),\qquad
 B=D+a+\sum_{m=1}^{n-2}B_m.
\]
Let $r_\omega$ be any Castelnuovo--Mumford regularity bound for
$i_*\omega_X$ on $\mathbf P^N$, and put
\begin{equation}\label{eq:explicit-d0}
 d_0=\max\left\{2n,1,r_\omega-1,
       \left\lfloor\frac{2^n n!B}{D}\right\rfloor+1\right\}.
\end{equation}
Every term is determined by the embedding and finitely many coherent
cohomology groups of $X$; this is a computable bound, not an unspecified
``sufficiently large'' constant.

\begin{theorem}\label{thm:all-X-explicit-positive-rank}
\status{Proved under stated hypotheses}
For every $d\geq d_0$, the generic smooth hypersurface section $Y$ of
degree $d$ has
$\dim E(Y,\mathcal O_Y)_{\mathrm{ss}}>0$.
More precisely, for every $d\geq\max\{1,r_\omega-1\}$, its generic
rank $f_{X,d}$ satisfies
\begin{equation}\label{eq:generic-rank-lower-bound}
 f_{X,d}\ \geq\
 D\binom dn-D\binom d{n-1}
 -\sum_{m=1}^{n-2}B_m\binom dm-a,
\end{equation}
where a binomial coefficient with lower index greater than $d$ is zero.
Consequently
\[
 f_{X,d}=\frac{D}{n!}d^n+O_X(d^{n-1})\quad(d\longrightarrow\infty).
\]
\end{theorem}

\begin{proof}
Take $d$ hyperplanes in simultaneous general position over the
algebraic closure and set $Y_0=\bigcup_{i=1}^d(X\cap H_i)$.
Every intersection $Z_I=X\cap\bigcap_{i\in I}H_i$ with $|I|\leq n$
is smooth of dimension $n-|I|$, and intersections of $n+1$ components
are empty.  Intersections with $|I|=n$ consist of $D$ reduced points.
The alternating normal-crossings resolution gives the Frobenius
equivariant spectral sequence
\[
 E_1^{s,t}=\bigoplus_{|I|=s+1}H^t(Z_I,\mathcal O_{Z_I})
 \ \Longrightarrow\ H^{s+t}(Y_0,\mathcal O_{Y_0}).
\]
The term $E_1^{n-1,0}$ has dimension $D\binom dn$ and bijective
Frobenius.  It has no outgoing differential on any page.  The image
of the incoming $d_1$ has dimension at most
$D\binom d{n-1}$: each smooth positive dimensional intersection has
at most $D$ connected components, since each has positive integral
degree.  For $r=2,\ldots,n-1$, the possible incoming $d_r$ has source
$E_r^{n-1-r,r-1}$.  Writing $m=n-r$, its dimension is at most
$\binom dm B_m$.  Indeed, the Koszul resolution of a transverse
$m$-fold hyperplane section bounds
\[
 h^{n-m-1}(Z_I,\mathcal O_{Z_I})
 \leq\sum_{j=0}^{m}\binom mj
 h^{n-m-1+j}(X,L^{-j})=B_m.
\]
Every quotient of the top term by an incoming image still has
bijective Frobenius: a Frobenius-stable subspace of a finite dimensional
bijective semilinear space over a perfect field itself has bijective
Frobenius.  Hence $E_\infty^{n-1,0}$ is a bijective-Frobenius subquotient
of $H^{n-1}(Y_0,\mathcal O_{Y_0})$ of dimension at least the first three
terms of the right side of \eqref{eq:generic-rank-lower-bound}.
Exactness of the stable-part functor gives the same lower bound for
$\dim H^{n-1}(Y_0,\mathcal O_{Y_0})_{\mathrm{ss}}$.

For $d\geq\max\{1,r_\omega-1\}$, regularity and Serre duality give
$H^{n-1}(X,L^{-d})=H^1(X,\omega_X\otimes L^d)^\vee=0$.
The divisor sequence, including for the normal-crossings divisor,
then makes $h^{n-1}(Y,\mathcal O_Y)$ independent of $Y$ in the space of
properly intersecting degree-$d$ sections: the terminal map
$H^n(X,L^{-d})\to H^n(X,\mathcal O_X)$ is surjective because
$H^n(Y,\mathcal O_Y)=0$.  The universal family is
flat, and cohomology and base change therefore gives a vector bundle
with Frobenius.  If its rank is $h$, the stable rank at a perfect
geometric point is the rank of $F^h$; thus its maximal value is attained
on an open, and the generic value is at least the value at $Y_0$.
The smooth locus is dense and meets this open.

For a smooth $Y$, R1 identifies the stable ambient image with
$H^{n-1}(X,\mathcal O_X)_{\mathrm{ss}}$, of dimension $a$; in fact the
coherent restriction is injective in the present range of $d$.
Exactness of the stable-part functor subtracts $a$ and proves
\eqref{eq:generic-rank-lower-bound}.
For $d\geq2n$ one has
$\binom dn\geq(d/2)^n/n!$ and $\binom dm\leq d^{n-1}$ for $m<n$.
The right side of \eqref{eq:generic-rank-lower-bound} is therefore at
least $d^{n-1}(Dd/(2^nn!)-B)$, which is positive under
\eqref{eq:explicit-d0}.  Finally, Serre duality, Riemann--Roch, and
Serre vanishing give
$\dim E(Y,\mathcal O_Y)=Dd^n/n!+O_X(d^{n-1})$; together with the lower
bound this proves the asymptotic formula.
\end{proof}

\begin{corollary}[The eventual parenthesis of 8.5.8]
\label{cor:eventual-hodge-parenthesis}
\status{Proved under stated hypotheses}
For the explicit $d_0$ of \eqref{eq:explicit-d0}, every $d\geq d_0$
has the following properties.  Every smooth degree-$d$ section $Y$
satisfies
\[
 \dim E(Y,\mathcal O_Y)
 =h^0(X,\omega_X\otimes L^d)-h^0(X,\omega_X)>0.
\]
For every $Y$ in the maximal-stable-rank open $U_d$, the number of
$p$-adic unit roots in vanishing cohomology is positive.  Thus, on
this open and in this stated eventual range, unit-root nonvanishing
and the maximal Hodge-level condition $E(Y,\mathcal O_Y)\ne0$ are
equivalent.  No assumption on crystalline $p$-torsion of $X$ is needed.
\end{corollary}

\begin{proof}
The coherent divisor sequence and the vanishings used in the preceding
proof identify $E(Y,\mathcal O_Y)$ with
$\ker(H^n(X,L^{-d})\to H^n(X,\mathcal O_X))$, whose dimension is the
displayed, $Y$-independent difference by Serre duality.  It is positive
because it is at least the positive generic stable rank of
Theorem~\ref{thm:all-X-explicit-positive-rank}.
For $Y\in U_d$, Theorems~\ref{thm:A-exact85} and
\ref{thm:A-unit-count} identify the unit-root count with
\[
 f_{X,d}+\dim H^n_{\mathrm{cris}}(X/W)[p]_{\mathrm{ss}},
\]
which is positive.  Both properties therefore hold for every point
under these quantifiers.
\end{proof}

\par\smallskip\noindent\textbf{Scope of the manuscript's parenthesis.}
The supplied transcription of cote~4, scan~8, first specifies a
section for which the semisimple rank has its maximum value.  Scan~9,
continuing 8.5.8, then says ``pour tous ces $Y$'' and, equivalently,
``pour $Y$ g\'en\'erique''.  Corollary~\ref{cor:eventual-hodge-parenthesis}
answers this eventual generic assertion, including varieties with
$p$-torsion.  It does not assert that an arbitrary smooth section
outside $U_d$ has a unit root whenever its coherent vanishing space
is nonzero.  For example, a supersingular smooth plane cubic has
one-dimensional $H^1(\mathcal O)$ and no unit root.  That example
distinguishes the stronger pointwise assertion from the quantifiers
actually written in 8.5.8; it is not substituted for the requested
eventual assertion.

\par\smallskip\noindent\textbf{Relation to Koblitz.}
Koblitz (1975), I.\S5, Theorems~1--2, pp.~139--141, proves the
asymptotic invertibility estimate with an error of order $d^{n-1}$.
The argument above gives explicit constants in terms of the embedding
and the displayed cohomology groups.  Koblitz's I.\S6, Theorem~3,
p.~142, concerns generic complete intersections, not every hypersurface
section of an arbitrary fixed ambient variety.  His revised conjecture
in II.\S3, pp.~149--151, must not be cited as an established exact rank
formula.

\subsection{All residues for hypersurfaces in every projective space}

\begin{theorem}\label{thm:projective-all-residues}
\status{Proved under stated hypotheses}
Let $k$ be finite of characteristic $p$, let $n\geq2$, and let
$d\geq pn+1$.  Every nonempty open in the degree-$d$ hypersurface
parameter space of $\mathbf P^n_k$ contains, over a finite extension,
a smooth hypersurface with any prescribed point-count residue
modulo $p$.
\end{theorem}

\begin{proof}
Choose the integer representative $t\in\{0,\ldots,p-1\}$ of the
desired residue.  Over a sufficiently large finite extension
$\mathbf F_Q/k$, choose a general-position collection of $d$
geometric hyperplanes on which Frobenius has $t$ orbits of length $n$
and one orbit of length $d-nt>n$.  Existence follows by taking a
rational point in the open set of general-position configurations
in the product of Weil restrictions of the dual projective space
from the indicated extension degrees.  This parameter space is
geometrically integral, since over the algebraic closure it is a
product of projective spaces; the general-position open is nonempty.
Lang--Weil supplies rational points over arbitrarily large fields.
One may first form these Weil restrictions over a fixed finite base
field and then take extension degrees congruent to one modulo the
least common multiple of the prescribed orbit lengths; this preserves
their field factors while making the field arbitrarily large.
Distinctness of every geometric hyperplane forces the specified
orbit lengths to be exact.

The union $Y_0$ of these hyperplanes is defined over $\mathbf F_Q$.
A rational point lying on one hyperplane lies on its entire Frobenius
orbit.  An orbit of length $n$ has exactly one common intersection
point in $\mathbf P^n$, and that point is rational.  An orbit of
length greater than $n$ has empty intersection.  The $t$ points from
the shorter orbits are distinct, because any $n+1$ distinct
hyperplanes have empty intersection.  Therefore
\[
 \#Y_0(\mathbf F_Q)=t
\]
as an equality of integers.
The general-position hyperplane-union argument proves that
$H^{n-1}(Y_0,\mathcal O_{Y_0})$ has bijective Frobenius, of the full
generic rank $\binom{d-1}{n}$.  Thus $Y_0$ belongs to the same open
$W_d$ on which the entire Lang local system is defined.
The coherent trace congruence, together with the standard vanishing
of intermediate coherent cohomology for a projective hypersurface,
reads
\[
 \#Y(\mathbf F_{Q^e})\equiv
 1+(-1)^{n-1}\operatorname{Tr}(\operatorname{Frob}_{Q^e}^{\mathrm{geom}},V_Y)
 \pmod p.
\]
Apply finite Chebotarev on $W_d$ to the Frobenius class of $Y_0$ and
the given smooth open.  The resulting smooth hypersurface has the
same residue $t$.  No exact monodromy-group computation is required.
\end{proof}

\subsection{Relation to the subsequent classifications}
The odd-degree characteristic-two obstruction refutes the universal
full-$\mathrm{GL}$ expectation, even for geometrically connected
$X=\mathbf P^2$.  The all-degree characteristic-two theorem identifies
the exact geometric and arithmetic images: full $\mathrm{GL}_g(2)$
for even $d\ge4$, and the full stabilizer
$2^{g-1}\!:\mathrm{GL}_{g-1}(2)$ for odd $d\ge5$.
The subsequent projective-space theorems settle every remaining
plane degree and every dimension, in all characteristics.
For arbitrary $X$, the explicit asymptotic positivity bound above is
proved using a normal-crossings union of hyperplane sections.
The later all-ambient theorem also proves every point-count residue
with an explicit degree bound. An effective finite generic-rank
calculation and exact determinant image are provided below.
A structural generic-rank formula and the complete nonabelian
monodromy group for arbitrary $X$ remain unresolved here.

% /root/monodromy, 19 November 2026.
% Complete mathematical construction; independently audited by /root.
\subsection{A smooth ambient surface whose relative Hasse--Witt map is always singular}
\label{subsec:relative-ordinarity-counterexample}
\status{Proved; independently audited by /root.  The finite Hopf-algebra checks
were executed; the final invariant-ring/Bertini search was not executed.}

The following construction refutes eventual bijectivity on
$E(Y,\mathcal O_Y)$ for arbitrary smooth ambient varieties.  The
ambient variety is smooth.  In particular, this does not reuse the
singular cones in Koblitz's counterexamples.

\begin{lemma}\label{lem:picard-finite-quotient}
\status{Proved under stated hypotheses}
Let $G$ be a finite commutative group scheme over a field $k$, and
let $Z\to X$ be a $G$-torsor, where $Z$ is a positive dimensional
projective complete intersection of dimension at least two and $X$
is smooth and geometrically connected.  If $G^D$ is connected, then
\[
 \operatorname{Pic}_{X/k}^0=G^D.
\]
\end{lemma}

\begin{proof}
The Koszul resolution gives $H^0(Z,\mathcal O_Z)=k$ and
$H^1(Z,\mathcal O_Z)=0$, compatibly with every extension of scalars;
smoothness or reducedness of $Z$ is unnecessary.  Therefore the
identity component of $\operatorname{Pic}_{Z/k}$ is trivial: its
Zariski tangent space is zero.  The kernel, as an fppf sheaf, of
\[
 \operatorname{Pic}_{X/k}\longrightarrow\operatorname{Pic}_{Z/k}
\]
is $G^D$.  Indeed, after trivializing the pullback of a line bundle,
its descent datum is a $G$-linearization of $\mathcal O_Z$.
All functions on $G\times Z$ come from $G$, because
$H^0(Z_R,\mathcal O)=R$ for every $k$-algebra $R$.  The descent
cocycle condition is precisely the condition that this unit be a
character $G_R\to\mathbf G_{m,R}$.  This proves the kernel assertion
without a surjectivity claim about the entire Picard scheme.
The connected group $\operatorname{Pic}_{X/k}^0$ maps to the trivial
identity component of $\operatorname{Pic}_{Z/k}$ and hence lies in
$G^D$.  Conversely the connected kernel $G^D$ lies in
$\operatorname{Pic}_{X/k}^0$.
\end{proof}

\par\smallskip\noindent\textbf{A specified finite group scheme and linear action.}
Work first over $\mathbf F_2$ and put
\[
 G=\operatorname{Spec}\mathbf F_2[a,b]/(a^2,b^2),\qquad
 \Delta(a)=a\otimes1+1\otimes a,\quad
 \Delta(b)=b\otimes1+1\otimes b+a\otimes a.
\]
Its group law is
$(a,b)(c,e)=(a+c,b+e+ac)$.  The basis $1,a,b,ab$ is the basis dual
to $1,t,t^2,t^3$ in the Hopf algebra of $\alpha_4$; the multiplication
rule is the divided-power rule
$f_if_j=\binom{i+j}{i}f_{i+j}$, truncated at index four.  Thus
$G=(\alpha_4)^D$.  Its Frobenius is zero, so
$V_{\alpha_4}=D(F_G)=0$, by Cartier duality for Verschiebung
\cite[Proposition~5.19]{EdixhovenVanDerGeerMoonenAV}.

Let $J$ be the direct sum of two nilpotent Jordan blocks of size
four on $\mathbf F_2^8$, with $J e_i=e_{i+1}$ in each block.
The formula
\begin{equation}\label{eq:alpha4-regular-action}
 (a,b)\longmapsto I+aJ+bJ^2+abJ^3
\end{equation}
is a faithful representation on $\mathbf A^8$ and acts on
$\mathbf P^7$.  Multiplication of these matrices gives exactly the
displayed group law.
The only subgroup of $G$ of order two is
$H=\{a=0\}\simeq\alpha_2$, and every nontrivial subgroup contains it.
One may see this directly from the restricted Lie algebra, with
basis $J,J^2$ and square operation $J^{[2]}=J^2$, $(J^2)^{[2]}=0$:
its only one-dimensional restricted subalgebra is $\langle J^2\rangle$.
Since $G$ has height one, its subgroup schemes correspond to these
restricted Lie subalgebras.  A projective line $\langle v\rangle$
is fixed by $H$ precisely when $J^2v=0$; a nilpotent operator has
no nonzero eigenvalue.  Consequently the nonfree locus is exactly
\[
 Z_{\rm bad}=\mathbf P(\ker J^2)=\mathbf P^3\subset\mathbf P^7.
\]
In particular, it has codimension four.

\par\smallskip\noindent\textbf{The surface and its fixed embedding.}
Let $Q=\mathbf P^7/G$ be the projective quotient and let
$Q^\circ=(\mathbf P^7\setminus Z_{\rm bad})/G$.  The quotient exists,
the map $\mathbf P^7\to Q$ is finite, and $Q^\circ$ is smooth;
these are Kothari's Theorem~3.3, p.~6
\cite{Kothari2026}.  The image of $Z_{\rm bad}$ has dimension three.
Fix a projective embedding $Q\hookrightarrow\mathbf P^M$ supplied
by a sufficiently divisible invariant Veronese ring.  Five general
hyperplanes of $\mathbf P^M$ cut out a smooth surface
\[
 X=Q\cap H_1\cap\cdots\cap H_5\subset Q^\circ.
\]
Their intersection avoids the three-dimensional bad locus and is
smooth by Bertini on $Q^\circ$.  The five pulled-back equations cut
out a surface $Z\subset\mathbf P^7$ which is a complete intersection;
the equality of the number of equations and the codimension makes
them a regular sequence in the polynomial ring.  The action on $Z$
is free and $Z/G=X$.  Geometric connectedness of $Z$, and then of
$X$, follows from the Koszul computation $H^0(Z,\mathcal O_Z)=k$.
This is also the particular dimension-two case of Kothari's
Proposition~3.4, pp.~6--7, and Corollary~3.7, p.~9; those statements
explicitly require smoothness of the quotient, not of $Z$.

This construction can be made unambiguous over a specified finite
field without leaving an unspecified ``general'' choice.  Use the
given ordered coordinates of $\mathbf P^7$ and their graded
lexicographic monomial order.  Compute the homogeneous invariant
ring by solving the coaction equations.  Every coordinate square is
invariant, because $a^2=b^2=0$, so this invariant ring is a finite
module over $\mathbf F_2[x_0^2,\ldots,x_7^2]$.  Its computation is
therefore a finite module-kernel calculation over an explicit
polynomial ring, rather than an unbounded search through degrees.
Take the least positive
Veronese degree which is generated in degree one.  Use its reduced
row-echelon monomial basis to embed $Q$.  Enumerate finite extensions
of $\mathbf F_2$ in increasing degree, each presented by the first
monic irreducible polynomial in lexicographic order, and then
enumerate ordered five-tuples of hyperplanes.  Retain the first
tuple for which the intersection is two dimensional, smooth, and
disjoint from the bad locus.  These conditions are decidable by
elimination and the Jacobian criterion.  The preceding nonempty
Bertini open proves termination.  Fix this output field $k$, this
surface $X$, and the resulting very ample bundle
$L=\mathcal O_{\mathbf P^M}(1)|_X$.
The description therefore defines one embedded surface by a finite
algorithm; the proof applies to every tuple passing these tests.
We do not claim to have numerically executed the invariant-ring and
Bertini search to identify the first tuple.

Lemma~\ref{lem:picard-finite-quotient} gives
\begin{equation}\label{eq:alpha4-picard-surface}
 \operatorname{Pic}_{X/k}^0=\alpha_4,\qquad
 h^1(X,\mathcal O_X)=1.
\end{equation}
Frobenius on this one-dimensional coherent cohomology is zero.
For a direct verification, the Cartan--Leray spectral sequence for
$Z\to X$ and $H^1(Z,\mathcal O_Z)=0$ identifies
$H^1(X,\mathcal O_X)$ with $H^1(G,k)$, the primitive elements of
$\mathcal O(G)$.  Those primitives are exactly $k a$; $b$ has the
additional term $a\otimes a$ in its coproduct, and $ab$ has the
terms $a\otimes b+b\otimes a$.  Absolute Frobenius sends $a$ to zero.

\begin{theorem}\label{thm:relative-ordinarity-fails-smooth}
\status{Proved under stated hypotheses}
For the fixed smooth projective geometrically connected surface
$(X,L)/k$ just constructed, let $r_\omega$ be a regularity bound for
its embedded dualizing sheaf, and put
$d_* =\max\{1,r_\omega-1\}$.
For every $d\geq d_*$ and every smooth section $Y\in|L^d|$, of
genus $g_Y$, one has
\[
 \dim E(Y,\mathcal O_Y)=g_Y-1,
 \qquad
 \dim E(Y,\mathcal O_Y)_{\mathrm{ss}}
       =f_Y\leq g_Y-2,
\]
where $f_Y$ is the $2$-rank of its Jacobian.  In particular the
Frobenius on $E(Y,\mathcal O_Y)$ is never bijective in this entire
range, including at the generic point.
\end{theorem}

\begin{proof}
Regularity and Serre duality give
$H^1(X,L^{-d})=H^1(X,\omega_X\otimes L^d)^\vee=0$.
Thus $H^1(X,\mathcal O_X)\to H^1(Y,\mathcal O_Y)$ is injective.
The induced homomorphism of group schemes
$\alpha_4=\operatorname{Pic}_{X/k}^0\to J_Y$ has zero Lie kernel.
Its finite kernel is therefore etale, but a subgroup of the local
group scheme $\alpha_4$ is local.  The kernel is trivial, and this
finite monomorphism is a closed immersion.

Naturality of Verschiebung and $V_{\alpha_4}=0$ put
$\alpha_4^{(2)}$ inside $\ker(V_{J_Y}:J_Y^{(2)}\to J_Y)$.
For an abelian variety of dimension $g_Y$ and $2$-rank $f_Y$, this
kernel has order $2^{g_Y}$ and its etale quotient has order $2^{f_Y}$;
its connected component therefore has order $2^{g_Y-f_Y}$.
Here the degree and point-count assertions are respectively
Propositions~5.20 and~5.22 of
\cite{EdixhovenVanDerGeerMoonenAV}: Frobenius is bijective on geometric
points, so $V\circ F=[2]$ identifies the geometric points of $\ker V$
with $J_Y[2](\bar k)$.
The connected subgroup $\alpha_4^{(2)}$ has order four, so
$g_Y-f_Y\geq2$.
Equation~\eqref{eq:alpha4-picard-surface} and coherent injectivity
give the dimension $g_Y-1$.  Since the ambient Frobenius is zero,
exactness of stable parts gives
$\dim E(Y,\mathcal O_Y)_{\mathrm{ss}}=f_Y$.
This proves the strict inequality and failure of bijectivity.
\end{proof}

\par\smallskip\noindent\textbf{What is and is not determined.}
This is a smooth counterexample to the universal eventual-relative-
ordinarity assertion.  It is compatible with the explicit positive
generic-rank theorem: for still larger degrees its generic rank is
positive and at least one less than its coherent dimension.  The
argument gives an exact structural obstruction and the displayed
upper bound; it does not assert an exact generic $2$-rank for this
surface.  The exact generic value for arbitrary $X$ remains a
separate determination.

\par\smallskip\noindent\textbf{Retained finite verification.}
The program \texttt{code/odd\_plane\_and\_alpha4\_checks.py} verifies
the Hopf group law, associativity, the matrix representation, its
divided-power multiplication table, the one-dimensional primitive
space and its zero Frobenius, and $\operatorname{rank}(J^2)=4$ on
the eight-dimensional representation.  The output is retained in
\texttt{outputs/odd\_plane\_and\_alpha4\_checks.json}.  These are
checks of the finite algebraic construction; existence and
smoothness of the surface are proved by the specified quotient and
Bertini construction, not inferred from a sample point calculation.

% Exact quintic case; independently audited by parent /root.
\subsection{The exact group for ordinary plane quintics in characteristic two}
\label{subsec:quintic-exact}

\begin{theorem}\label{thm:quintic-char2-monodromy}
\status{Proved under stated hypotheses}
For the universal ordinary smooth plane quintic in characteristic two,
the geometric and arithmetic monodromy groups on
$V=H^1_{\mathrm{et}}(C_{\bar\eta},\mathbf F_2)$ are both the full
stabilizer of the canonical nonzero covector:
\[
 G_{\mathrm{geom}}=G_{\mathrm{arith}}
 =\operatorname{Stab}_{\mathrm{GL}_6(\mathbf F_2)}(\ell)
 \simeq\mathbf F_2^5\rtimes\mathrm{GL}_5(\mathbf F_2).
\]
In particular their order is $319979520$, the representation is
reducible with invariant hyperplane of dimension five, and the
determinant character is trivial.
\end{theorem}
\begin{proof}

\par\smallskip\noindent\textbf{Theta-characteristic input, with a proof.}
For a smooth curve over an algebraically closed field of characteristic
two, let $B=F_*\mathcal O_C/\mathcal O_C$ for absolute Frobenius.
This is a line bundle with $B^{\otimes2}\simeq\omega_C$.
Locally a separating parameter $t$ gives the basis $1,t$ of
$F_*\mathcal O_C$ over $\mathcal O_C$; a change of separating parameter
$u=a^2+tb^2$ sends the generator of $B$ to $b$ times that generator,
whereas $du=b^2dt$.  This proves the claimed square relation.
The long exact sequence of
$0\to\mathcal O_C\to F_*\mathcal O_C\to B\to0$ shows
\[
 H^0(C,B)=\ker\bigl(F:H^1(C,\mathcal O_C)\to H^1(C,\mathcal O_C)\bigr).
\]
Thus $B$ is non-effective for an ordinary curve.
Every other theta characteristic is $B\otimes T$ for a nontrivial
two-torsion line bundle $T$.  Tensoring this sequence by $T$ and using
the projection formula and $F^*T=T^{\otimes2}\simeq\mathcal O_C$
gives an injection
\[
 H^0(C,\mathcal O_C)=H^0(C,F_*F^*T)
 \lhook\joinrel\longrightarrow H^0(C,B\otimes T),
\]
because $H^0(C,T)=0$.  Hence every theta characteristic other than $B$
is effective.  This is the statement recorded as St\"ohr--Voloch's
theorem in Ishitsuka--Ito, Theorem~3.1; Mumford (1971), p.~191,
describes the canonical characteristic-two theta characteristic.
For the geometric generic curve, the same reasoning may be performed
after perfecting its coefficient field.  Perfection does not change
the absolute Galois group, and $B$ is functorial under that group.

An ordinary curve of genus six has exactly $2^6=64$ geometric theta
characteristics.  On a plane quintic, $L=\mathcal O_C(1)$ is one of
them and $h^0(L)=3$, while $B$ is the unique non-effective one.
Under the identification of their set with the geometric points of
$J[2]$ by $T\mapsto L\otimes T$, the two distinguished points are
$0$ and $v=B\otimes L^{-1}\ne0$.  This is the dual representation
$V^\vee$, and $v$ is fixed by monodromy.

\par\smallskip\noindent\textbf{An irreducible incidence correspondence of contact conics.}
Let $S=\mathbf P(H^0(\mathbf P^2,\mathcal O(5)))$, of dimension twenty.
Let $I$ parametrize pairs $(C,Q)$ where $Q$ is a smooth plane conic
and the divisor cut out by $C$ on $Q$ is twice a divisor of degree five.
For a fixed conic $Q\simeq\mathbf P^1$, this means that the restriction
of its quintic equation, a binary form of degree ten, is a square.
The square forms constitute a vector subspace of dimension six inside
$H^0(Q,\mathcal O_Q(5))$, of dimension eleven.  The restriction map
from ternary quintics is surjective, with kernel
$Q\cdot H^0(\mathbf P^2,\mathcal O(3))$ of dimension ten.
Consequently the allowed quintic equations form a vector space of
dimension sixteen, or a projective space of dimension fifteen.
The smooth-conic parameter space has dimension five, so $\dim I=20$.
It is irreducible: fix $Q_0=(XZ+Y^2=0)$, take the irreducible product
$\mathrm{PGL}_3\times\mathbf P^{15}$ for its allowed quintics, and
map it surjectively to $I$ by changing projective coordinates.

The locus of pairs with a reduced reducible conic has dimension at
most nineteen.  Indeed, write $Q=L_1\cup L_2$ and $P=L_1\cap L_2$.
If $C$ is smooth and $C\cap Q$ is even, then $P\in C$: away from $P$
each of the two degree-five line sections must have even multiplicity,
so its multiplicity at $P$ is odd.  On each line the restriction is
therefore its linear equation for $P$ times the square of a binary
quadratic.  Each such restriction space has dimension three, and the
two restrictions automatically agree with value zero at $P$.
The combined restriction space has dimension six, so the quintic
equations again form a sixteen-dimensional vector space.
Pairs of distinct lines have dimension four, proving the bound
$4+15=19$.  Double-line conics cut exactly the divisors corresponding
to $L=\mathcal O_C(1)$.

Every effective theta characteristic $\theta$ and nonzero section $s$
gives a contact conic: $s^2\in H^0(C,\omega_C)=H^0(C,\mathcal O_C(2))$
lifts uniquely to a quadratic form on $\mathbf P^2$.
If that conic is a double line, then $\theta=L$.
Outside the dimension-at-most-nineteen image just described, each of
the other sixty-two theta characteristics therefore supplies a smooth
contact conic.  Thus $I\to S$ is dominant.  Since both dimensions are
twenty, its geometric generic fibre is zero dimensional.
It follows also that every such theta characteristic has $h^0=1$:
otherwise its projective space of sections would give a positive
dimensional family of conics in the same fibre.  Distinct theta
characteristics give distinct conics, because the half of an even
divisor on a smooth curve is unique.  The geometric generic fibre
therefore has exactly sixty-two points.

Irreducibility of $I$ implies transitivity of the geometric Galois
group on these sixty-two points.  This remains valid if the incidence
map has a purely inseparable part: the number of geometric generic
points is its separable degree, and a field's embeddings into a
separable closure form one orbit.  The Galois-equivariant
theta-characteristic interpretation proves that
$G_{\mathrm{geom}}$ on $V^\vee\setminus\{0,v\}$ is transitive.

\par\smallskip\noindent\textbf{The factor seven supplied by a quartic boundary.}
Take an ordinary smooth plane quartic $C_4$ and a line meeting it
transversely in four distinct points.  Their union $C_0$ is an ordinary
nodal quintic: the normalization sequence gives bijective Frobenius
on its $H^1(\mathcal O)$, of dimension six.  Over the parameter space
of such pairs, the etale normalization sequence gives a surjection
of local systems
\[
 H^1_{\mathrm{et}}(C_0,\mathbf F_2)
 \longrightarrow H^1_{\mathrm{et}}(C_4,\mathbf F_2).
\]
Forgetting the line is a smooth surjection with geometrically
connected fibres (nonempty opens of the dual plane), hence is
surjective on fundamental groups.  The quartic local system has
geometric monodromy $\mathrm{GL}_3(\mathbf F_2)$ by the already proved
quartic case.  Its order is divisible by seven, so the monodromy on
the nodal quintic family, and consequently the generic quintic group,
has order divisible by seven.  This specialization is legitimate
because the entire rank-six Lang system extends across $C_0$ on the
invertible-Hasse--Witt locus $W_5$.  The dense smooth open and $W_5$
have the same monodromy image for that extended finite etale system.

\par\smallskip\noindent\textbf{Finite group conclusion.}
Write $M=\operatorname{Stab}(v)\simeq2^5\!:\!\mathrm{GL}_5(2)$
in its action on $V^\vee$, and let $G\subset M$ be the geometric image.
Its transitive orbit of size sixty-two implies $31\mid|G|$; the
boundary argument gives $7\mid|G|$.
The complete maximal-subgroup list for $M$ is given in
Basheer--Moori (2021), Table~1, p.~3:
\[
\begin{array}{c|r|r}
\text{maximal subgroup}&\text{order}&\text{index in }M\\\hline
2_+^{1+8}:\mathrm{GL}_4(2)&10321920&31\\
2^{4+5}:\mathrm{GL}_4(2)&10321920&31\\
\mathrm{GL}_5(2)&9999360&32\\
2^{2+9}:(\mathrm{GL}_3(2)\times S_3)&2064384&155\\
2^9:(\mathrm{GL}_3(2)\times S_4)&2064384&155\\
(2^5:31):5&4960&64512.
\end{array}
\]
The first, second, fourth, and fifth orders are not divisible by
thirty-one.  The last is not divisible by seven.  The remaining
maximal subgroup is a complement $\mathrm{GL}_5(2)$; its action on
$V^\vee\setminus\{0,v\}$ has two orbits of length thirty-one, so it
cannot contain a group transitive on all sixty-two points.
No proper maximal subgroup can contain $G$, and therefore $G=M$.
The arithmetic image is contained in $M$ and contains $G$, hence
equals it.
\end{proof}

\par\smallskip\noindent\textbf{Independent elementary check of the complement issue.}
The same conclusion can be obtained from the maximal-subgroup table
for $\mathrm{GL}_5(2)$ in the ATLAS and the elementary fact
$H^1(\mathrm{GL}_m(2),\mathbf F_2^m)=0$ for $m\geq4$.
For completeness, let $c$ be a cocycle and put
$t_{ij}=1+E_{ij}$.  The relation $t_{ij}^2=1$ says that the $j$th
coordinate of $c(t_{ij})$ is zero.  Commuting $t_{ij}$ with
$t_{k\ell}$ for an auxiliary index $k$ distinct from $i,j,\ell$
shows that every coordinate except the $i$th is zero.
Thus $c(t_{ij})=a_{ij}e_i$.  The relation
$t_{ij}t_{jk}t_{ij}=t_{jk}t_{ik}$, for distinct indices,
gives $a_{ik}=a_{jk}$, so $a_{ij}=b_j$ depends only on $j$.
The coboundary associated with $b=(b_j)$ agrees with $c$ on all
elementary transvections, which generate the group.  Hence $c$ is
a coboundary.  The same result holds for the dual module by the
inverse-transpose automorphism.  Complements in $2^5:\mathrm{GL}_5(2)$
are therefore conjugate to the standard one, justifying the asserted
two-orbit description independently of the affine maximal-subgroup
classification.

% Exact characteristic-two plane sextic case; proof for independent audit.
\subsection{The exact group for ordinary plane sextics in characteristic two}
\label{subsec:sextic-exact}

\begin{theorem}\label{thm:sextic-full-monodromy}
\status{Proved under stated hypotheses}
For the universal ordinary smooth plane sextic in characteristic two,
the geometric and arithmetic monodromy images on
$V=H^1_{\mathrm{et}}(C_{\bar\eta},\mathbf F_2)$ are both
$\mathrm{GL}_{10}(\mathbf F_2)$.  Thus $V$ is irreducible and its
determinant character is trivial.
\end{theorem}

\begin{proof}
The generic genus is ten, and ordinarity gives $2^{10}=1024$
geometric theta characteristics.  The canonical characteristic $B_C$
is the unique non-effective one.  It is invariant under the
geometric Galois group, also when a purely inseparable extension is
needed to define it as a line bundle.  Identifying the set of theta
characteristics with $V^\vee$ using $B_C$ as origin, the remaining
1023 effective theta characteristics correspond to the nonzero
vectors.  We first prove that monodromy is transitive on them.

\par\smallskip\noindent\textbf{The contact reciprocity calculation.}
Let $C$ be a smooth plane curve of degree $d$, let $Q$ be a smooth
plane curve of degree $e=d-3$, and suppose $C\cap Q=2D$.
The half-intersection divisor defines the same zero-dimensional
subscheme $D$ on both curves: locally, in coordinates in which one
curve is $y=0$, the other has contact $y=u x^{2m}+\cdots$, and both
half divisors have ideal $(y,x^m)$.
Put
\[
 \theta_C=\mathcal O_C(D),\qquad
 \theta_Q=\mathcal O_Q(D-3H).
\]
These are theta characteristics.  The two ideal-sheaf sequences,
twisted by $e$, are
\[
 0\to\mathcal O_{\mathbf P^2}(-3)\to\mathcal I_D(e)
       \to\theta_C\to0,
 \qquad
 0\to\mathcal O_{\mathbf P^2}\to\mathcal I_D(e)
       \to\theta_Q\to0.
\]
Since $H^1(\mathbf P^2,\mathcal O(-3))=H^1(\mathbf P^2,\mathcal O)=0$,
they give the exact numerical relation
\begin{equation}\label{eq:contact-reciprocity}
 h^0(\theta_C)=1+h^0(\theta_Q).
\end{equation}

\par\smallskip\noindent\textbf{Smooth contact cubics.}
For $d=6$, every effective $\theta_C$ with section $s$ supplies a
contact cubic, because $s^2\in H^0(C,\omega_C)=H^0(C,\mathcal O_C(3))$
lifts uniquely to a ternary cubic.
If $Q$ is smooth, its theta characteristics are degree-zero line
bundles.  For an ordinary cubic, the only non-effective one is $B_Q$;
the other is $\mathcal O_Q$.  For a supersingular cubic, the sole
geometric theta characteristic is $\mathcal O_Q$.

The incidence corresponding to $\theta_Q=B_Q$ over the ordinary
smooth-cubic parameter space is irreducible of dimension 27.
To see this without a descent assumption on $B_Q$, pull the
coefficient space back along Frobenius (and work in an affine
coefficient chart).  The quotient of the relative Frobenius pushforward
by the structure sheaf supplies $B_Q$ on the Frobenius-twisted universal
curve.  This purely inseparable base change is irreducible and changes
no geometric Galois group.  The spaces
$H^0(Q,B_Q(3))$ form a vector bundle of rank nine, since its degree is
nine and the genus is one.  Squaring their sections gives the allowed
restrictions in $H^0(Q,\mathcal O_Q(6))$.
The restriction of ternary sextics is surjective, with kernel
$Q\cdot H^0(\mathbf P^2,\mathcal O(3))$ of dimension ten.
Thus the allowed sextics form, after this purely inseparable cover,
a nineteen-dimensional affine space, or a projective space of
dimension eighteen.  Smooth cubics have parameter dimension nine,
so the total is $9+18=27$, equal to the sextic parameter dimension.
The image of an irreducible space is irreducible.

The smooth-contact-cubic loci with effective $\theta_Q$ do not
dominate the sextic parameter space.  Their incidence dimensions are
at most 27 by the same degree-nine Riemann--Roch calculation; by
\eqref{eq:contact-reciprocity}, every such fibre over a sextic contains
the positive dimensional projective linear system of its
$\theta_C$, which has at least two sections.  Their images therefore
have dimension at most 26.  This remains true for the supersingular
cubic locus, whose base dimension is smaller.

\par\smallskip\noindent\textbf{Singular contact cubics.}
We show that these also occur only for sextics in a proper closed
subset of the 27-dimensional parameter space.
For every fixed cubic, the kernel of restriction of sextic forms has
dimension ten.  The following bounds are therefore bounds on the
dimension of the allowed restriction space, before adding ten and
projectivizing.

An irreducible nodal cubic has normalization $\mathbf P^1$ and two
points $P_+,P_-$ above the node.  Its stratum has dimension eight.
The pullback of $\mathcal O_Q(6)$ has degree eighteen.
Away from $P_+,P_-$, the orders of the restricted sextic are even.
At these two points the parities are equal.  If both are even, the
binary form is a square of a degree-nine form; the node gluing
condition gives one independent linear condition on that square root,
leaving dimension nine.  If both are odd, it is a fixed binary
quadratic with zeros $P_+,P_-$ times a square of degree eight, again
dimension nine; the gluing values are then both zero.  The incidence
dimension is consequently at most $8+9+10-1=26$.

An irreducible cuspidal cubic has normalization $\mathbf P^1$, one
point above its cusp, and parameter dimension at most seven.  All
orders of the pulled-back sextic are even, since its total degree is
eighteen and all orders away from that point are even.  It is therefore
a square of a degree-nine binary form, a space of dimension ten.
Any descent condition can only decrease the dimension.  The resulting
bound is $7+10+10-1=26$.  The dimensions eight and at most seven of
these strata follow directly by fixing a singular point: singularity
imposes three independent linear conditions on cubic forms, and a
repeated tangent cone imposes at least one further condition.

If the cubic is a transverse union of a line and a smooth conic, its
parameter dimension is seven.  On the two normalization components,
the sextic has degrees six and twelve.  Its odd orders, if any, must
occur at both of their two intersection points.  With no odd orders,
the square-root spaces have dimensions four and seven, subject to
two independent gluing conditions, giving nine.  With odd orders at
both points, extracting their fixed divisors leaves square-root
spaces of dimensions three and six, also giving nine.  The total
incidence dimension is at most 25.  If the line is tangent to the
conic, the cubic stratum has dimension at most six; each component's
restriction must have even order at the unique intersection point,
and the two square-root spaces satisfy at least their common-value
condition.  Their combined dimension is at most ten, again giving
at most 25.

For three distinct lines in a triangle, the parameter dimension is
six.  The possible odd-node patterns are either no nodes or all three
nodes.  In the first case the three four-dimensional square-root
spaces have three independent common-value conditions; in the second,
the three spaces have dimension three each.  Both bounds are nine,
giving total dimension at most 24.  For three concurrent lines the
parameter dimension is five.  Each line restriction has even order
at their sole common point, and the three square-root spaces satisfy
two common-value conditions.  Their combined dimension is at most
ten, again giving at most 24.

Finally a nonreduced cubic is $L^2M$, allowing $L=M$.
Its restriction to a smooth sextic can be even only if the line
section $M|C$ is even.  Sextics with such a line lie in dimension at
most 26: for a fixed line, binary sextics that are squares form a
four-dimensional subspace of a seven-dimensional restriction space;
the ternary restriction kernel has dimension 21.  Projectivizing and
allowing the line to vary gives $(4+21-1)+2=26$.
All possible singular cubics have now been covered.

Outside the proper loci just bounded, every one of the 1023 effective
theta characteristics of an ordinary sextic has a smooth contact
cubic, and this cubic corresponds to $B_Q$ by
\eqref{eq:contact-reciprocity}.  The irreducible incidence therefore
dominates the sextic space.  Its dimension is 27, so its generic
fibre is finite; every effective theta has $h^0=1$, and the fibre has
exactly 1023 geometric points.  As in the quintic case, irreducibility
implies geometric Galois transitivity on these points even if the
incidence has a purely inseparable part.  Thus $G_{\mathrm{geom}}$
is transitive on $V^\vee\setminus\{0\}$.

\par\smallskip\noindent\textbf{A boundary element that excludes the smaller transitive groups.}
Degenerate a sextic to an ordinary smooth quintic and a transverse
line, meeting in five distinct points.  This is an ordinary nodal
sextic.  Its rank-ten local system has a four-dimensional graph
submodule and the rank-six quintic local system as quotient.
The node permutation group is a subgroup of $S_5$; the graph
representation is its augmentation representation.
The exact quintic monodromy theorem supplies a quotient with order
divisible by seven and an element of order thirty-one acting on the
six-dimensional quotient as a one-dimensional identity plus an
irreducible five-dimensional summand.  Hence $7\mid|G_{\mathrm{geom}}|$.
Take the inverse image of the cyclic order-thirty-one subgroup in
the boundary monodromy group.  A Sylow thirty-one subgroup of this
inverse image maps onto the cyclic quotient.  There are no elements
of order $31^2$ in $\mathrm{GL}_{10}(2)$, since the multiplicative
order of $2$ modulo $31^2$ is 155.  Every nonidentity element of this
Sylow subgroup therefore has order thirty-one; choose one with
nonidentity quotient image and denote it by $g$.
It acts trivially on the graph submodule, since $31\nmid|S_5|$.
Its characteristic polynomial consequently has the form
\begin{equation}\label{eq:sextic-order31}
 \chi_g(T)=(T-1)^5 f(T),
\end{equation}
where $f$ is an irreducible polynomial of degree five whose roots
are primitive thirty-first roots of unity.
This polynomial is not reciprocal: inversion does not belong to
the orbit of squaring on primitive thirty-first roots, since
$\langle2\rangle\subset(\mathbf Z/31)^\times$ has odd order five.
As before, the specialization is inside the full-rank Lang locus,
so these are elements of the same generic monodromy image.

Apply Hering's classification of transitive linear groups in the
form of Churikov--Vasil'ev--Zvezdina (2022), Lemma~3.3, arXiv p.~6,
which was opened and checked.  In dimension ten over $\mathbf F_2$,
its alternatives, unless the group contains $\mathrm{SL}_{10}(2)$,
place the group in one of
\[
 N(\mathrm{Sp}_{10}(2)),\qquad
 \Gamma\mathrm L_5(4),\qquad
 \Gamma\mathrm L_2(32),\qquad
 \Gamma\mathrm L_1(1024).
\]
The exceptional alternatives have dimensions four or six or odd
characteristic and do not occur here; the $G_2$ alternative has
dimension divisible by six.
The symplectic normalizer preserves its unique invariant alternating
form up to a scalar, and over $\mathbf F_2$ the scalar is one.
Every one of its characteristic polynomials is reciprocal, contrary
to \eqref{eq:sextic-order31}.
An element of order thirty-one in $\Gamma\mathrm L_5(4)$ lies in
$\mathrm{GL}_5(4)$, so its fixed space has even dimension over
$\mathbf F_2$, contrary to the fixed dimension five in
\eqref{eq:sextic-order31}.
Finally, neither $\Gamma\mathrm L_2(32)$ nor
$\Gamma\mathrm L_1(1024)$ has order divisible by seven, which has
already been forced by the quintic quotient.
All proper alternatives are excluded.  Therefore
$G_{\mathrm{geom}}=\mathrm{SL}_{10}(2)=\mathrm{GL}_{10}(2)$.
The arithmetic image contains this group and is contained in it,
so it is the same group.
\end{proof}

% Characteristic-two plane septic case: separate pending independent audit.
\subsection{The exact group for ordinary plane septics in characteristic two}
\label{subsec:septic-exact}

\begin{theorem}\label{thm:septic-stabilizer-monodromy}
\status{Proved under stated hypotheses}
For the universal ordinary smooth plane septic in characteristic two,
the geometric and arithmetic monodromy images on
$V=H^1_{\mathrm{et}}(C_{\bar\eta},\mathbf F_2)$ are the full stabilizer
of the distinguished nonzero covector supplied by
$\mathcal O_C(2)$ and the canonical theta characteristic:
\[
 G_{\mathrm{geom}}=G_{\mathrm{arith}}
 =\operatorname{Stab}(\ell)
 \simeq 2^{14}\!:\mathrm{GL}_{14}(\mathbf F_2).
\]
In particular, the rank-fifteen representation is reducible, with its
distinguished invariant hyperplane, and its determinant is trivial.
\end{theorem}

\begin{proof}
We work first over an algebraic closure.  An ordinary genus-fifteen
curve has $2^{15}$ geometric theta characteristics.  Its canonical
theta characteristic $B_C$ is the unique non-effective one.  The
plane theta characteristic $L_C=\mathcal O_C(2)$ has six sections.
The class $v=L_C\otimes B_C^{-1}$ is the distinguished nonzero element
of $V^\vee$, fixed by arithmetic as well as geometric monodromy.
We prove transitivity on the other $2^{15}-2=32766$ elements of
$V^\vee$, then identify the whole stabilizer using a boundary subgroup.

\par\smallskip\noindent\textbf{Smooth contact quartics.}
Every nonzero section $s$ of an effective theta characteristic
$\theta_C$ supplies a unique quartic $Q$ with
$Q\cap C=2\operatorname{div}(s)$: the section $s^2$ belongs to
$H^0(C,\omega_C)=H^0(C,\mathcal O_C(4))$, and restriction from plane
quartics is an isomorphism.  If $Q$ is smooth, set
$D=\operatorname{div}(s)$ and
$\theta_Q=\mathcal O_Q(D-3H)$.  The contact reciprocity calculation
\eqref{eq:contact-reciprocity} gives
\[
 \theta_Q^2\simeq\omega_Q,\qquad
 h^0(\theta_C)=1+h^0(\theta_Q).
\]
For every smooth quartic and every theta characteristic on it,
$\theta_Q(3)$ has degree fourteen and exactly twelve sections, by
Riemann--Roch and Serre duality.  Restriction of septic forms to a
quartic is surjective and has kernel
$Q\cdot H^0(\mathbf P^2,\mathcal O(3))$, of dimension ten.

Consider first ordinary smooth quartics with $\theta_Q=B_Q$.
After a purely inseparable Frobenius base change, $B_Q$ is supplied
by the relative Frobenius sequence, so its twelve-dimensional spaces
of sections form a vector bundle.  Squaring their sections and lifting
through restriction gives an irreducible incidence of dimension
\[
 14+(12+10-1)=35,
\]
the dimension of the projective space of septics.  This incidence is
geometrically irreducible: the base is a nonempty open in the
geometrically irreducible quartic parameter space, the Frobenius
base change is a universal homeomorphism, and the section and lifting
spaces have geometrically irreducible affine fibres before
projectivizing.

For all smooth quartics and all theta characteristics, the incidence
dimension is at most thirty-five, since the scheme of roots is
finite over the quartic parameter space and the same twelve-section
calculation applies.  Those with effective $\theta_Q$ have images
of dimension at most thirty-four.  Indeed, reciprocity gives
$h^0(\theta_C)\ge2$, and the open part of
$\mathbf P H^0(C,\theta_C)$ near a smooth contact quartic is a
positive-dimensional fibre of this incidence over $C$.  This covers
the nonordinary smooth quartics as well, since all their geometric
theta characteristics are effective.  Thus a generic septic can
have a smooth contact quartic only through the canonical
non-effective theta characteristic of an ordinary quartic.

We next exclude singular contact quartics for the generic septic,
apart from the squares that give $L_C$.

\par\smallskip\noindent\textbf{Nonreduced quartics.}
A nonreduced quartic is either a square of a conic, or has the form
$Q=L^2R$ with $R$ a reduced conic.  A square gives exactly
$\theta_C=\mathcal O_C(2)$.  In the other case, evenness of $Q|C$
implies evenness of $R|C$.
For a fixed smooth conic, the restrictions of septics have binary
degree fourteen.  Those with even zero divisor are squares and form
an eight-dimensional space; the restriction kernel has dimension
$h^0(\mathbf P^2,\mathcal O(5))=21$.
Allowing the conic to vary gives dimension
$5+8+21-1=33$ for the image in the septic parameter space.
For a reduced pair of lines, each restricted binary form has odd
degree seven.  Its sole possible odd-order zero must be the node,
so each restriction is the fixed linear node factor times the
square of a cubic.  This leaves four parameters on each line;
both node values vanish, so their combined dimension is eight.
The conic parameter dimension is four, giving the smaller bound
$4+8+21-1=32$.  Consequently all non-square nonreduced contact
quartics occur over a proper subset of the septic parameter space.

\par\smallskip\noindent\textbf{Reduced reducible quartics.}
We give upper bounds that allow all coincidences and higher contacts
among components.  For a fixed reduced quartic, normalize each
component.  Away from the preimages of singular points, the divisor
of the restricted septic is even.  For each choice of the finite
odd-parity divisor $\epsilon$ on these preimages, its restriction
is a fixed section with divisor $\epsilon$ times a square.
The square-root line bundles form a finite set.  On a normalization
component of degree $a$ and genus $h$, the root has degree
$(7a-|\epsilon_i|)/2$ and, in the cases below, vanishing $H^1$;
hence its section dimension is
$(7a-|\epsilon_i|)/2+1-h$.
These degrees exceed $2h-2$: the only positive-genus component is
a smooth cubic, where the root degree is at least nine.
On a line the parity divisor has at most three points, and on a
conic at most four, so the genus-zero assertion is immediate.

For a smooth cubic and a line, the base dimension is eleven.
Both components have odd total intersection degree, so
$|\epsilon|\ge2$.  The sum of root dimensions is at most
$14+2-1-1=14$.  Adding the ten-dimensional restriction kernel and
projectivizing gives at most $11+14+10-1=34$.
If the cubic is irreducibly singular, its base dimension decreases
by at least one while the normalization genus decreases by one;
the same upper bound thirty-four follows.  Cases where the cubic
itself is reducible are included below.

For two smooth conics the base dimension is ten and the unrestricted
sum of root dimensions is sixteen.  If $\epsilon\ne0$, its total
degree is at least two, giving at most fifteen.  If
$\epsilon=0$, choose any intersection point of the two conics.
Descent of the septic restriction imposes equality of its two values,
and, after taking square roots over the algebraically closed field,
this is one independent linear condition on the root sections.
Independence follows from global generation on both normalization
components.  Thus the root dimension is again at most fifteen,
and the total bound is $10+15+10-1=34$.

For a smooth conic and two distinct lines the base dimension is
nine.  The two odd-degree line restrictions force
$|\epsilon|\ge2$, so the sum of root dimensions is at most
$14+3-1=16$.  The total is at most $9+16+10-1=34$.
For four distinct lines the base dimension is eight.  Each line
forces an odd parity divisor, so $|\epsilon|\ge4$ and the sum of
root dimensions is at most $14+4-2=16$; the total is at most
$8+16+10-1=33$.  No transversality assumption was used in these
upper bounds.

\par\smallskip\noindent\textbf{Integral singular quartics: normalization parameter spaces.}
An integral singular quartic has normalization of genus
$h\in\{0,1,2\}$.  We need two parameter bounds, and establish them
directly in characteristic two instead of using a Severi-variety
dimension assertion:
\begin{align*}
 \dim\{Q:\text{normalization genus }h\}&\le11+h,\\
 \dim\{Q:\text{normalization genus }h,
          \text{all singularities unibranch}\}&\le10+h.
\end{align*}
For $h=0$, a normalization map is given by a base-point-free
three-dimensional subspace of $H^0(\mathbf P^1,\mathcal O(4))$,
together with a projective frame.  The parameter dimension is
$\dim\operatorname{Gr}(3,5)+\dim\operatorname{PGL}_3
 -\dim\operatorname{PGL}_2=6+8-3=11$.
For $h=1$, temporarily mark an origin on the normalization.  Its
moduli have dimension one, its degree-four Picard variety dimension
one, the three-dimensional subspace of its four-dimensional section
space varies in dimension three, and a projective frame contributes
eight.  Forgetting the arbitrary origin subtracts one, giving
$1+1+3+8-1=12$.
For $h=2$, smooth-curve moduli have dimension three, the degree-four
Picard variety dimension two, and every degree-four line bundle has
three sections, giving $3+2+8=13$.
The dimension of genus-two moduli follows directly from curve
deformations: $H^0(T)=0$, $h^1(T)=3$ by Riemann--Roch, and $H^2(T)=0$.
The images under consideration are birational normalization maps,
so quotienting by reparametrization introduces no positive-dimensional
stabilizer.  These constructions cover all integral quartics of the
indicated normalization genus.

If all singularities are unibranch, the normalization map is
nonimmersive somewhere.  Otherwise at the preimage of each singular
point, one plane coordinate would be a uniformizer, making that
branch smooth.  For $h=0$ or $1$, and a degree-four line bundle $A$,
the map
$H^0(A)\to H^0(A|_{2P})$ is onto for every $P$, by Serre duality.
Among three-dimensional subspaces of $H^0(A)$, failure to have rank
two under this map has codimension two: for $h=1$ the rank-one locus
has dimension one in $\operatorname{Gr}(3,4)$ of dimension three;
for $h=0$ it has dimension four in
$\operatorname{Gr}(3,5)$ of dimension six.
Allowing $P$ to vary restores at most one dimension, so the
nonimmersive locus has codimension at least one in the parameter
space above.  This argument is invariant under reparametrization,
and the arbitrary elliptic origin remains a one-dimensional fibre.
For $h=2$, the system is the complete system $|A|$.
Failure to separate $2P$ is possible exactly when
$H^1(A(-2P))\ne0$, or $A\simeq\omega(2P)$.
As $P$ varies this is a subset of dimension at most one in the
two-dimensional degree-four Picard variety.  This remains a
dimension bound in characteristic two even if the map
$P\mapsto\omega(2P)$ is inseparable.  Thus the second bound follows
for all three genera.

Return to the restricted septic on such a quartic.  Its pullback
to the normalization has degree twenty-eight.  The odd-parity
divisor $\epsilon$ is supported above the finitely many singular
points.  There are at most six such branches, since the arithmetic
genus defect is at most three and the number of branches is at
most twice that defect.  A root after removing $\epsilon$ therefore
has degree at least eleven and vanishing $H^1$.  Its section
dimension is exactly
\[
 15-h-\frac{|\epsilon|}{2}.
\]
If $\epsilon\ne0$, then $|\epsilon|\ge2$; together with the
parameter bound $11+h$ this gives incidence dimension at most
$(11+h)+(14-h)+10-1=34$.
If $\epsilon=0$ and a singular point has two different branches,
descent imposes equality of the values at two distinct preimages.
The root line bundle separates these points, again by its degree,
so this is one independent linear condition after taking square
roots; the same bound thirty-four follows.
If all singularities are unibranch, the improved parameter bound
$10+h$ instead gives
$(10+h)+(15-h)+10-1=34$ without using any descent condition.
This excludes every integral singular contact quartic over a
generic septic.

\par\smallskip\noindent\textbf{Transitivity.}
Outside the proper closed subsets just bounded, every effective
theta characteristic other than $L_C$ has a smooth contact quartic,
and that contact quartic arises from its canonical non-effective
theta characteristic.  All $32766$ such theta characteristics
therefore occur in the geometrically irreducible canonical incidence
of dimension thirty-five.  This proves dominance.  Its generic
fibre is finite; in particular each of these theta characteristics
has exactly one section up to scalar, since otherwise its projective
section space would give a positive-dimensional fibre.  Consequently
the generic geometric fibre has precisely $32766$ points.
Irreducibility gives geometric Galois transitivity on these points;
a purely inseparable part of the incidence does not affect this
conclusion.  Hence $G_{\mathrm{geom}}$ is transitive on
$V^\vee\setminus\{0,v\}$.

\par\smallskip\noindent\textbf{The full quotient on the fourteen-dimensional space.}
Let $W=V^\vee/\langle v\rangle$, and let $H$ be the image on $W$.
Transitivity just proved implies transitivity of $H$ on
$W\setminus\{0\}$, of size $2^{14}-1$.
Degenerate to an ordinary sextic and a transverse line meeting it
in six points.  The rank-fifteen local system is still defined on
this nodal septic because Frobenius is bijective; its graph part
has rank five and its curve quotient has rank ten.  The projection
from this boundary family to the ordinary sextic parameter space
has geometrically connected fibres, so the quotient monodromy is
the full $\mathrm{GL}_{10}(2)$ established above.
In particular $11\mid|G_{\mathrm{geom}}|$.
The kernel of $\operatorname{Stab}(v)\to\mathrm{GL}(W)$ is a
two-group; therefore $11\mid|H|$.

Choose in the sextic quotient an order-thirty-one element acting
as a primitive five-dimensional block plus an identity of rank five.
An order-thirty-one lift exists by the Sylow argument used for
sextics: $\mathrm{GL}_{15}(2)$ has no element of order $31^2$, since
$\operatorname{ord}_{31^2}(2)=155$.
Its action on the graph part is trivial, since the node action
factors through $S_6$.  On $V^\vee$ it therefore has characteristic
polynomial $(T-1)^{10}f(T)$ for a primitive irreducible
degree-five factor $f$.  Quotienting by the fixed $v$ gives on $W$
\[
 (T-1)^9 f(T).
\]
This polynomial is not reciprocal, and the fixed-space dimension
is nine.

Hering's classification, in the opened form of
Churikov--Vasil'ev--Zvezdina (2022), Lemma~3.3, gives, in dimension
fourteen over $\mathbf F_2$, the alternatives
$\mathrm{GL}_{14}(2)$ or containment in
\[
 N(\mathrm{Sp}_{14}(2)),\quad
 \Gamma\mathrm L_7(4),\quad
 \Gamma\mathrm L_2(128),\quad
 \Gamma\mathrm L_1(16384).
\]
The symplectic alternative forces a reciprocal polynomial and is
excluded.  In $\Gamma\mathrm L_7(4)$ an order-thirty-one element
is $\mathbf F_4$-linear, so its fixed dimension over $\mathbf F_2$
is even, excluding nine.  The orders of the last two groups are
$1864254336$ and $229362$, respectively, neither divisible by eleven.
Thus $H=\mathrm{GL}_{14}(2)$.

\par\smallskip\noindent\textbf{The unipotent kernel and arithmetic equality.}
The intersection of $G_{\mathrm{geom}}$ with the kernel
$2^{14}$ of $\operatorname{Stab}(v)\to\mathrm{GL}(W)$ is an invariant
subspace for the full natural dual $\mathrm{GL}_{14}(2)$ module.
It is therefore zero or the whole kernel.  In the former case,
$G_{\mathrm{geom}}$ is a complement.  The elementary transvection
calculation $H^1(\mathrm{GL}_m(2),\mathbf F_2^m)=0$ for $m\ge4$,
proved in the quintic section (and equally valid for the dual),
shows that every such complement is conjugate to the standard one.
It has two orbits of size $2^{14}-1$ on
$V^\vee\setminus\{0,v\}$, contradicting transitivity.
Hence the whole kernel occurs, and the geometric group is the
full stabilizer.  Arithmetic monodromy also fixes $v$ and contains
the geometric group, so the two groups are equal.
\end{proof}

\par\smallskip\noindent\textbf{Finite certificate and scope.}
The retained program \texttt{sextic\_septic\_group\_checks.py}
checks the order of $2$ modulo $31^2$, the two distinct inverse
Frobenius orbits of primitive thirty-first roots, the semilinear
group orders and their missing prime divisors, and the two orbit
sizes used above.  It does not recompute Hering's classification
or replace the incidence argument by experimental Frobenius data.
This theorem settles degree seven in characteristic two only;
it does not assert the same stabilizer in every odd degree.

% All-degree characteristic-two contact argument; independent audit requested.
\subsection{Contact curves in every degree in characteristic two}
\label{subsec:all-degree-char2-contact}

The following argument removes the degree restriction from the
contact-incidence method.  Its global dimension input is
Christ--He--Tyomkin (2023), Proposition~2.7 (p.~8) and
Corollary~5.3 (p.~27), in the corrected publisher version of
\emph{On the Severi problem in arbitrary characteristic}.
Both statements apply over every algebraically closed field,
including characteristic two.  If a reduced plane curve of degree
$s$ has $c$ normalization components of genera $h_1,\ldots,h_c$,
their geometric genus convention is $1-c+\sum h_i$.  Consequently
the locus of such curves has dimension at most
\begin{equation}\label{eq:contact-severi-bound}
 3s+\sum_i h_i-c.
\end{equation}
If a family attains this dimension, its generic member is nodal.
The publisher's correction concerns the article history and missing
references; the opened corrected text contains these statements.

\begin{lemma}[The conductor inclusion for an even contact]
\label{lem:contact-conductor}
\status{Proved under stated hypotheses}
Let $C$ be a smooth plane curve of degree $d$ over an algebraically
closed field of characteristic two, and let $Q$ be a reduced plane
curve of degree $s<d$ with $Q\cap C=2D$ as divisors on $C$.
Write $\nu:\widetilde Q=\coprod Q_i\to Q$ for the normalization,
and write
\[
 \operatorname{div}(f|_{Q_i})=2A_i+\epsilon_i,
 \qquad M_i=\mathcal O_{Q_i}(A_i),
\]
where $f$ defines $C$ and $\epsilon_i$ is the reduced odd-order
divisor.  It is supported above the singular points of $Q$.
Set
\[
 J=\mathcal I_{D/Q}(s),\qquad
 L_C=\mathcal O_C(D),\qquad
 b=d-s.
\]
Then
\begin{equation}\label{eq:general-contact-reciprocity}
 h^0(C,L_C)=1+h^0(Q,J),
\end{equation}
and there is an inclusion
\begin{equation}\label{eq:conductor-section-inclusion}
 \nu_*\left(\bigoplus_i
       \omega_{Q_i}\otimes M_i^{-1}(3)\right)\hookrightarrow J.
\end{equation}
In particular, with $t=\sum_i h^1(Q_i,M_i)$,
\begin{equation}\label{eq:contact-h1-cancellation}
 h^0(Q,J)\ge t,
 \qquad t>0\Longrightarrow h^0(Q,J)\ge t+1.
\end{equation}
\end{lemma}

\begin{proof}
The same subscheme $D\subset\mathbf P^2$ is used on both curves.
Its ideal-sheaf sequences, twisted by $s$, are
\[
 0\to\mathcal O_{\mathbf P^2}(s-d)
  \to\mathcal I_D(s)\to L_C\to0,
 \qquad
 0\to\mathcal O_{\mathbf P^2}
  \to\mathcal I_D(s)\to J\to0.
\]
Here $\mathcal O_C(s)(-D)\simeq\mathcal O_C(D)$ because
$2D=Q|C$.  Since $s<d$ and the first cohomology of every line
bundle on $\mathbf P^2$ vanishes, these sequences prove
\eqref{eq:general-contact-reciprocity}.

We prove the inclusion locally and then express it globally.
At a point of $D$, complete the local ring of the plane.  Since
$C$ is smooth, choose its equation $f$ as one parameter and let
$b_0$ cut out the half divisor on $C$.  After multiplying the local
equation of $Q$ by a unit, that equation is
$b_0^2-fa_0$.  In the reduced one-dimensional ring
$A=\widehat{\mathcal O}_{Q,P}$, the ideal of $D$ is
$I=(f,b_0)$ and $f$ is a non-zero-divisor.
Reduction modulo $f$ gives
$A/(f)\simeq k[[x]]/(b_0^2)$, where the annihilator of $b_0$ is
$(b_0)$.  Hence
\[
 I^*:=\operatorname{Hom}_A(I,A)
 =\{z:zI\subset A\}=f^{-1}I,
 \qquad I^{**}=I.
\]
These equalities can be read directly in the total quotient ring:
$zf\in A$ and $zb_0\in A$ say that $zf$ belongs to
$(f):(b_0)=I$.

Let $v_i$ be the valuation on a normalization branch.  Put
\[
 r_i=v_i(f),\qquad j_i=\min\{v_i(f),v_i(b_0)\},
 \qquad\gamma_i=2j_i-r_i.
\]
The relation $b_0^2=fa_0$ gives $\gamma_i\ge0$; moreover
$\gamma_i\equiv r_i\pmod2$.  Thus
$\gamma_i=\epsilon_i(P_i)+2k_i$ for an integer $k_i\ge0$.
Write $c_i$ for the conductor exponent on this branch, so that
the conductor ideal on the normalization is generated there by
a parameter to the power $c_i$.
If a tuple $z$ of branch functions satisfies
\[
 v_i(z)\ge c_i+r_i-j_i\quad\hbox{for every branch},
\]
then its product with every member of $I^*=I/f$ lies in the
conductor, hence in $A$.  It follows that $z\in I^{**}=I$.
No separate assumption on a tame singularity, on its number of
branches, or on its analytic type is used here.

Let $\Delta$ be the conductor divisor on $\widetilde Q$, let
$\Gamma=\sum\gamma_iP_i$, and let $K_0=\sum k_iP_i$.
If $N=\nu^*J/\text{torsion}$, the valuation calculation gives
\[
 N|_{Q_i}\simeq M_i(-b)(-K_0|_{Q_i}),\qquad
 \nu_*N(\Gamma-\Delta)\subset J.
\]
Indeed the normalized ideal has vanishing order $j_i$, and
$\gamma_i=2j_i-r_i$ is precisely the adjustment from its square
to the divisor of $f$.

For a reduced plane curve, finite duality (Stacks Project,
Remark~48.20.12, Tag~0AX4, opened and checked) and adjunction give
\[
 \omega_{\widetilde Q}
 \simeq\nu^*\omega_Q(-\Delta)
 =\mathcal O_{\widetilde Q}(s-3)(-\Delta).
\]
For completeness, the conductor occurs because
$\operatorname{Hom}_{\mathcal O_Q}
(\nu_*\mathcal O_{\widetilde Q},\mathcal O_Q)$ is the conductor
ideal; finite duality identifies its tensor product with
$\omega_Q$ with $\nu_*\omega_{\widetilde Q}$.
Using $M_i^2\simeq\mathcal O_{Q_i}(d)(-\epsilon_i)$ now yields
\[
 \omega_{Q_i}\otimes M_i^{-1}(3)
 \simeq M_i(-b)(\epsilon_i-\Delta_i)
 =N|_{Q_i}(K_{0,i}+\epsilon_i-\Delta_i).
\]
Since $\Gamma=\epsilon+2K_0$, this line bundle is a subsheaf
of $N(\Gamma-\Delta)$, proving
\eqref{eq:conductor-section-inclusion}.

By Serre duality, $h^1(M_i)=h^0(\omega_{Q_i}M_i^{-1})$.
Twisting by the positive line bundle $\mathcal O_{Q_i}(3)$
does not decrease this dimension.  If the initial space is nonzero,
the increase is strict: choose a point where some section does not
vanish, and choose a plane cubic through its image whose pullback
is a nonzero divisor.  Multiplication by that cubic has image
consisting of sections vanishing there, whereas multiplying the
chosen initial section by a cubic not through the point gives a
section outside the image.  Summing over the components proves
\eqref{eq:contact-h1-cancellation}.
\end{proof}

\begin{lemma}[No lower-degree reduced contact for a generic curve]
\label{lem:no-low-degree-contact}
\status{Proved under stated hypotheses}
Fix $d\ge4$ and $0<s<d-3$ in characteristic two.
A generic smooth plane curve $C$ of degree $d$ has no reduced
plane curve $Q$ of degree $s$ with $Q|C$ even.
\end{lemma}

\begin{proof}
Let $N_d=\binom{d+2}{2}-1$.
For fixed normalization genera, component degrees, and odd-parity
divisors $\epsilon_i$, the line bundles
$M_i^2=\mathcal O_{Q_i}(d)(-\epsilon_i)$ have finite choices.
They can be parameterized after a finite base change of the
singularity stratum; inseparable base changes do not affect
dimension.  The normalized restrictions of $f$ are squares of
sections of these $M_i$, times the fixed odd-parity factors.
Ignoring descent conditions only increases the dimension.
Riemann--Roch gives
\[
 \sum_i h^0(M_i)
 =\frac{ds}{2}-\frac{|\epsilon|}{2}
   +c-\sum_i h_i+t,
 \qquad t=\sum_i h^1(M_i).
\]
The kernel of restricting degree-$d$ forms to $Q$ has dimension
$\binom{d-s+2}{2}$.  Combining this with
\eqref{eq:contact-severi-bound}, and projectivizing, gives the
incidence bound
\begin{equation}\label{eq:general-contact-dimension}
 \dim\mathcal I_{d,s}
 \le N_d-\frac{s(d-s-3)}2-\frac{|\epsilon|}2+t
\end{equation}
on each stratum with fixed $t$.

Here is the precise fibre argument used with this bound.
Suppose a component of the even-contact incidence dominates the
ordinary smooth degree-$d$ parameter space.  Over its generic
curve, roots $L_C^2=\mathcal O_C(s)$ form a finite scheme.
For each geometric root the contact forms are the projective
space of its nonzero sections, followed by squaring.
Restriction of degree-$s$ plane forms to $C$ is an isomorphism,
so no further choices of $Q$ occur.  The reduced-contact condition
is open in this projective section space.  The generic fibre of
our component therefore has dimension
$h^0(L_C)-1=h^0(J)$.
Its generic point belongs to a single stratum of the finite
stratification by the data above, and an open subset of the same
generic fibre remains in that stratum.  Thus the fibre dimension
may be subtracted in \eqref{eq:general-contact-dimension}; we are
not subtracting a family of contact curves from a special stratum
that it need not preserve.

The conductor lemma gives $h^0(J)\ge t$.  Dominance would imply
\[
 N_d\le N_d-\frac{s(d-s-3)}2-\frac{|\epsilon|}2<N_d,
\]
a contradiction.  There are finitely many relevant component
degrees, genera, parity choices in finite covers, and cohomology
dimensions.  Thus their images have proper closure, proving the
claim for a generic curve.
\end{proof}

\begin{theorem}[The theta orbits of generic plane curves in characteristic two]
\label{thm:all-char2-theta-orbits}
\status{Proved under stated hypotheses}
Let $d\ge4$, $g=(d-1)(d-2)/2$, and let $C$ be the geometric
generic ordinary smooth plane curve of degree $d$ in characteristic
two.
If $d$ is even, every theta characteristic other than the canonical
one $B_C$ has exactly one section, and geometric monodromy is
transitive on these $2^g-1$ characteristics.
If $d$ is odd, every theta characteristic other than $B_C$ and
$\mathcal O_C((d-3)/2)$ has exactly one section, and geometric
monodromy is transitive on these $2^g-2$ characteristics.
Equivalently, on $V^\vee$ geometric monodromy is transitive on
all nonzero vectors in even degree, and on
$V^\vee\setminus\{0,v\}$ in odd degree, where $v$ is the
distinguished nonzero invariant vector.
\end{theorem}

\begin{proof}
Put $e=d-3$.  A section of an effective theta characteristic
$\theta_C$ gives a contact form $Q$ of degree $e$, and conversely
such a form gives an effective theta characteristic and its
section up to scalar.  Work throughout over the ordinary smooth
degree-$d$ locus, so theta characteristics have finitely many
geometric choices, with constant count $2^g$.  Passing to a
purely inseparable cover if necessary changes neither dimensions
nor geometric Galois orbits.

\par\smallskip\noindent\textbf{The unique dominant smooth-contact incidence.}
For a smooth $Q$, contact reciprocity gives
$h^0(\theta_C)=1+h^0(\theta_Q)$, where
$\theta_Q=\mathcal O_Q(D-3)$.
Riemann--Roch gives
\[
 h^0(\theta_Q(3))=3e+h^0(\theta_Q(-3)).
\]
Thus the incidence dimension, for a stratum with these dimensions
constant, is at most $N_d+h^0(\theta_Q(-3))$.
If $\theta_Q$ is effective, then
$h^0(\theta_Q(-3))<h^0(\theta_Q)$: multiplication by a suitable
cubic gives a proper subspace, as in the final paragraph of the
conductor lemma.  The fibre dimension $h^0(\theta_C)-1$
is $h^0(\theta_Q)$, so such an incidence cannot dominate.
For an ordinary smooth $Q$, its canonical $B_Q$ is the unique
non-effective theta characteristic.  It has
$h^0(B_Q(3))=3e$.  For a nonordinary smooth $Q$, all its geometric
theta characteristics are effective.  Consequently the only
possible dominant smooth-contact incidence is the canonical one
over ordinary smooth $Q$.

That canonical incidence is geometrically irreducible of dimension
$N_d$: the ordinary degree-$e$ base is geometrically irreducible;
after Frobenius base change the relative canonical theta
characteristic exists; its $3e$ sections form a vector bundle;
and lifts of their squares through restriction form affine spaces
with the constant ten-dimensional kernel
$Q\cdot H^0(\mathbf P^2,\mathcal O(3))$.
Projectivization gives total dimension
$\binom{e+2}{2}-1+3e+10-1=N_d$.
For $e=1,2$ the genus-zero canonical theta characteristic is used,
and the same computation applies.

\par\smallskip\noindent\textbf{Reduced singular contact forms cannot dominate.}
Apply the notation and incidence bound of the preceding lemmas
with $s=e$.  A hypothetical dominating component has dimension
at most $N_d-|\epsilon|/2+t$, and generic fibre dimension
$h^0(J)$.  If $t>0$, the strict part of
\eqref{eq:contact-h1-cancellation} gives a contradiction.
If $|\epsilon|>0$, its total degree is even and at least two,
so the non-strict inequality also gives a contradiction.
If $h^0(J)>0$ while $t=|\epsilon|=0$, the fibre dimension again
gives a contradiction.  Only
\[
 t=0,\qquad \epsilon=0,\qquad h^0(J)=0
\]
could remain for a dominant component.

In this case its image in the degree-$e$ curve parameter space
must attain the maximal Severi dimension
$3e+\sum h_i-c$; otherwise the incidence bound loses at least
one further dimension.  By Christ--He--Tyomkin Corollary~5.3,
the generic $Q$ is therefore nodal.  Since it is singular, choose
one node.  The conductor inclusion and $h^0(J)=0$ imply
\[
 H^0(Q_i,\omega_{Q_i}M_i^{-1}(3))=0
 \quad\hbox{on every component}.
\]
If the two preimages $P_+,P_-$ of the node lie on one component,
there exists a plane cubic whose pullback contains both points.
It follows that
$H^0(\omega M_i^{-1}(P_++P_-))=0$, hence
$H^1(M_i(-P_+-P_-))=0$.  The evaluation map at the two points
is therefore onto.  If they lie on different components, the
same argument with one point proves surjectivity of each
one-point evaluation.  In either case, equality of the two values
needed for descent of the squared restriction imposes one
nonzero linear condition on the root sections after taking square
roots of its scalar coefficients.  This lowers the incidence
dimension by at least one, the final contradiction.

\par\smallskip\noindent\textbf{Nonreduced contact forms.}
Every nonreduced form of degree $e$ has a unique expression up to
scalars $Q=R^2S$, with $S$ reduced and $\deg S=s<e$.
If $s>0$, evenness of $Q|C$ implies evenness of $S|C$, excluded
for a generic $C$ by Lemma~\ref{lem:no-low-degree-contact}.
If $s=0$, then $e$ is even, and the half divisor of $Q|C$ is
the divisor of $R|C$.  It gives exactly the plane theta
characteristic $\mathcal O_C(e/2)$.  This case occurs precisely
when $d$ is odd and is the sole nonreduced contact case retained.

\par\smallskip\noindent\textbf{Counting the generic fibre and obtaining the orbit.}
Every noncanonical theta characteristic is effective.
All except the plane characteristic in odd degree must therefore
occur in the geometrically irreducible canonical smooth-contact
incidence.  This proves its dominance.  Since it has dimension
$N_d$, its generic fibre is finite, so each of these theta
characteristics has exactly one section up to scalar.
The generic geometric fibre has respectively $2^g-1$ or
$2^g-2$ points.  Irreducibility makes geometric Galois monodromy
transitive on them; purely inseparable factors of the incidence
do not change this orbit statement.  Translating theta
characteristics by the canonical $B_C$ identifies their set
with $V^\vee$, proving the asserted vector orbits.
\end{proof}


\subsection{Exact characteristic-two plane monodromy in every degree}
\label{subsec:char2-uniform}

\begin{lemma}[Two elementary group-theoretic inputs]
\label{lem:char2-group-inputs}
\status{Proved under stated hypotheses}
For every integer $a\ge8$, there is an odd prime $\ell_a$ with
$\operatorname{ord}_{\ell_a}(2)=a$.
For $N\ge4$, the first cohomology of $\mathrm{GL}_N(\mathbf F_2)$
with coefficients in its natural module or its dual is zero.
\end{lemma}
\begin{proof}
We include a proof of the needed special case of the primitive-divisor
theorem. If an odd prime $\ell$ divides $\Phi_a(2)$ without having
order $a$, the cyclotomic factorization and the lifting-the-exponent
identity show that
$a=\operatorname{ord}_{\ell}(2)\ell^j$ with $j\ge1$, and that
$v_\ell(\Phi_a(2))=1$. Indeed, substitute
$v_\ell(2^{fm}-1)=v_\ell(2^f-1)+v_\ell(m)$, where
$f=\operatorname{ord}_\ell(2)$, in the M\"obius formula for
$v_\ell(\Phi_a(2))$; the sum is zero unless $a/f$ is a power of
$\ell$, when it is one. Since $f\mid\ell-1$, such an $\ell$ is the
largest prime factor $P$ of $a$. Thus, if no primitive prime occurs,
$\Phi_a(2)$ is either $1$ or $P$.
On the other hand the cyclotomic product gives
\[
 \Phi_a(2)>2^{\varphi(a)}\prod_{j\ge1}(1-2^{-j})
 >2^{\varphi(a)-2}.
\]
For the last bound, the first three factors have product $21/64$,
and the remaining product is at least
$1-\sum_{j\ge4}2^{-j}=7/8$; their product exceeds $1/4$.
If $P\ge7$, then $\varphi(a)\ge P-1$ and $2^{P-3}>P$.
If $P\le3$ and $a\ge8$, then $\varphi(a)\ge4$ and the same
conclusion holds. If $P=5$, either $\varphi(a)\ge8$ or $a=10$;
the exceptional calculation is $\Phi_{10}(2)=11$.
This contradiction proves the first assertion. The cyclotomic
valuation calculation also agrees with
\cite[Lemma~3.1(b)--(c), pp.~125--126]{GlasbyEtAl2017}, whose
statement and proof were opened.

For the second assertion let $z$ be a $1$-cocycle with values in
$\mathbf F_2^N$ and put $t_{ij}=1+E_{ij}$.
The relation $t_{ij}^2=1$ forces the $j$th coordinate of
$z(t_{ij})$ to vanish. For $l\ne i,j$, choose
$k\notin\{i,j,l\}$. Comparing the cocycle on the commuting pair
$t_{ij},t_{kl}$ shows that its $l$th coordinate also vanishes.
Thus $z(t_{ij})=a_{ij}e_i$.
For distinct $i,j,k$, apply the cocycle to
$[t_{ij},t_{jk}]=t_{ik}$; it gives $a_{ik}=a_{jk}$.
Consequently $a_{ij}=b_j$ for a single vector $b$, and
$z(t_{ij})=(t_{ij}-1)b$. Since these transvections generate the
group, $z$ is a coboundary. Composing with inverse transpose
gives the dual-module assertion.
\end{proof}

\begin{theorem}[All plane degrees in characteristic two]
\label{thm:plane-char2-all-groups}
\status{Proved under stated hypotheses}
Let $k$ be any perfect field of characteristic two, let $U_d$ be
the ordinary smooth locus in the projective space of degree-$d$
plane forms, and choose a geometric generic base point $\bar\eta$.
Write
\[
 g=(d-1)(d-2)/2,\qquad
 V=H^1_{\mathrm{et}}(C_{\bar\eta},\mathbf F_2).
\]
For $d\ge4$, the images of $\pi_1(U_d\otimes_k\bar k,\bar\eta)$
and $\pi_1(U_d,\bar\eta)$ on $V$ are equal and are
\[
 \begin{cases}
 \mathrm{GL}_g(\mathbf F_2),&d\text{ even},\\
 \operatorname{Stab}(\ell)
   \simeq(\mathbf F_2)^{g-1}\rtimes
                  \mathrm{GL}_{g-1}(\mathbf F_2),&d\text{ odd},
 \end{cases}
\]
where $\ell\ne0$ is the canonical covector determined by
$\mathcal O_C((d-3)/2)\otimes B_C^{-1}$ in odd degree.
The nonzero representation is irreducible exactly in even degree;
its determinant is always trivial. In degrees one and two $V=0$;
in degree three it has rank one and trivial, hence irreducible,
monodromy. Thus there is no threshold after which all plane
degrees have full linear or irreducible monodromy.
\end{theorem}
\begin{proof}
The degrees at most four, the quintic stabilizer, and the sextic
full group have been proved above. We induct on $d\ge7$.
Work first over $\bar k$ and on the dual space $W_0=V^\vee$.
Theorem~\ref{thm:all-char2-theta-orbits} proves transitivity on
$W_0\setminus\{0\}$ when $d$ is even and on
$W_0\setminus\{0,\ell\}$ when $d$ is odd.
In the odd case put $W=W_0/\langle\ell\rangle$; in the even
case put $W=W_0$. The induced group on $W$ is therefore transitive
on its nonzero vectors. Its dimension is
\[
 N=\begin{cases}g,&d\text{ even},\\g-1,&d\text{ odd}.\end{cases}
\]

We specify the boundary subgroup used in the induction. The
universal plane-curve coherent $H^1$ is a vector bundle of rank
$g$, including singular fibres. The locus where its Frobenius is
bijective is open, and the Lang construction gives an etale
$\mathbf F_2$-local system on this whole locus. It contains the
nodal curves $C'\cup L$, where $C'$ is ordinary smooth of degree
$d-1$ and $L$ meets it transversely in $d-1$ distinct points.
The smooth ordinary locus is dense in this normal parameter
space. Restriction of connected finite etale covers to a dense
open remains connected, so its fundamental group maps onto the
fundamental group of the larger locus. Boundary monodromy thus
lies in the group being determined.

The normalization sequence for $C'\cup L$ gives
\[
 0\longrightarrow H^1(\Gamma,\mathbf F_2)
 \longrightarrow H^1_{\mathrm{et}}(C'\cup L,\mathbf F_2)
 \longrightarrow H^1_{\mathrm{et}}(C',\mathbf F_2)
 \longrightarrow0.
\]
Here the graph space has dimension $d-2$ and its monodromy factors
through the permutation group $S_{d-1}$ of the nodes. The
projection from the parameter space of transverse pairs $(C',L)$
to that of $C'$ is smooth, surjective, and has geometrically
connected fibres, namely nonempty opens of the dual plane.
It therefore surjects on fundamental groups and supplies the
entire already determined monodromy of $C'$.

Write $r=(d-2)(d-3)/2$. If $d$ is even, the preceding odd-degree
group has a quotient $\mathrm{GL}_{r-1}(2)$; the kernel of this
quotient in the boundary group has prime divisors only $2$ and
those dividing $(d-1)!$. If $d$ is odd, the dual boundary
sequence has an invariant subspace of dimension $r$ with full
$\mathrm{GL}_r(2)$ action. The nonzero invariant vector $\ell$
cannot lie in that subspace, since it has no nonzero invariant
vector. It remains nonzero at the boundary: an invariant section
of a finite etale sheaf on the dense open extends uniquely over
the normal larger locus, and the zero and nonzero sections are
disjoint. Thus the rank-$r$ subspace injects into $W$, with quotient
a subquotient of the graph permutation space. In this case also,
the kernel of the boundary quotient onto $\mathrm{GL}_r(2)$ has
prime divisors only $2$ and those dividing $(d-1)!$. These kernel
claims follow by writing the representation in block triangular
form: its off-diagonal part is a $2$-group, the preceding odd-degree
unipotent radical is a $2$-group, and its remaining diagonal
action factors through $S_{d-1}$.

Set $a=\lfloor N/2\rfloor+1$. The elementary inequalities
\[
 8\le a<a+1<N-1,\qquad a>d-1,\qquad
 a+1\le\begin{cases}r-1,&d\text{ even},\\r,&d\text{ odd}
 \end{cases}
\]
hold for even $d\ge8$ and odd $d\ge7$. For example the two initial
cases have $(N,r,a)=(21,15,11)$ and $(14,10,8)$.
For the upper bound, using $\lfloor N/2\rfloor\le N/2$
reduces it to $g\ge2d+2$ in even degree and $g\ge2d-1$ in odd
degree; both hold at the indicated initial degree and their
left-minus-right differences increase. The other inequalities
follow directly from the same genus formula.
For $j=a,a+1$, choose the prime $\ell_j$ in
Lemma~\ref{lem:char2-group-inputs}. We have $\ell_j>j>d-1$.
Also $\ell_j\nmid N$: otherwise $\ell_j>N/2$ forces $N=\ell_j$,
whereas $j\mid N-1$ and $N/2<j<N-1$, an impossibility.

In the boundary general linear quotient choose an element of
order $\ell_j$ with one irreducible $j$-dimensional block and
identity on its complement. The kernel has order prime to
$\ell_j$, so a Sylow subgroup of its inverse image supplies an
element of order $\ell_j$ with that image. It acts trivially on
the graph quotient, because $\ell_j\nmid(d-1)!$.
Semisimplicity of this odd-order action shows that its
characteristic polynomial on $W$ is
\[
                    (T-1)^{N-j}f_j(T),
 \qquad f_j\text{ irreducible of degree }j.
\]

Apply the opened classification in
\cite[Lemma~3.3, p.~6]{ChurikovVasilevZvezdina2022}.
Every proper extension-field alternative lies inside
$\Gamma\mathrm L_{N/b}(2^b)$ for some divisor $b>1$ of $N$.
Our order-$\ell_j$ elements belong to its linear subgroup,
since $\ell_j\nmid b$. Their fixed-space dimensions are therefore
divisible by $b$. The consecutive dimensions $N-a$ and $N-a-1$
force $b=1$, excluding all these alternatives, including the
one-dimensional semilinear case. Over the original field the
remaining symplectic alternative lies in $\mathrm{Sp}_N(2)$:
its normalizer preserves the unique invariant alternating form,
and the only nonzero scalar is one. One of $a,a+1$ is odd.
For this $j$, the irreducible factor $f_j$ is not reciprocal,
since inversion is not in the odd-order subgroup generated by
$2$ modulo $\ell_j$. Its characteristic polynomial contradicts
preservation of a nondegenerate alternating form. The $G_2$
alternative over a larger field was already excluded, and its
field-two exceptions have dimension six. The other exceptional
field-two cases have dimension four. Here $N\ge14$.
The classification therefore forces the image on $W$ to be
$\mathrm{SL}_N(2)=\mathrm{GL}_N(2)$.

For even $d$ this finishes the geometric assertion. For odd $d$
the kernel of the image in $\mathrm{GL}(W)$ is an invariant
subgroup of $\operatorname{Hom}(W,\langle\ell\rangle)$.
This is the irreducible dual natural module, so the kernel is
zero or the whole module. If it were zero, the group would be a
complement. Lemma~\ref{lem:char2-group-inputs} makes every such
complement conjugate to the block diagonal group. That group has
two separate orbits on $W_0\setminus\{0,\ell\}$, contradicting
the theta-orbit theorem. Thus the kernel is the whole module,
giving the full stabilizer and completing the induction.

The distinguished odd-degree covector is defined over the original
perfect field, as proved in the explicit Hasse--Witt construction.
Arithmetic monodromy is contained in its stabilizer and contains
the geometric group just determined; in even degree it is
contained in the full general linear group. Thus the two images
are equal in both cases. Irreducibility and the determinant
assertions follow immediately; $\mathbf F_2^\times$ is trivial.
\end{proof}

% Agent /root/monodromy; complete argument, 19 November 2026.
% Independently audited in full by /root, including the opened classification source.
\subsection{Plane curves in odd characteristic: the full group in every degree}
\label{sec:odd-plane-all-degrees}
\status{Proved; independently audited by /root, including the primary classification statement.}

Throughout this section $p$ is odd, $k$ is an algebraic closure of
$\mathbf F_p$, $d\geq4$, and $g=(d-1)(d-2)/2$.  Let $W_d$ be the
Hasse--Witt invertibility open in the projective space of degree-$d$
plane equations, including singular equations, and let $U_d\subset W_d$
be its smooth open.  The vector bundle $R^1\pi_*\mathcal O$ has constant
rank $g$ on the entire projective parameter space.  Its Lang fixed
sheaf $V$ is finite locally constant of rank $g$ on $W_d$.
The map $\pi_1(U_d)\to\pi_1(W_d)$ is surjective.  Thus computing its
image on $W_d$ computes the desired geometric image on $U_d$.
Choose as base point a union of $d$ lines with no triple intersection.
All groups below act on this same fibre, after choosing paths in the
specified connected parameter spaces.  This is permissible because
every boundary family used below is contained in $W_d$.

\subsubsection{The line-arrangement representation}

Put $P=\mathbf F_p^d$, with its permutation basis $e_1,\ldots,e_d$,
$u=e_1+\cdots+e_d$, and $W=P/\langle u\rangle$.  The graph of the
line arrangement is the complete graph $K_d$.  Its first cochain space
is $\bigwedge^2P$, with $e_i\wedge e_j$ representing an oriented edge.
The coboundaries are $u\wedge P$.  Consequently
\begin{equation}\label{eq:odd-arrangement-module}
 V=H^1(K_d,\mathbf F_p)
   =\bigwedge^2P/(u\wedge P)=\bigwedge^2W.
\end{equation}
This is cohomology, rather than the dual cycle space.  The coherent
normalization sequence gives the same identification after tensoring
with $k$, with bijective Frobenius on all the graph terms.

The space of ordered general-position lines is irreducible; its
quotient by the freely acting group $S_d$ is the unordered arrangement
space.  Its ordered cover is connected, so its permutation monodromy
is $S_d$.  Over the ordered cover the graph local system is constant.
Thus the image $G\subseteq\mathrm{GL}(V)$ contains exactly the
permutation action in \eqref{eq:odd-arrangement-module}.

Let $U=\{(x_i)\in P:\sum x_i=0\}$, let $D$ be its image in $W$, and
put $A=\bigwedge^2D\subseteq\bigwedge^2W$.

\begin{lemma}\label{lem:odd-wedge-submodules}
\status{Proved under stated hypotheses}
Every nonzero $S_d$-submodule of $\bigwedge^2W$ contains $A$.
If $p\nmid d$, then $A=V$ and $V$ is simple.  If $p\mid d$, then
$A$ is the unique proper nonzero submodule, and $V/A\simeq D$ is
simple of dimension $d-2$.
\end{lemma}

\begin{proof}
For four distinct indices $i,j,k,l$, put $r=e_i-e_j$ and $s=e_k-e_l$.
The group-algebra operator
\[
 (1-(ij))(1-(kl))
\]
on $\bigwedge^2W$ has image $\langle r\wedge s\rangle$ and is the
rank-one map obtained by evaluating the functional
$(x_i-x_j)\wedge(x_k-x_l)$ and multiplying by $r\wedge s$.
The latter vector is nonzero: the line $\langle u\rangle$ does not
meet $\langle r,s\rangle$, including when $d=4$, since $p\ne2$.

The vectors $r\wedge s$ for disjoint pairs span $\bigwedge^2U$ before
passage to $W$, and hence span $A$ afterwards.  Explicitly, on four
indices take $f_i=e_i-e_l$, $f_j=e_j-e_l$, $f_k=e_k-e_l$.  The three
perfect matchings give the columns
\[
 \begin{pmatrix}0&1&1\\1&0&-1\\-1&1&0\end{pmatrix}
\]
in the basis $f_i\wedge f_j,f_i\wedge f_k,f_j\wedge f_k$; their
determinant is $2$.  Every pair of difference-basis vectors belongs to
such a four-index set.  The same argument in the dual permutation
space shows that the corresponding functionals span
$\bigwedge^2W^*$, because $W^*$ is the sum-zero subspace of $P^*$.
For every nonzero vector of $V$, therefore, one of these rank-one
operators has nonzero value.  Its entire $S_d$-orbit spans $A$.
This proves the first assertion and also proves that $A$ is simple.

If $p\nmid d$, then $U\to W$ is an isomorphism.  If $p\mid d$,
$D=U/\langle u\rangle$ has dimension $d-2$, and the sum functional
$\sigma$ descends to $W$ with kernel $D$.  Contraction with $\sigma$
gives the exact sequence
\[
 0\longrightarrow\bigwedge^2D\longrightarrow\bigwedge^2W
 \xrightarrow{\ \iota_\sigma\ }D\longrightarrow0.
\]
The module $D$ is simple: a nonzero class has a representative whose
coordinates are not all equal; applying $1-(ij)$ for two unequal
coordinates produces a nonzero multiple of $e_i-e_j$, and the
translates of this difference span $D$.  Every proper nonzero
submodule of $V$ contains $A$ and has zero image in the simple
quotient unless it is all of $V$.  This proves uniqueness.
\end{proof}

\begin{proposition}\label{prop:odd-plane-irreducible}
\status{Proved under stated hypotheses}
The geometric monodromy group $G$ is irreducible on $V$.
\end{proposition}

\begin{proof}
Only the case $p\mid d$ needs another family.  Take a smooth ordinary
curve of degree $d-2$ and a smooth conic meeting it transversely in
$m=2d-4$ points.  The graph contribution is a sub-local-system
\[
 B=\mathbf F_p^m/\langle(1,\ldots,1)\rangle\subset V
\]
of rank $m-1=2d-5$.  Its node monodromy is the full symmetric group
$S_m$.  To verify this, fix the conic.  Restriction of degree-$(d-2)$
plane forms to it is surjective onto binary forms of degree $m$.
Over the squarefree binary forms, the ordered-root cover is connected.
Pulling it back gives an affine bundle after choosing a coefficient
chart (equivalently, an open of a projective bundle).  Imposing
smoothness and ordinarity of the plane curve removes an open subset
from this irreducible total space, and the open is nonempty.  The
ordered-root cover remains connected.  Thus the induced graph image
is $S_m$.

Here $p\nmid m$, since $p\mid d$ and $p$ is odd.  The deleted
permutation module $B$ is simple, by the same difference argument
used above.  If the only possible proper $G$-invariant subspace $A$
were invariant, the map $B\to V/A$ would be zero: its source is a
simple module of dimension $2d-5>d-2=\dim(V/A)$.  Hence $B\subset A$.

This inclusion can be tested at the common line arrangement.  Allow
the conic to become the union of the first two lines and the other
curve to become the remaining $d-2$ lines.  Keep the intersections
between the two subcurves transverse, disjoint from their internal
nodes.  Their ordinary coherent cohomology and the graph sequence
extend over an open containing both this point and the smooth
subcurve locus.  At this arrangement $B$ is the kernel of restriction
to the complete subgraph on vertices $3,\ldots,d$.  In particular the
edge cochain $e_1\wedge e_3$ represents an element of $B$.  But
\[
 \iota_\sigma(e_1\wedge e_3)=e_3-e_1\ne0\quad\hbox{in }D,
\]
so it does not belong to $A$.  This contradiction proves the claim.
No specialization through a drop of the Hasse--Witt rank has been
used.
\end{proof}

\subsubsection{Quartic boundaries and an elementary cocycle calculation}

\begin{lemma}\label{lem:sl3-natural-h1-odd}
\status{Proved under stated hypotheses}
For every odd prime $p$, the natural three-dimensional module $Q$
and its dual satisfy
$H^1(\mathrm{SL}_3(\mathbf F_p),Q)=
H^1(\mathrm{SL}_3(\mathbf F_p),Q^*)=0$.
\end{lemma}

\begin{proof}
Let $c$ be a cocycle.  The diagonal sign subgroup of determinant one
has order four, prime to $p$, so subtracting a coboundary makes $c$
zero on it.  Each $t_{ij}(a)=1+aE_{ij}$ commutes with the sign matrix
which is $-1$ in positions $i,j$ and $+1$ in the remaining position
$k$.  The cocycle relation then forces
$c(t_{ij}(a))\in\langle e_k\rangle$.
Write $x=t_{12}(a)$, $y=t_{23}(b)$, with $a,b\ne0$, and
$c(x)=Ae_3$, $c(y)=Be_1$.  Since $[x,y]=t_{13}(ab)$, the cocycle
identity gives
\[
 c([x,y])=-bA e_2-abA e_1.
\]
Its left side belongs to $\langle e_2\rangle$, hence $A=0$.
Permuting the three indices proves vanishing on every elementary
transvection.  They generate $\mathrm{SL}_3(\mathbf F_p)$, so the
normalized cocycle is zero.  Applying the inverse-transpose
automorphism proves the dual assertion.
\end{proof}

For any chosen four of the $d$ lines, fix the other $r=d-4$ lines and
vary an ordinary quartic, transverse to the fixed lines and avoiding
their pairwise intersections.  Include ordinary singular quartics
whenever the total degree-$d$ curve stays in $W_d$.  This connected
open contains the original arrangement.  Mayer--Vietoris for the
quartic and the fixed lines gives
\begin{equation}\label{eq:odd-quartic-boundary}
 0\longrightarrow T\longrightarrow V\longrightarrow Q
 \longrightarrow0,\qquad \dim Q=3,
\end{equation}
where $Q$ is the quartic Hasse--Witt local system.  Its image is
$\mathrm{GL}_3(\mathbf F_p)$ by the canonical quartic moduli argument
and Achter--Pries, Proposition~4.8, already checked above.  The image
on $T$ is solvable: its only moving graph edges are the four
intersections with each fixed line, so it factors through a subgroup
of $(S_4)^r$.  The nodes between the fixed lines are fixed.

Let $B$ be the image of this boundary family and let $D_T$ be its
image on $T$.  The image of $B$ in
$\mathrm{GL}(Q)\times D_T$ is a subdirect product.  Goursat's lemma,
or the elementary description of a subdirect product as equality in
a common quotient, shows that it contains
$\mathrm{SL}(Q)\times1$: every common quotient is solvable, whereas
$\mathrm{SL}_3(\mathbf F_p)$ is perfect.  Let $B_0$ be the inverse
image of $\mathrm{SL}(Q)\times1$.  It is normal in $B$, fixes $T$
pointwise, and maps onto $\mathrm{SL}(Q)$.
Its kernel is an additive submodule
\begin{equation}\label{eq:odd-quartic-kernel}
 K\subseteq\operatorname{Hom}(Q,T)=Q^*\otimes T.
\end{equation}
The natural dual is absolutely simple, with endomorphism ring
$\mathbf F_p$.  A submodule of this direct sum of copies of $Q^*$
therefore has the form $K=Q^*\otimes T'$ for a subspace $T'\subseteq T$.
The subspace $T'$ is $B$-invariant because $K$ is normal in $B$.
If $K\ne0$, it contains rank-one maps $w\otimes\phi$ and hence
transvections $1+w\otimes\phi$ of $V$.  If $K=0$, $B_0$ is a
complement to $\operatorname{Hom}(Q,T)$ in the indicated semidirect
product.  Lemma~\ref{lem:sl3-natural-h1-odd}, applied to each copy of
$Q^*$, makes that complement conjugate to the block-diagonal
$\mathrm{SL}(Q)$.  In this case too it contains transvections of $V$.
For $d=4$ the statement simply recovers the full quartic group.

\subsubsection{The normal subgroup generated by transvections}

\begin{lemma}\label{lem:transvection-clifford-blocks}
\status{Proved under stated hypotheses}
Let an irreducible finite group $G\subseteq\mathrm{GL}(V)$ contain a
transvection, and let $N\triangleleft G$ be the subgroup generated by
all transvections in $G$.  Then $V$ is a direct sum of pairwise
nonisomorphic simple $N$-modules, and $G$ permutes these summands
transitively.  Every transvection has its image line in exactly one
summand and fixes all the others pointwise.
\end{lemma}

\begin{proof}
The sum of the simple $N$-submodules, $\operatorname{soc}_N(V)$, is
nonzero and $G$-stable, since $N$ is normal.  Irreducibility of $G$
therefore makes it all of $V$.  The canonical isotypic summands are
permuted transitively by $G$.  Write any one of them as $S^{\oplus e}$.
For an element of $N$, the rank of its difference from the identity
on this summand is a multiple of $e$.  Some transvection has nonzero
action on some summand, and its total rank is one, so $e=1$ there;
transitivity gives $e=1$ for all summands.  Thus the canonical
isotypic summands are themselves simple.  Each transvection is
block diagonal because it belongs to $N$, and its rank-one
difference is nonzero on only one block.  This proves the assertion.
(The same argument with the finite division endomorphism ring of a
block shows it is absolutely simple, but that strengthening is not
needed below.)
\end{proof}

For a four-element subset $I\subset\{1,\ldots,d\}$, let
$U_I\subset P$ be its three-dimensional sum-zero permutation
subspace, supported on $I$.  Because $4\ne0$ in $\mathbf F_p$,
$U_I$ embeds in $W$.  Set
\[
 S_I=\bigwedge^2 U_I\subset V.
\]
The subgroup $S_I^{\rm perm}\simeq S_4$ permuting these four lines
lies in the corresponding quartic boundary group $B$.  Restricted
to it, the decomposition of the permutation space gives
\begin{equation}\label{eq:s4-quartic-decomposition}
 W\simeq U_I\oplus\mathbf F_p^{d-4},\qquad
 V\simeq S_I\oplus U_I^{\oplus(d-4)}
             \oplus\mathbf F_p^{\binom{d-4}{2}}.
\end{equation}
Both $U_I$ and $S_I$ are simple.  For the former use the elementary
deleted-permutation argument with $p\nmid4$; for the latter use
Lemma~\ref{lem:odd-wedge-submodules} with four indices.  They are
nonisomorphic: a transposition has trace $1$ on $U_I$ and trace $-1$
on $\bigwedge^2U_I$, distinct in every odd characteristic.  Neither
is trivial.  Thus \eqref{eq:s4-quartic-decomposition} is a semisimple
decomposition even when $p=3$, and $S_I$ occurs with multiplicity one.
The quotient $Q$ in \eqref{eq:odd-quartic-boundary} is isomorphic to
$S_I$ as an $S_4$-module, and $T$ contains no constituent of this type.

\begin{proposition}\label{prop:odd-transvection-normal-irreducible}
\status{Proved under stated hypotheses}
The normal subgroup $N\triangleleft G$ generated by all transvections
is irreducible on $V$.
\end{proposition}

\begin{proof}
Use the canonical $N$-block decomposition of
Lemma~\ref{lem:transvection-clifford-blocks}.  We first prove that for
each quartet $I$ its subspace $S_I$ lies entirely in one block.

If the kernel $K$ of its boundary group $B_0$ is zero, the cocycle
calculation makes $B_0$ a conjugate block-diagonal $\mathrm{SL}_3$.
Its commutator subspace $P_I=[B_0,V]$ is therefore three dimensional,
simple as a $B_0$-module, and maps isomorphically to $Q$.  Normality
of $B_0$ in $B$ makes $P_I$ stable under the quartet's $S_4$.
The unique-copy statement in \eqref{eq:s4-quartic-decomposition}
forces $P_I=S_I$.  This copy of $\mathrm{SL}_3$ is generated by
rank-one transvections of $V$, so $B_0\subseteq N$.  Consequently
each $N$-block is $B_0$-stable and
$[B_0,V]=\bigoplus_j[B_0,V_j]$.  Simplicity of $P_I$ puts it in one
$N$-block.

If $K\ne0$, write it as $Q^*\otimes T'$, with $T'\ne0$, and fix
one nonzero $\phi\in Q^*$.  For every nonzero $w\in T'$, the
transvection $1+w\otimes\phi$ is supported on a single $N$-block.
The functional here means the pullback of $\phi$ along $V\to Q$.
It is nonzero on that block and zero on every other block.  Since
the same nonzero functional is used for all $w$, that block must
be the same block $L$ for every $w$.  Thus $T'\subset L$.
Now fix one nonzero $w$ and vary $\phi$.  Every $\phi\in Q^*$
vanishes on every other $N$-block.  It follows that $L\to Q$ is
surjective.  Since $T'$ is stable under the quartet's $S_4$, this
$S_4$ must preserve $L$: its translate is another $N$-block sharing
the nonzero subspace $T'$.  As an $S_4$-submodule of the semisimple
module \eqref{eq:s4-quartic-decomposition}, $L$ is semisimple.
Its surjection onto $Q\simeq S_I$ forces it to contain the unique
copy $S_I$ in $V$.  This proves the claim in both cases.

When two quartets have three common indices, their subspaces $S_I$
intersect in the nonzero line $\bigwedge^2U_{I\cap J}$.
The graph of four-element subsets joined by such overlaps is
connected.  Hence all $S_I$ belong to the same $N$-block $L$.
Their span is $A=\bigwedge^2D$: the spanning disjoint-difference
wedges in Lemma~\ref{lem:odd-wedge-submodules} are supported on
quartets.  Thus $A\subset L$.

The subspace $A$ is $S_d$-stable.  Every translate of $L$ by $S_d$
is an $N$-block containing $A$, hence is $L$ itself.  Thus $S_d$
preserves $L$ and the sum $L'$ of the other blocks.  If $L'\ne0$,
Lemma~\ref{lem:odd-wedge-submodules} forces $A\subset L'$, which
contradicts $A\subset L$ and $L\cap L'=0$.  There are no other
blocks, and $N$ is irreducible.  The base case $d=4$ was already
settled by the quartic monodromy theorem.
\end{proof}

\subsubsection{Classification, exclusion of the alternating form, and determinant}

McLaughlin's transvection theorem, in the precise form of
Kantor--Seress (2001), Theorem~3.6, p.~22 \cite{KantorSeress2001},
states that for $n\geq3$
an irreducible subgroup of $\mathrm{SL}_n(\mathbf F_q)$ generated by
full transvection groups is $\mathrm{SL}_n(\mathbf F_q)$ or, when
$n$ is even, $\mathrm{Sp}_n(\mathbf F_q)$; the additional alternatives
in that theorem occur only when $q=2$.  Here $q=p$ is prime, and
the powers of $1+v\phi$ are precisely $1+a v\phi$, $a\in\mathbf F_p$.
Thus every generating transvection supplies its full transvection
group.  The theorem applies to $N$, which has determinant one and
is irreducible by Proposition~\ref{prop:odd-transvection-normal-irreducible}.
It leaves $N=\mathrm{SL}_g(\mathbf F_p)$ or $N=\mathrm{Sp}_g(\mathbf F_p)$.

The symplectic possibility is impossible.  If $N$ is symplectic,
normality puts $G$ in its group of similitudes: the invariant
alternating form is unique up to scalar.  For any quartet, the
multiplier on $B_0$ is trivial.  Indeed it is trivial on its
$p$-group kernel $K$ and on its perfect quotient
$\mathrm{SL}_3(\mathbf F_p)$.  Therefore $B_0$ preserves the
nondegenerate alternating form.  Since it fixes $T$ pointwise and
has fixed-point-free natural quotient $Q$, its fixed space is
$V^{B_0}=T$.  For a group preserving a nondegenerate bilinear form,
\[
 [B_0,V]=(V^{B_0})^\perp=T^\perp.
\]
The latter space has dimension three.  Its image in $Q$ is
$[B_0,Q]=Q$, so it maps isomorphically to $Q$ and meets $T$ trivially.
The alternating form is consequently nondegenerate on this
three-dimensional space $T^\perp$, a contradiction.  We have proved
\begin{equation}\label{eq:odd-plane-sl-contained}
 \mathrm{SL}_g(\mathbf F_p)\subseteq G.
\end{equation}

It remains to show that the geometric determinant is all of
$\mathbf F_p^\times$.  Take a Legendre elliptic cubic $E_t$, near a
supersingular parameter $t_0$, and add $d-3$ fixed lines chosen in
general position relative to $E_{t_0}$.  The Hasse polynomial of
the Legendre family has a simple root at $t_0\notin\{0,1\}$, as
proved in the elliptic calculation in
\S\ref{subsec:plane-residues}.  Over the strict local
trait at $t_0$, all graph contributions have bijective Frobenius
and determinant a unit.  The normalization sequence therefore
makes the total Hasse determinant a unit times the elliptic Hasse
invariant.  It has valuation exactly one at $t_0$.
On the punctured trait the total curve is ordinary, hence belongs
to $W_d$.  The determinant Lang torsor is
$z^{p-1}=h^{-1}$ by \eqref{eq:determinant-kummer}.  A valuation-one
$h$ gives tame inertia image of exact order $p-1$.  This is
geometric inertia, since the residue field is algebraically closed.
Thus $\det G=\mathbf F_p^\times$.  Combining this with
\eqref{eq:odd-plane-sl-contained} proves the theorem.

\begin{theorem}\label{thm:plane-odd-all-groups}
\status{Proved under stated hypotheses}
Let $p$ be odd and let $U_d$ be the ordinary smooth plane-curve locus
of degree $d$ over any perfect field $k$ of characteristic $p$.  At a
geometric generic point put
$V=H^1_{\mathrm{et}}(C_{\bar\eta},\mathbf F_p)$, of dimension
$g=(d-1)(d-2)/2$.  For every $d\geq3$, both the geometric and the
arithmetic monodromy images on $V$ are
$\mathrm{GL}_g(\mathbf F_p)$, with full determinant character.
For $d=1,2$ both images are trivial on the zero vector space.
For positive rank the representation is irreducible.
\end{theorem}

\begin{proof}
The preceding argument proves the geometric assertion for $d\geq4$
over $\overline{\mathbf F}_p$.  The connected finite etale frame cover
which it identifies is geometrically connected and remains connected
after every algebraically closed field extension.  Since the family
and its Frobenius local system are defined over $\mathbf F_p$, this
gives the same geometric image over the algebraic closure of any
perfect field of characteristic $p$.
The arithmetic image contains that image and is a subgroup of the
same finite general linear group, so it is equal to it.  The cubic
case follows from the same Legendre determinant calculation, since
its rank is one.  Degrees one and two have rank zero.
\end{proof}

% /root/monodromy, 19 November 2026.
% Independently audited in full by /root, including the determinant seeds.
\subsection{Hypersurfaces in projective space in odd characteristic}
\label{subsec:projective-space-odd-monodromy}
\status{Proved; independently audited by /root, including both determinant seeds.}

Fix an odd prime $p$ and an algebraically closed field $k$ of
characteristic $p$.  For degree $d$ hypersurfaces in $\mathbf P^n$,
let $W_{n,d}$ be the Hasse--Witt invertibility open in the full
projective space of equations, and let
\[
 h_{n,d}=\binom{d-1}{n},\qquad
 V_{n,d}=\left(H^{n-1}(\mathcal O_Y),F=1\right).
\]
The binomial coefficient is zero for $d\leq n$.
The full coherent bundle has this constant rank for all nonzero
equations, including singular equations.  A general-position union
of $d$ hyperplanes belongs to $W_{n,d}$, so this open is nonempty.
The desired smooth ordinary locus is a dense open in $W_{n,d}$;
passing to it does not change the finite monodromy image.
All the comparison families below therefore take place inside
one explicit constant-rank Lang local system, except for a
valuation-one Hasse degeneration used solely to compute its tame
determinant inertia.

\subsubsection{A connectedness lemma for linear projection}

\begin{lemma}\label{lem:linear-projection-monodromy}
\status{Proved under stated hypotheses}
Let $U$ be a nonempty open in $\mathbf P^N_k$, and let
$T\to U$ be a connected finite etale cover.  Let
$f:\mathbf P^N\dashrightarrow\mathbf P^M$ be a linear projection
whose center $Z$ meets $U$.  The restriction of $T$ to the
geometric generic fibre of $f|_{U\setminus Z}$ is connected.
Consequently a finite local system on $U$ retains its full
monodromy after restriction to that geometric generic fibre.
\end{lemma}

\begin{proof}
The open $U$ is smooth and integral, so $T$ is smooth and integral.
Let $B=\operatorname{Bl}_Z(\mathbf P^N)$ and put
$T'=T\times_U B_U$, where $B_U$ is the inverse image of $U$.
This is the blowup of $T$ in the smooth etale inverse image of
$Z\cap U$; in particular it is smooth and integral.  The resolved
projection gives a morphism $T'\to\mathbf P^M$.
Choose $z\in Z\cap U$.  The exceptional fibre above $z$ is
$E_z\simeq\mathbf P^M$, mapping isomorphically to the target.
Its pullback to $T'$ is the disjoint union of copies of $E_z$,
since it is the product of $E_z$ with the finite reduced set $T_z$.

Put $K=k(\mathbf P^M)$, and let $L$ be the algebraic closure of $K$
inside $k(T')$.  It is a finite extension of $K$.  Let $C$ be the
normalization of $\mathbf P^M$ in $L$; it is finite and integral.
There is a morphism $T'\to C$ over $\mathbf P^M$.  Indeed, on an
affine open of the target, each element of its integral closure
in $L$ is integral over the local ring at every point of the
normal scheme $T'$ lying above this open.  It therefore belongs
to each such local ring and defines a regular function there.
These morphisms glue.  One copy of $E_z$ in $T'$ consequently
gives a section $\mathbf P^M\to C$ of the finite map.
Its image is closed, has dimension $M$, and lies in the integral
scheme $C$ of dimension $M$; it is all of $C$.  Thus $C\simeq
\mathbf P^M$ and $L=K$.

On the nonempty open $T|_{U\setminus Z}$, the map to $\mathbf P^M$
is a composition of a finite etale map and a smooth linear
projection.  Hence $k(T')/K$ is separably generated.  The equality
$L=K$ implies that this extension is regular, so its geometric
generic fibre is integral.  Removing the exceptional divisor
does not alter that assertion.  For the local-system consequence,
apply the proved assertion to a connected component of its finite
frame cover, which is a torsor for its full monodromy image.
\end{proof}

\begin{corollary}\label{cor:restriction-independent-monodromy}
\status{Proved under stated hypotheses}
For $e\geq n+1$, fix a hyperplane $H\subset\mathbf P^n$.
The local system $V_{n,e}$ retains its full geometric monodromy
when the restriction of the equation to $H$ is prescribed to be
geometric generic.  On the parameter space where both the
equation and its restriction have invertible Hasse--Witt maps,
the joint monodromy image on $V_{n,e}$ and $V_{n-1,e}$ is the
product of their two full images.
\end{corollary}

\begin{proof}
Restriction of equations is a linear projection.  Its center
consists of equations divisible by the fixed equation of $H$.
It meets $W_{n,e}$: take $H$ and $e-1$ other general-position
hyperplanes.  Lemma~\ref{lem:linear-projection-monodromy} shows
that the joint image contains the first full group times $1$.
Its projection to the second group is surjective, since the
restriction map has geometrically integral generic fibre, an
open of an affine space.  The joint image is therefore the
product.  Removing further nonempty opens does not change
these finite images.
\end{proof}

\subsubsection{The arrangement module and independent quadratic reflections}

Take a general-position union of $d$ hyperplanes as base point.
Put $P=\mathbf F_p^d$, $u=\sum_i e_i$, $W=P/\langle u\rangle$,
$U=\ker(\sigma:P\to\mathbf F_p)$ with $\sigma(x)=\sum_i x_i$,
and $D=\operatorname{im}(U\to W)$.  The Cech resolution of the
union, whose strata are projective spaces, identifies its top
coherent cohomology and its Lang fibre with
\begin{equation}\label{eq:pn-arrangement-module}
 V_{n,d}=\bigwedge^nP/(u\wedge\bigwedge^{n-1}P)
       =\bigwedge^nW.
\end{equation}
The connected ordered-hyperplane cover of the unordered
arrangement locus gives the full permutation action of $S_d$
on this space.  Put $A=\bigwedge^nD$.

\begin{lemma}\label{lem:odd-general-hook-module}
\status{Proved under stated hypotheses}
Assume $1\leq n\leq d-2$.  If $p\nmid d$, then
$\bigwedge^nW$ is simple as an $S_d$-module.  If $p\mid d$, then
$A=\bigwedge^nD$ is its unique proper nonzero submodule; both
$A$ and the quotient $\bigwedge^{n-1}D$ are simple.
\end{lemma}

\begin{proof}
First suppose $2n\leq d$.  For $n$ pairwise disjoint transpositions
$\tau_j=(a_jb_j)$, the group-algebra operator
$\prod_j(1-\tau_j)$ has rank one on $\bigwedge^nW$.  Its image
is the wedge of the $n$ differences $e_{a_j}-e_{b_j}$, and its
functional is the wedge of the corresponding coordinate
differences.  This follows by simultaneously decomposing $W$
into the $n$ negative eigenlines and their common positive
eigenspace; $2$ is invertible.

These matching wedges span $\bigwedge^nU$.  To check this without
a modular representation-theory citation, use the difference
basis $f_i=e_i-e_o$ and the identity, for four distinct indices,
\[
 2f_a\wedge f_b
 =(e_a-e_b)\wedge(e_c-e_o)
  +(e_a-e_c)\wedge(e_b-e_o)
  +(e_a-e_o)\wedge(e_b-e_c).
\]
Beginning with a wedge of $n$ differences sharing $o$, repeatedly
introduce one unused index $c$.  Each application leaves one
fewer difference sharing $o$ and produces one disjoint pair.
There are $n-1$ available new indices because $d\geq2n$.
The resulting wedges are perfect matchings on $2n$ indices.
The same proof in the dual permutation space shows that the
matching functionals span $\bigwedge^nW^*$.
Thus every nonzero submodule contains a matching image and hence
all of $A$.  This proves simplicity of $A$, and proves simplicity
of the whole space when $p\nmid d$.

When $p\mid d$, the dot product gives a nondegenerate invariant
form on $D=U/\langle u\rangle$.  The same matching argument for
$n-1$ shows that $\bigwedge^{n-1}D$ is simple: matching images
span it, and their functionals span its dual under this form.
Contraction with the sum functional gives an exact sequence
\[
 0\longrightarrow\bigwedge^nD\longrightarrow\bigwedge^nW
 \longrightarrow\bigwedge^{n-1}D\longrightarrow0.
 \tag{*}
\]
Every nonzero submodule contains its kernel, and a submodule
which also surjects to the simple quotient is the whole space.
This proves the unique-submodule assertion in the stated range.

For $2n>d$, put $r=d-1-n$, so $1\leq r<d/2$.
The natural exterior pairing gives
$(\bigwedge^nW)^*\simeq\bigwedge^rW\otimes\operatorname{sgn}$.
If $p\nmid d$, the already proved simplicity for $r$ gives the
claim.  If $p\mid d$, duality preserves the property of having
exactly one proper nonzero submodule.  The displayed subspace
$\bigwedge^nD$ is nonzero and proper, so it must be that unique
submodule.  Its simplicity and that of the quotient follow
from the corresponding two simple factors for $r$.
\end{proof}

\begin{lemma}\label{lem:quadric-reflections}
\status{Proved under stated hypotheses}
For $n\geq2$ and $d\geq n+2$, the geometric group on
$V_{n,d}$ contains commuting rank-one reflections
$R_I$, indexed by $(n-1)$-subsets
$I\subset\{3,\ldots,d\}$.  At the arrangement base point,
their negative eigenlines are
\[
 \left\langle (e_1-e_2)\wedge\bigwedge_{i\in I}e_i\right\rangle,
\]
and their product is the permutation transposition $(12)$.
Here a rank-one reflection has eigenvalue $-1$ once and
eigenvalue $1$ on a codimension-one subspace.
\end{lemma}

\begin{proof}
Fix the last $d-2$ hyperplanes and replace the first two by a
variable quadric $Q$.  Require every positive-dimensional
quadric stratum to be connected with no higher coherent
cohomology, and every intersection with a line
$L_I=\bigcap_{i\in I}H_i$ to consist of two distinct points.
This is a nonempty open containing a smooth general quadric
and the reducible quadric $H_1H_2$.  A positive-dimensional
quadric hypersurface has $H^0(\mathcal O)=k$ and
$H^{>0}(\mathcal O)=0$, so the Cech complex computes all the
cohomology from its $H^0$ terms.  Only its top term has
nonconstant point sets: the two points on each $L_I$.
Swapping one pair acts on cohomology as a rank-one reflection,
since every coboundary has equal values on that pair.

All these swaps occur independently.  The restriction from
quadratic forms to $L_I\simeq\mathbf P^1$ is surjective.  The
discriminant of a binary quadratic is an irreducible polynomial
in its three coefficients in odd characteristic.  Its pullback
to the space of all quadratic forms is still irreducible, and
the discriminants for distinct lines are nonassociate.  To see
the last assertion, choose a point of $L_I\setminus L_J$ and
two hyperplanes through that point which meet $L_J$ in distinct
points.  Their product is tangent to $L_I$ and transverse to
$L_J$.  Their classes in the rational function field
modulo squares are therefore independent, by valuations at
these distinct prime divisors.  The ordered-pair cover has
Galois group $(\mathbf Z/2)^{\binom{d-2}{n-1}}$.
Removing the indicated open conditions leaves it connected.

The pair differences at $Q=H_1H_2$ are precisely the displayed
cochains, using the refinement of the cover with component $Q$
by its two hyperplanes.  Interchanging these two hyperplanes
interchanges every pair and fixes every other Cech term.
The transposition $(12)$ is thus the product of all the $R_I$.
Every fibre here has bijective Frobenius on the full top
coherent cohomology, because all terms are finite sets with
ordinary constant Frobenius.  The entire comparison is inside
$W_{n,d}$.
\end{proof}

\begin{proposition}\label{prop:pn-odd-irreducible}
\status{Proved under stated hypotheses}
For $n\geq2$ and $d\geq n+2$, the geometric monodromy group
$G_{n,d}$ is irreducible.
\end{proposition}

\begin{proof}
Only $p\mid d$ needs consideration, by
Lemma~\ref{lem:odd-general-hook-module}.  In that case any
proper invariant subspace must be $A=\bigwedge^nD$.
The negative eigenline of every $R_I$ lies outside $A$:
contraction with $\sigma$ sends its displayed generator to
\[
 -(e_1-e_2)\wedge\iota_\sigma
                  \left(\bigwedge_{i\in I}e_i\right)\ne0.
 \tag{**}
\]
Indeed the involved $n$ vectors remain linearly independent
modulo $u$: at least one index is outside $\{1,2\}\cup I$.
If $A$ were invariant, each $R_I$ would consequently fix $A$
pointwise.  Their product $(12)$ would do so too.  But on $D$
this transposition has a one-dimensional negative eigenspace
and a positive eigenspace of dimension $d-3$; on
$\bigwedge^nD$ its negative eigenspace has dimension
$\binom{d-3}{n-1}>0$.  This contradicts pointwise fixation.
\end{proof}


\subsubsection{The degree $n+2$ seed}

We use the following elementary observation about the Cech
computations, also in the determinant calculation below.  A family
consisting of coordinate hyperplanes and one further hypersurface
has the usual exact alternating resolution by its intersections
provided the latter equation restricts to a nonzero equation on
every coordinate stratum.  Locally its vanishing equation and any
of the coordinate equations form a regular sequence.  The
alternating union complex for a regular sequence is exact, by
induction using $(f)\cap(g)=(fg)$ and its versions modulo the
other equations.  Thus smoothness of the further hypersurface is
unnecessary.  Higher coherent cohomology of any hypersurface
stratum is given by its usual projective-space exact sequence.

\begin{proposition}\label{prop:pn-odd-seed-sl}
\status{Proved under stated hypotheses}
For $n\geq3$ and $d=n+2$, the group $G_{n,d}$ contains
$\mathrm{SL}_{n+1}(\mathbf F_p)$.
\end{proposition}

\begin{proof}
Let $N$ be the normal subgroup of $G_{n,d}$ generated by all its
transvections.  Replace any four of the arrangement hyperplanes,
indexed by $I$, by a varying quartic $Q$, and fix the other
$n-2$ hyperplanes.  The Cech spectral sequence has precisely
the following positive-cohomology terms contributing to total
degree $n-1$:
the quartic plane curve obtained by intersecting all the fixed
hyperplanes, of rank three, and the $n-2$ quartic surface strata,
each of rank one.  Higher-dimensional quartic strata have no
positive coherent cohomology.  The $H^0$ row has no positive
cohomology because the nerve is a full simplex.  There are no
nonzero incoming or outgoing differentials for the displayed
terms.  We consequently obtain an exact sequence of Frobenius
modules, and of Lang local systems on the invertibility open,
\[
 0\longrightarrow Q_I\longrightarrow V_{n,n+2}
      \longrightarrow T_I\longrightarrow0,
 \qquad \dim Q_I=3,\quad\dim T_I=n-2.
 \tag{1}
\]
Here $T_I$ is the direct sum of the rank-one surface local systems.
The open where all these terms are ordinary contains the original
arrangement.  Its image on $Q_I$ is $\mathrm{GL}_3(\mathbf F_p)$:
restriction of quartics to the fixed plane is a linear projection
with affine generic fibre, and the ordinary plane-quartic group
is already proved.  The inverse image of its connected frame
cover remains integral; imposing the other nonempty open
conditions preserves connectedness.  The image on $T_I$ is
abelian, being a subgroup of $(\mathbf F_p^\times)^{n-2}$.

Apply the quartic-boundary argument of the preceding plane
section to the dual of (1).  Goursat's lemma gives a normal
subgroup $B_0$ fixing $T_I^*$ pointwise and mapping onto
$\mathrm{SL}(Q_I^*)$.  Its kernel is
$Q_I\otimes T'$ for a subspace $T'\subseteq T_I^*$.
If the kernel is nonzero, it contains rank-one transvections.
If it is zero, Lemma~\ref{lem:sl3-natural-h1-odd} splits the
extension and again gives transvections.  Thus $N$ is nontrivial.
By Proposition~\ref{prop:pn-odd-irreducible} and
Lemma~\ref{lem:transvection-clifford-blocks}, $V_{n,n+2}^*$ is a
direct sum of canonical simple $N$-blocks permuted transitively
by $G_{n,n+2}$.

In the arrangement fibre there is the $S_d$-equivariant
identification
\[
 V_{n,n+2}^* \simeq W\otimes\operatorname{sgn},
 \qquad d=n+2.
 \tag{2}
\]
For a quartet $I$, let $U_I$ denote its three-dimensional
sum-zero subspace in $P$.  Since $4$ is invertible, it embeds
in $W$.  The restriction of (2) to its $S_4$ is the direct sum
of the unique simple module
$S_I=U_I\otimes\operatorname{sgn}$, of dimension three, and
$d-4$ copies of the sign character.  This is an actual
semisimple decomposition in every odd characteristic.
The quotient $Q_I^*$ in the dual of (1) is this unique
three-dimensional constituent.

The two cases in the proof of
Proposition~\ref{prop:odd-transvection-normal-irreducible} now
apply verbatim, with this specified decomposition: if the kernel
is zero, the commutator space of the split $\mathrm{SL}_3$ is
$S_I$ and belongs to one $N$-block; if it is nonzero, the
rank-one maps with a common quotient functional put $T'$ in one
block, and varying the functional makes that block surject to
$Q_I^*$.  The block is $S_4$-stable because $T'$ is stable, so
it contains the unique constituent $S_I$.

Quartets with three common indices give $S_I$'s with nonzero
intersection.  Their connected overlap graph puts all these
subspaces in one block $L$.  They span
$D\otimes\operatorname{sgn}$.  Every translate of $L$ by $S_d$
contains this same nonzero subspace and must equal $L$.
Thus $S_d$ also preserves the sum of the remaining blocks.
If that sum were nonzero, the degree-one case of
Lemma~\ref{lem:odd-general-hook-module}, with the sign twist,
would force it to contain $D\otimes\operatorname{sgn}$ too.
This contradicts the direct sum.  Therefore $N$ is irreducible.

Kantor--Seress, Theorem~3.6 \cite{KantorSeress2001}, applies:
over the prime field, powers of each transvection supply its
full transvection group.  It leaves
$N=\mathrm{SL}_{n+1}(\mathbf F_p)$ or a symplectic group.
The latter is excluded by
Lemma~\ref{lem:pn-reflection-excludes-symplectic} below.
\end{proof}

\begin{lemma}\label{lem:pn-reflection-excludes-symplectic}
\status{Proved under stated hypotheses}
If an irreducible normal subgroup $N\triangleleft G$ is the
natural symplectic group in dimension $h\geq3$ over an odd prime
field, then $G$ cannot contain a rank-one reflection.
\end{lemma}

\begin{proof}
The space of $N$-invariant alternating forms is one dimensional,
so normality puts $G$ in the symplectic similitude group.
Let $R$ be such a reflection, with positive eigenspace of
dimension $h-1$ and negative eigenline.  If its multiplier
were different from one, its positive eigenspace would be
totally isotropic, contradicting $h-1>h/2$.
Its multiplier is therefore one.  The negative eigenline is
then orthogonal to the positive eigenspace, since $p\ne2$,
and to itself because the form is alternating.  It is in the
radical of the form, a contradiction.
\end{proof}

\subsubsection{Induction in the degree}

\begin{proposition}\label{prop:pn-odd-degree-induction}
\status{Proved under stated hypotheses}
Suppose the groups for degree $d-1$ in $\mathbf P^n$ and
$\mathbf P^{n-1}$ contain their full special linear groups,
where $n\geq3$ and $d\geq n+3$.  Then
$G_{n,d}$ contains $\mathrm{SL}_{h_{n,d}}(\mathbf F_p)$.
If in addition the first of these two groups has full
determinant, then so does $G_{n,d}$.
\end{proposition}

\begin{proof}
Use the boundary $Y=A\cup H$, with $A$ a degree-$(d-1)$
hypersurface and $H$ a fixed hyperplane.  The usual exact
sequence of the union gives, on the simultaneous ordinary open,
\begin{equation}\label{eq:pn-two-piece-boundary}
 0\longrightarrow S\longrightarrow V_{n,d}
       \longrightarrow Q\longrightarrow0,
 \qquad
 S=V_{n-1,d-1},\quad Q=V_{n,d-1}.
\end{equation}
There are no adjacent positive coherent groups, since a
hypersurface in projective space has no intermediate coherent
cohomology.  Put
\[
 r=\dim Q=\binom{d-2}{n},\qquad
 s=\dim S=\binom{d-2}{n-1}.
\]
Both are at least three in the present range.
Corollary~\ref{cor:restriction-independent-monodromy} shows
that the graded image is the product of the full groups.
Take the inverse image $B_0$ of
$L=\mathrm{SL}(S)\times\mathrm{SL}(Q)$.
Its kernel $K$ is an additive $L$-submodule of
$M=\operatorname{Hom}(Q,S)$.  This outer tensor product of
the two absolutely simple natural modules is absolutely simple,
so $K=0$ or $K=M$.

If $K=M$, its rank-one elements are transvections of $V$.
Let $N$ be their full normal closure in $G_{n,d}$ and use
the canonical $N$-blocks, justified by irreducibility of $G$.
For every $w\in S$ and $\phi\in Q^*$, the operator
$1+w\otimes\phi$ is a transvection when both are nonzero.
Fixing $\phi$ puts all of $S$ in one block; fixing a nonzero
$w$ and varying $\phi$ makes that block surject to $Q$.
It contains $S$ and surjects to $V/S$, so it is all of $V$.
Thus $N$ is irreducible.

If $K=0$, the complement is conjugate to the split block
diagonal group $L$.  Here is the needed cohomology computation,
which requires no vanishing theorem for either natural module:
\[
 H^1(\mathrm{SL}(S)\times\mathrm{SL}(Q),S\otimes Q^*)=0.
\]
The first factor has no invariants in $S$.  In the low-degree
Hochschild--Serre sequence the remaining term is
$H^1(\mathrm{SL}(S),S)\otimes Q^*$, whose invariants under the
second factor are zero.  This proves the displayed vanishing
and hence the conjugacy of complements.
The split $L$ is generated by rank-one transvections on $V$,
so is contained in $N$.  Every canonical $N$-block is therefore
an $L$-submodule of the semisimple, multiplicity-free module
$S\oplus Q$.  If there is more than one block, they are exactly
these two summands.  Transitivity of $G$ on the blocks forces
$r=s$, equivalently $d=2n+1$.

Even this last possibility is impossible.  A rank-one reflection
cannot interchange two subspaces of common dimension $r>1$:
its difference from the identity would have rank at least $r$
on either one.  Thus every reflection of
Lemma~\ref{lem:quadric-reflections} preserves both blocks.
Their product $(12)$ does so too.  Conjugating within the
arrangement $S_d$ shows that every transposition, hence all
of $S_d$, preserves both blocks.  This contradicts
Lemma~\ref{lem:odd-general-hook-module}: a simple module, or a
module with a unique proper nonzero submodule, has no such
nontrivial direct sum decomposition.  We have proved that
$N$ is irreducible also when $K=0$.

The transvection classification again gives $N=\mathrm{SL}(V)$
or a symplectic group.  The latter is excluded by
Lemma~\ref{lem:pn-reflection-excludes-symplectic}, using the
quadric reflection already present in $G$.
Finally the full product assertion for the graded image
contains the old group on $Q$ times $1$ on $S$.  Its lifts
have the same determinant as their action on $Q$, because
the extension kernel is unipotent.  Full determinant for
that old group therefore implies full determinant for $G$.
\end{proof}

\subsubsection{Determinants in the two initial degrees}

\begin{lemma}\label{lem:elliptic-fixed-three-points}
\status{Proved under stated hypotheses}
There is a family of plane cubics over a strictly henselian
smooth trait over $k$ whose special fibre is smooth supersingular,
whose Hasse invariant has valuation one, and whose restriction
to a fixed line is the constant squarefree binary cubic
$g=x_0x_1(x_0-x_1)$.  There are also smooth ordinary plane
cubics with this same restriction.
\end{lemma}

\begin{proof}
Take the Legendre family at a simple supersingular zero of its
Hasse polynomial, as proved in the earlier elliptic calculation.
Choose a line transverse to the special fibre and identify it
with $\mathbf P^1$.  Transversality persists on the trait, and
the three intersection points form a finite etale scheme.
They split after strict henselization.  The unique projective
linear transformation on the line taking their chosen order
to $0,1,\infty$ varies regularly with the trait parameter.
Lift this transformation to the ambient projective plane and
rescale the cubic equation by a unit.  Its restriction to the
line is then exactly $g$.  Isomorphisms and unit changes of
cohomology trivialization multiply the Hasse invariant by a
unit, so its valuation is still one.  The same construction
with any ordinary elliptic plane cubic gives the last assertion.
\end{proof}

\begin{proposition}\label{prop:pn-seed-full-determinant}
\status{Proved under stated hypotheses}
For every $n\geq2$, the geometric determinant images in degrees
$n+1$ and $n+2$ are the full group $\mathbf F_p^\times$.
\end{proposition}

\begin{proof}
For degree $n+1$, fix $n-2$ coordinate hyperplanes and let their
intersection be a plane.  Lift the cubic family of
Lemma~\ref{lem:elliptic-fixed-three-points} from this plane to
$\mathbf P^n$, for instance by its same equation independent
of the other coordinates, and take its union with the fixed
hyperplanes.  Their intersections are proper.  All
positive-dimensional cubic strata of ambient dimension at
least three have no positive coherent cohomology; the only
nonzero term in total degree $n-1$ is the elliptic curve term.
The $H^0$ nerve is a full simplex and has no higher cohomology.
Thus the rank-one total Hasse invariant is a unit times the
elliptic one and has valuation exactly one.
For $n=2$ this is simply the elliptic family itself.

For degree $n+2$, use coordinates
$x_0,x_1,y_1,\ldots,y_{n-1}$ and hyperplanes $H_i=(y_i=0)$.
Their common line is $\ell=\mathbf P(x_0,x_1)$.
Let $f_1(x_0,x_1,y_1)$ be the family in the lemma, and for
$i>1$ let $f_i(x_0,x_1,y_i)$ be a fixed ordinary plane cubic,
all restricting to $g$ on $\ell$.  Define explicitly
\[
 q=g+\sum_{i=1}^{n-1}(f_i-g),\qquad
 Y=(q\,y_1\cdots y_{n-1}=0)\subset\mathbf P^n.
 \tag{3}
\]
Restriction to the plane where all $y_j$, $j\ne i$, vanish
is exactly $f_i$.  Every coordinate stratum has a nonzero
restricted cubic equation, because its restriction to $\ell$
is $g$.  The alternating union resolution described above
is therefore exact, even if the higher-dimensional cubic
strata happen to be singular.

Those higher-dimensional cubic strata have no positive
coherent cohomology, and every positive-dimensional stratum
has $H^0(\mathcal O)=k$.  The common zero-dimensional stratum
consists of the three constant, distinct roots of $g$.
The $H^0$ row is the full-simplex complex except that its last
copy of $k$ is replaced by $k^3$; it contributes an ordinary
rank-two piece in degree $n-1$.  In the positive-cohomology
rows the $n-1$ elliptic curves contribute their rank-one
cohomology, again in total degree $n-1$.  There are no
possible nonzero differentials.  This gives a Frobenius-stable
filtration on the total cohomology with ordinary graph piece
of rank two and the $n-1$ elliptic pieces.  Its total rank is
$2+(n-1)=n+1=\binom{n+1}{n}$.
The Hasse determinant is a unit times the product of their
Hasse invariants.  Exactly $f_1$ has valuation one and all the
other factors are units, so the total determinant has
valuation one.

In both constructions the punctured trait lies in $W_{n,d}$.
The determinant Lang torsor is
$z^{p-1}=h^{-1}$, as in \eqref{eq:determinant-kummer}.
A valuation-one $h$ gives tame geometric inertia of order
$p-1$.  This proves the assertion.  The special fibre is
allowed to have smaller Hasse--Witt rank: only its valuation
is used, and no stable local system is extended across it.
\end{proof}

\begin{theorem}\label{thm:pn-odd-full-groups}
\status{Proved under stated hypotheses}
Let $p$ be odd, let $k$ be any perfect field of characteristic
$p$, and let $n\geq2$.  On the ordinary smooth degree-$d$
hypersurface locus in $\mathbf P^n_k$, the geometric and
arithmetic monodromy images on
$H^{n-1}_{\mathrm{et}}(Y_{\bar\eta},\mathbf F_p)$ are both
\[
 \mathrm{GL}_{\binom{d-1}{n}}(\mathbf F_p).
\]
For $d\leq n$ the vector space is zero and the image is trivial.
For $d=n+1$ it has rank one and full determinant.
For every $d\geq n+2$ it is irreducible with full determinant.
\end{theorem}

\begin{proof}
Dimension two is the already proved plane theorem.
Induct on $n\geq3$.  The zero-rank degrees need no argument,
and Proposition~\ref{prop:pn-seed-full-determinant} proves the
rank-one degree $n+1$.  In degree $n+2$,
Proposition~\ref{prop:pn-odd-seed-sl} and the determinant
proposition give the full general linear group.
Now induct on $d\geq n+3$ using
Proposition~\ref{prop:pn-odd-degree-induction}.
The two smaller groups are known by the dimension and degree
inductions, and the old determinant propagates as proved there.
This proves the geometric image over $\overline{\mathbf F}_p$.
The connected finite frame cover remains geometrically
connected under algebraically closed extension; since the
family is defined over $\mathbf F_p$, the result holds over
the algebraic closure of any perfect $k$.  The arithmetic
image contains the full geometric general linear group, so
is equal to it.
\end{proof}

\begin{corollary}\label{cor:pn-odd-all-residues-sharp}
\status{Proved under stated hypotheses}
Let $k$ be finite of odd characteristic $p$ and $n\geq2$.
For every $d\geq n+2$, every nonempty open in the smooth
degree-$d$ hypersurface locus of $\mathbf P^n_k$ contains,
over suitable finite extensions of $k$, hypersurfaces with
each of the $p$ point-count residues.
For $d\leq n$ every smooth hypersurface has residue $1$.
For $d=n+1$ the ordinary open misses residue $1$.
Thus $d=n+2$ is the exact threshold for the assertion involving
every nonempty open.
\end{corollary}

\begin{proof}
The finite-field coherent trace formula gives
\[
 \#Y(k')\equiv 1+(-1)^{n-1}
       \operatorname{Tr}(F_{k'}\mid H^{n-1}(Y,\mathcal O_Y)_{\rm ss})
       \pmod p,
\]
because a projective hypersurface has no other positive
coherent cohomology.  In the zero-rank range the trace is zero.
In rank one on the ordinary open it is nonzero.
In rank $h\geq2$, traces of invertible matrices take every
value in $\mathbf F_p$: the block
$\begin{pmatrix}0&-1\\1&a\end{pmatrix}$ has determinant one
and trace $a$, and adjoining an identity block allows any
specified total trace.  By the theorem the finite geometric
and arithmetic image is the full group.  The finite
Chebotarev theorem already cited in the plane calculation
therefore realizes each such conjugacy class in every
nonempty open after suitable finite extension of $k$.
The intersection of that open with the ordinary locus is
nonempty because the latter is dense.
\end{proof}

% /root/monodromy, 19 November 2026.
% Elementary all-dimension characteristic-two computation; root independently audited.
\subsection{Hypersurfaces in characteristic two: an elementary exact group calculation}
\label{subsec:pn-char2-exact}
\status{Proved; independently audited.}

Fix an algebraically closed field $k$ of characteristic two,
$n\geq2$, and $d\geq n+2$.  Put $m=d-1$ and
$h=\binom{m}{n}$.  Work on the full Hasse--Witt invertibility
open $W_{n,d}$ as in the preceding section.  Its Lang fibre
at a general-position hyperplane arrangement is
$V=\bigwedge^n W$, where
$W=\mathbf F_2^d/\langle e_0+\cdots+e_m\rangle$.
Use $e_1,\ldots,e_m$ as a basis of $W$, so
$e_0=e_1+\cdots+e_m$.
Let $b_A=\bigwedge_{i\in A}e_i$, indexed by the $n$-subsets
$A\subset\{1,\ldots,m\}$, be the exterior basis of $V$.
Signs disappear in this characteristic.

\begin{lemma}\label{lem:char2-independent-quadric-swaps}
\status{Proved under stated hypotheses}
The geometric monodromy contains, for each
$(n-1)$-subset $I\subset\{2,\ldots,m\}$, the transvection
\begin{equation}\label{eq:char2-quadric-transvection}
 R_I=1+w_I\phi_I,\qquad
 w_I=\sum_{\substack{a\in\{2,\ldots,m\}\\a\notin I}}b_{I\cup\{a\}},
 \quad
 \phi_I=b_{\{1\}\cup I}^*.
\end{equation}
It fixes every exterior basis vector except
$b_{\{1\}\cup I}$, which it sends to
$\sum_{a\notin I}b_{I\cup\{a\}}$.
The swaps for different $I$ are independent monodromy elements.
\end{lemma}

\begin{proof}
Replace $H_0H_1$ by a varying quadric and keep the remaining
hyperplanes fixed.  The same alternating resolution as in
Lemma~\ref{lem:quadric-reflections} has constant lower terms
and a pair of points for every line
$L_I=\bigcap_{i\in I}H_i$.  On an open containing the original
arrangement these pairs are distinct, all positive-dimensional
quadric strata are connected and have no higher coherent
cohomology, and the entire top coherent Frobenius is bijective.

We check independence of the point-pair covers in characteristic
two, working on the affine space of quadric equations.  This
cover is the pullback of the projective parameter cover, so
its full image is already a sufficient geometric lower bound.
On a chosen basis for a line $L$, the restricted quadratic
form is $a_LX^2+b_LXY+c_LY^2$.  Its ordered roots give the
Artin--Schreier extension
\[
 z^2+z=a_Lc_L/b_L^2
\]
of the rational function field of the space of quadrics.
The restriction map to any line is surjective, so
$a_L,b_L,c_L$ are independent linear functions.
Moreover $b_L$ is evaluation of the polar alternating form
on the Pluecker vector of $L$.  For distinct lines these
linear functions are nonassociate, by injectivity of the
Pluecker embedding.

At the prime divisor $b_{L_I}=0$, the displayed function for
$I$ has pole order two and leading residue $a_{L_I}c_{L_I}$.
This residue is not a square in the divisor's rational
function field, since $a_{L_I},c_{L_I}$ remain independent
variables there.  The corresponding functions for all
$J\ne I$ have no pole at this divisor.
No Artin--Schreier coboundary $q^2+q$ can have that leading
term: a pole of order two would require a pole of order one
for $q$ and would have square leading residue.  Thus no
nonempty sum of these functions is an Artin--Schreier
coboundary.  Their classes are independent, and the compositum
has Galois group $(\mathbf Z/2)^{\binom{d-2}{n-1}}$.
Passing to the nonempty ordinary and transverse open preserves
connectedness of this cover and hence all the independent swaps.

At the arrangement, swapping the two points on $L_I$ exchanges
the cochains $e_0\wedge e_I$ and $e_1\wedge e_I$, and fixes
all other top cochains not involving both $0$ and $1$.
These latter cochains span the arrangement cohomology; those
involving both are eliminated by the coboundary relation
$u\wedge\bigwedge^{n-1}P=0$.
In the chosen exterior basis the only changed column is
$b_{\{1\}\cup I}$, and its difference is
$(e_0+e_1)\wedge e_I=w_I$.
This proves \eqref{eq:char2-quadric-transvection}.
The vector $w_I$ is nonzero because
$d-n-1\geq1$, and its $\phi_I$ value is zero.
Thus the operator is a transvection.
\end{proof}

\begin{lemma}\label{lem:char2-coordinate-generation}
\status{Proved under stated hypotheses}
Let $G_0$ be generated by the arrangement permutation group
and the swaps of Lemma~\ref{lem:char2-independent-quadric-swaps}.
Writing $\ell=\sum_A b_A^*$, one has
\[
 G_0=
 \begin{cases}
 \mathrm{GL}_{h}(\mathbf F_2),&d-n-1\ \text{odd},\\
 \operatorname{Stab}_{\mathrm{GL}_h(\mathbf F_2)}(\ell),
       &d-n-1\ \text{even}.
 \end{cases}
\]
In the second case the stabilizer is
$\mathbf F_2^{h-1}\rtimes\mathrm{GL}_{h-1}(\mathbf F_2)$.
\end{lemma}

\begin{proof}
Permuting the labels $1,\ldots,m$ and fixing label $0$ supplies
the following operators for every $(n-1)$-subset $I$ and
every $j\notin I$:
\[
 T^I_j(b_{I\cup\{j\}})=S_I,\qquad
 S_I=\sum_{a\notin I}b_{I\cup\{a\}},
\]
with every other exterior basis vector fixed.  This operator
exchanges the two vectors $b_{I\cup\{j\}}$ and $S_I$ and
fixes each of the other vectors in this sum.
For $j,t\notin I$, $j\ne t$, the conjugate
$T^I_jT^I_tT^I_j$ is therefore the pure coordinate
transposition exchanging $b_{I\cup\{j\}}$ and
$b_{I\cup\{t\}}$ and fixing all other exterior basis vectors.
The graph of $n$-subsets joined when they share $n-1$ elements
is connected.  Its edge transpositions generate the full
coordinate symmetric group $S_h$, so $G_0$ contains that group.

Put $l=d-n-1$.  Each $T^I_j$ is a transvection whose image
vector has exactly $l$ nonzero exterior coordinates, all
different from its input coordinate.  Conjugating by $S_h$
gives every transvection
$1+v b_a^*$ with $v$ supported on any $l$ coordinates other
than $a$.  If $l=1$, these are all elementary transvections
and generate $\mathrm{GL}_h(\mathbf F_2)$.
If $l\geq2$, choose two such supports with $l-1$ common
coordinates and differing in any desired pair $b,c$, both
different from $a$.  This is possible because
$h=\binom{d-1}{n}\geq d-1\geq d-n+1=l+2$
for $n\geq2,d\geq n+2$.
The product, with the same input functional, is
\[
 1+(b_b+b_c)b_a^*.
 \tag{4}
\]
If $l$ is odd, adding pairs from an $l$-element support reduces
it to a singleton.  Thus once again all elementary
transvections occur and $G_0=\mathrm{GL}_h(\mathbf F_2)$.

Suppose $l$ is even.  Every generator fixes $\ell$:
the coordinate transvections change a column by an even-weight
vector, and an arrangement permutation fixes $\ell$ as well.
For the latter assertion, permutations of $1,\ldots,m$ act
by coordinate permutations, while $(01)$ is the product of
the commuting swaps in
Lemma~\ref{lem:char2-independent-quadric-swaps}.
These permutations generate the full arrangement $S_d$.
Consequently $G_0\subseteq\operatorname{Stab}(\ell)$.

Relabel the exterior basis as $b_1,\ldots,b_h$ and put
$f_i=b_i+b_h$ for $1\leq i<h$.
The $f_i$ form a basis of $\ker\ell$, with complement
$\langle b_h\rangle$.  Operators (4), with image $f_i$ and
input functional $b_j^*$ for $i\ne j<h$, give the elementary
matrices on $\ker\ell$ and fix $b_h$.  They generate its
full $\mathrm{GL}_{h-1}(\mathbf F_2)$.
The coordinate transposition $(i\,h)$ is
$1+f_i(b_i^*+b_h^*)$.  Multiplying it by the commuting
operators $1+f_i b_j^*$ for all $j\ne i,h$ gives
$1+f_i\ell$.  These generate all the translations of the
complement by $\ker\ell$.  Hence the full stabilizer is
contained in $G_0$, proving equality.
\end{proof}

\begin{lemma}\label{lem:pn-char2-canonical-covector}
\status{Proved under stated hypotheses}
If $d\equiv n+1\pmod2$ and $d\geq n+1$, the universal ordinary
Hasse--Witt local system has a nowhere-zero invariant
$\mathbf F_2$-valued covector, defined over $\mathbf F_2$.
\end{lemma}

\begin{proof}
Write a degree-$d$ equation as $F=\sum_{|\alpha|=d}c_\alpha x^\alpha$.
Use the standard top-cohomology basis indexed by positive
$(n+1)$-tuples $u$ with $|u|=d$.  The Hasse--Witt matrix is
$H_{uv}=c_{2v-u}$, with an out-of-range coefficient understood
as zero.  Put
\[
 w_u=\begin{cases}c_u,&\text{every coordinate of }u\text{ is odd},\\
                   0,&\text{otherwise}.\end{cases}
\]
Then
\[
 H^{\mathsf T}w=w^{[2]}.
 \tag{5}
\]
Indeed, in the $v$ coordinate the sum is
$\sum_{u\ {\rm odd}}c_{2v-u}c_u$.  Whenever the coefficient
is nonzero, $2v-u$ is also a positive all-odd tuple.
The involution $u\leftrightarrow2v-u$ cancels all paired
terms.  Its only possible fixed term is $u=v$, giving
$c_v^2$ when $v$ is all odd and zero otherwise.

There exists a positive all-odd tuple of sum $d$ under the
assumed parity and bound.  If every $w_u$ vanished, the row
of $H$ indexed by any such tuple would vanish: every
possibly nonzero coefficient in that row is itself indexed
by an all-odd tuple.  Thus $w\ne0$ wherever $H$ is invertible.
The map $x\mapsto w^{\mathsf T}x$ is Frobenius equivariant
by (5), and restricts on Lang fixed vectors to an
$\mathbf F_2$-valued covector.  It is nonzero because those
fixed vectors span the coherent fibre.
This construction is global on projective parameter space:
the coherent bundle is $\mathcal O(-1)$ tensored with the
fixed top-cohomology vector space, and the entries $w_u$
are degree-one functions of the equation.  They define an
honest morphism of this bundle to $\mathcal O$, not merely
a covector up to scalar.  Everything is defined over
$\mathbf F_2$.
\end{proof}

\begin{theorem}\label{thm:pn-char2-all-groups}
\status{Proved under stated hypotheses}
Let $k$ be any perfect field of characteristic two and
$n\geq2$.  Put $h=\binom{d-1}{n}$.  The geometric and
arithmetic monodromy groups on the ordinary smooth
degree-$d$ hypersurface locus of $\mathbf P^n_k$ are:
\[
 \begin{array}{c|c}
 \text{degree range}&\text{both monodromy images}\\ \hline
 d\leq n&1\text{ on the zero vector space}\\
 d=n+1&\mathrm{GL}_1(\mathbf F_2)=1\\
 d\geq n+2,\ d\equiv n\pmod2&\mathrm{GL}_h(\mathbf F_2)\\
 d\geq n+3,\ d\equiv n+1\pmod2&
    \operatorname{Stab}(\lambda)
       \simeq\mathbf F_2^{h-1}\rtimes
                         \mathrm{GL}_{h-1}(\mathbf F_2),
 \end{array}
\]
where $\lambda$ is the nonzero covector in
Lemma~\ref{lem:pn-char2-canonical-covector}.
In the last case $\ker\lambda$ is a proper nonzero invariant
hyperplane, so the representation is reducible.
\end{theorem}

\begin{proof}
The two initial ranks are zero and one as stated.
In all higher degrees
Lemma~\ref{lem:char2-coordinate-generation} gives the required
geometric lower bound.  In the parity with a canonical
covector, Lemma~\ref{lem:pn-char2-canonical-covector} gives
the matching upper bound, even for the arithmetic image.
For precision in comparing covectors: the subgroup $G_0$
has exactly the one nonzero invariant covector $\ell$.
The canonical $\lambda$ is nonzero and $G_0$-invariant,
so its value on the arrangement fibre is $\ell$.
Thus the two stabilizers are the same.
In the other parity the lower bound is the entire general
linear group.  The result over arbitrary perfect $k$
follows by geometric connectedness of the corresponding
finite cover under algebraically closed extension, exactly
as in the odd-characteristic theorem.
\end{proof}

\begin{corollary}\label{cor:pn-all-p-sharp-residue-threshold}
\status{Proved under stated hypotheses}
For every finite field $k$, every $n\geq2$, and every
$d\geq n+2$, each nonempty open in the smooth degree-$d$
hypersurface parameter space of $\mathbf P^n_k$ contains
hypersurfaces, over suitable finite extensions of $k$,
with every point-count residue modulo the characteristic.
The threshold $n+2$ is exact: for $d\leq n$ the only
residue is $1$, and for $d=n+1$ the ordinary open misses $1$.
\end{corollary}

\begin{proof}
Odd characteristic was settled in
Corollary~\ref{cor:pn-odd-all-residues-sharp}.  In characteristic
two, the full general linear group in rank at least two
has both trace values.  In the covector-stabilizer case,
the trace is $1+\operatorname{Tr}(B)$ with
$B\in\mathrm{GL}_{h-1}(\mathbf F_2)$; here $h-1\geq2$,
so again both values occur.  The same finite Chebotarev
and coherent trace argument gives both point-count
residues in every open.  The zero- and one-rank failures
are as in the preceding corollary.
\end{proof}

\par\smallskip\noindent\textbf{Finite checks.}
The retained program
\path{code/projective_quadric_generators.py} constructs
these exact exterior matrices over $\mathbf F_2$.
It checks that every swap has rank-one difference and
order two, that their product is the hyperplane
transposition, and that the common fixed covector space
has dimension one exactly in the indicated parity.
For the tested ranks at most fifteen it additionally
enumerates all nonzero vector orbits, obtaining either
$2^h-1$ or $2^{h-1}-1$ and $2^{h-1}$ as predicted.
These finite checks support the formulas; the uniform
group-generation proof above is independent of sampling.

% /root/monodromy with independent suggestions and audit by /root.
% General-X A.5 proof; independently audited by /root on 19 November 2026.
\subsection{All point-count residues for arbitrary smooth ambient varieties}
\label{subsec:general-X-residues}
\status{Proved under the stated hypotheses; independently audited.}

The full generic monodromy group is unnecessary for this
consequence.  We isolate one alternating permutation factor
whose action on every other contribution has trivial trace
variation.  The argument permits a drop of the stable rank
at the reducible boundary, and includes its missing generic
unit roots in an explicit error representation.

\subsubsection{Independent node covers}

\begin{lemma}\label{lem:all-X-independent-nodes}
\status{Proved under stated hypotheses}
Let $X\subset\mathbf P^r_k$ be smooth projective of pure
dimension $n\geq2$ over an algebraically closed field.
Its connected components are denoted $X_1,\ldots,X_c$ and
have degrees $D_1,\ldots,D_c$.  Fix $e\geq3$ and $m\geq n+1$.
Take $m-1$ ordered variable ambient degree-$e$ hypersurfaces
and one of degree $e+\rho$, where $0\leq\rho<e$.
There is a nonempty open parameter space $S$ on which their
restrictions $D_i$ form a simple normal crossing divisor,
all positive-dimensional intersections on each $X_j$ are
smooth and connected, and all zero-dimensional intersections
are reduced.  On this space the joint monodromy of the
node sets
\[
 X_j\cap\bigcap_{i\in I}D_i,\qquad
 1\leq j\leq c,\quad
 I\subset\{1,\ldots,m-1\},\quad |I|=n,
\]
is the full product
\[
 \prod_{j=1}^c\prod_{|I|=n} S_{N_j},
 \qquad N_j=D_j e^n.
 \tag{6}
\]
It remains this product after restriction to any nonempty
open of $S$.
\end{lemma}

\begin{proof}
The required simultaneous smoothness and proper intersection
conditions hold on a nonempty open by Bertini, applied to
the finitely many subsets of the ordered divisors.
Positive-dimensional ample complete intersections in each
$X_j$ are geometrically connected.  This also follows by
iterating the degree-zero stable Lefschetz assertion R1.

We check the precise conditions of Entin's
Proposition~3.1, arXiv version~3, pp.~7--8
\cite{Entin2021}, for the disjoint union of these node covers
over a larger open on which they are finite.
The incidence cover for one pair $(j,I)$ is irreducible and
smooth: over $X_j$ it is a product of projective bundles
given by the linear conditions that the $n$ specified
equations vanish at the point, together with the unrestricted
other parameters.  Its fibre square off the diagonal is
irreducible for the same reason over
$X_j\times X_j\setminus\Delta$; evaluation at two distinct
points is surjective for ambient degree-$e$ forms.
Thus condition (i) of that proposition holds, and the
single-factor point action is doubly transitive.

For condition (ii), fix $(j,I)$.  Choose $n-1$ of these
divisors generally so that their intersection on $X_j$ is
a smooth connected curve $C$.  Choose a point $P\in C$.
There is a degree-$e$ equation whose restriction to $C$
has a zero of order exactly two at $P$ and only simple
zeros elsewhere, while all its intersections on the other
components are transverse.  Here are the required
jet-separation details.  In local coordinates from the
given projective embedding, the functions $1,t,t^2$ occur
as restrictions of degree-$e$ forms after multiplying by a
nonvanishing linear form to make the degrees equal.
Thus the second-order jet at $P$ can be prescribed.
The linear subsystem vanishing twice at $P$ separates
first-order jets at every other point $Q$: multiply a
degree-$(e-2)$ form by the square of a hyperplane equation
which vanishes at $P$ but not at $Q$.  Since $e-2\geq1$,
these residual forms separate tangent directions.
Bertini for this subsystem therefore gives just the one
double zero and no other nontransverse intersection.
All equations not indexed by $I$ may now be chosen generally
and not to vanish at $P$.  Every other node set stays
transverse.  The same imposed-jet Bertini argument makes
all proper partial intersections smooth.  At $P$, choose
the differential of the last equation as a linear combination
of the first $n-1$ independent differentials with every
coefficient nonzero, and prescribe its quadratic coefficient
along the tangent direction to $C$ to be nonzero.  Ambient
forms of degree at least two realize these jets.  Every
proper subtuple is then transverse at $P$, and only the
full $n$-tuple loses transversality.

The base parameter space is smooth.  The point-incidence
source is smooth at its point above $P$, hence formally
irreducible there.  Its fibre has exactly $N_j-1$ geometric
points, with all but the one double point etale.
All other node covers are etale over this parameter.
These are exactly Entin's condition (ii), valid also in
characteristic two.  Concretely, the colliding length-two
algebra has the local quadratic deformation
$z^2+az+b$; the independent constant and linear jet
parameters give the transposition even in characteristic
two.  The proposition consequently gives the full product
(6).  Removing nonempty opens preserves connectedness of
each component of its finite Galois cover, proving the
last assertion.
\end{proof}

\subsubsection{The finite part of the Lang kernel at a boundary}

\begin{lemma}\label{lem:boundary-Lang-finite-part}
\status{Proved under stated hypotheses}
Let $B$ be a smooth integral scheme over an algebraically
closed field of characteristic $p$, let $\mathcal E$ be a
vector bundle on $B$ with a Frobenius-semilinear endomorphism,
and let $U\subset B$ be a dense open on which its stable
rank is the generic rank $r$.
Let $S$ be a smooth connected affine scheme with a morphism
$f:S\to B$, such that the pulled-back stable rank is a
constant $r_0$.  After shrinking $S$ if necessary, the
special Lang local system $W$ of rank $r_0$ and any specified
finite collection of finite local systems on $S$ occur as
representations of one group $H$ which also acts on a
generic fibre $V$ of rank $r$.  There is an $H$-equivariant
inclusion $W\hookrightarrow V$, the image on the specified
finite local systems is their full joint image on $S$,
and the image on $V$ is a subgroup of the geometric
monodromy image of $\mathcal E|_U$.
\end{lemma}

\begin{proof}
Shrink to affine opens so that $f$ takes values in a smooth
affine open $B_0\subset B$.  Use the graph of $f$ in
$S\times B_0$ and henselize this smooth affine scheme along
the graph.  Write $T=\operatorname{Spec}(A)$ for the result,
with closed subscheme $\operatorname{Spec}(A/I)=S$.
The ring is Noetherian and normal, and is connected, hence integral:
henselization of a Noetherian ring is Noetherian and is a
filtered etale limit of the smooth ambient ring.  Every
connected component meets the connected closed subscheme
because $I$ is in the Jacobson radical.  The nonzero flat
map from the ambient domain to its henselization is injective,
so $T$ dominates $S\times B_0$.  The inverse image $T_U$
of $U$ is consequently a nonempty dense open.

The kernel scheme of $F-1$ on the additive vector bundle
is affine, separated, of finite presentation, and etale:
in a local trivialization its equations are
$x-Ax^{[p]}=0$, whose Jacobian is the identity.
Its fibres are finite, so it is quasi-finite.
Its restriction to $S$ is finite etale after shrinking
$S$, because the stable rank is constant.
The finite-part theorem for henselian pairs,
Stacks Lemma~41.20.5, Tag~0BTN \cite{Stacks0BTN},
splits this kernel over $T$ as a finite etale part and an
open-and-closed part disjoint from $S$.
The finite part is a subgroup: sums and inverses of its
sections stay in it. Indeed the inverse image of the away part
under multiplication or inversion is open and closed in a finite
etale scheme over the henselian pair and has no closed-fibre points;
the Jacobson-radical property makes it empty. Its restriction to
$T_U$ is a subgroup of the generic Lang local system $V$,
and its restriction to $S$ is exactly $W$.

By Stacks Lemma~58.8.2, Tag~09ZS \cite{Stacks09ZS},
finite etale covers, and their morphisms and group laws,
extend uniquely from $S$ to $T$.
Thus all the specified finite local systems extend.
The map $\pi_1(T_U)\to\pi_1(T)\simeq\pi_1(S)$ is
surjective: a connected finite etale cover of the normal
integral $T$ remains connected after restriction to a
nonempty open.  Take the joint finite image of
$\pi_1(T_U)$ on $V$ and these extended systems to be $H$.
All schemes were taken over an algebraically closed field,
so its image on $V$ lies in the desired geometric image.
No full stable-rank local system has been extended across
the rank drop; only its finite boundary part was extended.
\end{proof}

\begin{lemma}\label{lem:alternating-factor-isolation}
\status{Proved under stated hypotheses}
Let a finite group $H$ act on $\mathbf F_p$-vector spaces
$P,V,W,R,Q_0,Q_1$, with exact sequences
\[
 0\to R\to P\to M\to0,\qquad
 0\to M\to W\to Q_0\to0,\qquad
 0\to W\to V\to Q_1\to0.
 \tag{7}
\]
Suppose the image on $P$ is the product
$A=\prod_{\alpha=1}^t A_{N_\alpha}$ in its direct sum
of natural permutation representations, with every
$N_\alpha\geq\max(5,p+2)$.
Set $c=\dim R+\dim Q_0+\dim Q_1$ and
\[
 K_p=
 \begin{cases}
 2,&p=2,\\
 v_2(p^2-1)+1,&p\text{ odd}.
 \end{cases}
\]
If $t>K_pc$, then traces of elements of the image of $H$
on $V$ take every value in $\mathbf F_p$.
\end{lemma}

\begin{proof}
Let $K$ be the kernel of the action on
$R\oplus Q_0\oplus Q_1$.  Its image in the product $A$
is a normal subgroup.  Since each alternating factor is
nonabelian simple, it is the product of a subset of the
factors.  The product of the omitted factors is a quotient
of $H/K$, a subgroup of $\mathrm{GL}_c(\mathbf F_p)$.
Choose $\ell=3$ for $p=2$ and $\ell=2$ otherwise.
Each omitted factor has order divisible by $\ell$.
The number of omitted factors is therefore at most
$v_\ell(|\mathrm{GL}_c(\mathbf F_p)|)\leq K_pc$.

For completeness, the last estimate is elementary.
For $p=2$,
\[
 v_3(|\mathrm{GL}_c(\mathbf F_2)|)
 =\lfloor c/2\rfloor+
             v_3(\lfloor c/2\rfloor!)\leq2c.
\]
For odd $p$, the lifting-the-exponent formulas give
\[
 v_2(|\mathrm{GL}_c(\mathbf F_p)|)
 =c\,v_2(p-1)+\lfloor c/2\rfloor v_2(p+1)
                 +v_2(\lfloor c/2\rfloor!)
 \leq K_pc.
\]
The estimates also hold when $c=0$.
At least one factor $A_{N_\alpha}$ thus survives in the
image of $K$.  For any element $a$ of that single factor,
choose a lift $g\in K$ acting as $a$ on its permutation
summand of $P$ and identically on every other summand.
Trace additivity in (7) gives
\[
 \operatorname{Tr}(g\mid V)
  =\dim V-N_\alpha+\#\operatorname{Fix}(a)
       \quad\text{in }\mathbf F_p.
 \tag{8}
\]
Every integer $b=0,\ldots,p-1$ is the fixed-point count of
an even permutation on $N_\alpha$ letters.  Fix $b$ letters
and derange the remaining $q=N_\alpha-b\geq3$.
For odd $q$ use one $q$-cycle, and for even $q\geq4$ use
disjoint cycles of lengths $2$ and $q-2$.
These permutations are even and have exactly the required
fixed points.  Equation (8) proves the claim.
\end{proof}

\subsubsection{Explicit bounds for the error representation}

Return to a smooth projective connected
$X\subset\mathbf P^r_k$ over a finite field of characteristic
$p$, and put $L=\mathcal O_X(1)$.
The following numerical data are all computable from this
fixed embedded variety.  Cohomology dimensions are unchanged
after passage to $\bar k$ and include all its geometric
components.
Let
\[
 D=\deg_L(X),\quad A=h^{n-1}(X,\mathcal O_X),\quad
 P_\omega(t)=\chi(X,\omega_X\otimes L^t).
\]
Choose a regularity bound $r_\omega$ for the embedded
dualizing sheaf and put $d_{\rm coh}=\max(1,r_\omega-1)$.
For every $d\geq d_{\rm coh}$ the coherent relative bundle
has rank
\begin{equation}\label{eq:all-X-relative-coherent-rank}
 b_d=P_\omega(d)-h^0(X,\omega_X).
\end{equation}
Indeed Serre duality and regularity give
$H^{n-1}(X,L^{-d})=0$, and the divisor sequence identifies
the coherent cokernel with the kernel of the surjection
$H^n(X,L^{-d})\to H^n(X,\mathcal O_X)$.
This last map is onto because $H^n(Y,\mathcal O_Y)=0$.
Serre duality and the higher-cohomology vanishing for
$\omega_X\otimes L^d$ give \eqref{eq:all-X-relative-coherent-rank}.
The same construction works in families of arbitrary
effective Cartier divisors, including the boundary unions:
the dual map is multiplication by the nonzerodivisor
defining the divisor, so the displayed kernel is a vector
bundle and carries its natural semilinear Frobenius.

Choose an integer
\[
 e\geq3,\qquad e^n\geq\max(5,p+2),
 \qquad N=De^n .
 \tag{9}
\]
This guarantees that the node count on every geometric
component is at least $\max(5,p+2)$.
For $1\leq a\leq n-1$ and $t\geq0$ define the finite
cohomological upper bounds
\begin{equation}\label{eq:all-X-B-at}
 B_{a,t}=
 \sum_{j=0}^a\binom aj
 \max_{je\leq v\leq j(2e-1)}
       h^{t+j}(X,L^{-v}),
\end{equation}
where $v$ is integral and a cohomology group in degree
greater than $n$ is zero.
The Koszul resolution shows that $B_{a,t}$ bounds
$h^t(Z,\mathcal O_Z)$ for every proper complete intersection
$Z$ of $a$ divisors whose degrees lie in $[e,2e-1]$.
Its hypercohomology spectral sequence has terms
$H^{t+j}(X,L^{-\sum d_i})$ for subsets of $j$ equations,
which gives precisely the bound
\eqref{eq:all-X-B-at}.  The dimension of
$H^t_{\mathrm{et}}(Z_{\bar k},\mathbf F_p)$ is at most this
coherent dimension, by Artin--Schreier/Lang.

Set
\[
 R_0=A+\sum_{a=1}^{n-1}B_{a,n-a-1},
 \qquad
 Q_0=\sum_{a=1}^{n-1}B_{a,n-a}
           +D e^{n-1}(2e-1).
 \tag{10}
\]
For each $\rho=0,\ldots,e-1$, form the polynomial
\[
 q_\rho(z)=P_\omega(ez+\rho)-h^0(X,\omega_X)
                           -N\binom{z-1}{n}.
 \tag{11}
\]
Its degree is at most $n-1$, since the leading terms cancel.
Let $H_0$ be the maximum, over these $e$ polynomials, of
the sum of the absolute values of their coefficients.
Put
\begin{equation}\label{eq:all-X-C-bound}
\begin{aligned}
 C&=\left\lceil2R_0+Q_0+H_0\right\rceil,
       \qquad d_{\rm res}=eM_0,\\
 M_0&=\max\left\{2(n+1),\left\lceil d_{\rm coh}/e\right\rceil,
       \left\lfloor2\cdot3^{n-1}n!K_pC\right\rfloor+1\right\}.
\end{aligned}
\end{equation}
The constants can be very large; no optimality is asserted.

\begin{lemma}\label{lem:all-X-boundary-errors}
\status{Proved under stated hypotheses}
For $d=em+\rho\geq d_{\rm coh}$, with $m\geq n+1$ and
$0\leq\rho<e$, the construction of
Lemma~\ref{lem:all-X-independent-nodes} and
Lemma~\ref{lem:boundary-Lang-finite-part} gives exact
sequences of the form (7).  Here $P$ is the direct sum of
the node permutation spaces excluding the one exceptional
degree-$(e+\rho)$ component, so
$\dim P=N\binom{m-1}{n}$; $V$ is the generic relative
Lang fibre; and
\[
 \dim R+\dim Q_0^{\rm err}+\dim Q_1^{\rm err}
       \leq C(m+1)^{n-1}.
 \tag{12}
\]
The superscript distinguishes these error spaces from the
numerical constant $Q_0$ in (10).
\end{lemma}

\begin{proof}
Work first over $\bar k$ and shrink the arrangement
parameter space $S$ so that the relative stable rank is
constant and its Lang sheaf $W$ is finite locally constant.
The normal crossing cover of $Y=\bigcup_{i=1}^mD_i$
has an exact alternating resolution of its constant
$\mathbf F_p$-sheaf by the constant sheaves of its
intersections: exactness is checked on a geometric
stalk, where it is the augmented complex of a simplex.
Its cohomology spectral sequence is
\[
 E_1^{s,t}=
 \bigoplus_{|I|=s+1}
 H^t_{\mathrm{et}}(D_I,\mathbf F_p)
 \Longrightarrow H^{s+t}_{\mathrm{et}}(Y,\mathbf F_p).
 \tag{13}
\]
The last cochain term $E_1^{n-1,0}$ maps naturally to the
abutment in degree $n-1$.  Its summands from intersections
not involving the exceptional component form $P$.
Compose this map with the quotient by the ambient
$H^{n-1}_{\mathrm{et}}(X,\mathbf F_p)$ to obtain
$P\to W$, and write its kernel as $R$ and image as $M$.
This is a map of finite local systems on $S$ after the
chosen shrinking, and therefore its kernel and cokernel
are also finite local systems.

The kernel of the full last-cochain map to the abutment
is bounded by all incoming terms in the spectral sequence:
for $r=1,\ldots,n-1$ they originate at
$(n-1-r,r-1)$, of total degree $n-2$.
Taking the ambient quotient adds at most $A$ to the
kernel dimension.  Thus
\[
 \dim R\leq A+
       \sum_{a=1}^{n-1}\binom ma B_{a,n-a-1}
       \leq R_0(m+1)^{n-1}.
 \tag{14}
\]
There is no assumption that these higher differentials
vanish.

The quotient of the abutment by the full last-cochain
image is bounded by its other degree-$(n-1)$ graded terms,
namely the terms with $t>0$ in (13).  Removing from the
last-cochain space the intersections involving the
exceptional component adds at most
$D e^{n-1}(e+\rho)\binom{m-1}{n-1}$.
Passing to the relative ambient quotient cannot increase
this cokernel.  Consequently
\[
 \dim(W/M)\leq
 \sum_{a=1}^{n-1}\binom ma B_{a,n-a}
 +D e^{n-1}(2e-1)\binom{m-1}{n-1}
 \leq Q_0(m+1)^{n-1}.
 \tag{15}
\]

The finite-boundary-part lemma supplies $W\subset V$,
while $\dim V\leq b_d$.  Since
$\dim W\geq\dim P-\dim R$, we obtain
\[
 \dim(V/W)
 \leq b_d-N\binom{m-1}{n}+\dim R
 \leq(H_0+R_0)(m+1)^{n-1}.
 \tag{16}
\]
The last inequality is the coefficient bound from (11).
Equations (14)--(16) prove (12).
The exact sequences (7) now hold with
$Q_0^{\rm err}=W/M$ and $Q_1^{\rm err}=V/W$.
All node covers and these special finite systems are
included in the joint group of the boundary-part lemma;
their full independent node image is retained.
\end{proof}

\begin{theorem}\label{thm:all-X-residue-density}
\status{Proved under stated hypotheses}
Let $k$ be a finite field of characteristic $p>0$, let
$X\subset\mathbf P^r_k$ be smooth projective connected of
dimension $n\geq2$, and define the explicit integer
$d_{\rm res}$ by \eqref{eq:all-X-C-bound}.
For every $d\geq d_{\rm res}$ and every nonempty open
$\mathcal U$ in the degree-$d$ ambient hypersurface
parameter space, each $a\in\mathbf F_p$ occurs as
\[
 \#(X\cap H)(k')\equiv a\pmod p
\]
for some finite extension $k'/k$ and some
$H\in\mathcal U(k')$ with $X\cap H$ smooth.
Moreover, traces of the geometric generic relative
Hasse--Witt monodromy group take every value in
$\mathbf F_p$ in this degree range.
Geometric connectedness of $X$ is not required.
\end{theorem}

\begin{proof}
Write $d=em+\rho$ with $0\leq\rho<e$, so $m\geq M_0$.
The number of independent alternating factors is
$c_X\binom{m-1}{n}$, where $c_X\geq1$ is the number of
geometric connected components.  In particular it is
at least $\binom{m-1}{n}$.
Since $m\geq2(n+1)$,
\[
 \binom{m-1}{n}\geq\frac{m^n}{2^n n!},
 \qquad
 (m+1)^{n-1}\leq(3m/2)^{n-1}.
 \]
The strict inequality in the definition of $M_0$ therefore
gives
\[
 \binom{m-1}{n}>K_pC(m+1)^{n-1}.
 \tag{17}
\]
Restrict the joint group of the preceding lemma to the
inverse image of the product of alternating node groups.
It still surjects to that product.  Lemma
\ref{lem:alternating-factor-isolation}, with the bound (12),
shows that its traces on $V$ take every value in
$\mathbf F_p$.  This group maps into geometric monodromy
on the full generic smooth maximal-stable-rank open, so
the latter has the same trace-surjectivity.

It remains to pass from traces to the stated rational
points, with the exact extension-of-constants direction.
Replace $k$ by a finite extension splitting the geometric
components of $X$ and killing the finite Frobenius actions
on all constant spaces
$H^i_{\mathrm{et}}(X_{\bar k},\mathbf F_p)$.
We may further kill the finite arithmetic quotient of
the generic finite monodromy image.  A rational point
over any further finite extension is also a point over a
finite extension of the original $k$, so this replacement
preserves the required existential conclusion.
Every original nonempty open remains nonempty after this
base change and meets the dense smooth maximal-rank open.

Apply the finite Chebotarev theorem used earlier
\cite[Theorem~9.7.13]{KatzSarnak1999} to the finite
representation over this open.  Each geometric conjugacy
class giving a specified trace occurs over a suitable
finite extension, inside the prescribed open.
The Artin--Schreier form of weak Lefschetz R1 identifies
the lower stable cohomology of the section with that of
$X$, and identifies the remaining top quotient with $V$.
The coherent trace congruence
\cite[Corollary~3.2]{Katz1973Congruence} consequently gives
\[
 \#Y(k')\equiv
  \sum_{i=0}^{n-1}(-1)^i
       \dim_{\mathbf F_p}H^i_{\mathrm{et}}(X_{\bar k},\mathbf F_p)
  +(-1)^{n-1}\operatorname{Tr}(F_{k'}\mid V)\pmod p.
 \tag{18}
\]
Its first term is constant by the chosen field extension.
The second takes every value, which proves the claim.
If one uses arithmetic rather than geometric Frobenius
in defining the monodromy representation, replace the
chosen group element by its inverse; all conjugacy classes
are still available.  This keeps the coherent-Frobenius
trace convention in (18) unchanged.
\end{proof}

\par\smallskip\noindent\textbf{Scope of the conclusion.}
This proves the residue consequence A.5 with a specified
bound for every ambient variety in the question, including
the non-geometrically-connected case.  It does not identify
the full monodromy group in A.4 or the exact generic stable
rank in A.3.  In particular it never assumes generic
relative ordinarity, which the preceding smooth
$\alpha_4$ example disproves.



% /root/monodromy, 19 November 2026; independent parent audit passed.
\subsection{Effective generic-rank calculation and exact determinant images}
\label{subsec:effective-rank-determinant}
\status{Proved under stated hypotheses; independently audited.
These do not constitute a structural classification of general ambient monodromy.}

\begin{proposition}\label{prop:effective-generic-rank}
\status{Proved under stated hypotheses}
Let $k$ be a finite field of characteristic $p$ and let
$(X,L)$ be specified by homogeneous equations and its
projective embedding.  Fix a degree $d$ for which the
coherent relative group is the vector bundle of rank
$b=h^0(X,\omega_X\otimes L^d)-h^0(X,\omega_X)$ described
above; in particular every $d\geq d_{\rm coh}$ is allowed.
There is a finite exact algorithm to compute the generic
stable rank $r_d$ and equations for a dense open on
which that rank is attained.

More precisely, take a rational basis of the relative
coherent bundle over the parameter function field
$K=k(t_1,\ldots,t_M)$, and let $A\in M_b(K)$ be its
semilinear Frobenius matrix.  With $A^{[p^j]}$ denoting
entrywise $p^j$-th powers, one has
\begin{equation}\label{eq:generic-rank-finite-product}
 r_d=\operatorname{rank}_K
       \left(A A^{[p]}\cdots A^{[p^{b-1}]}\right).
\end{equation}
For $b=0$ this rank is zero.  The nonvanishing of the
maximal nonzero minors of this displayed product gives
the maximal-stable-rank open inside the chosen basis
chart, after also imposing smoothness.
\end{proposition}

\begin{proof}
There is no search through hypersurfaces of increasing
finite fields in this calculation.  The cohomological
matrix itself can be computed over the rational function
field from the input equations.  A finite graded free
resolution computes the coherent cohomology spaces, and
the standard affine Cech presentation computes the
induced maps by taking $p$-th powers and multiplying by
the universal section to the power $p-1$.
More intrinsically, the relative bundle is the kernel of
\[
 H^n(X,L^{-d})\otimes\mathcal O(-1)
       \longrightarrow H^n(X,\mathcal O_X)\otimes\mathcal O
\]
given by multiplication by the universal section.
The induced Frobenius on the source is the composite
\[
 H^n(X,L^{-d})
 \xrightarrow{\ F\ }H^n(X,L^{-pd})
 \xrightarrow{\ s^{p-1}\ }H^n(X,L^{-d});
\]
it preserves that kernel.  Kernel bases and the matrices
of these maps are finite linear-algebra computations
on the graded-resolution/Cech representatives; polynomial
ideal membership is decidable by Groebner bases.
Thus an explicit rational matrix $A$ is obtained by a
finite procedure.

For the stabilization bound, put
$B_0=1$ and
$B_j=A A^{[p]}\cdots A^{[p^{j-1}]}$ for $j>0$.
The spaces $I_j=\operatorname{im}_K B_j$ are decreasing.
If $\dim I_j=\dim I_{j+1}$, then they are equal and
the Frobenius-linearized map
$F^*I_j\to I_j$ is surjective, hence an isomorphism:
its image is exactly $I_{j+1}$ and both have the same
dimension.  All subsequent images are the same.
Every strict decrease lowers dimension by at least one,
so stabilization occurs by $j=b$.  Extension to the
perfect closure of $K$ preserves all these ranks, and
the stable part there has dimension $\operatorname{rank}B_b$.
This proves \eqref{eq:generic-rank-finite-product}.
The same argument applies in every geometric fibre
where the rational basis and its matrix are regular.
Nonvanishing of its generic maximal minors consequently
cuts out the asserted open.
\end{proof}

\begin{remark}\label{rem:algorithm-versus-generic-formula}
\status{Proved under stated hypotheses}
Proposition~\ref{prop:effective-generic-rank} is an exact
algorithm for each specified input.  We distinguish it
from a closed structural formula for $r_d$ in terms of
fixed invariants of $X$, or a uniform description of its
defect from coherent dimension as $d$ varies.
The latter determination for arbitrary $X$ remains
unsettled in this dossier.  No claim that the algorithm
was executed for every possible input is intended.
\end{remark}

\begin{proposition}\label{prop:exact-general-determinant-image}
\status{Proved under stated hypotheses}
Let $k$ be a perfect field of characteristic $p$.
Let $\mathcal V$ be the rank-$r$ Lang local system of an
invertible Frobenius bundle on a nonempty open
$U\subset\mathbf P^M_k$, as for the stable relative
Hasse--Witt bundle.  Choose a rational trivialization of
its determinant line and write its Frobenius as
$e\mapsto h e$, where
$h\in K^\times$ and $K=k(\mathbf P^M)$.
Put $N=p-1$ and
\begin{equation}\label{eq:det-valuation-gcd}
 g=\gcd\bigl(N,\ v_D(h)\text{ for every prime divisor }
                         D\subset\mathbf P^M_{\bar k}\bigr),
 \qquad m=N/g.
\end{equation}
Then $h=c f^g$ for some $f\in K^\times$ and $c\in k^\times$.
The geometric determinant image is the unique subgroup
of $\mathbf F_p^\times$ of order $m$.

The arithmetic determinant image is the unique subgroup
of order $m\delta$, where
$\delta$ is the order of the class of $c$ in
$k^\times/(k^\times)^g$.  In particular, for
$k=\mathbf F_q$,
\begin{equation}\label{eq:det-arithmetic-explicit}
 \delta=\operatorname{ord}_{\mu_g}
                   \left(c^{(q-1)/g}\right),
 \qquad
 |\det G_{\rm arith}|=\frac{p-1}{g}\,
                   \operatorname{ord}_{\mu_g}
                   \left(c^{(q-1)/g}\right).
\end{equation}
Both determinant images are trivial when $p=2$, or when
the local system has rank zero.
For finite $k$, factorization of the numerator and
denominator of $h$ computes these quantities exactly.
\end{proposition}

\begin{proof}
The determinant line $\mathcal L=\det\mathcal E_{\rm ss}$
has a nowhere-vanishing Frobenius section in
$\mathcal L^{1-p}$ on $U$.  In the chosen rational basis
it is $h$.  A Lang fixed basis $ze$ satisfies
$h z^p=z$, so the determinant frame torsor is
\begin{equation}\label{eq:det-kummer-general-image}
 z^N=h^{-1}.
\end{equation}
All $N$-th roots of unity lie in $\mathbf F_p^\times$.
The determinant image is therefore the Galois group
of the Kummer extension cut out by this equation.

Over $\bar k$, an element of the rational function field
is an $a$-th power, for $a\mid N$, exactly when all its
divisor valuations are divisible by $a$.
One direction is immediate.  Conversely, if
$\operatorname{div}(h)=aD$, then $a[D]=0$ in
$\operatorname{Pic}(\mathbf P^M_{\bar k})=\mathbf Z$.
Thus $D$ is principal, and the remaining constant has
an $a$-th root in $\bar k$.
The Kummer class of $h$ consequently has exact order
$N/g=m$ over the geometric function field, proving the
geometric assertion.

Because $k$ is perfect, its prime divisors become
geometrically reduced after algebraic extension.
The divisor $\operatorname{div}(h)/g$ is therefore an
integral divisor defined over $k$.
Its class in $\operatorname{Pic}(\mathbf P^M_k)=\mathbf Z$
is zero.  It is the divisor of a rational function
$f\in K^\times$, and $h/f^g=c\in k^\times$.
In the Kummer extension of \eqref{eq:det-kummer-general-image}
the element $w=z^m f$ satisfies $w^g=c^{-1}$.
Thus the extension contains the constant field
$k'=k(c^{1/g})$, whose degree is $\delta$ by elementary
Kummer theory.  Over $k'K$, the remaining equation has
degree exactly $m$: its divisor valuations, divided by
$g$, have gcd one with $m$ by the definition of $g$.
This degree remains $m$ over $\bar k K$, so $k'$ is
the entire algebraic constant extension.
The arithmetic extension therefore has degree
$m[k':k]=m\delta$, as claimed.

For a finite field, the map
$k^\times/(k^\times)^g\to\mu_g$ sending
$c$ to $c^{(q-1)/g}$ is an isomorphism of cyclic groups.
This proves \eqref{eq:det-arithmetic-explicit}.
Changing $f$ multiplies $c$ by a $g$-th power and does
not change $\delta$.  Changing the rational determinant
basis multiplies $h$ by an $N$-th power and does not
change any of the images.
For rank zero the determinant is $1$; for $p=2$ the
target $\mathbf F_2^\times$ is trivial.
\end{proof}

\begin{remark}\label{rem:det-frobenius-convention}
\status{Proved under stated hypotheses}
At a rational point over $\mathbf F_q$ where a determinant
basis is regular, the determinant of coherent $q$-Frobenius
is $h^{(q-1)/(p-1)}$.  On the Lang frame torsor,
arithmetic Frobenius acts by its inverse, since
$z^q/z=h^{-(q-1)/(p-1)}$.
Thus the geometric-Frobenius determinant character is
the coherent Frobenius determinant; the arithmetic
Frobenius character is its inverse.  The subgroup orders
in Proposition~\ref{prop:exact-general-determinant-image}
are unaffected by this convention.
\end{remark}


% /root/monodromy, 19 November 2026, 05:39 UTC.
\subsection{An ambient adjunction obstruction in characteristic two}
\label{subsec:general-adjunction-covector}
\status{Proved under the stated hypotheses; independently audited.}

\begin{theorem}\label{thm:general-adjunction-covector}
\status{Proved under stated hypotheses}
Let $X/k$ be smooth, projective and geometrically connected, of
dimension $n\geq2$, over a perfect field of characteristic two.
Let $L=\mathcal O_X(1)$ for the fixed embedding $X\subset\mathbf P^r$.
Write $R_I=\operatorname{reg}(\mathcal I_X)$.
Suppose that, for some
integer $d_*$, there is a line bundle $M_*$ defined over $k$ with
\[
 M_*^{\otimes2}\simeq\omega_X\otimes L^{d_*}.
\]
Put $d=d_*+2t$ and $M=M_*\otimes L^t$.
Write $F=F_{X/k}:X\to X^{(2)}$ and
$\mathcal B=\operatorname{coker}(\mathcal O_{X^{(2)}}\to F_*\mathcal O_X)$.
Choose Castelnuovo--Mumford regularities
\[
 R_M=\operatorname{reg}_L(M_*),\qquad
 R_B=\operatorname{reg}_{L^{(2)}}
       (\mathcal B^\vee\otimes M_*^{(2)}),
 \qquad
 \kappa=\frac{c_1(\omega_X)c_1(L)^{n-1}}{c_1(L)^n}.
\]
Assume
\begin{equation}\label{eq:adjunction-covector-threshold}
 t\geq\max\{0,R_M,R_B-1\},\qquad
 d\geq\max\{d_{\rm coh},R_I-1\},\qquad d>\kappa.
\end{equation}
On the maximal-stable-rank smooth-hypersurface open $U_d$,
the relative Hasse--Witt local system $\mathcal V_d$ has a
canonical nonzero morphism
\begin{equation}\label{eq:adjunction-invariant-covector}
 \lambda_M:\mathcal V_d\longrightarrow\mathbf F_2.
\end{equation}
Here ``canonical'' is relative to the given square root and its
isomorphism with $\omega_X L^d$; the construction is described below.
In particular, both the geometric and arithmetic images fix a
nonzero covector.  If $\operatorname{rank}\mathcal V_d\geq2$,
they are proper subgroups of $\GL(\mathcal V_{d,\bar\eta})$ and
the representation is reducible.

If the square root is defined only after an algebraic extension
of $k$, the conclusion about the geometric image still holds.
There are arbitrarily large degrees satisfying
\eqref{eq:adjunction-covector-threshold}; the rank is at least two
for all sufficiently large such degrees, by the previously proved
explicit stable-rank lower bound (or by the all-residue theorem
when $k$ is finite).
\end{theorem}

\begin{proof}
The morphism $F$ is finite flat of degree $2^n$.  Finite duality,
or the Cartier trace calculation of
\cite[Lemma~1.3.6, Proposition~1.3.7, Theorem~1.3.8, pp.~26--28]
{BrionKumar2005}, gives
\begin{align}
 F^!M^{(2)}
 &\simeq \omega_X\otimes F^*(\omega_{X^{(2)}}^{-1}\otimes M^{(2)})
       \simeq\omega_X^{-1}\otimes M^{\otimes2}
       \simeq L^d,\label{eq:adjunction-finite-duality}\\
 H^0(X,L^d)
 &\simeq\Hom_{X^{(2)}}(F_*\mathcal O_X,M^{(2)}).
 \label{eq:adjunction-section-map}
\end{align}
These are relative-Frobenius identifications and are $k$-linear.
Given $s\in H^0(X,L^d)$, let $\tau_s$ be the corresponding map.
Set
\begin{equation}\label{eq:adjunction-w-definition}
 v_s=\tau_s(1)\in H^0(X^{(2)},M^{(2)}),\qquad
 w_s=F^*v_s\in H^0(X,M^{\otimes2})
                  =H^0(X,\omega_X L^d).
\end{equation}
The map $s\mapsto w_s$ is $k$-linear.  Relative Frobenius is
essential in this assertion: its pullback fixes base-field
coefficients.  On a perfect geometric fibre one can express
$w_s$ as a square using the absolute Cartier operator, but
one must not square the universal parameter coefficients.

We check the Frobenius identity on geometric fibres.  Locally
choose a generator $m$ of $M$ and etale coordinates whose top
form is $\eta$.  The induced generator of $L^d$ is $m^2/\eta$;
write $s=fm^2/\eta$.  Write $T$ for the coefficient form of
the absolute Cartier trace, so that
\[
 C(f\eta)=T(f)\eta,\qquad T(a^2b)=aT(b).
\]
Formula~\eqref{eq:adjunction-w-definition} becomes
$w_s=T(f)^2m^2$.  The Cartier operator dual to
\[
 F_s:H^n(X,L^{-d})\longrightarrow H^n(X,L^{-d}),
       \quad \xi\longmapsto sF_X(\xi),
\]
sends $gm^2$ to $T(fg)m^2$.  Consequently
\begin{equation}\label{eq:adjunction-covector-Cartier}
 C_s(w_s)=T(fT(f)^2)m^2=T(f)^2m^2=w_s.
\end{equation}
Under Serre duality this is exactly
\[
 \langle F_s\xi,w_s\rangle=\langle\xi,w_s\rangle^2.
\]
Thus $w_s$ is a Frobenius-equivariant scalar-valued covector.
In matrix notation this is $A^tw_s=w_s^{[2]}$, consistently
with the explicit projective-space formula.

Because $w_s$ is linear in $s$, it defines a global covector
on $H^n(X,L^{-d})\otimes\mathcal O(-1)$ over
$\mathbf P H^0(X,L^d)$.  It restricts to the relative coherent
bundle
\[
 \mathcal E_d=\ker\bigl(H^n(X,L^{-d})\otimes\mathcal O(-1)
       \xrightarrow{\ s\ }H^n(X,\mathcal O_X)\otimes\mathcal O\bigr).
\]
The Frobenius identity shows that the restricted covector
kills the Frobenius-nilpotent part and induces a map from
the stable Lang local system to the constant sheaf
$\mathbf F_2$.

It remains to prove nonvanishing.  The sheaf $\mathcal B$ is
locally free: etale locally the standard square-free
monomials form a basis of $F_*\mathcal O_X$ containing $1$.
Dualizing its defining exact sequence and tensoring with
$M^{(2)}$ gives
\[
 0\longrightarrow\mathcal B^\vee\otimes M^{(2)}
 \longrightarrow\mathcal Hom(F_*\mathcal O_X,M^{(2)})
 \xrightarrow{\ \tau\mapsto\tau(1)\ } M^{(2)}
 \longrightarrow0.
\]
The regularity bounds in
\eqref{eq:adjunction-covector-threshold} give
\[
 H^1(X^{(2)},\mathcal B^\vee\otimes M^{(2)})=0,
       \qquad H^0(X^{(2)},M^{(2)})\ne0.
\]
Indeed, the first is the degree-one regularity vanishing,
and the second follows from global generation of the
nonzero line bundle $M$.  Thus $s\mapsto v_s$ is surjective
onto a nonzero space, and $w_s\ne0$ for generic $s$.

Take such an $s$ with $Y=(s=0)$ smooth.  If its covector
vanished on $E_s$, duality for the displayed kernel would
give $w_s=s\eta_0$ for some $\eta_0\in H^0(X,\omega_X)$.
Every valuation of $w_s=F^*v_s$ is even.  The divisor of
$s$ is reduced, so the valuation of $\eta_0$ along each
component of $Y$ would be odd and positive.  Hence
$\eta_0\in H^0(X,\omega_X L^{-d})$.  This last space is zero:
a nonzero section would give an effective divisor with
intersection number $(\kappa-d)c_1(L)^n<0$ against
$L^{n-1}$, an impossibility.  Thus the restricted covector
is nonzero generically.  A morphism of finite local
systems on the connected $U_d$ has locally constant rank;
it is therefore nonzero everywhere on the maximal-rank open
of the complete linear system.

For the prescribed ambient degree-$d$ hypersurface parameter
space, the additional bound $d\geq R_I-1$ gives
$H^1(\mathbf P^r,\mathcal I_X(d))=0$.  The restriction
$H^0(\mathbf P^r,\mathcal O(d))\to H^0(X,L^d)$ is therefore
surjective.  Its induced linear projection, away from its
center, pulls back the construction just given.  The smooth
ambient open avoids the center; its maximal stable rank is
the same as in the complete linear system, and the pulled-back
covector is nonzero on it.  This is the $U_d$ in the statement.

Everything was defined over $k$ if the square root was,
so arithmetic as well as geometric monodromy fixes it.
After extending the base field the same construction
gives the asserted geometric conclusion.  The kernel
of a nonzero covector is a nonzero proper invariant
subspace when the rank is at least two.
\end{proof}

\begin{remark}\label{rem:adjunction-square-Pn}
\status{Proved under stated hypotheses}
For $X=\mathbf P^n$ and $L=\mathcal O(1)$, the square-root
condition is $d-n-1\equiv0\pmod2$.  In homogeneous
monomial coordinates the construction recovers the
covector whose entry indexed by $u$ is the defining
coefficient $c_u$ when all coordinates of $u$ are odd,
and zero otherwise.  Thus this obstruction is an
ambient adjunction phenomenon, rather than a property
limited to plane theta characteristics.
The theorem gives a subgroup constraint for general
$X$, not equality with the full covector stabilizer.
\end{remark}


% /root/monodromy, 19 November 2026; independently audited by /root.
\subsection{Many fixed covectors for ordinary abelian ambient varieties}
\label{subsec:ordinary-abelian-covectors}
\status{Proved; independent parent proof and source audit passed.}

\begin{proposition}\label{prop:ordinary-abelian-fixed-quotient}
\status{Proved under stated hypotheses}
Let $A$ be an ordinary abelian variety of dimension $n\geq2$
over a perfect field of characteristic two, with fixed very
ample $L=\mathcal O_A(1)$.  For every sufficiently large even
degree $d=2t$, the geometric monodromy of $\mathcal V_d$
acts trivially on a quotient of dimension $2^n-1$.
More precisely, write $F:A\to A^{(2)}$ and, after a field
extension splitting $\ker F$, write
\[
 F_*\mathcal O_A=\bigoplus_{\chi\in H}\mathcal N_\chi,
 \qquad H=\operatorname{Hom}(\ker F,\mathbf G_m)
                 \simeq(\mathbf Z/2)^n.
\]
Set
\[
 R=\max_{\chi\in H}
         \operatorname{reg}_{L^{(2)}}(\mathcal N_\chi).
\]
It suffices that $t\geq\max\{1,R\}$ and
$2t\geq\max\{d_{\rm coh},\operatorname{reg}(\mathcal I_A)-1\}$.
After the same splitting extension, the arithmetic image
also fixes this quotient pointwise.
\end{proposition}

\begin{proof}
The character decomposition is the decomposition of the
algebra of a torsor under the diagonalizable group
$\ker F\simeq\mu_2^n$.  In particular,
$\mathcal N_0=\mathcal O_{A^{(2)}}$ and multiplication gives
$\mathcal N_\chi\otimes\mathcal N_\psi\simeq
\mathcal N_{\chi+\psi}$.
The line bundles are pairwise nonisomorphic two-torsion
bundles in $\Pic^0(A^{(2)})$.
This decomposition, including all the two-torsion line
bundles, is also stated and proved in
\cite[Lemma~4.5, p.~11]{SannaiTanaka2016}.

Choose the unique nonzero invariant top differential
$\omega$ normalized by $C(\omega)=\omega$, where $C$
is absolute Cartier.  Such a normalization exists because
$A$ is ordinary and the characteristic is two: if
$C(\omega_0)=c\omega_0$ with $c\ne0$, take
$\omega=c^2\omega_0$.
The coefficient trace $F_*\mathcal O_A\to\mathcal O_{A^{(2)}}$
defined by this differential is the projection to the
trivial character.  Indeed its restriction to every
nontrivial character line is zero, since a nontrivial
degree-zero line bundle has no global sections; on
$\mathcal N_0$ the normalization makes it the identity.

Because $L^{2t}=F^*(L^{(2)t})$, the projection formula
gives a direct sum decomposition
\begin{equation}\label{eq:abelian-section-character-decomposition}
 W:=H^0(A,L^{2t})
    =\bigoplus_{\chi\in H}
       H^0(A^{(2)},L^{(2)t}\otimes\mathcal N_\chi).
\end{equation}
Write $s=\sum_\chi s_\chi$ for the corresponding
character components, regarded again as sections on $A$.
Each summand in \eqref{eq:abelian-section-character-decomposition}
is nonzero under the stated regularity bound.  Thus for
generic $s$ all its $2^n$ components are nonzero, and
these components are linearly independent.

We claim that each
\[
 w_\chi(s)=s_\chi\omega\in H^0(A,\omega_A L^{2t})
\]
is a Frobenius-equivariant covector for $F_s=sF_A$.
The twisted Cartier operator dual to $F_s$ sends it to
$C(s s_\chi\omega)$.  In the sum
$s s_\chi=\sum_\psi s_\psi s_\chi$, every term with
$\psi\ne\chi$ has nontrivial character and is killed
by Cartier, whereas
\[
 C(s_\chi^2\omega)=s_\chi C(\omega)=s_\chi\omega.
\]
Consequently $C_s(w_\chi)=w_\chi$.
The character projections are linear in $s$, so these
are global covectors on the tautologically twisted
coherent bundle, just as in
Theorem~\ref{thm:general-adjunction-covector}.

By Serre duality the dual of the relative coherent
fibre is $W/k s$, since $H^0(A,\omega_A)=k\omega$.
The $2^n$ independent vectors $s_\chi$ span a subspace
whose intersection with $ks$ is exactly $ks$.
Their classes in $W/ks$ therefore have dimension
$2^n-1$, with the single relation
\begin{equation}\label{eq:abelian-covector-single-relation}
 \sum_{\chi\in H}w_\chi=0
       \quad\hbox{on the relative coherent fibre}.
\end{equation}
They kill the nilpotent part by Frobenius equivariance,
so remain independent on the stable part.  Lang descent
turns them into $2^n-1$ independent $\mathbf F_2$-valued
covectors; independence can be checked after tensoring
with the algebraic closure.  Their rank is constant on
the connected maximal-rank open.  The ambient restriction
bound in the statement permits pullback from the complete
linear system to the prescribed hypersurface parameter
space.  This proves the asserted fixed quotient.
\end{proof}

\begin{remark}\label{rem:abelian-obstruction-not-exact-group}
\status{Proved under stated hypotheses}
Already for an ordinary abelian surface there is a fixed
quotient of dimension three.  Thus even the proposed
replacement ``full stabilizer of one adjunction covector''
would be too large for general ambient varieties.
The proposition gives a lower bound for the full fixed
covector space; it does not assert that no further
invariants occur, nor determine the whole monodromy group.
\end{remark}

\begin{example}\label{ex:ordinary-abelian-surface-explicit}
\status{Proved under stated hypotheses}
Here is a completely specified ambient surface, with an executed
ordinary-hypersurface calculation.  Over $\mathbf F_2$, let
\[
 E:\quad Y^2Z+XYZ+X^3+Z^3=0,\qquad
 A=E\times E,\qquad L=\mathcal O_E(3O)\boxtimes\mathcal O_E(3O),
\]
where $O=[0:1:0]$, and embed $A$ in $\mathbf P^8$ by the
Segre embedding of these two plane cubics.  The derivatives of
the displayed cubic have no common projective zero on $E$;
its four $\mathbf F_2$-points are
$[0:1:1],[1:0:1],[1:1:1],O$.
Thus $E$ is smooth and ordinary (its Frobenius trace is $-1$).
The degree-one section
\begin{equation}\label{eq:abelian-degree-one-example}
 Z_1Z_2+Z_1Y_2+X_1Y_2+Y_1X_2=0
\end{equation}
cuts out a smooth ordinary genus-ten curve on $A$.

The exact programs \path{ordinary_abelian_surface.py} and
\path{ordinary_abelian_degree_one.py}, with their retained
JSON outputs, verify this as follows.  On $E\setminus\{O\}$
write $x=X/Z,y=Y/Z$ and $\omega=dx/x$.
The coefficient Cartier map satisfies
\[
 T(x^{2a})=x^a,\quad T(x^{2a+1})=0,\quad
 T(x^{2a}y)=x^{a+1},\quad T(x^{2a+1}y)=x^a(y+1).
\]
These formulas follow directly from
$((y+1)/x)^2=y/x+x$ and the decomposition
$k(E)=k(E)^2\oplus xk(E)^2$.
Multiplication uses $y^2=x^3+xy+1$.
On the tensor basis $(1,x,y)\otimes(1,x,y)$ the Cartier
map $g\mapsto T_1T_2(sg)$ has the following nine columns,
encoded as binary integers from the first basis vector:
\[
 (3,21,46,16,170,357,154,493,116).
\]
Its rank over $\mathbf F_2$ is nine.  Quotienting the fixed
line $ks$ gives relative rank eight.  Together with the
ordinary ambient $H^1(A,\mathcal O_A)$ of rank two, this
proves $p$-rank ten.

Smoothness is checked on all four product charts
$Z_i\ne0$ or $Y_i\ne0$.
The invariant tangent vector field has components
$(x,x^2+y)$ on the first chart and
$(1+u+v^2,u^2+v)$ on $u=X/Y,v=Z/Y$.
On each product chart the ideal generated by the two
elliptic equations, the section, and its two tangent
derivatives has Groebner basis $[1]$ over $\mathbf F_2$.
These charts cover $A$ and the vector fields never vanish,
so these certificates prove smoothness over the algebraic
closure.  The direct count is eight rational points,
consistent with the computed Cartier trace zero.

It follows that the generic degree-$d$ section on this
\emph{fixed} embedded surface is ordinary for every $d\geq1$.
Indeed choose $d$ general ordinary members of $|L|$;
they meet transversely, with no triple intersections.
Each has genus ten and each pair meets in $L^2=18$ points.
Their connected nodal union has
\[
 h^1(\mathcal O)=10d+18\binom d2-d+1=9d^2+1.
\]
The normalization sequence expresses this cohomology
as an extension of the ordinary component cohomology
by the Frobenius-bijective graph cohomology.  Its
Frobenius is therefore bijective.  Openness of that
condition and the existence of smooth members give
ordinary generic smooth sections.  Thus
\begin{equation}\label{eq:abelian-example-generic-rank}
 r_d=9d^2-1\qquad\text{for every }d\geq1.
\end{equation}
The ambient restriction map is surjective in every
degree: restriction of plane forms to each cubic is
surjective, and Segre monomials span all forms of
bidegree $(d,d)$.

For every even $d\geq2$, the four character summands
have nonzero spaces of sections (their factors on each
elliptic curve have degree $3d/2$), so the fixed quotient
of Proposition~\ref{prop:ordinary-abelian-fixed-quotient}
has dimension three with no additional degree bound.
The degree-two program also gives an explicit smooth
ordinary genus-37 section: its 36-dimensional Cartier
matrix has rank 36; the relative rank is 35, and its
four character covectors have rank four before taking
the single relation and rank three afterwards.
All defining coefficients, matrix columns, covectors,
chart certificates and the count of seven rational
points are retained in its JSON output.
\end{example}


% /root/monodromy, 19 November 2026; independently audited by /root.
\subsection{An odd-characteristic ambient obstruction in every degree}
\label{subsec:product-surface-obstruction}
\status{Proved; independent parent proof and source audit passed.}

\begin{proposition}\label{prop:product-surface-norm-obstruction}
\status{Proved under stated hypotheses}
Let $E_1,E_2$ be ordinary elliptic curves over a perfect
field of characteristic $p$, and let
\[
 A=E_1\times E_2,\qquad
 L=\mathcal O_{E_1}(aO_1)\boxtimes\mathcal O_{E_2}(bO_2),
       \qquad a,b\geq3.
\]
For every $d\geq1$ with $p\mid ad$ and $p\mid bd$, the
geometric relative Hasse--Witt monodromy of smooth members
of $|L^d|$ fixes pointwise a quotient of dimension two.
The same is true of arithmetic monodromy if both
$E_i[p]^{\et}$ are constant over the base field.
The statement applies equally to ambient degree-$d$
hypersurfaces for the product embedding by $L$.

In particular, over $\mathbf F_3$ take
\begin{equation}\label{eq:odd-product-explicit-elliptic}
 E:\quad y^2=x^3+x^2+1,
 \qquad A=E\times E,\qquad
 L=\mathcal O_E(3O)\boxtimes\mathcal O_E(3O).
\end{equation}
For this fixed smooth embedded surface, both monodromy
images have a fixed quotient of dimension two for
\emph{every} $d\geq1$.  In particular no eventual
full-linear-group assertion holds for this $X$, even
in odd characteristic.
\end{proposition}

\begin{proof}
Work first over the algebraic closure.  For a smooth
ample member $i:Y\hookrightarrow A$, put
$\pi_j=\operatorname{pr}_j|_Y$.
The restriction and Albanese maps are
\[
 j:\Pic^0(A)=E_1\times E_2\longrightarrow J_Y,
 \quad (P,Q)\longmapsto\pi_1^*P\otimes\pi_2^*Q,
\]
and
\[
 u:J_Y\longrightarrow A,
       \quad N\longmapsto
            (\operatorname{Nm}_{\pi_1}N,
             \operatorname{Nm}_{\pi_2}N),
\]
using the usual principal identifications on the elliptic
curves and the Jacobian.

The map $j$ is injective on the geometric points of
$\Pic^0(A)[p]$.  To see this directly, let
$P\boxtimes Q$ be a nontrivial degree-zero line bundle
on $A$.  It has no global sections, whereas
\[
 H^1\bigl(A,(P\boxtimes Q)\otimes L^{-d}\bigr)=0
\]
by the Kunneth formula: both factors have strictly
negative degree and hence no $H^0$.
The restriction exact sequence consequently gives
$H^0(Y,i^*(P\boxtimes Q))=0$.
Thus its restriction cannot be trivial.  Passing to
the etale $p$-torsion quotient of $J_Y$ preserves this
injection: over an algebraically closed field that
quotient induces a bijection on the geometric points
of the finite $p$-torsion group scheme.

We calculate the composite $u\circ j$.
The projection degrees are
$\deg\pi_1=bd$ and $\deg\pi_2=ad$.
The norm of a pulled-back line bundle is its power
by the degree, so the two diagonal entries are
$[bd]$ and $[ad]$.
The off-diagonal entries are zero.  For example,
represent $Q\in\Pic^0(E_2)$ as $[q]-[q_0]$.
Pushing its pullback to $Y$ along $\pi_1$ gives the
difference of the intersection divisors of $Y$ with
$E_1\times\{q\}$ and $E_1\times\{q_0\}$, viewed on
$E_1$.  Each divisor has line bundle
$\mathcal O_{E_1}(adO_1)$, because the section defining
$Y$ has bidegree $(ad,bd)$.  Their difference is
therefore trivial in $\Pic^0(E_1)$.
This calculation holds for generic $q,q_0$ and hence
for the entire homomorphism.  Thus
\begin{equation}\label{eq:product-norm-composition}
 u\circ j=\begin{pmatrix}[bd]&0\\0&[ad]\end{pmatrix}.
\end{equation}
Under the divisibility hypotheses it vanishes on
$\Pic^0(A)[p]^{\et}$.

The Artin--Schreier/Lang comparison identifies
\[
 \mathcal V_Y=\operatorname{coker}
       \bigl(H^1_{\et}(A,\mathbf F_p)
                 \longrightarrow H^1_{\et}(Y,\mathbf F_p)\bigr).
\]
For an abelian variety, the dual of $H^1_{\et}(-,\mathbf F_p)$
is its etale $p$-torsion; for a curve it is the etale
$p$-torsion of its Jacobian.  Functoriality identifies
the dual of restriction with $u[p]^{\et}$.
It follows that
\[
 \mathcal V_Y^\vee
   =\ker\bigl(J_Y[p]^{\et}\longrightarrow A[p]^{\et}\bigr)
\]
contains the injected copy of $\Pic^0(A)[p]^{\et}$
provided by $j$ and \eqref{eq:product-norm-composition}.
Both elliptic factors are ordinary, so this copy has
dimension two over $\mathbf F_p$.
These maps exist throughout the smooth family; on its
maximal-rank open they give an injection of the constant
ambient geometric local system into $\mathcal V_d^\vee$.
Thus geometric monodromy fixes a two-dimensional quotient
of $\mathcal V_d$.  If the ambient etale torsion is
constant over $k$, the same holds arithmetically.
Pulling the family back to ambient hypersurface
parameters preserves the constraint.  Smooth ambient
members form a nonempty open since $L^d$ is very ample.

For \eqref{eq:odd-product-explicit-elliptic}, the affine
partial derivative of $x^3+x^2+1$ is $2x$; at its only
zero the polynomial equals one, and the point at
infinity is smooth.  The numbers of affine points
above $x=0,1,2$ are respectively $2,1,2$.
Hence $\#E(\mathbf F_3)=6$ and the Frobenius trace
is $-2$, nonzero modulo three.  The curve is ordinary,
and its etale three-torsion is rational because the
group of six rational points contains all three
geometric points of that torsion.
Now $a=b=3$, so the divisibility conditions hold for
every $d$.  The two-dimensional fixed quotient is
incompatible with the full linear group in any
dimension at least two; its existence itself forces
the relative rank to be at least two.
\end{proof}


% /root/monodromy, 19 November 2026; independently audited by /root.
\subsection{The exact fixed quotient on products of ordinary elliptic curves}
\label{subsec:product-exact-fixed-quotient}
\status{Proved; independent parent proof and source audit passed.
This determines all invariant covectors, not the full group on their kernel.}

\begin{theorem}\label{thm:product-exact-fixed-covectors}
\status{Proved under stated hypotheses}
Let $E_1,\ldots,E_n$ be ordinary elliptic curves over a
perfect field of characteristic $p$, where $n\geq2$, and set
\[
 X=\prod_{i=1}^nE_i,\qquad
 L=\boxtimes_{i=1}^n\mathcal O_{E_i}(a_iO_i),\quad a_i\geq3,
 \qquad h_d=d^n\prod_i a_i.
\]
Use the product of the indicated complete embeddings and
the Segre embedding as the fixed projective embedding of $X$.
Then the generic relative Hasse--Witt rank is
\begin{equation}\label{eq:product-elliptic-generic-rank}
 r_d=h_d-1\qquad\text{for every }d\geq1.
\end{equation}
In particular the generic relative coherent Frobenius is
invertible in every degree.

Suppose in addition that $p\mid a_i d$ for every $i$.
For geometric monodromy on its maximal-rank open one has
\begin{equation}\label{eq:product-entire-fixed-space}
 \dim_{\mathbf F_p}(\mathcal V_{d,\bar\eta}^{\vee})^{G_{\rm geom}}
 =\begin{cases}
       2^n-1,&p=2,\\
       n,&p\text{ odd}.
   \end{cases}
\end{equation}
After a finite extension splitting the Frobenius kernel
and normalizing its invariant top Cartier differential,
the arithmetic invariant space is the same space.
Thus these numbers give the dimension of the entire
trivial quotient, rather than a lower bound for it.
\end{theorem}

\begin{proof}
First we prove generic ordinarity.  For each factor
choose a section of $\mathcal O_{E_i}(a_i dO_i)$ with
distinct zeros.  It can be chosen as a product of $d$
sections of $\mathcal O_{E_i}(a_iO_i)$ with disjoint
reduced zero sets.  The resulting product section on
$X$ cuts out an SNC union of coordinate fibres.
Every nonempty stratum is a product of ordinary elliptic
curves, so Frobenius is bijective on every coherent
cohomology group of that stratum.  In the structure-sheaf
Cech spectral sequence, Frobenius is consequently
bijective on every page and on the total cohomology.

For every nonzero defining section, the Kunneth formula
gives $H^{n-1}(X,L^{-d})=0$ and
$h^n(X,L^{-d})=h_d$.  Also $h^n(X,\mathcal O_X)=1$.
The multiplication map to $H^n(X,\mathcal O_X)$ is
surjective because its cokernel lies in
$H^n(Y,\mathcal O_Y)=0$.  Hence the relative coherent
rank is constantly $h_d-1$ on the whole projective
linear system.  Invertibility at the SNC section is
an open condition, and the smooth open is nonempty.
This proves \eqref{eq:product-elliptic-generic-rank}.

The ambient hypersurface family is the same complete
linear system up to a surjective linear projection.
For completeness, a line bundle $N$ of degree at least
three on an elliptic curve is normally generated.
For multiplication $H^0(N)\otimes H^0(N^j)\to H^0(N^{j+1})$
with $j\geq2$, a base-point-free pencil in $H^0(N)$
and $H^1(N^{j-1})=0$ prove surjectivity.
For $j=1$, choose a point $q$ and a base-point-free
pencil in $H^0(N(-q))$.  The pencil exact sequence
after tensoring with $N$ has kernel $\mathcal O(q)$,
whose $H^1$ vanishes, so its products span
$H^0(N^2(-q))$.  The square of a section nonzero at
$q$ spans the remaining one-dimensional quotient.
This proves normal generation.  Tensoring the factors
and grouping Segre monomials proves the required
surjectivity for $X$ in every degree.

We now work over the algebraic closure.  Put
$W=H^0(X,L^d)$ and $P=\mathbf P(W)$, and normalize an
invariant top form $\omega$ by $C(\omega)=\omega$.
The coefficient can be normalized since $X$ is ordinary
and the base field is algebraically closed.
The relative coherent bundle fits in the exact sequence
\begin{equation}\label{eq:product-global-bundle}
 0\longrightarrow\mathcal E_d\longrightarrow
       H^n(X,L^{-d})\otimes\mathcal O_P(-1)
       \longrightarrow\mathcal O_P\longrightarrow0.
\end{equation}
It is a vector bundle on all of $P$, not only on the
ordinary or smooth locus.  Frobenius extends there by
$F_s=s^{p-1}F_X$.

Any geometric monodromy-invariant covector corresponds,
on the invertibility open, to a Frobenius-equivariant
map $\lambda:\mathcal E_d\to\mathcal O$.
It extends over all of $P$.  Indeed, in a regular local
frame its equation is $A^t\lambda=\lambda^{[p]}$ with
$A$ regular.  If its largest pole along any prime
divisor had order $m>0$, the right side would have
a coordinate of valuation $-pm$, whereas every coordinate
on the left has valuation at least $-m$, a contradiction.
There are no divisorial poles; normality of $P$ and local
freeness extend the section also across codimension two.

Dualizing \eqref{eq:product-global-bundle} and taking
global sections, using $H^1(P,\mathcal O_P)=0$, shows
that such a covector is represented by
\begin{equation}\label{eq:product-linear-lift}
 w_s=D(s)\omega,\qquad D\in\End(W),
\end{equation}
modulo scalar multiples of the identity.
Its Frobenius defect lies in the constant one-dimensional
kernel in the Frobenius pullback of the dual sequence.
It is therefore $c s^{[p]}$ for a constant $c$.
Adding $\alpha\operatorname{id}_W$ to $D$ changes this
constant by $\alpha-\alpha^p$; an Artin--Schreier equation
over the algebraically closed field removes it.
Thus every invariant covector lifts to a $D$ satisfying
\begin{equation}\label{eq:product-exact-Cartier-operator}
 C\bigl(s^{p-1}D(s)\omega\bigr)=D(s)\omega
            \quad\text{for all }s\in W.
\end{equation}
Two such exact lifts differ by $c\operatorname{id}$
with $c\in\mathbf F_p$.

The divisibility assumptions give
$L^d=F^*M$ with
$M=\boxtimes_i\mathcal O_{E_i^{(p)}}((a_i d/p)O_i^{(p)})$.
Since $\ker F\simeq\mu_p^n$, its character algebra gives
\begin{equation}\label{eq:product-character-W}
 W=\bigoplus_{\chi\in H}W_\chi,
 \qquad H\simeq\mathbf F_p^n,
 \qquad W_\chi=H^0(X^{(p)},M\otimes\mathcal N_\chi).
\end{equation}
Every $W_\chi$ is nonzero: each factor is a positive
degree line bundle on an elliptic curve.
The normalized Cartier trace is projection to the
trivial character followed by taking the Frobenius
root.  Nontrivial character terms are killed.

Apply \eqref{eq:product-exact-Cartier-operator} to a
section $s\in W_\chi$, and put
$T=\operatorname{pr}_\chi\circ D|_{W_\chi}$.
Only this component contributes to Cartier.  Raising
the equation to the $p$-th power gives
\[
 D(s)^p=s^{p-1}T(s).
\]
We do not assume that the Cartier root preserves the
character space.  Instead choose a nonzero eigenvector
$s_0$ of the finite-dimensional operator $T$, say
$T(s_0)=a s_0$.  The displayed equation implies
$D(s_0)=a^{1/p}s_0$; projecting to $W_\chi$ then gives
$a=a^{1/p}$, so $a\in\mathbf F_p$.
For arbitrary $s\in W_\chi$, substitute $s_0+z s$
and compare the coefficient of $z$.
The left side has no such coefficient, and the right
side gives $s_0^{p-1}(T(s)-a s)=0$.
Since $X$ is integral, $T(s)=a s$.
The same equation now gives $D(s)^p=a s^p$, hence
$D(s)=a s$.  Thus $D$ does preserve every character
space and acts there by a scalar in $\mathbf F_p$:
\begin{equation}\label{eq:product-character-scalar}
 D|_{W_\chi}=c_\chi\operatorname{id},\qquad c_\chi\in\mathbf F_p.
\end{equation}

It remains to determine the allowed functions
$c:H\to\mathbf F_p$.
Choose one nonzero section in each character space
and vary its coefficient independently.
For a nonpure multi-index $(m_\chi)$ with
$\sum m_\chi=p$, the coefficient of the corresponding
monomial in $s^{p-1}D(s)$ is the nonzero factor
$(p-1)!/\prod m_\chi!$ times
$\sum m_\chi c_\chi$.
Cartier kills that monomial unless
$\sum m_\chi\chi=0$ in $H$; if this weight is zero,
the product of the chosen nonzero sections is nonzero
and its Frobenius root is nonzero.
Independence of the coefficient monomials therefore
shows that the exact conditions are
\begin{equation}\label{eq:product-scalar-relations}
 \sum_\chi m_\chi c_\chi=0
 \quad\text{whenever}\quad
 \sum m_\chi=p,\quad\sum m_\chi\chi=0,
 \quad(m_\chi)\text{ is not pure}.
\end{equation}
Pure monomials already give the right side of
\eqref{eq:product-exact-Cartier-operator}.

When $p=2$, every zero-weight pair consists of two
equal characters, so there are no conditions in
\eqref{eq:product-scalar-relations}.  The space of
exact lifts has dimension $2^n$ over $\mathbf F_2$.
When $p$ is odd, put $a_x=c_x-c_0$.
The multiset consisting of $p-2$ zeros, $x$, and $-x$
gives $a_{-x}=-a_x$.
The multiset consisting of $p-3$ zeros, $x$, $y$, and
$-x-y$ then gives $a_{x+y}=a_x+a_y$.
The only nontrivial pure exception is $p=3$, $x=y$;
there the same additivity follows already from
$a_{-x}=-a_x$, since $2x=-x$.
Thus $a$ is an $\mathbf F_p$-linear functional on $H$.
Conversely every affine-linear function
$c_x=c_0+\ell(x)$ satisfies all displayed relations.
The space of exact lifts has dimension $n+1$.

Finally quotient the exact lifts by the one-dimensional
space $\mathbf F_p\operatorname{id}_W$.
This gives precisely \eqref{eq:product-entire-fixed-space}.
The character projectors, their indicated combinations,
and the normalized top differential are defined over
a finite field extension when the original base is
finite; over any perfect base they descend after a
finite algebraic extension.  Over that extension the
whole geometric invariant space is already arithmetic.
\end{proof}

\begin{corollary}\label{cor:product-two-explicit-fixed-spaces}
\status{Proved under stated hypotheses}
For the explicit surface over $\mathbf F_3$ in
\eqref{eq:odd-product-explicit-elliptic}, the relative
rank is $9d^2-1$ and the full geometric and arithmetic
invariant covector spaces both have dimension two,
for every $d\geq1$.
For the surface over $\mathbf F_2$ in
Example~\ref{ex:ordinary-abelian-surface-explicit}, the
relative rank is also $9d^2-1$, and both invariant
covector spaces have dimension three for every even
$d\geq2$.
The exact full images inside the pointwise stabilizers
of these quotients are not determined by this theorem.
\end{corollary}


% /root/monodromy, 19 November 2026; independently audited by /root.
\subsection{A terminating finite-field algorithm for the entire monodromy image}
\label{subsec:effective-entire-monodromy}
\status{Proved; independent parent proof audit and exact quartic computation passed.
This is an effective answer for specified finite-field inputs, not a structural
classification for arbitrary perfect base fields.}

\begin{theorem}\label{thm:effective-monodromy-matrices}
\status{Proved under stated hypotheses}
Suppose $k=\mathbf F_q$ and the fixed smooth projective
embedding $X\subset\mathbf P^m_k$ is given by homogeneous
equations.  For every
\[
 d\geq d_{\rm coh}=\max\{1,\operatorname{reg}(\omega_X)-1\},
\]
there is an explicit terminating algorithm that outputs,
up to a simultaneous choice of basis, the complete finite
lists of matrices
\[
 G_{\rm geom},G_{\rm arith}\subset\GL_{r_d}(\mathbf F_p)
\]
for the local system in A.4.  It also decides geometric
and arithmetic irreducibility and outputs the determinant
images.  The algorithm uses only finite linear algebra,
Groebner-basis computations, construction of finite fields,
and bounded enumeration of polynomials and matrices.
It does not call an oracle for Galois groups or for
factorization over an algebraic closure.
\end{theorem}

\begin{proof}
We give the steps and their termination proofs.

\emph{Step 1: the stable matrix.}
Use Proposition~\ref{prop:effective-generic-rank} over
the rational function field $K=\mathbf F_q(t_1,\ldots,t_M)$
of an affine chart of the ambient degree-$d$ parameter
space.  Let $B$ be its coherent Frobenius matrix and
$b$ the coherent dimension.  Compute the image of
$B B^{[p]}\cdots B^{[p^{b-1}]}$ by Gaussian elimination.
The Frobenius-linearized map on that image is invertible;
a chosen basis gives an explicit $A\in\GL_r(K)$, where
$r=r_d$.  If $r=0$, output the trivial groups and the zero
representation (which is not called irreducible), and stop.
This construction is valid even though $K$ is imperfect.

\emph{Step 2: the frame algebra.}
Let $T=(T_{ij})$ be an $r\times r$ matrix of indeterminates
and introduce one more indeterminate $z$.  Form
\begin{equation}\label{eq:effective-frame-algebra}
 R=K[T_{ij},z]/
        (A T^{[p]}-T,\ z\det T-1),\qquad
 D=\prod_{i=0}^{r-1}(p^r-p^i)=|\GL_r(\mathbf F_p)|.
\end{equation}
This is a finite etale $K$-algebra of dimension $D$.
Indeed the equations for one column, $Av^{[p]}-v=0$,
have Jacobian $-1$.  Their leading homogeneous forms
are an invertible linear combination of the $p$-th
powers of the coordinates, so they have no common
projective zero at infinity and define a scheme of
length $p^r$.  Its geometric points form an
$\mathbf F_p$-space of dimension $r$ and span the
whole $r$-dimensional vector space: otherwise their
span, of dimension $s<r$, would itself have only
$p^s$ fixed points by the same argument.
The invertible frames are thus a torsor under
$\GL_r(\mathbf F_p)$ and number exactly $D$.
Compute a Groebner basis of the displayed ideal,
a $K$-basis of $R$, and its multiplication table.

\emph{Step 3: a primitive element from a finite list.}
Let $e_1,\ldots,e_D$ be the computed basis and put
$N=D(D-1)/2$.  Choose the finite set
$S=\{1,t_1,t_1^2,\ldots,t_1^N\}\subset K$.
Test the elements
\[
 \theta=\sum_{j=1}^D a_j e_j,\qquad (a_1,\ldots,a_D)\in S^D,
\]
until the vectors $1,\theta,\ldots,\theta^{D-1}$ are
linearly independent.  At least one test succeeds.
Over a separable closure, the algebra has $D$ distinct
embeddings.  The product of the pairwise differences
of their values at $\sum a_je_j$ is a nonzero polynomial
of degree $N$ in the $a_j$.  A nonzero polynomial of
degree at most $N$ cannot vanish on a grid having more
than $N$ distinct values in every coordinate, by
induction on the number of variables.  The successful
$\theta$ has a square-free monic minimal polynomial
$f\in K[Z]$ of degree $D$, computed by linear algebra,
and $R=K[\theta]\simeq K[Z]/(f)$.

\emph{Step 4: factorization by bounded enumeration.}
Clear the denominators of $f$ to obtain a nonzero
$F\in\mathbf F_q[t_1,\ldots,t_M,Z]$.
Enumerate all normalized polynomials over $\mathbf F_q$
of total degree at most $\deg F$, and test divisibility
by polynomial division or Groebner reduction.
This is a finite list.  The irreducible factors of
$F$ are among it, since total degrees add in a
polynomial domain.  Repeated exact division therefore
factors $F$ completely.  Discard factors independent
of $Z$, which are units over $K$, and normalize the
remaining ones to be monic in $Z$.
Gauss's lemma proves that these are precisely the
irreducible factors of $f$ over $K$.
Thus this step is a fully specified finite algorithm;
practical multivariate factorization can replace the
enumeration without changing its correctness.

Choose one factor $f_0$ by a fixed lexicographic rule.
It defines a field component $R_0$ of $R$ and the
corresponding maximal ideal
\[
 I_0=(A T^{[p]}-T,\ z\det T-1,\ f_0(\theta))
       \subset K[T_{ij},z].
\]

\emph{Step 5: the arithmetic matrices.}
Enumerate all $g\in\GL_r(\mathbf F_p)$.
Its action on the frame algebra is the substitution
\[
 T\longmapsto Tg,\qquad z\longmapsto(\det g)^{-1}z.
\]
Retain exactly those $g$ for which $f_0(\theta)$ after
this substitution reduces to zero modulo $I_0$.
This is a finite collection of Groebner reductions.
The retained list is the stabilizer of the chosen
field component, hence exactly $G_{\rm arith}$ in
the corresponding frame.  Preservation of the one
extra equation suffices because the frame ideal is
already preserved; inclusion between the resulting
maximal ideals is equality.

To justify the monodromy identification explicitly,
choose a geometric frame in the chosen component.
Its Galois orbit is the image of $\operatorname{Gal}(K)$
acting on the frames.  Under the simply transitive
right frame action, the stabilizer of its connected
component is the same subgroup of $\GL_r(\mathbf F_p)$,
up to the common initial choice of frame.
Restricting the original local system to any dense
open does not change its image: the fundamental group
of a dense open in a normal connected scheme surjects
onto that of the scheme.  Thus the generic calculation
computes the image over the stated $U_d$.

\emph{Step 6: the geometric matrices with a finite
constant-field bound.}
Construct the finite field $k_D=\mathbf F_{q^D}$ and
put $K_D=k_D(t_1,\ldots,t_M)$.
A finite field of this degree is constructed by
enumerating monic degree-$D$ polynomials and testing
irreducibility; the elementary existence theorem
for finite fields guarantees termination.
Factor $f_0$ over $K_D$ by Step 4, now over $k_D$,
choose one factor $f_{00}$, and repeat the matrix
preservation test with its component ideal.
The resulting subgroup is $G_{\rm geom}$, and it is
a subgroup of the already chosen arithmetic image.

Here is the reason the finite extension degree $D$
suffices.  The selected field component is a connected
torsor under its stabilizer $H$, so its degree over
$K$ is $|H|$, a divisor of $D$.  Its full algebraic
constant field is some $\mathbf F_{q^c}$, where
$c\mid |H|$, hence $c\mid D$.
After extending to $k_D$, every resulting field
component has its constant field algebraically closed
in it.  A finitely generated field over a perfect
field is separably generated; therefore each such
component is geometrically integral.
Consequently these are exactly the geometric
components, and their stabilizers give the geometric
image, not an upper approximation to it.

\emph{Step 7: irreducibility and determinant.}
For each of the two finite matrix lists, enumerate
all row-echelon matrices representing the proper
nonzero subspaces of $\mathbf F_p^r$ and test
invariance under every listed element.
There is an invariant subspace if and only if the
representation is reducible.  Taking determinants
of the listed matrices gives both determinant
images.  Proposition~\ref{prop:exact-general-determinant-image}
is a much smaller independent calculation of the
same images and provides a consistency check.

Every step has a finite bound or an explicitly proved
termination criterion.  This proves the theorem.
\end{proof}

\begin{remark}\label{rem:effective-monodromy-scope}
\status{Proved under stated hypotheses}
The theorem outputs the whole finite group for each
specified finite-field input, with a proved finite
algorithm.  Its worst-case complexity is enormous:
even the frame dimension is $|\GL_r(\mathbf F_p)|$.
It does not supply a uniform structural description
of the groups as $X$ and $d$ vary.  Nor does it supply
a uniform algorithm for arithmetic monodromy over
an arbitrary perfect field presented without effective
field operations and factorization data.
Those distinctions are retained in the result ledger.
\end{remark}

\begin{example}\label{ex:effective-quartic-frame-algebra}
\status{Proved under stated hypotheses}
The retained program \path{lang_frame_quartic.py}
executes a nontrivial small frame calculation for a genuine
smooth plane-quartic family.  Over $\mathbf F_2(t)$ take
\[
 f_t=x^4+y^4+z^4+x^2yz+xz^3+txy^3+xyz^2.
\]
In the standard genus-three monomial basis its Hasse matrix is
\[
 A_t=\begin{pmatrix}1&0&0\\0&0&1\\0&t&1\end{pmatrix}.
\]
The program extracts these coefficients from $f_t$ and checks
that all three affine-chart singularity ideals have Groebner
basis $[1]$ over $\mathbf F_2(t)$.
Thus the generic curve is smooth and ordinary, since
$\det A_t=t\ne0$.

The first coordinate is a constant summand.  A frame of the
remaining rank-two Lang module has columns $(c^2,c)^t$ and
$(d^2,d)^t$, with algebra
\[
 R=\mathbf F_2(t)[c,d]/
 \left(c^3+\frac{c+1}{t},\quad
              d^2+cd+c^2+\frac1t\right).
\]
Its basis is $1,c,c^2,d,cd,c^2d$, and its frame determinant
is $cd(c+d)=1/t$.
The program verifies all six substitutions from
$\GL_2(\mathbf F_2)$ preserve this algebra.
It also finds the primitive element $\theta=c+td$,
whose minimal polynomial, after clearing denominators, is
\[
 t^2 Z^6+(t^3+t^2)Z^3+(t^4+t^2+1)Z^2
       +(t^4+t)Z+t^6+t^5+1.
\]
Linear independence of $1,\theta,\ldots,\theta^5$ and
annihilation by this polynomial are checked exactly.

Its geometric irreducibility has a direct certificate.
The first equation gives $t=(c+1)/c^3$, a separable
degree-three map of rational function fields.
Setting $u=d/c$ transforms the second equation into
\[
 u^2+u=\frac1{c+1}.
\]
The right side has a simple pole and is not an
Artin--Schreier coboundary over $\overline{\mathbf F}_2(c)$:
every pole of $v^2+v$ has even order.
Thus the frame algebra remains a field of degree six
after extending constants to the algebraic closure.
Both monodromy images of this particular curve family
are exactly
$\operatorname{diag}(1,\GL_2(\mathbf F_2))\subset\GL_3(\mathbf F_2)$.
This is a restricted one-parameter family, not the full
space of plane quartics, whose image was determined
earlier.  Its exact coefficient matrix, basis, six group
elements, minimal polynomial and smoothness certificates
are retained in \texttt{lang\_frame\_quartic.json}.
\end{example}


% /root/monodromy with suggestions by /root, 19 November 2026.
\subsection{Finite constant data for an arbitrary perfect base field}
\label{subsec:perfect-field-monodromy-data}
\status{Proved; independent parent proof and original-source audit passed.
The output is finite polynomial data and matrix lists, rather than a
uniform structural classification of the resulting finite groups.}

\begin{lemma}\label{lem:effective-separable-factorization}
\status{Proved under stated hypotheses}
Let $k_0$ be a finitely generated field over $\mathbf F_p$,
given by a finite field presentation.  Arithmetic and
factorization of separable univariate polynomials over
$k_0(t_1,\ldots,t_M)$ are terminating exact procedures.
For a polynomial that defines a finite separable field
extension of this rational function field, its absolutely
irreducible factors over $\bar k_0(t_1,\ldots,t_M)$ can
also be found by a terminating procedure, with their
coefficients in explicitly presented finite separable
extensions of $k_0$.
\end{lemma}

\begin{proof}
We give details so that no absolute-factorization or
Galois-group oracle is concealed in the assertion.

A separating transcendence basis of $k_0/\mathbf F_p$
and a finite separable primitive presentation over the
corresponding rational function field can be found
by elimination, linear algebra and enumeration.
Explicitly, enumerate candidate rational functions
for a transcendence basis, check algebraic independence
by elimination, and check that the remaining algebra
is finite by a zero-dimensional Groebner computation.
Separability is checked by the rank of the differentials,
or by the nonzero derivative of a primitive minimal
polynomial.  Such a candidate exists because the
ground field is perfect.  For a finite separable
extension, the primitive element is found by testing
linear combinations of a finite basis, as in the
finite-grid argument of the previous theorem.
Thus we may write
\[
 k_0(t)=R(\alpha),\qquad
 R=\mathbf F_q(u_1,\ldots,u_s,t_1,\ldots,t_M),
\]
with a separable minimal polynomial for $\alpha$.
All arithmetic is then polynomial reduction and linear
algebra over a rational function field.

Here is a direct factorization procedure over $R(\alpha)$.
For a monic separable polynomial $f(Z)$, try translates
$g_c(Z)=f(Z+c\alpha)$ for successively enumerated
$c\in R$, until
\[
 N_c(Z)=\operatorname{Norm}_{R(\alpha)/R}(g_c(Z))
\]
is square-free.  Its norm is computed by a resultant.
Only finitely many $c$ are bad: over an algebraic
closure its roots are $\beta-c\sigma(\alpha)$, where
$\beta$ ranges over the roots of $f^\sigma$ and
$\sigma$ over the distinct embeddings of $R(\alpha)$.
Within one embedding they are distinct; across two
different embeddings an equality excludes at most
one value of $c$.  The rational function field $R$
is infinite here.  If there are no transcendental
variables, the field is finite and bounded polynomial
enumeration directly factors $f$ instead.

Factor $N_c$ over $R$ by clearing denominators and
the bounded finite-field polynomial enumeration in
Step 4 of Theorem~\ref{thm:effective-monodromy-matrices}.
For each irreducible factor $h$ of $N_c$, compute
$\gcd(g_c,h)$ over $R(\alpha)$ and discard the units.
These gcds are irreducible and give the factors of
$g_c$.  To verify irreducibility, two roots of $g_c$
in the same $R$-Galois orbit are carried to one another
by an automorphism fixing $R(\alpha)$: otherwise the
target root would occur in two different conjugates
of $g_c$, contradicting square-freeness of $N_c$.
Undoing the translate factors $f$.  Thus all separable
factorization steps used below are effective.

For absolute factorization, clear denominators of the
given separable polynomial and remove its content
in $k_0[t]$, obtaining a primitive polynomial
$F(t,Z)$.  Its total degree bounds the total degree
of every polynomial factor over $\bar k_0$.
For each possible pair of leading monomials and
factor degrees, write two polynomials with unknown
coefficients, normalize the leading coefficient of
one to one, and equate their product with $F$.
Impose the selected $Z$-degree of the proposed factor
to be positive and require the selected leading
coefficients to be nonzero (an inverse variable
imposes this as an equation).  Permit the complementary
factor to be the unit polynomial; in particular the
full-degree factor is always one of the cases.
There are finitely many cases.  The coefficient
ideal is proper if and only if that case has a
factorization over the algebraic closure, by the
Nullstellensatz.  Properness is decided by whether
its Groebner basis contains one.

Select the least positive $Z$-degree for which a
factor exists.  A factor of that degree is absolutely
irreducible: a nontrivial splitting would yield a
smaller positive $Z$-degree.  There are no factors
independent of $Z$, by primitivity and Gauss's lemma.
Thus a finite list of coefficient tests proves what
degree to seek; it is not a search whose negative
outcome is undecidable.

We can find the factor's coefficients explicitly.
The field defined by the original irreducible factor
is separably generated over $k_0$.  Its algebraic
constant field is finite separable, and after adjoining
it its components are geometrically integral.
Indeed finiteness and geometric irreducibility are
\cite[Lemmas~10.47.8 and 10.47.13]{Stacks037P,Stacks037Q};
geometric reducedness follows from separability by
\cite[Lemmas~10.44.2 and 10.44.3]{Stacks030W,Stacks030X}.
Consequently the normalized absolute factors have
coefficients in finite separable extensions of $k_0$.
Enumerate separable irreducible polynomials over
$k_0$, testing irreducibility by the procedure above,
and enumerate coefficient tuples in their finite
extensions.  Dovetail these enumerations and test
the finitely many coefficient equations.  A solution
exists in one such extension, so the search terminates.
Repeat division and the finite coefficient tests
until all factors have been found; their positive
$Z$-degrees sum to the original degree.
This proves the absolute-factorization assertion
without assuming a preexisting algebraic closure
computing oracle.
\end{proof}

\begin{theorem}\label{thm:perfect-field-exact-monodromy-data}
\status{Proved under stated hypotheses}
Let $k$ be any perfect field of characteristic $p$,
and let the fixed smooth projective embedding $X/k$
be defined over a finitely generated subfield
$k_0\subset k$.  Fix $d\geq d_{\rm coh}$.
There is a terminating construction from a finite
presentation of this field of definition and the
equations of $X$ of the following finite data:
\[
 H_{\rm geom}\triangleleft H_0\subset\GL_{r_d}(\mathbf F_p),
 \qquad E/k_0\text{ finite Galois},\qquad
 \rho:H_0\twoheadrightarrow\operatorname{Gal}(E/k_0),
 \quad\ker\rho=H_{\rm geom}.
\]
The two groups are output as complete lists of matrices,
$E$ by a finite field presentation, and $\rho$ by its
action on generators of $E$.
In a compatible geometric frame, the requested images
over $k$ are exactly
\begin{equation}\label{eq:perfect-field-monodromy-formula}
 G_{\rm geom}=H_{\rm geom},\qquad
 G_{\rm arith}=\rho^{-1}
       \bigl(\operatorname{Gal}(E/(E\cap k))\bigr).
\end{equation}
The intersection is taken using the chosen embedding
of $E$ into a separable algebraic closure of $k_0$
inside an algebraic closure of $k$.
Changing the initial frame conjugates all the displayed
matrix groups simultaneously.

Geometric irreducibility is a finite subspace test on
$H_{\rm geom}$.  Arithmetic irreducibility is the same
test on the explicit preimage in
\eqref{eq:perfect-field-monodromy-formula}; the only
additional base-field information is the finite
intersection $E\cap k$.
\end{theorem}

\begin{proof}
The cohomological matrix and stable-image algorithm
of Proposition~\ref{prop:effective-generic-rank} work
over $k_0(t)$ as well: they use rational linear
algebra and Frobenius linearization, not perfection
of this intermediate field.  Form the frame algebra
and find a primitive element as before.
Use Lemma~\ref{lem:effective-separable-factorization}
to factor its polynomial over $k_0(t)$ and choose
one field component $L$.
The finite list of matrix substitutions preserving
that component is $H_0$.
It acts as the full Galois group of $L/k_0(t)$:
the component is a connected torsor under its
stabilizer, whose order equals its degree.

Let $E$ be the algebraic closure of $k_0$ in $L$.
It is finite separable by the field facts cited
in the preceding lemma, and it is Galois over $k_0$
because $L/k_0(t)$ is Galois.
The field $L/E$ is separably generated with $E$
algebraically closed in it, hence geometrically
integral by the same cited lemmas.
It follows that the components after extending
constants to $\bar k_0$ are indexed by the embeddings
of $E$; their stabilizer in $H_0$ is
\[
 H_{\rm geom}=\operatorname{Gal}(L/E(t)).
\]
The restriction map to $E$ is $\rho$.
Both matrix lists are obtained by the finite ideal
preservation tests, now using one absolute factor
for $H_{\rm geom}$.

We explain how the field $E$ itself is recovered,
rather than leaving it as an unspecified relative
algebraic closure.  Let $f_0$ be the primitive
polynomial of $L$, and choose a monic absolutely
irreducible factor $f_{00}$ over $\bar k_0(t)$.
Normalize every rational-function coefficient by
coprime numerator and denominator with one leading
coefficient equal to one.
Let $E_0$ be the field over $k_0$ generated by all
the scalar coefficients in these normalized expressions.
The factorization algorithm gives them in a finite
explicit extension, so $E_0$ is computed by adjoining
them and closing their vector span under multiplication.
This terminates inside that finite extension.

In fact $E_0=E$ with the selected embedding.
Since $L/E$ is geometrically integral, the geometric
components are indexed by the embeddings of $E$.
The stabilizer of one component under
$\operatorname{Gal}(\bar k_0/k_0)$ is exactly
$\operatorname{Gal}(\bar k_0/E)$.
On the other hand, that automorphism fixes the
chosen factor precisely when it fixes all of its
normalized coefficients, so this stabilizer is
$\operatorname{Gal}(\bar k_0/E_0)$.
The finite separable fields are equal.
Thus $E$ has been explicitly obtained.

Its embedding into $L$ and the restriction map
$\rho$ can also be calculated by linear algebra.
The selected factor realizes $L$ inside
$E(t)[Z]/(f_{00})$.
This latter field has degree
$[E:k_0]\deg f_{00}=[L:k_0(t)]$ over $k_0(t)$,
so it is the same field $L$.
Express the computed generators of $E$ in the
basis $1,\theta,\ldots,\theta^{[L:k_0(t)]-1}$
by solving a linear system.  Apply every retained
matrix substitution to these expressions and
reduce.  This outputs $\rho$ explicitly.

Finally extend the constants to $k$.
Put $F=E\cap k$.  Since $E/k_0$ is finite Galois,
$\operatorname{Gal}(Ek/k)=\operatorname{Gal}(E/F)$.
Since $L/E$ is geometrically integral, extending
$E$ to $Ek$ does not change its geometric degree
or create new components.  The component over
$k(t)$ therefore has degree
\[
 |H_{\rm geom}|[E:F],
\]
and its automorphism group is exactly the preimage
of $\operatorname{Gal}(E/F)$ under $\rho$.
This proves \eqref{eq:perfect-field-monodromy-formula}.
The geometric image is unchanged by any extension
of the algebraically closed constant field, giving
the first formula.  Subspace enumeration gives the
irreducibility assertions exactly as before.
\end{proof}

\begin{remark}\label{rem:perfect-field-data-scope}
\status{Proved under stated hypotheses}
The theorem is an exact finite-data description for
an arbitrary perfect field, and an algorithm whenever
the finite field-of-definition presentation and the
finite constant intersection can be evaluated.
It makes no claim that membership in an arbitrary
unpresented, possibly noncomputable perfect field
can be decided by a universal program.
For a finite base field, the intersection is computed
by finite-field arithmetic and the previous theorem
already gives a complete unconditional algorithm.
The structural question of identifying these outputs
uniformly as familiar finite groups is separate;
for projective space it was answered above, whereas
for general ambient varieties it is not inferred
from this construction.
\end{remark}


% /root/monodromy, 19 November 2026; independently audited by /root.
\subsection{Two exact full monodromy groups on a product surface}
\label{subsec:product-full-groups}
\status{Proved; independent parent proof, primary-source audit and exact-program rerun passed.}

The following computation goes beyond determining invariant covectors.
It determines the entire images in two degrees for one fixed embedded
ordinary abelian surface.  Its use of a projective parameter space is
essential: it permits a short, checkable certificate for all divisorial
inertia, including divisors whose general point parametrizes a singular
curve.

\begin{lemma}[A line certificate for generic stable defect one]
\label{lem:hasse-line-certificate}
\status{Proved under stated hypotheses}
Let $M(s)$ be a $b\times b$ matrix of homogeneous linear forms over
$\mathbf F_2$ in homogeneous coordinates on $P=\mathbf P^m$.
Suppose $h(s)=\det M(s)$ is nonzero.  Let $a,b_0$ be coefficient
vectors with $\det M(b_0)=1$.  If, at every root $\alpha$ of
$\det M(a+t b_0)$ over $\overline{\mathbf F}_2$, one has
\[
 \rank\bigl(M(a+\alpha b_0)^{[2]}M(a+\alpha b_0)\bigr)=b-1,
\]
then at the geometric generic point of every irreducible component
of $h=0$ the semilinear map with matrix $M(s)^t$ has stable rank
$b-1$.
\end{lemma}

\begin{proof}
The determinant is homogeneous of degree $b$, and its restriction to
the indicated affine line has degree exactly $b$.  For every geometric
irreducible homogeneous factor $q$ of $h$, one has $q(b_0)\ne0$.
Consequently $q(a+t b_0)$ has its full positive degree and has a root
over the algebraic closure.  If the rank of $M^{[2]}M$ were below
$b-1$ generically on $q=0$, then $q$ would divide every $(b-1)$-minor
of that matrix.  All these minors would vanish at every root of
$q(a+t b_0)$, contrary to the hypothesis.  Thus the first and second
Frobenius images both have dimension $b-1$ at that generic point.
Over a perfect geometric field their equality makes this image
Frobenius-stable and bijective.  This is the asserted stable rank.
The argument does not require the determinant divisor to be reduced.
\end{proof}

\begin{lemma}[Rank-one loss and transvection generation]
\label{lem:char2-defect-one-inertia}
\status{Proved under stated hypotheses}
Let $\mathcal E$ be a vector bundle of rank $r$ on a projective space
$P$ over an algebraically closed field of characteristic two, with
a Frobenius-linear endomorphism.  Suppose this map is generically
bijective, and at the generic point of every component of its
noninvertibility divisor its stable rank is $r-1$.
Then the geometric image of the Lang local system on the
invertibility open is generated by transvections.
\end{lemma}

\begin{proof}
At such a divisor, pass to the strict henselization of the discrete
valuation ring.  The Lang scheme $x-F(x)=0$ is separated, affine,
quasi-finite and etale: its Jacobian is the identity, and its
geometric fibres are finite.  Its finite etale part has rank
$2^{r-1}$ and is a subgroup, by
Lemma~\ref{lem:boundary-Lang-finite-part}, or directly by
\cite[Lemma~41.20.5]{Stacks0BTN} and
\cite[Lemma~58.8.2]{Stacks09ZS}.
Since the local ring is strictly henselian, inertia fixes its
$(r-1)$-dimensional generic fibre pointwise.  The quotient has
dimension one over $\mathbf F_2$, so inertia also acts trivially on
that quotient.  Every nonidentity inertia element therefore has
the form $1+u\otimes\lambda$, where $u$ lies in the fixed hyperplane
and $\lambda$ annihilates it.  Its difference from the identity has
rank one and square zero: it is a transvection.

Let $G$ be the geometric image and divide it by the normal subgroup
generated by these divisorial inertia groups.  Normalize $P$ in the
resulting finite Galois extension.  This finite normal cover is
unramified in codimension one and hence is etale everywhere by
purity, \cite[Lemma~58.21.4]{Stacks0BMB}.  Projective space over an
algebraically closed field has no nontrivial connected finite etale
cover, so the quotient is trivial.

For clarity, the last fact can also be proved without a fundamental
group calculation.  For $\mathbf P^1$ it is Riemann--Hurwitz.
For higher-dimensional projective space, pull a cover back to the
blowup of a point.  This blowup is a $\mathbf P^1$-bundle over
$\mathbf P^{m-1}$.  On each fibre the pulled-back cover is trivial;
cohomology and base change applied to its structure sheaf identifies
it with the pullback of a finite etale cover of $\mathbf P^{m-1}$.
The exceptional section is mapped to the original point, so the
restriction of the cover to that section is trivial.  Thus the
cover of $\mathbf P^{m-1}$, and then the pulled-back cover, is
trivial.  The blowup is birational, so the original connected
function-field extension also has degree one.
\end{proof}

\begin{lemma}[A small Frobenius certificate for a transvection group]
\label{lem:product-frobenius-group-test}
\status{Proved under stated hypotheses}
Let $N$ be either $8$ or $32$, and let
$H\triangleleft G\subseteq\GL_N(\mathbf F_2)$.
Suppose $H$ is nontrivial and generated by transvections, and
$G$ contains an element whose characteristic polynomial is
irreducible of degree $N$ and is not reciprocal.  Then
$H=G=\GL_N(\mathbf F_2)$.
\end{lemma}

\begin{proof}
The irreducible polynomial makes $G$ irreducible.  The socle of its
restriction to $H$ is nonzero and $G$-stable, so it is the whole
space.  The isotypic $H$-summands are permuted transitively by $G$.
If a simple constituent occurred with multiplicity $e$, the rank
of $h-1$ on its isotypic summand would be a multiple of $e$.
A nontrivial generating transvection has total rank one, so
$e=1$ on one summand and, by transitivity, on all of them.
Thus the space is a direct sum of simple $H$-modules of a common
dimension $c$, and $G$ embeds in
$\GL_c(\mathbf F_2)\wr S_{N/c}$ if there is more than one summand.
This is precisely the elementary argument of
Lemma~\ref{lem:transvection-clifford-blocks}; it includes the
possibility of imprimitive groups rather than assuming it away.

For $N=8$, an irreducible element has order divisible by $17$:
otherwise its order divides $15=2^4-1$, contradicting its degree.
For $N=32$, its order is divisible by $65537$: otherwise it divides
$65535=2^{16}-1$.  Here $65537$ is prime, as direct trial division
up to $257$ verifies, and
\[
 \operatorname{ord}_{17}(2)=8,\qquad
 \operatorname{ord}_{65537}(2)=32.
\]
The indicated prime does not divide $|\GL_c(\mathbf F_2)|$ for
$c<N$, nor $|S_{N/c}|$.  It therefore cannot divide the order of
the displayed proper wreath product.  Thus $H$ is irreducible.

We now use exactly one external finite-group classification:
\cite[Theorem~3.6, printed p.~22]{KantorSeress2001}.
Over $\mathbf F_2$, an irreducible group of dimension at least three
generated by transvection groups is the full special linear group,
or a symplectic group, or one of the listed orthogonal and symmetric
groups in their even-dimensional natural representations.
Each non-linear alternative preserves a nondegenerate alternating
bilinear form.  This form is unique: the endomorphism ring of the
irreducible $H$-module is a finite field, whose degree over
$\mathbf F_2$ divides the rank one of a transvection, and hence is
one.  Schur's lemma then gives a one-dimensional space of invariant
forms.  Every normalizing element of $G$ preserves that form, since
$\mathbf F_2^\times=1$.  Its characteristic polynomial would be
reciprocal, contrary to the hypothesis.  Only
$H=\mathrm{SL}_N(\mathbf F_2)=\GL_N(\mathbf F_2)$ remains.
\end{proof}

\begin{theorem}[Exact groups in degrees one and two on one fixed surface]
\label{thm:product-char2-two-full-groups}
\status{Proved under stated hypotheses}
Let
\[
 E:\quad Y^2Z+XYZ+X^3+Z^3=0\quad\text{over }\mathbf F_2,
 \qquad X=E\times E,
 \qquad L=\mathcal O_E(3O)\boxtimes\mathcal O_E(3O),
\]
with the fixed Segre embedding of the two plane cubics.
For the relative Hasse--Witt local system on the ordinary smooth
degree-$d$ hypersurface-section open, both geometric and arithmetic
images are, up to a common choice of basis,
\[
 \begin{array}{c|c|c}
 d&\rank\mathcal V_d&G_{\rm geom}=G_{\rm arith}\\ \hline
 1&8&\GL_8(\mathbf F_2)\\
 2&35&\left\{
       \begin{pmatrix}A&B\\0&I_3\end{pmatrix}:
       A\in\GL_{32}(\mathbf F_2),\ B\in M_{32\times3}(\mathbf F_2)
       \right\}.
 \end{array}
\]
In particular the second image is the full pointwise stabilizer of
the three independent covectors already constructed, of order
$2^{96}|\GL_{32}(\mathbf F_2)|$.  The same statements hold after
base change to any perfect field of characteristic two.
\end{theorem}

\begin{proof}
The generic ranks and the equality between the ambient hypersurface
family and the complete linear system have already been established
in Theorem~\ref{thm:product-exact-fixed-covectors}.  We recall the
matrix model, to make the finite computations reproducible.
Write $x=X/Z$, $y=Y/Z$ and $\omega=dx/x$, normalized by
$C(\omega)=\omega$.  Put $C(f\omega)=T(f)\omega$.
The equation $y^2+xy=x^3+1$ gives
\[
 T(x^{2a})=x^a,\quad T(x^{2a+1})=0,\quad
 T(x^{2a}y)=x^{a+1},\quad
 T(x^{2a+1}y)=x^a(y+1).
\]
For $d=1$ use the ordered basis $(1,x,y)$ of
$H^0(E,\mathcal O(3O))$; for $d=2$ use
$(1,x,x^2,x^3,y,xy)$ of $H^0(E,\mathcal O(6O))$.
Use the lexicographically ordered tensor basis on $X$.
An integer $s=\sum s_j2^j$ encodes the section with coefficient
$s_j\in\mathbf F_2$ on the $j$th tensor basis vector, beginning
with index zero.  Write $W=H^0(X,L^d)$ and $b=\dim W=9d^2$.

For a universal coefficient vector $s$, let $M(s)$ be the linear
matrix whose $j$th column is the coefficient vector of
$T\otimes T(s e_j)$, with the scalar Frobenius root removed.
Thus the actual Cartier map is
$v\mapsto M(s)^{[1/2]}v^{[1/2]}$, and the dual coherent
Frobenius matrix is $M(s)^t$.
The full coherent bundle is $W^*\otimes\mathcal O(-1)$, and its
kernel under evaluation at $s$ is the relative bundle
$\mathcal E_d=\Omega^1_{\mathbf P(W)}$.
Its quotient is $H^2(X,\mathcal O_X)\otimes\mathcal O$, with
bijective Frobenius.  The identity
\[
                       M(s)s=s^{[2]}
\]
expresses this quotient.  Therefore $\det M(s)$ has exactly the
same invertibility open as the relative Frobenius.

This open is defined on the whole projective parameter space,
including singular sections.  The Lang local system extends over
all of it.  Removing the smooth-discriminant locus does not change
its image: a dense open in a normal integral scheme surjects onto
its etale fundamental group.  In particular every ordinary
$\mathbf F_2$ coefficient vector below supplies an actual
arithmetic Frobenius class for the image in question; it need not
itself define a smooth curve.  On the dual representation its
characteristic polynomial is that of the induced Cartier matrix;
on the original representation it is reciprocal.  Our irreducibility,
nonreciprocity, order and fixed-space tests are unchanged by this
duality convention.

\par\smallskip\noindent\textbf{Divisorial inertia certificates.}
The retained exact program
\path{code/product_char2_group_certificates.py}
constructs all matrices from the four displayed Cartier formulas.
Its full output, including every binary column and every finite-field
rank, is
\path{outputs/product_char2_group_certificates.json}.
The following two affine lines have ordinary leading coefficient:
\[
 \begin{array}{c|r|r|l}
 d&a&b_0&\det M(a+t b_0)\\ \hline
 1&265&303&t^2(t^2+t+1)(t^5+t^4+t^2+t+1)\\
 2&56866400210&54262944944&(t+1)^4q_7(t)q_9(t)q_{16}(t),
 \end{array}
\]
where
\[
\begin{split}
 q_7&=t^7+t^6+t^5+t^3+t^2+t+1,\\
 q_9&=t^9+t^8+t^6+t^3+1,\\
 q_{16}&=t^{16}+t^{15}+t^{13}+t^{10}+t^6+t^3+t^2+t+1.
\end{split}
\]
The listed factors are irreducible over $\mathbf F_2$.
At a root of each listed factor, both $M$ and $M^{[2]}M$
have rank $8$ for $d=1$, respectively $35$ for $d=2$.
The program verifies irreducibility by the finite-field Frobenius
criterion, multiplies the factors back to the determinant, and
performs Gaussian elimination in the respective finite fields.
The determinants have full degrees $9$ and $36$.
Lemma~\ref{lem:hasse-line-certificate} therefore gives stable rank
$b-1$ for the full coherent module at every geometric generic
Hasse divisor.  Removing its bijective rank-one quotient gives
stable relative rank $b-2$, a loss of exactly one.
Lemma~\ref{lem:char2-defect-one-inertia} proves that both geometric
images are generated by transvections.  This argument includes
all Hasse components without asserting their irreducibility or
reducedness.

\par\smallskip\noindent\textbf{Degree one.}
For the ordinary section encoded by $258$, the matrix on $W$
has characteristic polynomial
\[
 (T+1)(T^8+T^4+T^3+T+1).
\]
The degree-eight factor is irreducible and not reciprocal.
For the ordinary section encoded by $27$, the characteristic
polynomial is $(T+1)^9$.
Thus the geometric image is not trivial: if it were trivial,
the local system would be constant over the algebraic closure,
and all rational points would have the same arithmetic Frobenius
characteristic polynomial.  The arithmetic image is irreducible
because of the first section.  Apply
Lemma~\ref{lem:product-frobenius-group-test} with $N=8$.
It gives both images equal to $\GL_8(\mathbf F_2)$.

\par\smallskip\noindent\textbf{The action on the common kernel in degree two.}
The four character projections $s_\chi$ of $s$ give four fixed
vectors in $W$, with sum $s$; in $W/\langle s\rangle$ they give
three independent invariant covectors.  All are defined over
$\mathbf F_2$, by the explicit Cartier calculation above and
Theorem~\ref{thm:product-exact-fixed-covectors}.
Let $K\subset V$ be their common kernel.  Then $\dim K=32$ and
the quotient $Q=V/K$ is the constant three-dimensional module
for both geometric and arithmetic monodromy.

At the ordinary section encoded by $17707720254$, the characteristic
polynomial on $W$ is $(T+1)^4f_{32}(T)$, where
\[
\begin{split}
 f_{32}(T)={}&T^{32}+T^{28}+T^{27}+T^{26}+T^{25}+T^{24}
             +T^{23}+T^{22}+T^{20}+T^{17}\\
            &+T^{14}+T^{11}+T^8+T^6+T^5+T^3+T^2+T+1.
\end{split}
\]
The program verifies that this polynomial is irreducible and not
reciprocal.  Consequently the arithmetic action on $K$ is
irreducible.  At the second ordinary section, encoded by
$65866883812$, the polynomial on $K$ has a root at $1$.
Thus the geometric image on $K$ is nontrivial, by the same
constant-local-system argument as in degree one.

The restriction to $K$ of a transvection of $V$ is either the
identity or a transvection: every image line lies in $K$, since
the three quotient covectors are fixed.  Hence the geometric
image on $K$ is generated by transvections.  Applying
Lemma~\ref{lem:product-frobenius-group-test} with $N=32$ shows
that this image is the entire $\GL(K)$.

\par\smallskip\noindent\textbf{The translation subgroup in degree two.}
It remains to prove that every map $Q\to K$ occurs; surjectivity
onto $\GL(K)$ alone does not prove this.
Choose a vector-space splitting $V=K\oplus Q$ and let $G_0$ be
the geometric image.  Its kernel $R$ in
\[
 1\longrightarrow R\longrightarrow G_0
          \longrightarrow\GL(K)\longrightarrow1
\]
is a $\GL(K)$-submodule of $\Hom(Q,K)=K\otimes Q^*$.
The natural module is absolutely simple with endomorphism ring
$\mathbf F_2$.  Thus $R=K\otimes U$ for a subspace $U\subseteq Q^*$.
Modulo $R$, the subgroup is the graph of a cocycle in
$H^1(\GL(K),K\otimes(Q^*/U))$.
This cohomology is zero by the elementary transvection calculation
in Lemma~\ref{lem:char2-group-inputs}, valid in every dimension
at least four.  Conjugating the splitting therefore makes
$G_0=R\rtimes\GL(K)$, and its invariant-vector space projects
isomorphically onto $U^\perp\subseteq Q$.

That invariant-vector space is stable under arithmetic monodromy.
Its projection to the constant quotient $Q$ is injective, so
arithmetic monodromy acts trivially on it as well.
We now use the additional exact certificate at $s=65866883812$:
if $C=M(s)$ and the four character vectors are $s_\chi$, then
\begin{equation}\label{eq:product-no-complement-certificate}
 \dim\bigl(\im(C-I)+\langle s\rangle\bigr)
 =\dim\bigl(\im(C-I)+\langle s,s_\chi:\chi\rangle\bigr)=33.
\end{equation}
The full columns and the four vectors are retained in the output.
On the relative coherent fibre $V\otimes\bar k=\operatorname{ann}(s)$,
a vector fixed by this rational-point Frobenius lies in
$\ker(C^t-I)\cap\operatorname{ann}(s)$.
Equation~\eqref{eq:product-no-complement-certificate} forces it
to annihilate every $s_\chi$ and hence to lie in $K$.
Every geometric invariant vector is arithmetic invariant, so this
applies to it.  But $K^{\GL(K)}=0$.  Therefore $V^{G_0}=0$,
$U^\perp=0$, and $U=Q^*$.
The kernel is the whole $\Hom(Q,K)$, proving the asserted full
pointwise stabilizer.  The arithmetic image is contained in that
same stabilizer and contains the geometric image, so the two agree.

Finally the geometric covers remain connected under algebraically
closed extension.  Both displayed groups are already defined over
$\mathbf F_2$, and the three degree-two covectors are rational
over that field.  After any perfect base extension the arithmetic
image is therefore still between the same equal geometric image
and the same upper bound.  This proves the final assertion.
\end{proof}

\begin{remark}[Verification scope]\label{auto:A_monodromy:9178c750cd75}
\status{Proved under stated hypotheses}

The group conclusions use the lemmas proved here, the elementary
cocycle and socle arguments already proved, and the single opened
finite-group classification specified above.  The retained program
provides exact polynomial, rank and Frobenius certificates.  It does
not infer monodromy from a random sample.  The exploratory search
chooses certificates, after which every statement needed from each
certificate is verified exactly.  The high-dimensional characteristic
polynomials are also recomputed independently by SymPy's integer
characteristic-polynomial routine and reduced modulo two.
No assertion for degree $d\ge3$ on this surface follows just from
these two degree calculations.
\end{remark}


% /root/monodromy, 19 November 2026; independently audited by /root and /root/reduction_descent.
\subsection{Uniform structure on the fixed ordinary product surface}
\label{subsec:product-uniform-structure}
\status{Proved; parent proof, primary-source and code audits passed.
The uniform jet, incidence and union argument was also independently audited by the reduction/descent agent.}

We first record the general finite-group form of the prime-order
test used above.  We then use one further computed base degree to
obtain a statement in every degree; the latter statement does not
identify the entire group in every degree.

\begin{lemma}\label{lem:primitive-prime-transvection-test}
\status{Proved under stated hypotheses}
Let $N\ge4$ and $H\triangleleft G\subseteq\GL_N(\mathbf F_2)$.
Suppose that $H$ is nontrivial and generated by transvections.
Assume $G$ contains an element $g$ whose characteristic polynomial
is irreducible and not reciprocal, and whose order is divisible
by a prime $\ell$ with $\operatorname{ord}_{\ell}(2)=N$.
Then $H=G=\GL_N(\mathbf F_2)$.
\end{lemma}

\begin{proof}
Use the socle and multiplicity argument in
Lemma~\ref{lem:product-frobenius-group-test}.
If $H$ were reducible, $G$ would be contained in
$\GL_c(\mathbf F_2)\wr S_{N/c}$ for a proper divisor $c$ of $N$.
The order condition gives $\ell\equiv1\pmod N$, hence $\ell>N$;
and $\ell$ divides no $2^i-1$ with $i<N$.
It therefore divides neither factor of this wreath-product order.
This contradicts the order of $g$.
Thus $H$ is irreducible.  The same opened
\cite[Theorem~3.6, p.~22]{KantorSeress2001}, uniqueness of the
invariant alternating form in its non-linear alternatives, and
nonreciprocity exclude those alternatives exactly as before.
\end{proof}

\begin{corollary}\label{cor:product-degree-three-full-group}
\status{Proved under stated hypotheses}
For the fixed surface and embedding in
Theorem~\ref{thm:product-char2-two-full-groups}, the geometric and
arithmetic groups in degree three are both
$\GL_{80}(\mathbf F_2)$, over every perfect extension of
$\mathbf F_2$.
\end{corollary}

\begin{proof}
Here $W$ has dimension $81$.  The one-factor ordered basis is
\[
                 (1,x,x^2,x^3,x^4,y,xy,x^2y,x^3y).
\]
Use the same tensor and binary coefficient conventions as before.
The ordinary section encoded by $1284714998244677440413117$
has characteristic polynomial $(T+1)f_{80}(T)$, where the binary
coefficient encoding of $f_{80}$ is
\[
                         1592084315789834604401113.
\]
The retained program verifies that it is irreducible and not
reciprocal.  Its root order is divisible by
$\ell=4278255361$, and
$\operatorname{ord}_{\ell}(2)=80$.
The independent verifier checks primality by trial division
through $\lfloor\sqrt\ell\rfloor$, checks the order by the
exponents $80,40,16$, and checks the root-order assertion by
modular polynomial exponentiation.  The ordinary section
$810021528320842360421666$ has a different relative characteristic
polynomial, so the geometric image is nontrivial.

For the line with
\[
 a=998130104707720808648124,
 \qquad b_0=2032072755533295013746953,
\]
the determinant is monic of degree $81$.  Its six irreducible
factors have degrees $5,10,10,11,11,34$, and at a root of every
factor both $M$ and $M^{[2]}M$ have rank $80$.
Thus Lemmas~\ref{lem:hasse-line-certificate} and
\ref{lem:char2-defect-one-inertia} give transvection generation.
Apply Lemma~\ref{lem:primitive-prime-transvection-test}.

The complete matrices, factors, ranks and section coefficients
are in \path{outputs/product_degree3_certificates.json};
\path{code/verify_product_degree3.py} reconstructs them, independently
computes both $81\times81$ characteristic polynomials over the
integers before reduction modulo two, and repeats all the line
rank calculations.  Its exact output is
\path{outputs/verify_product_degree3.json}.
\end{proof}

\begin{lemma}[The required jet separation]
\label{lem:cubic-power-jets}
\status{Proved under stated hypotheses}
If $L$ is very ample on a smooth variety, $L^3$ separates second
jets at one point and first jets at two distinct points.
On a smooth projective surface its complete linear system has an irreducible
discriminant whose generic member has exactly one ordinary node;
the nonreduced locus has codimension at least two.
\end{lemma}

\begin{proof}
Ratios of sections of $L$ give local coordinates, so products of
at most two such ratios span the local ring modulo the cube of
the maximal ideal.  Multiplication by a section nonzero at the
point identifies this with second-jet generation by $L^3$.
For distinct points $P,Q$, take a section $m$ of $L$ vanishing
at $Q$ and nonzero at $P$.  Sections $m^2t$ with $t\in H^0(L)$
vanish to order at least two at $Q$ and give arbitrary first
jets at $P$.  Interchanging the points proves the assertion.

The incidence of a section with a specified singular point is
a projective bundle over the surface, since value and first
derivative impose three independent linear conditions.  It is
irreducible of dimension one less than the parameter space.
Second-jet separation makes an ordinary quadratic node a nonempty
open condition in this incidence, including in characteristic two.
Two specified distinct singular points impose six independent
conditions; varying the two points adds four dimensions.
Thus the locus with two distinct singular points has codimension
at least two in the parameter space.  The generic singular member
has exactly one ordinary node, and the discriminant is an
irreducible divisor.  A nonreduced divisor is singular along a
positive-dimensional locus, so belongs to the codimension-two
locus just described.
\end{proof}

\begin{theorem}[Transvection generation in every degree]
\label{thm:product-all-degree-transvections}
\status{Proved under stated hypotheses}
For the fixed $X=E\times E$ and $L$ of
Theorem~\ref{thm:product-char2-two-full-groups}, every geometric
monodromy group $G_{{\rm geom},d}$, $d\ge1$, is generated by
transvections.  Equivalently for the proof, every geometric
irreducible component of the full degree-$d$ Hasse divisor has
generic stable rank exactly $9d^2-2$ on the relative coherent
bundle of rank $9d^2-1$.
\end{theorem}

\begin{proof}
The divisor statement has been checked in degrees one, two and
three by the exact line certificates above.  Fix $d\ge4$, put
$e=d-3$, and work over the algebraic closure.
Choose a general smooth ordinary curve $D\in|L^e|$.
Such curves exist for every $e\ge1$ by
Theorem~\ref{thm:product-exact-fixed-covectors}.
Write $P_3=\mathbf P H^0(X,L^3)$ and let $C$ vary in this space.

We describe the bad loci for a general pencil in $P_3$.
The base Hasse divisor has generic stable rank loss one, and its
exceptional stable-rank locus has codimension at least two in
$P_3$.  By Lemma~\ref{lem:cubic-power-jets}, nonreduced curves and
curves worse than the generic one-node member form loci of
codimension at least two in the discriminant analysis relevant
here.  In particular every divisorial Hasse component has a
reduced generic member with at most one node.

For fixed smooth $D$, the incidence
\[
 \{(P,C): P\in D,\ s_C(P)=0,\ ds_C|_{T_PD}=0\}
\]
is a projective bundle over $D$ and is irreducible.  Its image
$\Delta_D\subset P_3$ is an irreducible divisor.  At its generic
point, $C$ is smooth and $C\cap D$ has one point of length two
and is otherwise reduced.  Indeed, contact of order at least
three imposes three conditions at one point of $D$, giving
codimension at least two after that point varies.  Two distinct
tangencies impose four independent conditions and the two points
vary in dimension two, again giving codimension at least two.
A singularity of $C$ at the tangency point adds a further linear
condition.  A singularity elsewhere, together with the tangency,
imposes five independent conditions at two points, which vary in
dimension three.  All these independence claims follow from
the two jet-separation assertions of
Lemma~\ref{lem:cubic-power-jets}.

For a general ordinary $D$, the generic $C$ on $\Delta_D$ is
ordinary.  Here is a dominance check to justify this assertion,
rather than a transversality assumption on Hasse strata.
Consider the joint incidence $(P,C,D)$ with both curves smooth
at $P$ and with equal tangent lines there.  Value and first-jet
generation for $L^3$ and $L^e$ identify it, locally over $P\in X$,
with the irreducible rank-one locus of two nonzero proportional
cotangent vectors, times affine kernel spaces.  Hence the joint
incidence is irreducible.  It dominates each of the two complete
linear systems: for a smooth curve in either system choose a
point on it and prescribe a proportional nonzero first jet for
the other section.  The inverse images of the two ordinary opens
are therefore nonempty dense opens of this same irreducible
incidence, so their intersection is dense.  Passing to a general
$D$ gives the asserted ordinary generic member of $\Delta_D$.
Consequently $\Delta_D$ is not a component of the base Hasse
divisor.  Their intersections and their exceptional loci have
codimension at least two in $P_3$.

The locus of $C$ containing $D$ is also of codimension at least
two.  For $e\le3$ it is the linear subspace
$s_D H^0(L^{3-e})$, whose codimension is respectively $45,72,80$;
for $e>3$ it is empty.  We may therefore choose a projective
pencil $\Lambda\subset P_3$ avoiding all the codimension-two
loci above.  Its point at infinity is chosen ordinary and
transverse to $D$.

For every member of this pencil, put $Y=C\cup D$.
There are no common components.  The exact sequence
\[
 0\longrightarrow\mathcal O_Y
 \longrightarrow\mathcal O_C\oplus\mathcal O_D
 \longrightarrow\mathcal O_{C\cap D}\longrightarrow0
\]
and exactness of stable parts give
\begin{equation}\label{eq:product-union-stable-rank}
 f(Y)=f(C)+f(D)+\#(C\cap D)(\bar k)-1,
\end{equation}
where $f$ denotes the stable dimension of $H^1(\mathcal O)$.
To see the finite-scheme term explicitly, Frobenius kills the
nilradical of $H^0(\mathcal O_{C\cap D})$ and is bijective on
its reduced geometric quotient.  Also $H^0(\mathcal O_C)=
H^0(\mathcal O_D)=\bar k$.  The coherent analogue uses the
scheme length instead of the number of points:
\[
 p_a(Y)=p_a(C)+p_a(D)+\operatorname{length}(C\cap D)-1.
\]
At a Hasse point of $C$, the intersection is reduced and $C$
loses exactly one stable dimension.  At a point of $\Delta_D$,
both components are ordinary and the number of intersection
points is one less than their scheme length.  These are distinct
sets of pencil parameters.  At every other parameter the union
is ordinary.  Thus every nonordinary union on the pencil has
stable rank loss exactly one.

Multiplication by the fixed defining section $s_D$ embeds this
pencil as a genuine projective line in $\mathbf P H^0(X,L^d)$.
Its point at infinity is ordinary, so every geometric component
of the full Hasse divisor meets the line at a finite parameter.
If such a component had generic stable rank loss at least two,
the closed stable-rank locus defined by the appropriate minors
of an iterated Frobenius matrix would contain the component and
all of those intersection points.  This contradicts
\eqref{eq:product-union-stable-rank}.  Hence every full-system
Hasse component has generic loss one.  The relative and absolute
losses agree, since $H^1(X,\mathcal O_X)$ has bijective Frobenius.
Lemma~\ref{lem:char2-defect-one-inertia} now proves transvection
generation in every degree.
\end{proof}

\begin{theorem}[All characteristic-two invariant covectors on products]
\label{thm:product-char2-all-parities-fixed-quotient}
\status{Proved under stated hypotheses}
Let $X=\prod_{i=1}^n E_i$ be a product of ordinary elliptic
curves in characteristic two, $n\ge2$, and
$L=\boxtimes_i\mathcal O_{E_i}(a_iO_i)$ with $a_i\ge3$.
For any $d\ge1$, let
$e=\#\{i: a_i d\text{ is even}\}$.
Then the entire geometric invariant covector space has dimension
\[
             \dim(V_d^\vee)^{G_{{\rm geom},d}}=2^e-1.
\]
After a finite extension splitting the ordinary Frobenius kernels
and normalizing the invariant top Cartier differential, the
arithmetic invariant space is the same space.
In particular, on the fixed $E\times E$ with $a_1=a_2=3$,
this dimension is zero for odd $d$ and three for even $d$.
\end{theorem}

\begin{proof}
Work over an algebraic closure and normalize the invariant top
form by Cartier.  The proof of
Theorem~\ref{thm:product-exact-fixed-covectors}, through
\eqref{eq:product-exact-Cartier-operator}, used no divisibility
assumption.  It identifies invariant covectors with exact lifts
$D\in\End(W)$, modulo $\mathbf F_2\operatorname{id}$, satisfying
\begin{equation}\label{eq:char2-projector-exact}
        \beta(s,D(s))=D(s)^{[2]}\quad\text{for every }s\in W,
\end{equation}
where $\beta$ is the symmetric relative Cartier pairing on
global sections and $\beta(s,s)=s^{[2]}$.

Every such $D$ is an idempotent.  An eigenvector with eigenvalue
$a$ in \eqref{eq:char2-projector-exact} gives $a=a^2$.
If a nontrivial Jordan block remained, subtract its eigenvalue
times the identity and choose $v,w\ne0$ with $D(v)=w$ and
$D(w)=0$.  Substituting $v+tw$ in the same equation makes the
coefficient of $t$ on the left equal to $\beta(w,w)=w^{[2]}$,
whereas the right side is independent of $t$, a contradiction.
Thus $W=A\oplus B$ is its $1$- and $0$-eigenspace decomposition.
Substituting $a+tb$ shows that
\begin{equation}\label{eq:char2-projector-orthogonality}
                    \beta(A,B)=0.
\end{equation}

Let $F:X\to X^{(2)}$ be relative Frobenius and
$\mathcal F=F_*L^d$.  Finite duality gives a perfect symmetric
pairing
\[
 \mathcal F\otimes\mathcal F\longrightarrow L^{d,(2)},
\]
obtained by multiplying sections and applying the normalized
Cartier trace.  Its pairing on global sections is $\beta$.
The bundle $\mathcal F$ is globally generated.  On a factor
$E_i$, evaluation on a Frobenius fibre is evaluation on a
length-two subscheme; $H^1(\mathcal O_{E_i}(a_idO_i)(-Z))=0$
because its degree is $a_i d-2>0$.  Tensoring these surjections
proves global generation on the product.

The projector $D$ preserves every evaluation kernel.
Indeed, write a global section whose value in a fibre is zero
as $a+b$ with $a\in A,b\in B$.  Its pairing with every section
is zero in that fibre.  By
\eqref{eq:char2-projector-orthogonality}, the value of $a$ pairs
to zero with $A$ and with $B$, hence with the whole fibre by
global generation.  Perfectness makes that value zero, and the
same holds for $b$.  The evaluation kernel is a vector bundle;
vanishing on all closed fibres therefore gives a descended
bundle endomorphism $\mathcal D\in\End(\mathcal F)$.
Its action on global sections is the original $D$.

We compute this endomorphism algebra directly.  For an ordinary
elliptic curve and a line bundle $N$ of odd degree $m$, the
rank-two bundle $F_*N$ has degree $m$ and is stable.
If a line subbundle $A\subset F_*N$ existed with degree at least
$m/2$, adjunction would give a nonzero map $F^*A\to N$ and
$2\deg A\le m$, impossible for the proposed integral degree.
Thus every line subbundle has degree strictly below the slope.
A noninvertible nonzero endomorphism would have a rank-one
kernel and image, both of degree below the slope (saturating the
image only raises its degree), contradicting additivity of
degree.  Over the algebraically closed field it follows that
\[
                     \End(F_*N)=\bar k.
\]

For even degree, $N=F^*N'$ and the ordinary splitting gives
$F_*N=N'\oplus(N'\otimes T)$, where $T$ is the nontrivial
two-torsion character line bundle.  The two summands are
nonisomorphic and have no homomorphisms between them.
This is the dimension-one case of the opened ordinary-Frobenius
splitting in \cite[Lemma~4.5, p.~11]{SannaiTanaka2016}.
Their cross Cartier pairing is zero, since the trace kills the
nontrivial character.

Relative Frobenius on the product is the product of the relative
Frobenii.  The bundle $\mathcal F$ is therefore a direct sum of
$2^e$ external products, one for each choice of a character on
the even-degree factors.  Each external product has endomorphism
algebra $\bar k$, by the degree-zero Kunneth formula.  Distinct
summands have no homomorphisms, because on one even-degree factor
the corresponding degree-zero nontrivial character line has no
sections.  Thus
\[
                      \End(\mathcal F)=\bar k^{\,2^e}.
\]
Every exact $D$ is consequently a sum of some of these $2^e$
projectors.  Conversely the summands are orthogonal for the
Cartier pairing; every such projector satisfies
\eqref{eq:char2-projector-exact}.  The exact-lift space therefore
has dimension $2^e$ over $\mathbf F_2$.
Quotienting by the identity gives the asserted $2^e-1$.
All splittings and the normalized top form descend after a
finite extension, which also proves the arithmetic statement.
\end{proof}

\begin{remark}[What the uniform results leave open]\label{auto:A_monodromy:1386149fe4fc}
\status{Proved under stated hypotheses}

For the fixed surface, we now know the exact generic rank, the
entire invariant quotient in every degree, and that geometric
monodromy is generated by transvections in every degree.
The complete images are identified above in degrees one, two
and three.  A transvection-generated group with the stated
invariant quotient can still have further invariant subspaces
or a proper translation subgroup.  The two uniform theorems do
not remove those possibilities for all $d$, and no all-degree
full-stabilizer assertion is made here.
\end{remark}


% /root/monodromy, 19 November 2026; complete conditional reduction independently audited by /root.
\subsection{A single geometric gate for the full product groups}
\label{subsec:product-irreducible-hasse-reduction}
\status{Reducedness and the complete reduction proved and independently audited by the parent. The geometric hypothesis is discharged for every degree at least four in Theorem~\ref{thm:product-Hasse-irred}; the resulting unconditional all-degree images are Corollary~\ref{thm:product-char2-all-full-groups}.}

\begin{proposition}\label{prop:product-hasse-reduced}
\status{Proved under stated hypotheses}
For the fixed surface $E\times E$ and embedding considered above,
the full Hasse polynomial $h_d(s)=\det M_d(s)$ is geometrically
reduced for every $d\ge1$.
\end{proposition}

\begin{proof}
At the geometric generic point of any irreducible Hasse component,
Theorem~\ref{thm:product-all-degree-transvections} gives stable rank
$b-1$ for the full coherent matrix, where $b=9d^2$.
Thus $M$ and $M^{[2]}M$ both have rank $b-1$.
Let $u$ span $\ker M$ and let $\ell$ span the left kernel.
Since $\ker M^{[2]}$ is spanned by $u^{[2]}$, equality of the two
ranks says that $u^{[2]}$ is not in $\im M$, or equivalently
$\ell(u^{[2]})\ne0$.

The symmetric Cartier multiplication pairing gives
$M(s)w=M(w)s$ and $M(w)w=w^{[2]}$.
The adjugate at this rank-$b-1$ point is a nonzero scalar multiple
of $u\ell$.  Differentiating the determinant in the coefficient
direction $u$ therefore gives a nonzero scalar multiple of
\[
                     \ell M(u)u=\ell(u^{[2]})\ne0.
\]
Each component is consequently generically smooth and occurs
with multiplicity one.  The repeated factors in some retained
pencil determinants reflect tangency of those pencils, not a
nonreduced global Hasse divisor.
\end{proof}

\begin{lemma}[Frobenius-periodic reflexive sheaves]
\label{lem:periodic-reflexive-projective-trivial}
\status{Proved under stated hypotheses}
Let $P$ be projective space over an algebraically closed field of
characteristic two.  If a coherent reflexive sheaf $\mathcal S$
has an isomorphism $F^*\mathcal S\simeq\mathcal S$, then it is a
trivial vector bundle with its constant unit-Frobenius structure.
\end{lemma}

\begin{proof}
Fitting ideals commute with base change.  In a local ring of $P$,
each Fitting ideal $I$ therefore satisfies $I=I^{[2]}$.
If $I$ is proper, this gives $I\subseteq\mathfrak m I$, so
Nakayama makes $I=0$.  The usual Fitting-ideal criterion, or a
minimal presentation, now shows that the sheaf is locally free
of its generic rank.  Its Frobenius isomorphism gives a finite
etale Lang frame cover.  Projective space has trivial geometric
fundamental group, as proved in
Lemma~\ref{lem:char2-defect-one-inertia}, so this frame cover has
a section.  The resulting frame is Frobenius-fixed.
\end{proof}

\begin{lemma}[Determinant allocation]
\label{lem:irreducible-hasse-quotients-trivial}
\status{Proved under stated hypotheses}
Let $P=\mathbf P^m$ over an algebraically closed field of
characteristic two, and let $\mathcal E=\Omega_P^1$ carry a
generically invertible Frobenius-linear map.  Suppose its Hasse
divisor is reduced and irreducible.
Every nonzero proper invariant subspace of its Lang representation
has trivial quotient representation.
\end{lemma}

\begin{proof}
Let $U$ be the invertibility open and extend the unit-Frobenius
subbundle of an invariant subspace by saturation in $\mathcal E$.
This gives a coherent exact sequence
\[
 0\longrightarrow\mathcal S\longrightarrow\mathcal E
       \longrightarrow\mathcal Q\longrightarrow0,
\]
where $\mathcal Q$ is torsion-free and $\mathcal S$ is reflexive.
The latter assertion follows locally from the depth lemma:
the kernel of a vector bundle mapping to a torsion-free sheaf
on a regular scheme satisfies the depth-two condition.
Frobenius preserves the saturated subsheaf, because it does so
at the generic point.  Flatness of Frobenius on the regular
scheme preserves reflexivity (regular sequences may equivalently
be replaced by their squares).  It therefore induces regular
maps on $\mathcal S$ and, after reflexive extension, on
$\mathcal Q^{**}$.

In codimension one all terms are vector bundles, and determinants
of these maps give effective divisors satisfying
\[
                     D_{\mathcal E}=D_{\mathcal S}+D_{\mathcal Q}.
\]
If $D_{\mathcal S}=0$, the map
$F^*\mathcal S\to\mathcal S$ is an isomorphism in codimension
one.  Both sheaves are reflexive, so it is an isomorphism
everywhere.  Lemma~\ref{lem:periodic-reflexive-projective-trivial}
makes $\mathcal S$ a nonzero trivial bundle.  Its inclusion in
$\Omega_P^1$ contradicts $H^0(P,\Omega_P^1)=0$, which follows
at once from the Euler sequence.

The reduced irreducible divisor must consequently be allocated
entirely to $\mathcal S$, and $D_{\mathcal Q}=0$.
The same argument makes $\mathcal Q^{**}$ a trivial unit-Frobenius
bundle.  On $U$ it is the original quotient bundle, because
that quotient is already locally free there.  The composite
$\mathcal E\to\mathcal Q\to\mathcal Q^{**}$ is Frobenius-equivariant
and exhibits the quotient local system as constant.
\end{proof}

\begin{lemma}[A fixed natural block in every higher degree]
\label{lem:product-degree-three-block}
\status{Proved under stated hypotheses}
For every $d\ge3$, the geometric group on the fixed surface
contains a subgroup acting as the natural
$\GL_{80}(\mathbf F_2)$ on an $80$-dimensional direct summand
and trivially on a complementary summand.
\end{lemma}

\begin{proof}
The assertion for $d=3$ has already been proved.  For $d>3$ fix
$d-3$ ordinary smooth members $D_i\in|L|$ in general transverse
position, with no triple points, and put $D=\bigcup D_i$.
This fixed connected nodal curve is ordinary.
Vary an ordinary degree-three curve $C$, transverse to every
$D_i$ and avoiding all nodes of $D$.
The resulting curves $C\cup D$ are ordinary by the union sequence,
so their Lang sheaf is a restriction of the full ordinary
coherent family in degree $d$.  Its image is a subgroup of the
degree-$d$ geometric group.  The allowed $C$ form a dense open
of the degree-three ordinary parameter space; hence their
relative degree-three image remains $\GL_{80}(\mathbf F_2)$.

The union sequence and the constant ambient
$H^1(X,\mathbf F_2)$ give a filtration on this boundary
representation $V$,
\[
 0\subset T\subset T_0\subset V,
 \qquad T_0/T\text{ constant},\qquad V/T_0=Q,
 \quad\dim Q=80.
\]
Here $T$ is the graph term from $C\cap D$ and $Q$ is the relative
degree-three local system.  More explicitly, the quotient by
$T$ is
$(H^1(C,\mathbf F_2)\oplus H^1(D,\mathbf F_2))/H^1(X,\mathbf F_2)$;
the inverse image of the constant $H^1(D,\mathbf F_2)$ is $T_0$.
The ambient maps are injective, so these identifications are exact.

Each $C\cap D_i$ consists of $3L^2=54$ points.
Thus the image on $T$ is a subgroup of a product of $S_{54}$'s.
The image on $T_0$ is an extension of such a subgroup by a
subgroup of the additive two-group
$\Hom(T_0/T,T)$.  Every prime in its order is at most $53$.
In contrast $\GL_{80}(\mathbf F_2)$ is simple and its order is
divisible by $4278255361$.

We recall an elementary simplicity proof to specify this input.
The group $\GL_N(\mathbf F_2)$ for $N\ge3$ is perfect, since
$1+E_{ij}=[1+E_{ik},1+E_{kj}]$ for distinct indices, and the
elementary matrices generate it.  It acts faithfully and doubly
transitively on nonzero vectors.  A nontrivial normal subgroup
is therefore transitive.  The stabilizer of a vector $v$ has
the abelian normal subgroup
$U_v=\{1+v\otimes\phi:\phi(v)=0\}$, and the conjugates of $U_v$
generate the whole group.  Modulo a transitive normal subgroup,
every conjugate of $U_v$ has the same image: move $v$ by an
element of the normal subgroup and then use its stabilizer.
The quotient is thus abelian, and perfectness makes it trivial.

Goursat's elementary common-quotient description now says that
the boundary image on $Q\oplus T_0$ contains
$\GL(Q)\times1$: a nontrivial common quotient would be the
simple group $\GL(Q)$, whose large prime cannot occur on $T_0$.
Let $B$ be the inverse image of $\GL(Q)\times1$.
It fixes $T_0$ pointwise and maps onto $\GL(Q)$.
Its kernel $R$ is a submodule of
$\Hom(Q,T_0)=Q^*\otimes T_0$, hence
$R=Q^*\otimes T'$ for a subspace $T'\subseteq T_0$.
Modulo $R$, the subgroup is a graph of a cocycle with values in
$Q^*\otimes(T_0/T')$.  The elementary vanishing
$H^1(\GL_{80}(\mathbf F_2),Q^*)=0$ from
Lemma~\ref{lem:char2-group-inputs} makes it conjugate to the
zero cocycle.  Its inverse image is then
$R\rtimes\GL(Q)$, and in particular contains the required
block-diagonal $\GL(Q)$.  A quotient surjection alone would
not have supplied this last subgroup; the cocycle step does.
\end{proof}

\begin{theorem}[Conditional exact uniform product groups]
\label{thm:product-full-group-if-hasse-irreducible}
\status{Proved under stated hypotheses}
Fix $d\ge3$ on the surface $E\times E$ above.
If the degree-$d$ Hasse divisor is geometrically irreducible,
then both geometric and arithmetic groups are
\[
 \begin{cases}
 \GL_{9d^2-1}(\mathbf F_2),&d\text{ odd},\\
 \mathbf F_2^{\,3(9d^2-4)}\rtimes\GL_{9d^2-4}(\mathbf F_2),
       &d\text{ even},
 \end{cases}
\]
where in the second line the group is the full pointwise
stabilizer of the three known independent covectors.
Theorem~\ref{thm:product-Hasse-irred} below proves the required
geometric irreducibility for every $d\ge4$, thereby completing
the entire fixed-surface full-image problem together with the
three certified initial degrees.
\end{theorem}

\begin{proof}
Put $q=0$ for odd $d$ and $q=3$ for even $d$; these are the
entire invariant-covector dimensions by
Theorem~\ref{thm:product-char2-all-parities-fixed-quotient}.
Let $K$ be the common kernel and $N=\dim K=9d^2-1-q$.
Proposition~\ref{prop:product-hasse-reduced} and the hypothesis
permit Lemma~\ref{lem:irreducible-hasse-quotients-trivial}.
Every nonzero proper invariant subspace therefore contains $K$.
It follows that $K$ is irreducible (or that $V=K$ is irreducible
when $q=0$), and that $V$ has no invariant vector: an invariant
line could not contain $K$, whose dimension is at least $80$.

The image on $K$ is irreducible and generated by transvections,
by Theorem~\ref{thm:product-all-degree-transvections}.
The natural $80$-dimensional block of
Lemma~\ref{lem:product-degree-three-block} lies in $K$, because
a nontrivial natural module has no equivariant map to a trivial
module.  On $K$ that subgroup acts naturally on this block and
trivially on its complement.
No nondegenerate alternating form on $K$ can be invariant under
it.  The natural module of $\GL_{80}(\mathbf F_2)$ has no
invariant bilinear form: elementary transvections first kill
every off-diagonal matrix entry, using a third index, and then
every diagonal entry.  Its pairing with a trivial module is
also zero.  Hence the entire natural block would lie in the
radical of any such form.

The alternatives in the opened
\cite[Theorem~3.6, p.~22]{KantorSeress2001} are therefore excluded,
and the image on $K$ is $\GL(K)$.
If $q=0$ this is the desired full group.
If $q=3$, the translation-kernel and $H^1$ argument in
Theorem~\ref{thm:product-char2-two-full-groups} applies verbatim.
The kernel is $K\otimes U$ for some $U\subseteq Q^*$, where
$Q=V/K$ is constant; vanishing of the natural-module first
cohomology conjugates the group to
$(K\otimes U)\rtimes\GL(K)$.
Its invariant-vector space is $U^\perp$.
The already proved vanishing $V^G=0$ forces $U=Q^*$, so all
translations occur.  The arithmetic image is between this
geometric group and the same pointwise stabilizer, because the
three covectors are defined over $\mathbf F_2$.
It is consequently the same group.
\end{proof}

\begin{remark}[Why reducedness alone was insufficient]\label{auto:A_monodromy:24716fc9b916}
\status{Proved under stated hypotheses}

Generic stable loss one proves reducedness, but does not by
itself exclude two reduced components meeting in codimension
two. The subsequent proof of
Theorem~\ref{thm:product-Hasse-irred} supplies the missing
irreducibility by mixed tacnodal unions and smooth points
joining all restriction factors. It does not assume a generic
symmetric-matrix rank bound for the special Cartier matrices.
\end{remark}


% /root/monodromy, 2026-11-19 08:15 UTC. Complete proof independently audited by root and RD.
\subsection{Irreducible product Hasse divisors and all full images}
\label{subsec:product-hasse-tacnode}
\status{Proved. The complete Cartier, tacnode, small-degree jet and component-merging arguments were independently audited by the parent and the reduction-descent agent. This completes every degree for the fixed surface; it is not a classification for arbitrary ambient varieties.}

\begin{lemma}[A pointed Cartier calculation]
\label{lem:pointed-Cartier}
\status{Proved under stated hypotheses}
Let $C$ be a smooth projective curve in characteristic two with
one-dimensional Cartier kernel, generated by $\omega$.
If $P\in C$ satisfies $\omega(P)\ne0$, then Cartier induces a
surjection
\[
 H^0(C,\omega_C(2P))\longrightarrow H^0(C,\omega_C),
\]
and its kernel is $k\omega$. If $D$ is ordinary, the corresponding
map for $D$ is surjective with a one-dimensional kernel, whose
nonzero members have a nonzero double-pole coefficient at $P$.
\end{lemma}

\begin{proof}
There is no residue for a meromorphic differential with its only
possible pole at $P$, by the residue theorem. Cartier kills its
double-pole term, so the displayed target is correct.
The source has dimension $g(C)+1$ by Riemann--Roch.
For an ordinary curve, Cartier is already bijective on its
$g(C)$-dimensional holomorphic subspace, which proves the second
assertion.

For the first assertion, Serre duality and finite Frobenius
duality identify the displayed Cartier map with Frobenius
\[
 H^1(C,\mathcal O_C(-P))\longrightarrow
 H^1(C,\mathcal O_C(-2P)).
\]
The duality is the curve case of the local Cartier trace and
finite-duality formula in
\cite[Lemma~1.3.6 and Proposition~1.3.7, pp.~26--27]{BrionKumar2005};
its twists are exactly $F^*\mathcal O_C(-P)=\mathcal O_C(-2P)$.
The first space maps isomorphically to $H^1(C,\mathcal O_C)$,
while the kernel of the second space mapping to
$H^1(C,\mathcal O_C)$ is
$\mathfrak m_P/\mathfrak m_P^2$, the first-jet line.

Here is the boundary coefficient explicitly. Use an affine cover
with only one member $U_0$ containing $P$. A Cech representative
$c_{ij}$ of a nonzero Frobenius-kernel class satisfies
$c_{ij}^2=b_i-b_j$. The differentials $db_i$ glue, because
$d(c_{ij}^2)=0$, and give a nonzero global Cartier-killed
form $\omega$, up to scale. This also follows directly from
the local exact sequence
$0\to\mathcal O_C\xrightarrow{F}F_*\mathcal O_C
\xrightarrow{d} B^1_C\to0$: in a uniformizer, the exact
differentials are exactly those killed by Cartier.
The map on global constants is bijective, so the connecting
map identifies $H^0(B^1_C)$ with the Frobenius kernel on
$H^1(\mathcal O_C)$. In particular the glued differential is
nonzero.

Subtract the same constant from all the $b_i$ so that $b_0(P)=0$.
The Cech intersections avoid $P$, so the original cocycle also
represents its unique lift to $H^1(\mathcal O_C(-P))$.
Its Frobenius image is the boundary of
$b_0\bmod\mathfrak m_P^2\in\mathfrak m_P/\mathfrak m_P^2$.
The coefficient of a local parameter in this class is
$db_0(P)=\omega(P)$, up to the harmless sign of the connecting
homomorphism. Hence the lift of the one-dimensional Frobenius
kernel has nonzero image in this jet line. Any kernel of the
displayed twisted Frobenius would already lie in that original
kernel; it is therefore zero. Duality proves the surjectivity,
and the dimension calculation gives its Cartier kernel.
\end{proof}

\begin{lemma}[Smoothness at a mixed tacnodal union]
\label{lem:mixed-tacnode-Hasse}
\status{Proved under stated hypotheses}
Let $X=E\times E$ be the fixed ordinary surface above.
Let $C\in|L^a|$ and $D\in|L^e|$ be smooth curves without common
components, ordinary except that $C$ has a one-dimensional
Cartier kernel $k\omega$. Suppose they have a single simple
tangency at $P$, are otherwise transverse, and $\omega(P)\ne0$.
Then the full degree-$(a+e)$ Hasse divisor is smooth at $C\cup D$.
\end{lemma}

\begin{proof}
The normalization condition can be read directly from the
hypersurface dualizing generator, so no assertion about a smooth
normalization being a smooth family is required. In completed
coordinates the two branches have equations $y=0$ and
$y=\phi(x)$, with $\phi(x)=c x^2+O(x^3)$ and $c\ne0$.
The normalization ring is the pair of power-series rings, and
its subring for the union consists of pairs congruent modulo
$\phi(x)$, equivalently modulo $x^2$.
The adjunction generator $dx/(\partial(y(y-\phi))/\partial y)$
restricts on each branch to $dx/\phi(x)$, up to a sign, which
is immaterial in characteristic two. Multiplying by the
normalization subring shows that dualizing forms have at most
double poles, with matching double-pole coefficients and
matching residues. Conversely these two equalities are precisely
the congruence modulo $x^2$. At an ordinary node the same
calculation gives simple poles and opposite residues.

Cartier on dualizing forms acts componentwise on these rational
differentials. This follows either from the same local trace
formula or its compatibility with residues: the residue of a
Cartier image is the square root of the original residue.
The local matching conditions are preserved; the double-pole
term is killed and the simple-pole term becomes its square root.
A Cartier-killed dualizing form therefore has zero residues,
so it consists of Cartier-killed elements of
$H^0(\omega_C(2P))$ and $H^0(\omega_D(2P))$ with matching
double-pole coefficients. By
Lemma~\ref{lem:pointed-Cartier}, the first is a multiple
of the holomorphic $\omega$, while a nonzero element of the
second has nonzero polar coefficient. The matching condition
forces the second to be zero. Thus the Cartier kernel on
$C\cup D$ is exactly $k(\omega,0)$.

The exact adjunction sequence
$0\to\mathcal O_X\xrightarrow{s}L^{a+e}\to\omega_{C\cup D}\to0$
has an ordinary two-dimensional ambient Cartier cokernel on
global sections, because $H^1(X,\mathcal O_X)$ is ordinary.
Its kernel is the Cartier-unit line $ks$. Consequently its
Cartier kernels identify canonically: the one-dimensional
kernel just computed lifts uniquely, after adjusting by $ks$,
to
$\ker M(s)$, where $s=f_Cg_D$. Let $u$ be its generator.
Since its differential vanishes on $D$, the section is divisible
by $g_D$: $u=g_Da_0$ with $a_0\in H^0(L^a)$ and $a_0(P)\ne0$.

Choose local coordinates so $C$ is $y=0$ and $D$ has equation
$y+x^2+$ terms of higher order, allowing a nonzero coefficient
of $x^2$ and multiplying by a unit. The section $s$ has quadratic
term proportional to $y^2$, with no $xy$ term. The local Cartier
trace of a top differential at $P$ extracts its $xy$ coefficient.
Consequently the Frobenius-twisted evaluation at $P$ annihilates
$M(s)$ and spans its left kernel. Explicitly, for the matrix
convention with the output scalar square root removed, this is
the functional with coefficient vector $(e_i(P)^2)_i$ in a
local frame: $\mathcal C(sv)(P)^2=\operatorname{ev}^{[2]}_P M(s)v$.
This twist does not affect vanishing or the following nonzero
local coefficient. The section $u$ has linear term a nonzero multiple
of $y$. Choose a global deformation $\dot s$ with zero value at
$P$ and first jet $x$; very ampleness supplies it.
Then
\[
       \operatorname{ev}^{[2]}_P\bigl(M(\dot s)u\bigr)\ne0.
\]
The rank-one adjugate formula makes this a nonzero scalar
multiple of $dh_{a+e}(\dot s)$. Hence the Hasse divisor is
smooth at this union.
\end{proof}

\begin{lemma}[The degree-one jet exception is small]
\label{lem:product-small-degree-jets}
\status{Proved under stated hypotheses}
On the fixed product, $L^e$ for $e\ge2$ separates second jets at
one point and first jets at any two distinct points.
For $L$ itself, second jets are separated at a point
$P=(P_1,P_2)$ when neither $P_i$ is a flex in its plane cubic.
For every $Q\ne P$, the joint first-jet evaluation has rank at
least five; for fixed $P$, rank six fails only on the two fibres
through $P$ and a finite set of other points.
\end{lemma}

\begin{proof}
An elliptic line bundle of degree at least six interpolates
any subscheme of length at most four, since the residual line
bundle has positive degree and hence vanishing first cohomology.
Tensoring the interpolation maps on the two factors gives the
first assertion. For degree three, the map to length-three
jets at $P_i$ is an isomorphism unless
$\mathcal O_E(3P_i)\simeq\mathcal O_E(3O)$, precisely the flex
condition. Tensoring again gives second jets on the surface.

Write $V_i=H^0(E_i,\mathcal O(3O))^*$, let $e_P$ be the
one-dimensional evaluation direction in its projective plane,
and let $S_P\subset V_i$ be the two-dimensional vector tangent
plane to the embedded cubic. The dual of the surface first-jet
map at $P$ has image
\[
 A_P=e_{P_1}\otimes S_{P_2}+S_{P_1}\otimes e_{P_2},
 \qquad\dim A_P=3.
\]
Distinct points of a smooth plane cubic have distinct tangent
lines: a common tangent line would have total intersection
multiplicity at least four with a cubic.
If $P_1\ne Q_1$ and $P_2\ne Q_2$, then
\[
 A_P\cap A_Q\subseteq
 (S_{P_1}\cap S_{Q_1})\otimes(S_{P_2}\cap S_{Q_2}),
\]
whose dimension is one. A nonzero decomposable tensor in $A_P$
has either its first factor $e_{P_1}$ or its second factor
$e_{P_2}$; applying the same observation at $Q$ leaves only
\[
 e_{P_1}\in S_{Q_1},\ e_{Q_2}\in S_{P_2},
 \quad\text{or}\quad
 e_{Q_1}\in S_{P_1},\ e_{P_2}\in S_{Q_2}.
\]
For fixed $P$ these conditions permit only finitely many $Q$:
a tangent meets a cubic in a finite subscheme. For completeness,
tangents through a fixed point are finite for the specific
cubic here without any separability assumption on a Gauss map.
For $F=Y^2Z+XYZ+X^3+Z^3$, its partial derivatives are
\[
 F_X=YZ+X^2,\qquad F_Y=XZ,\qquad F_Z=Y^2+XY+Z^2.
\]
No nonzero constant linear combination is zero: compare the
coefficients of $X^2$, $Y^2$, and $XZ$. A nonzero quadratic
cannot vanish identically on the irreducible cubic. Hence for
every projective point $(a:b:c)$ the equation
$aF_X+bF_Y+cF_Z=0$ cuts a finite subscheme on $E$. Its points
are exactly those whose tangent contains $(a:b:c)$.
If, for example, $P_1=Q_1$ and $P_2\ne Q_2$, expanding in a
basis of $S_{P_1}$ beginning with $e_{P_1}$ shows
\[
 A_P\cap A_Q=e_{P_1}\otimes(S_{P_2}\cap S_{Q_2}),
\]
again of dimension one. This proves all stated rank bounds.
\end{proof}

\begin{theorem}[Uniform Hasse irreducibility]
\label{thm:product-Hasse-irred}
\status{Proved under stated hypotheses}
For the fixed surface and embedding, the degree-$d$ Hasse divisor
is geometrically irreducible for every $d\ge4$.
\end{theorem}

\begin{proof}
Work over an algebraic closure and put $e=d-3\ge1$.
Let $H_1,\ldots,H_t$ be the irreducible components of the
already reduced degree-three Hasse divisor in
$P_3=\mathbf P H^0(X,L^3)$.
Its general member on each $H_i$ is smooth, has stable loss one,
and has a one-dimensional Cartier kernel. Smoothness follows
because the singular-curve discriminant is irreducible and its
general member is ordinary: it contains an ordinary transverse
union of an ordinary degree-one and an ordinary degree-two
curve, so its nonempty ordinary open is dense.

For each $H_i$, choose a smooth general point $C$ and a point
$P\in C$ where its unique Cartier-killed differential is nonzero.
The tangent linear functional $dh_3|_C$ fails to be independent
of the first-jet functionals at at most one point of $C$.
Indeed first-jet evaluations at two distinct points are jointly
surjective by Lemma~\ref{lem:cubic-power-jets}; their dual spans
have zero intersection. Since $dh_3|_C\ne0$, it cannot belong
to both spans. We may therefore also impose this independence
at $P$, and require its two elliptic coordinates to avoid the
finite flex sets. These exclusions remove only finitely many
points of the smooth irreducible ample curve $C$.

Consider the incidence of $(C,P,D)$ over this open of pairs
$(C,P)$, with $D\in|L^e|$ tangent to $C$ at $P$.
The universal smooth curve over the irreducible open of $H_i$
is irreducible, since its fibres are connected smooth curves.
The two conditions on $D$ are independent linear first-jet
conditions, so the incidence is a projective bundle over an
irreducible base before removing open conditions.

The projection to $|L^e|$ is submersive on the chosen locus.
For precision, work on the affine cone over $H_i$: its tangent
space is $\ker dh_3$, and independence of $dh_3$ from the
first-jet functionals makes its map to the three-dimensional
first-jet space surjective. Projective scaling is harmless.
For an arbitrary infinitesimal change of $D$, move $P$ to
cancel its value change, using $dD(P)\ne0$. Then choose the
first jet of the change of $C$ inside $\ker dh_3$ to preserve
$C(P)=0$ and equality of tangent directions. These are exactly
the three incidence equations. The projection is consequently
dominant.

We also check nonemptiness of its good open, including $e=1$.
Fix a good $(C,P)$. If $e\ge2$,
Lemma~\ref{lem:product-small-degree-jets} gives the
usual independent conditions: second jets permit a nonzero
normal derivative and contact of exactly two at $P$;
a singularity of $D$ at another $Q$ imposes three further
conditions, while $Q$ varies in dimension two; another contact
with $C$ imposes two further conditions, while $Q\in C$ varies
in dimension one. Thus both bad loci are proper.

For $e=1$, second jets are still separated at the chosen $P$.
Outside the exceptional two fibres and finite set of
Lemma~\ref{lem:product-small-degree-jets}, the same
counts apply. On that exceptional locus, joint first-jet rank
is at least five. The two contact conditions at $P$ together
with three singularity conditions at $Q$ therefore have total
rank at least four; here $Q$ varies in dimension at most one.
Relative to the two fixed conditions, the bad locus again has
codimension at least one. For an additional contact at $Q\in C$,
the four conditions have rank at least three. The exceptional
$Q$ on $C$ are finite, since an irreducible ample curve contains
neither fibre through $P$. This bad locus is likewise proper.
These estimates prove that smooth $D$ with just the chosen
simple tangency and all other intersections transverse form a
nonempty open in the same irreducible incidence.

A nonempty open in an irreducible dominant incidence remains
dominant. It intersects the inverse image of the ordinary open
of $|L^e|$. Its image therefore contains a nonempty open of that
linear system. There are finitely many $H_i$, so a single general
smooth ordinary $D$ works for every $i$.

Fix this $D$, with equation $g$. On the multiplication subspace
$gP_3\subset\mathbf P H^0(L^d)$, the support of the Hasse divisor
is precisely
\[
                      H_1\cup\cdots\cup H_t\cup\Delta_D,
\]
as a divisorial statement. Here $\Delta_D$ is the irreducible
simple-contact divisor of
Theorem~\ref{thm:product-all-degree-transvections}.
This follows from the union stable-rank formula; exceptional
nonreduced members and members containing $D$ have codimension
at least two. The generic points of these displayed divisors
are distinct, and there is an ordinary member outside them.

For each $i$ the incidence just constructed gives a point of
$H_i\cap\Delta_D$ where the full Hasse divisor is smooth, by
Lemma~\ref{lem:mixed-tacnode-Hasse}. Any global Hasse component whose restriction contains either
$H_i$ or $\Delta_D$ contains this chosen intersection point,
because it contains that entire closed prime divisor on $gP_3$.
At a smooth point of the full Hasse divisor there is only one
such global component. Thus all components assigned to either
prime are the same. This argument allows arbitrary pullback
multiplicities and does not assume unique assignments in
advance.
Finally, every global Hasse component restricts nontrivially
to $gP_3$: a nonzero homogeneous polynomial of positive degree
cannot become a nonzero constant on this projective linear
subspace, and the subspace contains an ordinary point.
Thus its restriction contains one of the displayed prime
divisors. All have just been assigned to a single component,
so the full Hasse divisor is irreducible.
\end{proof}


\begin{corollary}[All full groups on the fixed ordinary product]
\label{thm:product-char2-all-full-groups}
\status{Proved under stated hypotheses}
Let $k$ be any perfect field of characteristic two, let
\[
 E:\ Y^2Z+XYZ+X^3+Z^3=0,\qquad
 X=E\times E,\qquad
 L=\mathcal O_E(3O)\boxtimes\mathcal O_E(3O),
\]
and use the fixed Segre embedding. For every $d\ge1$, the
maximal relative Hasse--Witt rank is $r_d=9d^2-1$. On the smooth
ordinary degree-$d$ hypersurface-section open, the entire
geometric and arithmetic images coincide and are
\[
 G_{{\rm geom},d}=G_{{\rm arith},d}=
 \begin{cases}
 \GL_{9d^2-1}(\mathbf F_2),&d\text{ odd},\\[2mm]
 \left\{\begin{pmatrix}A&B\\0&I_3\end{pmatrix}:
 A\in\GL_{9d^2-4}(\mathbf F_2),\quad
 B\in M_{(9d^2-4)\times3}(\mathbf F_2)\right\},&d\text{ even}.
 \end{cases}
\]
In even degree this is exactly the pointwise stabilizer of all
three independent invariant covectors, and its order is
$2^{3(9d^2-4)}|\GL_{9d^2-4}(\mathbf F_2)|$.
The representation is irreducible exactly in odd degree.
\end{corollary}

\begin{proof}
The generic rank and exact invariant covectors were proved in
Theorem~\ref{thm:product-char2-all-parities-fixed-quotient} and
the preceding product rank calculation. Degrees one and two
are Theorem~\ref{thm:product-char2-two-full-groups}; degree three
is Corollary~\ref{cor:product-degree-three-full-group}.
For $d\ge4$, Theorem~\ref{thm:product-Hasse-irred} discharges the
geometric hypothesis in
Theorem~\ref{thm:product-full-group-if-hasse-irreducible}.
Its proof supplies the entire linear group in odd degree and
all translations as well as the full linear kernel group in
even degree. The arithmetic and geometric equality persists
under any perfect extension of $\mathbf F_2$: the finite
geometrically connected covers giving the geometric images
remain connected after algebraically closed field extension,
and the three covectors in even degree are already defined
over $\mathbf F_2$. Finally the natural full linear module is
irreducible, while the even-degree common kernel of the three
covectors is a nonzero proper invariant subspace.
\end{proof}


% /root/monodromy, 2026-11-19. Proved; independent root and reduction_descent audits passed.
\subsection{Uniform Hasse geometry for every ordinary elliptic product}
\label{subsec:ordinary-products-general-structure}
\status{Proved and checked. The parent and reduction/descent agent independently audited the complete proof. The universal polynomial symmetrizer calculation was independently rerun by both agents. Full group exclusion is treated separately.}

\begin{lemma}[A universal elliptic Cartier symmetrizer]
\label{lem:ordinary-elliptic-universal-symmetrizer}
\status{Proved under stated hypotheses}
Let $E:y^2+xy=x^3+A x^2+B$ over a perfect field of characteristic
two, with $B\ne0$, and put $N=\mathcal O_E(3O)$.
In the basis $(1,x,y)$ and with the Cartier-invariant differential
$\eta=dx/x$, the root-removed Cartier multiplication matrix is
\[
 M(s)=\begin{pmatrix}
 s_0&B s_2&B s_1\\
 s_2&s_1+A s_2&s_0+A s_1\\
 0&s_2&s_1+s_2
 \end{pmatrix}.
\]
The constant invertible matrix
\[
 R=\begin{pmatrix}1&0&B\\0&0&B\\0&B&(1+A)B\end{pmatrix},
 \qquad\det R=B^2,
\]
satisfies
\[
 RM(s)=\begin{pmatrix}
 s_0&0&B s_2\\
 0&B s_2&B(s_1+s_2)\\
 B s_2&B(s_1+s_2)&B s_0+B s_1+(1+A)B s_2
 \end{pmatrix}.
\]
Thus $RM(s)$ is symmetric and its diagonal is $R^t s$.
For any pair of ordinary elliptic curves with degree-three
line bundles, the analogous nine-dimensional Cartier pencil
has a constant invertible symmetrizer and an invertible
linear map to its nine diagonal entries.
\end{lemma}

\begin{proof}
The smooth ordinary Weierstrass equation can be put in the
displayed form by first normalizing its nonzero $xy$
coefficient, then translating $x$ to remove the $y$ term and
translating $y$ by a constant to remove the $x$ term.
These changes preserve $O$ and the degree-three line bundle.
The discriminant in the resulting equation is $B$, so $B\ne0$.
A degree-three line bundle is a translate of $\mathcal O(3O)$
over the algebraic closure, which suffices for the geometric
assertion.

Writing a rational function as $a^2+x b^2$ gives
$\mathcal C((a^2+x b^2)\eta)=a\eta$.
The Weierstrass relation consequently gives
\[
 \mathcal C(1)=1,\quad \mathcal C(x)=0,\quad
 \mathcal C(y)=x,\quad
 \mathcal C(xy)=\sqrt B+\sqrt A\,x+y.
\]
Together with $\mathcal C(x^2)=x$ and
$\mathcal C(y^2)=y$, this is precisely the matrix written above
when the output scalar square root is removed.
Direct multiplication proves the formulas for $R$.
The matrix $R$ itself need not be symmetric; only $RM(s)$ is.

On a product the Cartier multiplication coefficients are tensor
products of the two one-factor coefficient matrices. Therefore
$R_1\otimes R_2$ symmetrizes all nine coefficient matrices.
Its diagonal map is $R_1^t\otimes R_2^t$, and its determinant
is $B_1^6B_2^6\ne0$.
The universal polynomial identities over
$\mathbf F_2[A_1,B_1,A_2,B_2]$ are independently reproducible in
\path{agents/monodromy/code/ordinary_product_universal_symmetrizer.py};
the complete retained output is
\path{agents/monodromy/outputs/ordinary_product_universal_symmetrizer.json}.
\end{proof}

\begin{lemma}[Squarefreeness from independent symmetric diagonals]
\label{lem:symmetric-independent-diagonal-squarefree}
\status{Proved under stated hypotheses}
Let $S(z)$ be an $m\times m$ symmetric matrix of linear forms
in characteristic two, and suppose its diagonal entries are
independent linear coordinates $z_1,\ldots,z_m$.
Then $\det S$ is geometrically squarefree.
\end{lemma}

\begin{proof}
In the determinant expansion, permutations having a cycle of
length at least three cancel with the inverse of the first such
cycle. The surviving terms are indexed by matchings:
\[
 \det S=\sum_{\mathcal M}
 \prod_{\{i,j\}\in\mathcal M}S_{ij}^2
 \prod_{i\text{ unmatched}}z_i.
\]
Only the empty matching contributes to the monomial
$z_1\cdots z_m$, and its coefficient is one. Every term of a
nonempty matching includes a square of a linear form and hence
has some exponent at least two in each of its monomials.
If a nonconstant polynomial $q$ had $q^2\mid\det S$ over an
algebraic closure, write $\det S=q^2r$ with $\deg r<m$.
Every exponent in $q^2$ is even. A term contributing to
$z_1\cdots z_m$ would therefore require a monomial of $r$
with all $m$ exponents odd, whose degree is at least $m$.
This is impossible. Thus no irreducible factor is repeated.
\end{proof}

\begin{lemma}[Jet bounds for arbitrary ordinary cubic factors]
\label{lem:ordinary-product-uniform-jets}
\status{Proved under stated hypotheses}
Let $X=E_1\times E_2$ with both elliptic curves ordinary in
characteristic two and $L=\mathcal O_{E_1}(3O)\boxtimes
\mathcal O_{E_2}(3O)$. The conclusions of
Lemma~\ref{lem:product-small-degree-jets} hold on $X$.
In addition, the singular-curve discriminant of $|L|$ is
irreducible with a general one-node member, and its general
member is ordinary.
\end{lemma}

\begin{proof}
The tensor first-jet proof of that lemma uses only the plane
cubic embeddings. Its one finiteness point is also uniform:
for $F=Y^2Z+XYZ+X^3+A X^2Z+B Z^3$ the partials are
\[
 F_X=YZ+X^2,\qquad F_Y=XZ,\qquad
 F_Z=Y^2+XY+A X^2+B Z^2.
\]
Their constant linear combinations are linearly independent,
by comparing $Y^2,X^2,XZ$. No nonzero quadratic vanishes
identically on the cubic, so the tangents through any fixed
point form a finite set. Thus the exceptional rank-five pairs
are the pairs sharing one elliptic coordinate, of dimension
three in $X\times X$, together with the crossed incidence
locus, which is finite over the first point and has dimension
at most two. Elsewhere the two first jets have rank six.

A specified singular point imposes three independent conditions,
so its incidence is an irreducible projective bundle over $X$.
An ordinary node occurs in a nonempty open of every fibre:
products of sections vanishing simply in the two elliptic
directions provide an arbitrary nonzero mixed quadratic
coefficient. For two singular points, the rank-six locus gives
codimension $6-4=2$ after the points vary. The exceptional
rank-five loci give codimensions at least $5-3=2$ and $5-2=3$.
Hence the general discriminant member has exactly one node,
and nonreduced members have codimension at least two.

Finally $|L|$ contains an ordinary grid: take three distinct
vertical fibres with sum $3O$ on $E_1$ and three distinct
horizontal fibres with sum $3O$ on $E_2$.
Its normalization components are ordinary elliptic curves and
its graph term is Frobenius-bijective, so the grid is ordinary.
The irreducible discriminant thus contains an ordinary point,
and its ordinary open is nonempty and dense.
\end{proof}

\begin{theorem}[Uniform irreducible Hasse divisors on ordinary products]
\label{thm:ordinary-products-Hasse-irreducible}
\status{Proved under stated hypotheses}
Let $E_1,E_2$ be arbitrary ordinary elliptic curves over a
perfect field of characteristic two, and let
$L=\mathcal O_{E_1}(3O)\boxtimes\mathcal O_{E_2}(3O)$.
The full Hasse divisor in $|L|$ is geometrically reduced and
has generic Cartier corank one on every component.
For every $d\ge2$, the full Hasse divisor in $|L^d|$ is
geometrically irreducible and reduced, with generic stable
rank loss exactly one. Consequently the geometric relative
Hasse--Witt group in degree $d\ge2$ is generated by
transvections.
\end{theorem}

\begin{proof}
The two preceding algebraic lemmas make the degree-one Hasse
polynomial geometrically squarefree. Its derivative is nonzero
at each generic component point, so its matrix has corank one
there. By Lemma~\ref{lem:ordinary-product-uniform-jets}, no
component is the singular-curve discriminant, whose generic
point is ordinary. Thus the generic curve on each degree-one
Hasse component is smooth with a one-dimensional Cartier
kernel. Stable loss one in this degree is not assumed or
needed in the following argument.

Fix $d\ge2$, put $e=d-1$, and repeat the mixed-tacnode proof
with $C\in|L|$ and a general ordinary $D\in|L^e|$.
We spell out the differences caused by using $|L|$ in place
of $|L^3|$.
For a smooth general $C$ on a base Hasse component, the
nonzero tangent functional $dh_1$ belongs to the dual
first-jet space at only finitely many points of $C$.
Indeed, if there were infinitely many, choose one such $P$;
all but finitely many other points $Q$ on $C$ avoid the two
fibres and the finite crossed set through $P$.
The first-jet dual spaces at $P$ and such a $Q$ are disjoint,
contradicting that the same nonzero $dh_1$ lies in both.
We can therefore choose $P$ avoiding those finitely many
points, the zeros of the Cartier-killed differential, and the
flex fibres. The affine tangent cone $\ker dh_1$ maps
surjectively to first jets at this $P$.

The incidence with $D$ tangent to $C$ at $P$ consequently
projects dominantly to $|L^e|$, by the same three-equation
infinitesimal calculation in
Theorem~\ref{thm:product-Hasse-irred}. The good open of smooth
$D$ with exactly one simple contact is nonempty. For $e\ge2$
use full two-point first jets; for $e=1$ use rank at least five
and its exceptional one-dimensional locus through $P$.
The latter meets $C$ in only finitely many points. Exactly the
same dimension counts as in that theorem exclude a second
contact or a singularity of $D$.
Hence a single general ordinary $D$ admits such a mixed
contact with every base Hasse component.

For completeness the fixed-$D$ contact divisor is still
irreducible with generic ordinary $C$ despite the weaker
base jets. Its incidence over $P\in D$ is a projective bundle
cut out by two independent conditions. Two contacts impose
four conditions on the rank-six locus of pairs in $D^2$;
the rank-five exceptional pairs have dimension at most one
and impose at least three. Both give codimension at least
two in $|L|$. A singularity of $C$ at a second point imposes
five conditions generically, with the point pair varying in
dimension three; on the exceptional locus the rank is at
least four and the dimension is at most two. A singularity
at the contact point also has codimension at least two.
Order-three contact imposes three conditions at a nonflex
point; the flex fibres meet $D$ in finitely many points, so
even there the two contact conditions give codimension two.
The generic contact is therefore simple with $C$ smooth.
The joint tangency incidence for $(P,C,D)$ is irreducible
and dominates both complete linear systems by first jets.
Its two ordinary inverse images are nonempty dense opens,
so the generic contact $C$ is ordinary for a general ordinary
$D$.

On the multiplication subspace $g_D|L|$, the divisorial
nonordinary support is the union of the base Hasse components
and this contact divisor, by the exact union formula.
The exceptional nonreduced base locus has codimension at least
two; when $e=1$, $C=D$ is a single point of $\mathbf P^8$,
and when $e>1$ no base member contains $D$.
At each mixed contact, Lemmas~\ref{lem:pointed-Cartier} and
\ref{lem:mixed-tacnode-Hasse} make the full Hasse divisor
smooth. Their proofs use only ordinarity of the ambient product,
not the particular elliptic equation, and require only Cartier
corank one on $C$, not stable rank loss one on $C$. Therefore every base
Hasse component and the contact divisor are assigned to the
same global Hasse component. Every global component has a
positive-degree nonzero restriction to this linear subspace,
which contains an ordinary point. The component-merging
argument of Theorem~\ref{thm:product-Hasse-irred} proves that
there is exactly one global component. The smooth mixed
point also makes its multiplicity one.

There is an ordinary pair $(C,D)$ with a single simple
contact. The union formula gives stable loss one at that
point of the now irreducible Hasse divisor. The closed
iterated-Frobenius rank loci imply generic stable loss at
most one, and noninvertibility makes it exactly one.
Lemma~\ref{lem:char2-defect-one-inertia} proves the last assertion.
\end{proof}

\begin{corollary}[An irreducible kernel for every ordinary pair]
\label{cor:ordinary-products-irreducible-kernel}
\status{Proved under stated hypotheses}
In the preceding theorem let $d\ge2$, put $b=9d^2$, and set
$q=0$ for odd $d$ and $q=3$ for even $d$.
The entire invariant-covector space has dimension $q$.
Its common kernel $K\subset V_d$ is irreducible of dimension
$b-1-q$, and $V_d$ has no invariant vector.
The image on $K$ is irreducible and generated by transvections.
\end{corollary}

\begin{proof}
The fixed-covector calculation is already
Theorem~\ref{thm:product-char2-all-parities-fixed-quotient}.
Its statement and proof allow distinct ordinary factors: the
one-factor Frobenius pushforward calculations and their external
tensor product do not require the elliptic curves to be isomorphic.
The relative coherent bundle on $\mathbf P^{b-1}$ is
$\Omega^1_{\mathbf P^{b-1}}$, by the Euler/adjunction sequence
and the ordinary one-dimensional ambient $H^2(\mathcal O_X)$.
Theorem~\ref{thm:ordinary-products-Hasse-irreducible} therefore
permits Lemma~\ref{lem:irreducible-hasse-quotients-trivial}.
Every nonzero proper invariant subspace has trivial quotient,
so contains the common kernel. This makes $K$ irreducible and
rules out invariant lines, since $\dim K\ge32$.
The transvection assertion follows by restriction; a nontrivial
generator restricts either to a transvection or to the identity.
\end{proof}


% /root/monodromy, 2026-11-19. Proved; independent root and reduction_descent audits passed.
\subsection{An isogeny construction of a uniform linear block}
\label{subsec:ordinary-products-isogeny-full}
\status{Proved and checked. The parent and reduction/descent agent independently audited the geometric generic-fibre argument, nodal-image incidence, group extraction, arithmetic fixed covectors, and effective isogeny threshold. The original Chai theorem and Stacks generic-fibre lemma were opened by both auditors; the supporting exact computations were independently rerun.}

\begin{lemma}[An open locus of full degree-one images]
\label{lem:ordinary-pair-GL8-open}
\status{Proved under stated hypotheses}
There is a nonempty Zariski open subset $B^\circ$ of the space
$B=\mathbf G_m^2$ of pairs
\[
 E_{b_i}:y^2+xy=x^3+b_i
\]
on which the geometric relative degree-one image for
$\mathcal O(3O)\boxtimes\mathcal O(3O)$ is
$\GL_8(\mathbf F_2)$.
\end{lemma}

\begin{proof}
The nine sections $(1,x,y)\otimes(1,x,y)$ trivialize the
coefficient projective bundle over $B$.
Let $U\subset\mathbf P^8\times B$ be the smooth ordinary open.
It has nonempty geometrically integral fibres.
The relative rank-eight Lang frame cover $Z\to U$ is finite
etale of degree $|\GL_8(\mathbf F_2)|$.
At $b_1=b_2=1$ its geometric fibre is connected by
Theorem~\ref{thm:product-char2-two-full-groups}.

We justify passage to the geometric generic fibre, since $U/B$
is not proper. Suppose $Z_{\bar\eta}$ were disconnected.
After a finite separable extension of $k(B)$ all its components
are defined. Normalize $B$ in that extension; the resulting
integral normal scheme $B'$ is finite surjective over $B$: the
base is of finite type over a field, hence excellent and has
finite normalization in a finite field extension.
The scheme $U'=U\times_B B'$ is a nonempty open of a projective
bundle over a normal integral base, hence normal and integral.
The finite etale cover $Z'=Z\times_B B'$ is normal, so its
irreducible components are disjoint open-and-closed subschemes.
Each component is finite etale over $U'$ and has nonempty image,
which is both open and closed; it therefore surjects onto $U'$.
The distinct geometric generic components consequently all
meet the fibre over a point of $B'$ lying above $(1,1)$.
This contradicts geometric connectedness of that fibre.
Thus $Z_{\bar\eta}$ is connected. It is smooth, being finite
etale over an open of geometric projective space, so it is
normal; distinct irreducible components of a normal Noetherian
scheme are disjoint. It is therefore irreducible, as required.

By \cite[Lemma~37.27.5]{Stacks0559}, a finite-type morphism with
geometrically irreducible generic fibre has geometrically
irreducible fibres after restricting the integral base to a
nonempty open. Apply this to $Z\to B$.
Connectedness of a full frame cover is precisely transitivity
on all frames, hence equality of the monodromy image with
$\GL_8(\mathbf F_2)$. This gives $B^\circ$.
\end{proof}

\begin{lemma}[Ordinary nodal images under odd product isogenies]
\label{lem:ordinary-product-isogeny-nodal-images}
\status{Proved under stated hypotheses}
Let $\phi_i:E_i'\to E_i$ be odd-degree isogenies of ordinary
elliptic curves in characteristic two, of degrees $m_i$, and
put $\phi=\phi_1\times\phi_2$.
For a general smooth ordinary
$C'\in|\mathcal O_{E_1'}(3O)\boxtimes\mathcal O_{E_2'}(3O)|$,
its image $C=\phi(C')$ is an ordinary nodal integral curve,
$C'\to C$ is its normalization, and
\[
 \mathcal O_X(C)=
 \mathcal O_{E_1}(3m_2O)\boxtimes\mathcal O_{E_2}(3m_1O).
\]
Over a suitable dense open of the $C'$ parameter space, the
permutation image on the nodes of $C$ is a subquotient of a
product of copies of $S_{18}$.
\end{lemma}

\begin{proof}
Work over an algebraic closure, so the odd finite etale kernel
of $\phi$ is a constant group. Its translations act freely.
A general $C'$ is not preserved by any nonzero kernel
translation. If such a translation preserves the linear system,
its projective action on it is nontrivial: a scalar action
would act identically on the very ample embedding of the
surface, contradicting that the translation is nonzero.
Its fixed locus is therefore proper. If it does not preserve
the line bundle, it cannot preserve a member.
Thus $\phi|_{C'}$ is generically one-to-one.

For a nonzero kernel point $g$, the possible double points
are parametrized by $P\in C'\cap(C'-g)$, and this intersection
has degree $L'^2=18$. We show it is transverse for a general
$C'$. The two conditions $C'(P)=C'(P+g)=0$ are independent.
If both coordinates of $g$ are nonzero, the two first jets
have rank six except at finitely many $P$: the crossed tangent
conditions reduce to equations $3P_i=\text{constant}$.
On the rank-six locus, proportionality of the two nonzero
gradients is one further independent equation. The bad
incidence has codimension three with $P$ varying in dimension
two, and hence has proper image in $\mathbf P^8$.
The finitely many exceptional $P$ already impose two independent
value conditions, so their image is proper as well.

If exactly one coordinate of $g$ is zero, the joint first-jet
rank is five everywhere. Its sole relation involves only the
values and the two derivatives along the factor on which
$g$ translates; the transverse derivatives remain independent.
After the two values are set to zero, the relation on the two
parallel derivatives is nonzero, since two distinct evaluation
directions are independent. Thus the determinant of the two
gradients is not identically zero on their three-dimensional
allowed derivative space: one may choose a nonzero parallel
derivative and then an independent transverse derivative.
Tangency is again one extra equation. The same codimension
three, two-dimensional-point count excludes it generically.

Triple points are excluded by three-point values.
Three distinct Segre rank-one tensors are collinear only when
one factor is constant on all three and their other factor
points are collinear. For $P,P+g,P+h$ with distinct kernel
translations, this exceptional case requires both translations
to occur in the same elliptic factor, and then collinearity is
$3P_i+g_i+h_i=0$. Its locus in the surface has dimension one.
Outside it the three evaluations have rank three, while on it
they have rank two. The corresponding incidences have dimensions
at most $8-3+2=7$ and $8-2+1=7$, respectively.
Only finitely many kernel pairs are involved, so a general
$C'$ has no triple identifications. Etaleness of $\phi$ and
transversality now show that its image has only ordinary nodes.

The image class is the norm of $L'$.
Pushing the divisor $3(\{O\}\times E_2')+
3(E_1'\times\{O\})$ multiplies the first component by $m_2$
and the second by $m_1$, which gives the displayed exact
line bundle, with no unspecified degree-zero twist.
Ordinarity follows from the normalization exact sequence:
its graph term is Frobenius-bijective and $C'$ is ordinary.
Finally each $g$ gives an etale intersection cover of degree
$18$. The unordered nodes are quotients of the union of these
covers, with the constant identification of $g$ and $-g$.
Therefore the node permutation image is a subquotient of a
product of $S_{18}$'s, as claimed.
\end{proof}

\begin{lemma}[A natural eight-dimensional block after an isogeny]
\label{lem:ordinary-product-isogeny-GL8-block}
\status{Proved under stated hypotheses}
Suppose the degree-one relative image on $E_1'\times E_2'$
is $\GL_8(\mathbf F_2)$, and let $\phi_i$ be as above.
For every
\[
                    d\ge \max(m_1,m_2)+1,
\]
the degree-$d$ geometric group on $E_1\times E_2$ contains
an actual natural $\GL_8(\mathbf F_2)$ block fixing a
complementary summand pointwise.
\end{lemma}

\begin{proof}
Choose a fixed reduced vertical divisor of degree $3(d-m_2)$
linearly equivalent to $3(d-m_2)O$ on $E_1$, and a fixed
reduced horizontal divisor of degree $3(d-m_1)$ similarly.
These degrees are at least three, so such divisors can be
chosen in general position. Let $D$ be the resulting grid
of elliptic fibres on $X$.
It is ordinary and nodal. By the preceding line-bundle formula,
$Y=\phi(C')\cup D$ is a member of $|L^d|$.
For $C'$ in a dense open, it meets $D$ transversely, avoids its
nodes, and its own image nodes avoid $D$.
The last avoidance is also an elementary incidence condition:
for a fixed image fibre, the two value conditions at $P,P+g$
have $P$ varying in dimension one, and so have proper image
in the coefficient projective space.
All these $Y$ are ordinary nodal curves.

On this boundary family the relative Lang representation has
a filtration
\[
       0\subset T\subset T_0\subset V_d,
 \qquad T_0/T\text{ constant},\qquad
 V_d/T_0=Q,\quad\dim Q=8,
\]
where $T$ is the graph contribution and $Q$ is the degree-one
relative system of $C'$.
Indeed odd-degree isogenies identify the ambient
$H^1(X,\mathbf F_2)$ with $H^1(X',\mathbf F_2)$, and the
normalization exact sequence gives this filtration.
The image on $Q$ remains full after shrinking the parameter
space.

The image nodes of $C'$ give degree-$18$ permutation covers.
An intersection with a fixed image fibre splits upstairs into
fixed fibres of the corresponding odd isogeny; on each such
fibre, $C'$ cuts a divisor of degree three.
Their permutation images are therefore subgroups of products
of $S_3$. The nodes of $D$ itself are constant.
It follows that the image on $T$ has order with prime factors
at most $17$. The whole $T_0$ image is an extension of this
by a subgroup of the additive two-group
$\Hom(T_0/T,T)$, so the same prime bound holds.

The simple group $\GL_8(\mathbf F_2)$ has order divisible by
$31$ (the factor $2^5-1$ in its order).
Goursat therefore makes the kernel of the entire $T_0$ action
still surject onto $\GL(Q)$.
Inside its block parabolic, the translation kernel is
$Q^*\otimes T'$ for a subspace $T'\subset T_0$.
The elementary $H^1(\GL_8(\mathbf F_2),Q^*)=0$ calculation
of Lemma~\ref{lem:char2-group-inputs} splits the remaining
cocycle after conjugation, exactly as in
Lemma~\ref{lem:product-degree-three-block}.
Thus this subgroup contains the stated block-diagonal
$\GL_8$, not merely a quotient isomorphic to it.
The boundary lies in the full coherent ordinary open; its
image is a subgroup of the same monodromy image as the dense
smooth ordinary open.
\end{proof}

\begin{theorem}[Eventual full images for every ordinary pair]
\label{thm:ordinary-products-eventual-full}
\status{Proved under stated hypotheses}
Let $E_1,E_2$ be any ordinary elliptic curves over a perfect
field of characteristic two, with
$L=\mathcal O_{E_1}(3O)\boxtimes\mathcal O_{E_2}(3O)$.
There exist odd-degree isogenies
$\phi_i:E_i'\to E_i$, of degrees $m_i$, such that for every
$d\ge d_1=\max(m_1,m_2)+1$ the entire geometric and arithmetic
images are
\[
 \begin{cases}
 \GL_{9d^2-1}(\mathbf F_2),&d\text{ odd},\\
 \mathbf F_2^{3(9d^2-4)}\rtimes\GL_{9d^2-4}(\mathbf F_2),
       &d\text{ even},
 \end{cases}
\]
where the second line is the full pointwise stabilizer of the
three independent invariant covectors.
The isogenies, and hence the displayed integer $d_1$, can be
found by a terminating finite procedure over a finite field or
an effectively presented finitely generated field of definition.
No degree-one full-image assertion for an arbitrary pair is
being inferred from the generic pair.
\end{theorem}

\begin{proof}
By \cite[Theorem~2, printed p.~477]{Chai1995}, for any prime
$\ell\ne2$ the $\ell$-power Hecke orbit of an ordinary elliptic
curve is Zariski dense in its one-dimensional moduli space.
We use only the case $g=1$, with $\ell=3$.
The product of the two such dense orbits meets the nonempty
open $B^\circ$ of Lemma~\ref{lem:ordinary-pair-GL8-open}.
Choose a pair there; reversing an isogeny by its dual gives
isogenies $\phi_i:E_i'\to E_i$ of odd degrees $m_i$.
Geometric base extensions needed to define them do not change
the geometric image.

For $d\ge d_1$, Corollary~\ref{cor:ordinary-products-irreducible-kernel}
gives an irreducible common kernel $K$, no invariant vectors,
and transvection generation.
Lemma~\ref{lem:ordinary-product-isogeny-GL8-block} provides an
actual natural $\GL_8$ summand inside $K$, fixing its complement.
This excludes a nondegenerate invariant alternating form on
$K$: the natural $\GL_8$ has no invariant bilinear form and
no pairing with a trivial module, so this summand would lie
in the radical. The opened transvection classification
\cite[Theorem~3.6, printed p.~22]{KantorSeress2001} therefore
forces the full $\GL(K)$.
The translation-kernel and $H^1$ argument already used in
Theorem~\ref{thm:product-full-group-if-hasse-irreducible} gives
all translations when the fixed quotient has dimension three.
This proves the geometric images.

The three even-degree covectors are defined over the original
perfect field. Indeed the root line bundles
$\mathcal O_{E_i}(3dO/2)$ are defined there, and the ordinary
characteristic-two splitting $F_*\mathcal O_{E_i}=\mathcal O\oplus T_i$
is defined there: its nontrivial two-torsion character line is
unique. Normalize the invariant differential to $dx/x$ in the
ordinary Weierstrass form, for which Cartier is the identity.
The resulting four product projectors and their three quotient
covectors are then arithmetic invariants. The arithmetic image
is between the proved geometric group and its same pointwise
stabilizer, so the two images coincide.

For the effective assertion, enumerate finite chains of
three-isogenies from each input elliptic curve. This is a finite
branching enumeration: at each step solve the three-division
polynomial, enumerate its finite etale cyclic subgroups, and
form their quotients; reverse the resulting chains by duals.
Run the exact Lang-frame geometric-image algorithm of
Sections~\ref{subsec:effective-entire-monodromy} and
\ref{subsec:perfect-field-monodromy-data} on each resulting pair
in degree one, and stop at the first full $\GL_8$ image.
Dovetail the finite computations. Every individual test
terminates, and Hecke density plus the nonempty open proves
that some enumerated pair passes. Its explicitly retained
isogeny degrees give $d_1$. This is a terminating construction,
not an unspecified Galois-group oracle. For an arbitrary
perfect field with no effective coefficient presentation the
same argument proves existence; the calculation is performed
on a finitely generated field of definition when such a
presentation is supplied.
\end{proof}


\begin{remark}[Finite supporting checks]\label{auto:A_monodromy:3a04f9cdd890}
\status{Proved under stated hypotheses}

The exact program
\path{agents/monodromy/code/ordinary_product_isogeny_incidence_checks.py}
checks all $18424$ triples of distinct rank-one $3\times3$
matrices over $\mathbf F_2$, including all $98$ collinear
triples, against the common-factor criterion. It checks that
each of the three nonzero parallel-gradient relations admits
a transverse pair of gradients, and checks the prime $31$
separation from degree-$18$ node covers. The output is
\path{agents/monodromy/outputs/ordinary_product_isogeny_incidence_checks.json}.
These are finite supporting calculations; the universal
incidence assertions are proved in the preceding lemmas.
\end{remark}


\subsection{An exact parity formula for supersingular elliptic products}
\label{sec:DE-A3-ss-product}

The following family gives an explicit answer in every degree, and shows that
eventual relative ordinarity can fail even when the ambient Picard scheme is
smooth and its crystalline cohomology is torsion-free.  It does not purport to
determine the generic rank for an arbitrary fixed ambient variety.

Throughout this subsection let $k=\overline{\mathbf F}_2$, let
\[
 E_0:\quad y^2+y=x^3,
 \qquad X=E_0\times E_0,
 \qquad L=\OO_{E_0}(3O)\boxtimes\OO_{E_0}(3O),
\]
and embed $X\hookrightarrow\mathbf P^8_k$ by the product of the plane cubic
embeddings followed by the Segre embedding.  Write
$N_m=\OO_{E_0}(mO)\boxtimes\OO_{E_0}(mO)$, so $L^d=N_{3d}$.
For a smooth member $Y\in|N_m|$, put
\[
 Q_Y=\coker\bigl(H^1(X,\OO_X)\longrightarrow H^1(Y,\OO_Y)\bigr).
\]
Frobenius on these vector spaces is $a\mapsto a^2$-semilinear; ``stable rank''
means the dimension of the eventual image of its iterates over the perfect
field $k$.

\begin{lemma}[Explicit endomorphisms of every positive degree]
\label{lem:DE-A3-ss-norm}
\status{Proved}
For every integer $n\ge1$ there is an explicitly specified endomorphism
$b:E_0\to E_0$ with
$b^\dagger b=bb^\dagger=[n]$, where $\dagger$ is the Rosati involution for
$\OO_{E_0}(O)$.  In particular $\deg b=n$.
\end{lemma}
\begin{proof}
Choose $\omega\in\mathbf F_4$ with $\omega^2+\omega+1=0$, and define
\[
 i(x,y)=(x+1,y+x+\omega),\qquad
 j(x,y)=(x+\omega,y+\omega^2x+\omega).
\]
For $a\in\mathbf F_4^\times$ and $c^2+c=1$, substitution shows that
$(x,y)\mapsto(x+a,y+a^2x+c)$ preserves the equation, because $a^4=a$ and
$a^3=1$.  These maps extend to automorphisms fixing $O$.  Direct composition
gives $i^2=j^2=[-1]$ and $ij=-ji$, with $[-1](x,y)=(x,y+1)$.
Since they fix the polarizing divisor $O$, their Rosati adjoints are their
inverses.  Consequently
\[
 (a+bi+cj+dij)^\dagger(a+bi+cj+dij)
 =[a^2+b^2+c^2+d^2]
\]
for integers $a,b,c,d$; the same identity holds in the opposite order.

We recall a short proof that every positive integer is a sum of four squares,
so the construction does not depend on an unproved norm-surjectivity claim.
The quaternion norm identity makes products of such sums sums of four
squares.  For an odd prime $\ell$, the two sets of residues $a^2$ and
$-1-b^2$ each have $(\ell+1)/2$ elements.  They intersect.  Choosing
$|a|,|b|\le(\ell-1)/2$ gives
$a^2+b^2+1=k\ell$ with $0<k<\ell$.
Choose the smallest positive multiplier $k$ for which $k\ell$ is a sum of
four squares, and suppose $k>1$.  Write $k\ell=\sum a_i^2$ and choose
$x_i\equiv a_i\pmod k$ with $|x_i|\le k/2$.  Then
$\sum x_i^2=kt$ with $0\le t\le k$.  The case $t=0$ would imply
$k\mid\ell$, which is impossible.  If $t<k$, all four coefficients of the
quaternion product $a\overline x$ are divisible by $k$, since modulo $k$
this product is $x\overline x=kt$.  Dividing those coefficients by $k$
expresses $t\ell$ as a sum of four squares, contradicting minimality.
If $t=k$, equality in the four bounds forces $k$ even and every
$a_i\equiv k/2\pmod k$.  Hence $k\mid4\ell$, so $k=2$ or $4$.
For $k=2$, the four odd squares have sum $4\pmod8$, whereas $2\ell$ is
$2$ or $6\pmod8$.  For $k=4$, the four squares are each $4\pmod{16}$,
whereas $4\ell$ is $4$ or $12\pmod{16}$.  Both are impossible.  Thus
$k=1$.  The prime $2$ and the integer $1$ are immediate, and multiplicativity
finishes the proof.  To specify $b$ without choices, take the first
nonnegative quadruple in lexicographic order with sum of squares $n$ and
use the displayed endomorphism.  Finally
$(\deg b)^2=\deg[n]=n^2$, so $\deg b=n$.
\end{proof}

\begin{lemma}[The ambient quotient and the even-degree obstruction]
\label{lem:DE-A3-ss-prym}
\status{Proved}
For every smooth $Y\in|N_m|$, $m\ge3$, its genus is $g=m^2+1$ and
$\dim Q_Y=m^2-1$.  Its $2$-rank $f_Y$ satisfies
\[
 f_Y\le m^2-1,
 \qquad
 f_Y\le m^2-3\quad\hbox{if }2\mid m.
\]
Moreover $\dim(Q_Y)_{\ss}=f_Y$.
\end{lemma}
\begin{proof}
The plane cubic is smooth: on the affine chart the partial derivative with
respect to $y$ is $1$, and at its unique point at infinity the homogeneous
equation has a nonzero partial derivative.  Its inverse map is
$(x,y)\mapsto(x,y+1)$, so $E_0[2](k)=0$.  Thus its $2$-rank is zero and
Frobenius on $H^1(E_0,\OO_{E_0})$ is zero.  The latter also follows from
Cartier: $\omega_0=dx$ is a regular invariant differential and
$C(dx)=0$.  Hence Frobenius is zero on $H^1(X,\OO_X)$, a space of
dimension two.

K\"unneth gives $H^1(X,P\otimes N_m^{-1})=0$ for every
$P\in\Pic^0(X)$: on each elliptic factor the two line bundles in question
have negative degree, so both possible $H^1$ tensor summands have a zero
$H^0$ factor.  In particular $H^0(Y,\OO_Y)=k$ and
$H^1(X,\OO_X)\hookrightarrow H^1(Y,\OO_Y)$.  Adjunction, or the same
exact sequence and Serre duality, gives $g=m^2+1$.

Restriction
\[
 j:\Pic^0(X)=E_0^2\longrightarrow J_Y
\]
is a closed immersion of group schemes.  Indeed a nontrivial
$P\in\Pic^0(X)(k)$ has $H^0(X,P)=0$.  The preceding vanishing and the
restriction sequence then give $H^0(Y,P|_Y)=0$, so $P|_Y$ is not trivial.
The tangent map of $j$ is the already proved injection on $H^1(\OO)$.
Its proper kernel therefore has only the identity as a geometric point and
has zero tangent space.  It is a finite local group scheme with zero
cotangent space, and Nakayama's lemma makes it trivial.

Identify $X$ with its dual by the product principal polarization.  The map
dual to $j$ is
\[
 q:J_Y\longrightarrow X,\qquad
 M\longmapsto
 \bigl(\operatorname{Nm}_{\pi_1}M,
       \operatorname{Nm}_{\pi_2}M\bigr),
\]
where $\pi_i:Y\to E_0$ are the projections.  They are finite of degree
$m$: $Y$ is smooth and connected, hence integral, and cannot contain an
elliptic fibre.  The diagonal entries of $qj$ are $[m]$ by norm and
pullback.  Each off-diagonal entry is zero.  For example the intersection
divisor on the first $E_0$ obtained by fixing a point on the second factor
always belongs to $|\OO_{E_0}(mO)|$, so the norms of the pullbacks of two
point divisors on the second factor are linearly equivalent.  We obtain
\[
 qj=[m]_X.                                                     \tag{\ast}
\]
In particular $q$ is surjective.

For completeness its kernel is connected and smooth, not merely connected
up to isogeny.  Let $P=(\ker q)^0_{\mathrm{red}}$.  Over the perfect field
$k$ this is an abelian subvariety.  The quotient $J_Y/P$ is an abelian
variety and $q$ factors through an isogeny $\nu:J_Y/P\to X$.
Dualizing, the kernel of $\nu^t$ lies in the kernel of $q^t=j$, which is
trivial.  The isogeny duality theorem gives
$\ker(\nu^t)\simeq(\ker\nu)^D$, so $\nu$ is an isomorphism.
Thus $\ker q=P$.  Here the exact source statements are the quotient
existence theorem and its abelian-variety example, and the isogeny duality
theorem \cite[Theorem~4.38, Example~4.40, Theorem~7.5]{EdixhovenVanDerGeerMoonenAV}.

The $2$-rank is additive in an exact sequence of abelian varieties.
One can see this without choosing an isogeny complement: on $k$-points,
the sequence of $2^a$-torsion groups is exact, since a lift of a torsion
point can be adjusted by a point of $P(k)$ using the surjectivity of
$[2^a]:P(k)\to P(k)$.  Counting points gives additivity.
Since $f_X=0$, we have $f_Y=f_P$ and $\dim P=g-2$.  This proves the first
bound.

If $m$ is even, (\ast) restricts to a closed immersion of finite group
schemes $j_2:X[2]\hookrightarrow P[2]$.  Its Frobenius twist maps
$X^{(2)}[2]$ into $P^{(2)}[2]$.  The group $\ker V_X$ lies in
$X^{(2)}[2]$, and finite-group Verschiebung is functorial for $j_2$
(or use $V_{J_Y}j^{(2)}=jV_X$ and the restriction of $V_{J_Y}$ to
$P^{(2)}$).  Consequently it embeds
\[
 \ker(V_X:X^{(2)}\to X)\lhook\joinrel\longrightarrow
 \ker(V_P:P^{(2)}\to P).
\]
The source is connected of order $4$, since $\deg V_X=2^2$ and $f_X=0$.
For any abelian variety $A$ of dimension $t$, $\deg V_A=2^t$ and
$\#\ker V_A(k)=2^{f_A}$: relative Frobenius is bijective on $k$-points
and $V_AF_A=[2]$.  The identity component of $\ker V_A$ therefore has
order $2^{t-f_A}$.  The displayed injection yields
$\dim P-f_P\ge2$, hence $f_Y\le g-4=m^2-3$.
These degree and Frobenius identities are exactly
\cite[Propositions~5.19--5.22]{EdixhovenVanDerGeerMoonenAV}; the argument
uses the connected subgroup, so it is not just a bound on geometric
$2$-torsion points.

Finally Artin--Schreier and the semilinear Fitting decomposition identify
the stable rank of $H^1(Y,\OO_Y)$ with $f_Y$.  The stable-part functor on
finite-dimensional semilinear vector spaces over a perfect field is exact:
split each space into its nilpotent and bijective parts, or apply a
sufficiently high Frobenius iterate to an exact sequence.  Quotienting by
the Frobenius-zero space $H^1(X,\OO_X)$ therefore leaves the stable rank
unchanged.  This proves the last assertion.
\end{proof}

\begin{theorem}[Exact generic rank in every hypersurface degree]
\label{thm:DE-A3-ss-product-rank}
\status{Proved}
For the specified $X\subset\mathbf P^8_{\mathbf F_2}$, the geometric
generic stable rank in degree $d\ge1$ is
\[
 \boxed{\quad
 \dim E(Y,\OO_Y)_{\ss,\mathrm{gen}}
 =\begin{cases}
 9d^2-1,&d\ \text{odd},\\
 9d^2-3,&d\ \text{even}.
 \end{cases}\quad}                                           \tag{\dagger}
\]
The coherent relative space has dimension $9d^2-1$ in every degree.
Consequently its generic Frobenius is bijective exactly for odd $d$.
There is no integer $d_0$ making it generically bijective for every
$d\ge d_0$.  The generic stable rank is positive in every degree, so
for this example the positivity threshold is exactly $d_0=1$.
\end{theorem}
\begin{proof}
First let $m\ge3$ be odd.  Choose the endomorphism $b$ in
Lemma~\ref{lem:DE-A3-ss-norm} with $b^\dagger b=bb^\dagger=[m-1]$ and put
\[
 \phi=\begin{pmatrix}1&b\\-b^\dagger&1\end{pmatrix}:X\longrightarrow X.
\]
Its Rosati adjoint satisfies $\phi^\dagger\phi=[m]I$; hence $\phi$ is an
isogeny of degree $m^2$ with kernel killed by the odd integer $m$.
It is therefore \'etale.  The inverse images of the two coordinate
elliptic divisors are
\[
 C_1=\{x+by=0\},\qquad C_2=\{-b^\dagger x+y=0\}.
\]
Each is a graph isomorphic to $E_0$.  They meet transversely in
$\ker\phi$, which consists of exactly $m^2$ reduced points.

We need linear equivalence with the fixed $N_m$.  The polarization of
$\OO(C_1+C_2)=\phi^*(\OO_{E_0}(O)\boxtimes\OO_{E_0}(O))$ is $[m]I$,
by the pullback formula for polarizations
\cite[Proposition~7.6]{EdixhovenVanDerGeerMoonenAV}.
This is also the polarization of $N_m$, so their difference belongs to
$\Pic^0(X)$ by \cite[Corollary~7.22]{EdixhovenVanDerGeerMoonenAV}.  Both line bundles are symmetric; on $\Pic^0(X)$ inversion
pulls a line bundle back to its inverse.  Their difference is thus a
$2$-torsion $k$-point of $\Pic^0(X)$.  There are no such nontrivial points,
since $E_0[2](k)=0$.  Hence $C_1+C_2\in|N_m|$ exactly.

For a connected nodal curve $D$ with smooth components $D_i$, the
normalization sequence gives a Frobenius-equivariant exact sequence
\[
 0\longrightarrow H^1(\Gamma_D,k)
 \longrightarrow H^1(D,\OO_D)
 \longrightarrow\bigoplus_iH^1(D_i,\OO_{D_i})\longrightarrow0,
\]
where $\Gamma_D$ is its dual graph and Frobenius on the first term is
bijective.  Thus its stable rank is
$b_1(\Gamma_D)+\sum_i f_{D_i}$, and its genus minus stable rank is
$\sum_i(g(D_i)-f_{D_i})$.  Applied to $C_1\cup C_2$, this gives genus
$m^2+1$ and stable rank $m^2-1$.

The universal Cartier divisor over the projective space $|N_m|$ is proper
and flat.  Every nonzero section has connected divisor and
$h^1(\OO_D)=m^2+1$: this follows from
$H^1(N_m^{-1})=0$ and the surjection
$H^2(N_m^{-1})\to H^2(\OO_X)$ dual to the nonzero section.
Cohomology and base change therefore make $H^1(\OO_D)$ a vector bundle
on the parameter space.  Its relative Frobenius is a bundle map from the
Frobenius pullback.  For vector spaces of dimension $g=m^2+1$, the stable
rank is the rank of the $g$th iterate.  Rank at least a fixed integer is
an open condition.  The nodal member just constructed consequently
supplies a nonempty open locus of rank at least $m^2-1$.
Since $N_m$ is very ample for $m\ge3$, the smooth-member locus is also
nonempty open by Bertini.  These opens meet in the irreducible projective
space $|N_m|$.  Lemma~\ref{lem:DE-A3-ss-prym} proves that the generic
smooth member has rank exactly $m^2-1$.

Now let $m=3d$ be even, so $m\ge6$.  Write $m=3+(m-3)$, with both
summands odd and at least three.  Choose smooth members
$C\in|N_3|$ and $D\in|N_{m-3}|$ in the just proved maximal-rank opens.
They may be chosen transverse: both systems are very ample, so a general
member of the second system meets a fixed smooth member of the first
transversely.  Their union is a connected nodal member of $|N_m|$.
Each component has genus minus stable rank two.  The normalization
formula therefore gives
\[
 g(C\cup D)-f_{C\cup D}=4,\qquad f_{C\cup D}=m^2-3.
\]
The same open-rank and smoothness argument gives generic smooth rank at
least $m^2-3$.  The even-degree upper bound of
Lemma~\ref{lem:DE-A3-ss-prym} makes this an equality.

Finally these are precisely the complete linear systems cut out by
ambient hypersurfaces of degree $d$.  For the plane cubic,
$H^0(\mathbf P^2,\OO(d))\to H^0(E_0,\OO(3dO))$ is onto, by
$H^1(\mathbf P^2,\OO(d-3))=0$.  The Segre coordinate ring has all its
bihomogeneous pieces of bidegree $(d,d)$ generated by degree-$d$
polynomials in the pairwise products of coordinates.  Consequently
\[
 H^0(\mathbf P^8,\OO(d))\twoheadrightarrow H^0(X,L^d).
\]
The nonempty open rank and smoothness loci pull back to nonempty open
loci in the ambient hypersurface parameter space.  Taking $m=3d$ gives
(\dagger).  The entire construction is defined over an algebraic
extension of $\mathbf F_2$ and computes the geometric generic ranks for
the displayed variety over $\mathbf F_2$ itself.
\end{proof}

\begin{corollary}[An ambient Cartier-nilpotent dimension is insufficient]
\label{cor:DE-A3-ss-product-defect}
\status{Proved}
For this abelian surface, $\Pic^0(X)$ is smooth and
$H^*_{\cris}(X/W(k))$ is torsion-free.  All its global one-forms are
Cartier-killed, so the space of global forms eventually killed by Cartier
has dimension two.  Nevertheless the generic curve has
$g(Y)-f_Y=2$ in odd degree and $g(Y)-f_Y=4$ in even degree.
Thus neither eventual relative ordinarity nor a formula equating the
generic full defect with this ambient Cartier-nilpotent dimension holds
for every sufficiently large degree.
\end{corollary}
\begin{proof}
The Picard variety is the smooth abelian surface $E_0^2$.  Crystalline
cohomology of a product of elliptic curves is the exterior algebra on
their free first crystalline cohomology, hence is torsion-free
\cite[II, (7.1.1), p.~651]{Illusie1979DRW}.
The two pullbacks of $dx$ span the global one-forms and have Cartier zero.
The defect assertion is Theorem~\ref{thm:DE-A3-ss-product-rank} and
$g(Y)=9d^2+1$.
\end{proof}

\begin{remark}[Computations and exact scope]
\label{rem:DE-A3-ss-product-computation}
\status{Proved under stated hypotheses}
The exact finite-field probe
\texttt{A3\_supersingular\_product\_probe.py} uses the basis
$\{x^a:2a\le m\}\cup\{x^ay:2a+3\le m\}$ of
$H^0(E_0,\OO(mO))$ and the Cartier identities
\[
 C(x^{2a}dx)=0,\quad C(x^{2a+1}dx)=x^a dx,\quad
 C(x^{2a}y\,dx)=x^{a+1}dx,\quad
 C(x^{2a+1}y\,dx)=x^ay\,dx.
\]
Its recorded best stable matrix ranks for $d=1,\ldots,6$ are
$8,33,80,141,224,321$; elapsed time is
$29.180787964000046$ seconds.  It retains the exact section coefficients,
matrix columns, rank sequences and seeds.  These finite probes motivated
the parity formula; the nodal deformation and Prym arguments prove it
uniformly, without inferring either smoothness or an all-degree theorem
from random samples.  The additional certificate
\texttt{A3\_supersingular\_product\_checks.py} checks the explicit
automorphisms, both symbolic quaternion norm identities, four-square
witnesses for $1\le n\le4096$, the degeneration arithmetic for
$1\le d\le10000$, and all six retained matrix witnesses.  It passed
in $0.027462210000521736$ seconds; its complete code and output are retained.
The generic-rank problem for arbitrary $X$ in A.3 remains a separate
question.  No priority claim for this explicit formula is made here.
\end{remark}


% /root/monodromy, 2026-11-19. Proved; independent root and reduction_descent audits passed.
\subsection{A supersingular product: a double Hasse component and a full image}
\label{subsec:supersingular-product-degree-one}
\status{Proved and checked. The parent and reduction/descent agent independently
audited the complete determinant, local wild-inertia and full-group proofs.
Both independently read and reran the polynomial/Pfaffian certificate and
the separate finite-field group verifier; all checks passed.}

Let
\[
 E:\ y^2+y=x^3\quad\text{over }\mathbf F_2,
 \qquad X=E\times E,\qquad
 L=\mathcal O_E(3O)\boxtimes\mathcal O_E(3O).
\]
The indicated Segre embedding identifies ambient hyperplane sections with
the complete system $P=\mathbf P(W)$, where $W=H^0(X,L)$ has dimension nine.
Write $s=(s_{ij})_{0\le i,j\le2}$, in lexicographic order, in the tensor
basis of $(1,x,y)$, and write $S=(s_{ij})$ for the corresponding
$3\times3$ matrix. Put $J_3=(\delta_{i,2-j})$ and $J=J_3\otimes J_3$.

\begin{proposition}[An explicit nonreduced relative Hasse divisor]
\label{prop:ss-product-relative-Hasse}
\status{Proved under stated hypotheses}
The relative coherent Frobenius bundle is $\mathcal E=\Omega^1_P$.
Its determinant section, up to the constant trivialization implicit in
the chosen basis, is
\begin{equation}\label{eq:ss-product-Hasse-factorization}
 h(s)=D(s)^2Q(s),\qquad
 D(s)=\det S,\qquad Q(s)=s^{[2]t}Js.
\end{equation}
The cubic $D$ is geometrically irreducible; the cubic $Q$ is
geometrically smooth and irreducible; they are distinct. Thus the
relative Hasse divisor has a component of multiplicity two and a
different component of multiplicity one. The generic relative stable
rank is eight. A smooth generic section has genus ten and $2$-rank eight.
\end{proposition}

\begin{proof}
Use $\eta=dx$ on $E$. Writing a rational function as $a^2+xb^2$
gives $C((a^2+xb^2)\,dx)=b\,dx$. The curve equation gives
\[
 C(1)=0,\quad C(x)=1,\quad C(y)=x,\quad
 C(x^2)=C(y^2)=0,\quad C(xy)=y.
\]
Consequently the root-removed Cartier multiplication matrix on the
degree-three one-factor space is
\[
 M_E(u)=\begin{pmatrix}u_1&u_0&0\\u_2&0&u_0\\0&u_2&u_1\end{pmatrix},
 \qquad
 J_3M_E(u)=\begin{pmatrix}0&u_2&u_1\\u_2&0&u_0\\u_1&u_0&0\end{pmatrix}.
\]
If $M_i$ is the coefficient of $u_i$, then the full tensor matrix is
$M(s)=\sum s_{ij}M_i\otimes M_j$. Direct differentiation of the
six-term determinant shows
\begin{equation}\label{eq:ss-Hessian}
 A(s):=JM(s)=\operatorname{Hess}_s(D),\qquad
 A(s)^t=A(s),\quad A(s)_{ii}=0,\quad M(s)s=0.
\end{equation}
The eight-dimensional Pfaffians of the alternating $9\times9$
matrix $A(s)$ form the vector
\[
       \operatorname{Pf}_8(A(s))=D(s)s.
\]
This identity may be verified by expansion of the nine Pfaffians.
The retained exact polynomial program checks all nine over
$\mathbf F_2[s_0,\ldots,s_8]$, without specialization. The standard
Pfaffian adjugate identity, or the same expansion of cofactors, gives
\[
 \operatorname{adj}M(s)=D(s)^2s\,(Js)^t,
 \qquad
 \operatorname{adj}(M(s)^t)=D(s)^2(Js)s^t.
\]

Kunneth gives $H^1(X,L^{-1})=0$, and the hypersurface exact sequence
gives the Euler sequence
\[
 0\longrightarrow\mathcal E\longrightarrow
 W^*\otimes\mathcal O_P(-1)\longrightarrow\mathcal O_P\longrightarrow0.
\]
The quotient Frobenius is zero because $X$ is supersingular.
In affine coefficient coordinates the relative Frobenius is therefore
the restriction of $M(s)^t$ from $(s^{[2]})^\perp$ to $s^\perp$.
The cofactor formula identifies the determinant of this restriction
with $D(s)^2(s^{[2]})^tJs$: indeed the exterior-product vector of a
basis of the source hyperplane is its normal $s^{[2]}$, and the
cofactor matrix carries it to
$D(s)^2s\,(Js)^ts^{[2]}$, a scalar multiple of the target normal $s$.
The projective twists agree: $\det\mathcal E=\mathcal O_P(-9)$, so
the Frobenius determinant is a section of $\mathcal O_P(9)$.
The displayed homogeneous polynomial has degree nine, proving
\eqref{eq:ss-product-Hasse-factorization} globally.

The rank-at-most-two matrix locus is the image closure of the
irreducible space of products of a $3\times2$ and a $2\times3$
matrix, and is the determinant hypersurface. Thus $D$ is
geometrically irreducible. Moreover
\[
                  \nabla Q(s)=Js^{[2]},
\]
which never vanishes at a projective point. Thus $Q$ is smooth;
two distinct positive-degree projective hypersurfaces in $\mathbf P^8$
intersect, so smoothness also implies geometric irreducibility.
At $S=\operatorname{diag}(0,1,1)$ one has $D=0,Q=1$, proving that
the two cubics are distinct. At $S=I_3$ both $D$ and $Q$ are one.
Thus the relative Frobenius is generically invertible.
The smooth open is nonempty by the very ample embedding, so it meets
this invertibility open. Adjunction gives $g=1+L^2/2=10$.
The ambient $H^1(X,\mathcal O_X)$ has nilpotent Frobenius and the
relative stable rank is eight, so the smooth generic curve has
$2$-rank eight.
\end{proof}

\begin{lemma}[A first-derivative criterion for wild transvection inertia]
\label{lem:ss-simple-pole-transvection}
\status{Proved under stated hypotheses}
Let $R=k[[t]]$ with $k$ algebraically closed of characteristic two,
and let a rank-$r$ vector bundle over $R$ have regular semilinear
Frobenius $F$. Suppose that its special fibre has stable rank $r-1$,
its linearized kernel has dimension one, and its determinant has
valuation one. Let $u$ generate the semilinear kernel in the special
fibre, and let $v$ be the first derivative of $F(\widetilde u)$ at
$t=0$, for a regular lift $\widetilde u$. If $v\notin ku$, the
generic Lang local system has a nontrivial transvection in its inertia.
\end{lemma}

\begin{proof}
The finite part of the Lang scheme supplies $r-1$ regular fixed
sections, whose special values span the stable subspace. This is the
finite-part lemma already proved using Stacks~\cite{Stacks0BTN,Stacks09ZS}.
They form a Frobenius-fixed direct summand. Complete them by a lift of
$u$. In this frame
\[
 F=\begin{pmatrix}I_{r-1}&b(t)\\0&a(t)\end{pmatrix},
 \qquad a(t)=a_1t+O(t^2),\qquad b(t)=b_1t+O(t^2),\qquad a_1\ne0.
\]
The constant column is zero because $u$ is in the special kernel.
Changing the lift changes its Frobenius image only in order at least
two, so its first derivative is well defined in the special fibre.
The hypothesis $v\notin ku$ is precisely $b_1\ne0$.

The nonzero fixed vector in the one-dimensional generic quotient has
coordinate $z=a(t)^{-1}$. Lifting it requires solving
\[
                  y_i^2+y_i=b_i(t)/a(t)^2.
\]
At least one right side has a nonzero simple pole. Such an
Artin--Schreier equation is ramified: if $y$ lay in an unramified
extension, its negative valuation would give an even pole order
for $y^2+y$, a contradiction; a nonnegative valuation gives no pole.
Over the strictly henselian field all inertia fixes the regular
$(r-1)$-dimensional Lang subspace and its one-dimensional quotient.
A nonidentity element therefore differs from the identity by a rank-one
square-zero map, which is a transvection.
\end{proof}

\begin{theorem}[Full degree-one image on a supersingular product]
\label{thm:ss-product-degree-one-full-group}
\status{Proved under stated hypotheses}
For $X,L$ above, over every perfect field of characteristic two,
the degree-one relative Hasse--Witt local system on its maximal-rank
smooth open has
\[
                   G_{\rm geom}=G_{\rm arith}=\GL_8(\mathbf F_2).
\]
In particular its geometric representation is irreducible, despite
the double component of its coherent Hasse divisor.
\end{theorem}

\begin{proof}
We give two small exact certificates over
$k_8=\mathbf F_2[\alpha]/(\alpha^3+\alpha+1)$.
For readability an integer $c_0+2c_1+4c_2$ denotes the field element
$c_0+c_1\alpha+c_2\alpha^2$; all coefficient lists below use this
encoding and the fixed lexicographic tensor basis.

First take
\[
                s=(6,7,0,5,4,1,6,0,4).
\]
The relative Frobenius is invertible and the characteristic polynomial
of its third iterate, up to replacing it by its reciprocal according
to the chosen Frobenius convention, is
\begin{equation}\label{eq:ss-degree-one-Frobenius-polynomial}
                    T^8+T^7+T^5+T+1.
\end{equation}
The retained verifier checks its irreducibility over $\mathbf F_2$
by the Frobenius-power and gcd criterion and checks that it is not
reciprocal. Thus the arithmetic image contains the element required
in Lemma~\ref{lem:product-frobenius-group-test}.
The verifier uses an independent restriction-of-scalars calculation:
if $B=A A^{[2]}A^{[4]}$ is the full third Frobenius iterate, then
$B^{[2]}=A^{[2]}A^{[4]}A$ has the same characteristic polynomial,
because cyclic products of square matrices do. Thus that polynomial
has coefficients in $\mathbf F_2$. Restriction from $\mathbf F_8$ to
$\mathbf F_2$ cubes it. The directly computed $27\times27$ integer
characteristic polynomial, reduced modulo two, is
$(T(T^8+T^7+T^5+T+1))^3$; unique factorization, with monic factors,
therefore proves the asserted nine-dimensional polynomial and its
eight-dimensional relative factor.

For a geometric transvection, take
\[
 s_0=(0,6,4,5,5,5,3,3,4),\qquad\delta=(1,0,\ldots,0),
 \qquad s(t)=s_0+t\delta.
\]
Here $D(s_0)\ne0$, $Q(s_0)=0$, and $dQ_{s_0}(\delta)=6\ne0$.
The full linearized matrix $M(s_0)^t$ has rank eight; its semilinear
stable rank is seven. A generator of its semilinear kernel is
\[
 u=(2,7,7,3,3,3,2,4,0)=(Js_0)^{[1/2]}.
\]
The verifier checks $s_0^tu=0$, so this is also the relative kernel.
The derivative on this kernel along the displayed affine line is
\[
 v=M(\delta)^tu^{[2]}=(0,0,0,0,4,3,0,5,5).
\]
These two vectors are independent. The determinant factorization
shows that the relative determinant has valuation one along this
line at $t=0$.

For completeness, this ambient calculation is also the derivative
in a regular frame of the relative bundle. A lift of $u$ in
$s(t)^\perp$ has square $u^{[2]}+O(t^2)$; hence its image derivative
is exactly $M(\delta)^tu^{[2]}$. A change of output frame has no
first-order effect because its image at $t=0$ is zero. Passing
from the affine coefficient cone to a projective chart only subtracts
a radial tangent vector; that adds a multiple of
$M(s_0)^tu^{[2]}=0$ and changes neither derivative test.
Lemma~\ref{lem:ss-simple-pole-transvection} therefore supplies a
nontrivial geometric transvection.

Let $H$ be its normal closure in the arithmetic image $G$.
Then $H$ is nontrivial, generated by transvections, normal in $G$,
and lies in the geometric image. Applying
Lemma~\ref{lem:product-frobenius-group-test} with $N=8$ proves
$H=G=\GL_8(\mathbf F_2)$. That lemma explicitly handles Clifford
blocks using the primitive prime $17$ and excludes all form-preserving
alternatives by nonreciprocity; its only classification input is
\cite[Theorem~3.6, printed p.~22]{KantorSeress2001}.

All calculations were made on the invertibility open of the coherent
bundle over the full projective parameter space. Restriction to the
nonempty smooth open does not change the image: a connected finite
etale cover of a normal integral scheme remains connected on a
dense open. Thus the use of coefficient points or a local line does
not require a separate smoothness assertion for those particular
sections. Finally geometric connectedness of the finite frame cover
is preserved by algebraically closed base extension. The same
geometric full group, and hence the same arithmetic full group, holds
over every perfect base field of characteristic two.
\end{proof}

\begin{remark}[Retained exact calculations]\label{auto:A_monodromy:3502a4b6e144}
\status{Proved under stated hypotheses}

The polynomial certificate is
\path{agents/monodromy/code/supersingular_product_degree_one.py}, with
output \path{agents/monodromy/outputs/supersingular_product_degree_one.json}.
It verifies all nine Pfaffian identities and all $511$ projective
$\mathbf F_2$ coefficient vectors. The extension-field discovery
calculation is
\path{agents/monodromy/code/supersingular_product_frobenius_probe.py};
its output retains the field polynomials, coefficients, characteristic
polynomials and local derivative data. A separate certificate verifier
\path{agents/monodromy/code/verify_supersingular_degree_one_group.py}
checks only the two stated $\mathbf F_8$ points and does not depend
on the discovery search. Its output is
\path{agents/monodromy/outputs/verify_supersingular_degree_one_group.json}.
\end{remark}


% /root/monodromy, 2026-11-19. Proved; independent root, reduction_descent and obstructions_crystals proof audits passed.
\subsection{All invariant covectors on the supersingular product}
\label{subsec:ss-exact-covectors}
\status{Proved and checked. The parent and two other agents independently
audited the complete proof; the local Cartier certificate was independently
rerun by the parent. The persistent coherent nilpotent part is handled by a
bounded Frobenius-root extension and a homogeneous-coefficient argument.}

\begin{theorem}\label{thm:ss-product-exact-covectors}
\status{Proved under stated hypotheses}
Let $E_0:y^2+y=x^3$ over a perfect field of characteristic two, let
$X=E_0\times E_0$ and
$L=\mathcal O_{E_0}(3O)\boxtimes\mathcal O_{E_0}(3O)$.
For every $d\ge1$, the entire geometric and arithmetic invariant
covector spaces of the relative Hasse--Witt local system have dimension
\[
 \dim(V_d^\vee)^{G_{\rm geom}}
 =\dim(V_d^\vee)^{G_{\rm arith}}
 =\begin{cases}0,&d\text{ odd},\\1,&d\text{ even}.
 \end{cases}
\]
In even degree the unique nonzero covector is represented under
adjunction by $D(s)\omega$, where
$\omega=dx_1\wedge dx_2$, $D=\delta_1\delta_2$, and $\delta_i$
is the invariant derivation dual to $dx_i$, lifted to $L^d$ by its
canonical Frobenius-pullback connection. Thus every even-degree image
is contained in the pointwise stabilizer of one nonzero covector.
This theorem determines the entire invariant quotient, not the full group.
\end{theorem}

\begin{proof}
First work over an algebraic closure. Put $m=3d$,
$N=\mathcal O(mO)\boxtimes\mathcal O(mO)$,
$W=H^0(X,N)$ and $b=m^2$. On $P=\mathbf P(W)$ the full coherent
bundle is $\mathcal B=W^*\otimes\mathcal O_P(-1)$, with semilinear
Frobenius matrix $M(s)^t$. Here $M(s)$ is linear in the coefficient
vector $s$ and is the root-removed Cartier multiplication matrix.
Its relative subbundle is $\mathcal E=\Omega_P^1$; the quotient
$\mathcal B/\mathcal E=H^2(X,\mathcal O_X)\otimes\mathcal O_P$
has zero Frobenius. Thus the unit Lang local systems of the full
and relative bundles coincide.

\par\smallskip\noindent\textbf{The generic linearized image is already the stable image.}
The identity $C(\omega)=0$ gives $M(s)s=0$ in every degree.
For odd $m$, Theorem~\ref{thm:DE-A3-ss-product-rank} gives stable
rank $b-1$, so the generic linearized rank is also $b-1$.
For even $m$, $N$ is the relative Frobenius pullback of $N_{m/2}$.
The invariant derivations act on its sections, since they kill the
square transition functions of pullback frames. They commute and
satisfy $\delta_1^2=\delta_2^2=0$. A local $2$-basis calculation gives
\[
 C(s\delta_1(s)\omega)=C(s\delta_2(s)\omega)=0.
\]
Thus $s,\delta_1(s),\delta_2(s)$ lie in $\ker M(s)$.
They are generically independent: for the section $s=x_1x_2$
they are $x_1x_2,x_2,x_1$. Hence the generic linearized rank is
at most $b-3$. The exact stable-rank theorem gives $b-3$, so
equality holds. In both parities the image of the first linearized
Frobenius at the generic point is therefore its unit subspace.

\par\smallskip\noindent\textbf{A purely inseparable lift of an arbitrary invariant covector.}
Let $K=k(P)$ and let $I\subset\mathcal B_K$ be that unit image.
The restriction $F^*I\to I$ is invertible. The finite Lang-frame
equivalence identifies a geometric invariant covector with a
$K$-linear map $f:I\to K$ satisfying
$f(Fx)=f(x)^2$ on $I$. Define a functional on the full generic
bundle, with coefficients in $K^{1/2}$, by
\begin{equation}\label{eq:ss-covector-root-extension}
                  \lambda(x)^2=f(Fx).
\end{equation}
The right side is linear in the squares of the coordinates of $x$,
so this defines a linear functional over $K^{1/2}$. On $I$ it
equals $f$ by uniqueness of square roots. Since $Fx$ lies in $I$,
it satisfies $\lambda(Fx)=\lambda(x)^2$ on the whole bundle.
This extension is unique: any equivariant functional zero on $I$
has square zero by \eqref{eq:ss-covector-root-extension}.

In coordinates, $M(s)\lambda=\lambda^{[2]}$. The squared coefficient
vector $w=\lambda^{[2]}$ is $K$-rational and is a rational section of
$F_P^*\mathcal B^*=W\otimes\mathcal O_P(2)$. It satisfies
\begin{equation}\label{eq:ss-square-covector-equation}
                       M(s)^{[2]}w=w^{[2]}.
\end{equation}
It has no divisorial poles. In a regular local frame, if the largest
pole order among its coordinates were $a>0$, the right side would
have a coordinate of pole order $2a$, whereas the left side has
pole order at most $a$. Normality and local freeness extend the
section also across codimension two. Thus $w$ is a vector of
homogeneous quadratic polynomials in the coefficients $s_i$.

\par\smallskip\noindent\textbf{Cross terms cannot occur.}
Write $M(s)=\sum s_i M_i$. The common kernel of the $M_i$ is zero.
Indeed, relative finite-Frobenius duality, with the invariant top
form trivializing the canonical bundle, gives a perfect symmetric
Cartier pairing
\begin{equation}\label{eq:ss-perfect-Cartier-pairing}
 F_{X/k*}N\otimes F_{X/k*}N\longrightarrow N^{(2)}.
\end{equation}
The finite-duality and coefficient-trace formulas used here are
\cite[Lemma~1.3.6, Proposition~1.3.7, Theorem~1.3.8, pp.~26--28]
{BrionKumar2005}, already used in the general adjunction proof.
Its pairing on global sections is the root-removed bilinear map
$\beta(s,t)=M(s)t$. The bundle $F_{X/k*}N$ is globally generated.
On either elliptic factor, evaluation on a Frobenius fibre is
evaluation on a length-two subscheme, and the obstruction group is
$H^1(\mathcal O(mO)(-Z))=0$ because $m-2>0$.
Tensoring the two evaluation surjections proves the assertion on
the product. A global section pairing to zero with every global
section therefore vanishes in every fibre, by perfectness, and
is zero. This proves the common-kernel assertion, also after
squaring the matrices.

For $j<k$, let $c_{jk}$ be the coefficient of $s_js_k$ in $w$.
The right side of \eqref{eq:ss-square-covector-equation} has only
even exponents. The monomials with precisely $s_j,s_k$ occurring
to odd powers on its left side give
\[
              \sum_i s_i^2M_i^{[2]}c_{jk}=0.
\]
Thus every $M_i^{[2]}c_{jk}=0$, and $c_{jk}=0$.
All cross terms vanish. Since $k$ is perfect,
\[
                         w=D(s)^{[2]}
\]
for a $k$-linear endomorphism $D$ of $W$. Taking square roots of
\eqref{eq:ss-square-covector-equation} shows
\begin{equation}\label{eq:ss-exact-linear-operator}
                 \beta(s,D(s))=D(s)^{[2]}.
\end{equation}
The original functional is exactly $D(s)$, not merely its
purely inseparable extension. Conversely every $D$ satisfying
\eqref{eq:ss-exact-linear-operator} is a regular equivariant
covector. No quotient by scalar identities is taken here:
the full-to-relative quotient has Frobenius zero, so it has no
nonzero equivariant covectors.

\par\smallskip\noindent\textbf{Descent to a small endomorphism algebra.}
Polarizing \eqref{eq:ss-exact-linear-operator} gives
\[
                   \beta(Ds,t)=\beta(s,Dt).
\]
If a global section $s$ has zero value in a fibre of $F_*N$,
the right side pairs that zero value with $Dt$, so the left
side vanishes for every $t$. Global generation and the perfect
pairing force $Ds$ to have zero value there too.
Hence $D$ preserves every evaluation kernel and descends to an
$\mathcal O_{X^{(2)}}$-linear endomorphism of $F_*N$.
Formally the kernel of the global evaluation map is locally free;
the induced map from it to $F_*N$ vanishes on every closed fibre
and therefore vanishes identically. The descended endomorphism
acts on global sections as the original $D$.

For a line bundle $A$ of odd degree $m$ on an elliptic curve,
$F_*A$ is stable of rank two and degree $m$.
A line subbundle $T$ of degree at least $m/2$ would, by adjunction,
give a nonzero map $F^*T\to A$, hence $2\deg T\le m$,
contradicting the integral degree assumption. A noninvertible
nonzero endomorphism would have rank-one kernel and image; stability
puts both their degrees strictly below $m/2$, contradicting
additivity of degree. Every nonzero endomorphism is therefore
invertible, and over the algebraically closed field its endomorphism
algebra is $k$. Kunneth consequently gives
$\operatorname{End}(F_*N)=k$ when $m$ is odd.
Equation~\eqref{eq:ss-exact-linear-operator}, together with
$\beta(s,s)=0$, forces the scalar $D$ to be zero.

For even $m$, projection formula gives
$F_*\mathcal O(mO)=\mathcal O_{E_0^{(2)}}((m/2)O)\otimes U$,
where $U=F_*\mathcal O_{E_0}$. Its rank-one quotient in
$0\to\mathcal O\to U\to B\to0$ has degree zero. The cohomology
sequence and the zero Frobenius on $H^1(E_0,\mathcal O)$ give
$h^0(B)=1$, so $B\simeq\mathcal O$.
The extension is nonsplit because $h^0(U)=1$.
Every endomorphism preserves its unique global section and induces
scalars $a,b$ on the subline and quotient. Compatibility with its
nonzero extension class gives $a=b$. Subtracting $a$ times the
identity leaves a multiple of the quotient-followed-by-inclusion
map, which is nonzero and square zero. The invariant derivation
$\delta$ is precisely a nonzero multiple of this map. Thus
\[
 \operatorname{End}(U)=k[\delta]/(\delta^2),\qquad
 \operatorname{End}(F_*N)
       =k[\delta_1,\delta_2]/(\delta_1^2,\delta_2^2).
\]

It remains to impose the Cartier equation on
$D=a+b\delta_1+c\delta_2+e\delta_1\delta_2$.
In a pullback frame, write
$s=u^2+x_1v^2+x_2w^2+x_1x_2z^2$.
Taking the coefficient of $x_1x_2$ gives
\[
 C(s^2\omega)=C(s\delta_1s\,\omega)
 =C(s\delta_2s\,\omega)=0,
 \qquad C(s\delta_1\delta_2s\,\omega)=\delta_1\delta_2s\,\omega.
\]
Consequently $C(sD(s)\omega)=D(s)\omega$ is equivalent to
$a=b=c=0$ and $e^{1/2}=e$, or $e\in\mathbf F_2$.
The four operators are independent on global sections because
the evaluation map is surjective, so no further solutions arise.
The nonzero solution $D=\delta_1\delta_2$ exists and is nonzero:
$D(x_1x_2)=1$. This proves the geometric dimensions in both parities.

Finally $D$, the invariant differential, the line-bundle square root
and its Cartier connection are defined over $\mathbf F_2$.
The unique nonzero even-degree covector is therefore arithmetic
invariant over the stated perfect base field. In odd degree there
are no geometric invariant covectors and hence no arithmetic ones.
This proves the asserted arithmetic dimensions as well.
\end{proof}

\begin{remark}[Exact local checks and scope]\label{auto:A_monodromy:9cc9a63fa2c6}
\status{Proved under stated hypotheses}

The local four-operator identities are reproduced symbolically in
\path{agents/monodromy/code/supersingular_cartier_operator_checks.py},
with output
\path{agents/monodromy/outputs/supersingular_cartier_operator_checks.json}.
The same retained program performs a bounded check of constant
alternating symmetrizers for the one-factor Cartier pencils: the
space has dimension one in pole degree three and zero in the tested
odd pole degrees five, seven, nine, eleven and thirteen. This check
is not used in the proof and asserts nothing for untested degrees;
it records why the special degree-one Hessian calculation has not
been extrapolated to higher degrees.
\end{remark}


% /root/monodromy, 2026-11-19. Proved; independent root and reduction_descent audits passed.
\subsection{A natural eight-dimensional block in every supersingular degree}
\label{subsec:ss-uniform-GL8-block}
\status{Proved and checked. The parent and reduction/descent agent
independently audited both boundary constructions, the full node-cover
comparison and the extraction of an actual natural block. This determines
a uniform subgroup; the higher-degree full-image question remains separate.}

\begin{theorem}\label{thm:ss-uniform-GL8-block}
\status{Proved under stated hypotheses}
For $X=E_0\times E_0$, $E_0:y^2+y=x^3$, with
$L=\mathcal O(3O)\boxtimes\mathcal O(3O)$ in characteristic two,
the degree-$d$ geometric Hasse--Witt image contains, for every $d\ge1$,
a subgroup acting as the natural $\GL_8(\mathbf F_2)$ on an
eight-dimensional subspace and fixing a complementary subspace pointwise.
Consequently the whole representation admits no invariant nondegenerate
alternating bilinear form. In even degree the same conclusion holds on
the common kernel of the invariant covectors of
Theorem~\ref{thm:ss-product-exact-covectors}.
No irreducibility or full-image assertion in higher degree is made here.
\end{theorem}

\begin{proof}
We work over an algebraic closure. The degree-one case is
Theorem~\ref{thm:ss-product-degree-one-full-group}.
We use the exact ranks of Theorem~\ref{thm:DE-A3-ss-product-rank}.
The maximum stable-rank open of the coherent family includes singular
divisors; its unit Lang system restricts to the local system in the
question on the dense smooth open, with the same image.
Indeed an iterated Frobenius image of constant maximum rank is a
unit subbundle, and restriction of a finite etale cover on a normal
integral scheme to a dense open preserves its connected components.

\par\smallskip\noindent\textbf{Odd degree: an odd self-isogeny and nodal images.}
Suppose $d\ge3$ is odd. Lemma~\ref{lem:DE-A3-ss-norm} gives an
endomorphism $b$ of $E_0$ with $b^\dagger b=bb^\dagger=[d-1]$.
Put
\[
 \phi=\begin{pmatrix}1&b\\-b^\dagger&1\end{pmatrix}:X\longrightarrow X.
\]
Then $\phi^\dagger\phi=\phi\phi^\dagger=[d]I$, and $\phi$ is
an etale isogeny of degree $d^2$ with constant odd kernel $G$.
For a generic smooth degree-one member $C\in|L|$, the image
$Z=\phi(C)$ is integral and nodal, and $C$ is its normalization.
Here is why the nodal-image argument applies also to this nonproduct
isogeny and to supersingular elliptic factors.

Only translations by points of the constant kernel enter that argument.
A generic $C$ is not invariant under any nonzero translation: if a
translation preserves $L$, its projective action is nontrivial because
$L$ is very ample; otherwise it preserves no member of $|L|$.
Thus $C\to Z$ is generically one-to-one. For a nonzero kernel point
$g$, double identifications are parametrized by
$C\cap(C-g)$, of degree $L^2=18$.

The tensor first-jet calculation of
Lemma~\ref{lem:product-small-degree-jets} uses only the plane cubic
embeddings and the fact that tangents through a fixed point are finite.
For the present cubic $Y^2Z+YZ^2+X^3$, the partial derivatives are
$X^2,Z^2,Y^2$, which are linearly independent, and no nonzero quadratic
vanishes on the cubic. Thus that finiteness condition holds.
Two first jets have rank at least five; rank six fails only on the
two coordinate-fibre loci and the finite crossed-tangency locus
over either fixed first point. For pairs $P,P+g$ with both coordinates
of $g$ nonzero, the exceptional points satisfy equations
$3P_i=\text{constant}$ and are finite. Away from them, the two value
conditions and proportionality of the two nonzero gradients impose
three conditions with $P$ varying in dimension two.
At the finitely many exceptions, the two independent values already
give a proper incidence image.

If one coordinate of $g$ is zero, the joint first-jet rank is five.
After imposing the two values, its sole remaining relation involves
the two derivatives along the translated factor and is nonzero;
the two transverse derivatives remain independent. The determinant
of the two gradients is not identically zero on this three-dimensional
allowed derivative space. Tangency is again a third condition.
Thus every $C\cap(C-g)$ is transverse for generic $C$.

Three-point identifications are excluded by the same Segre value
calculation used in Lemma~\ref{lem:ordinary-product-isogeny-nodal-images}.
Three distinct rank-one Segre tensors are collinear only if one factor
is constant and the other three points are collinear. For
$P,P+g,P+h$, the exceptional locus has dimension at most one, since
it requires translations in the same factor and
$3P_i+g_i+h_i=0$. The value ranks are three off this locus and two
on it; the incidence dimensions in $|L|=\mathbf P^8$ are at most
$8-3+2=7$ and $8-2+1=7$. Etaleness of $\phi$ therefore gives only
ordinary nodes. Each $g$ contributes an etale intersection cover of
degree eighteen, with the constant identification of $g$ and $-g$.
The node permutation image is consequently a subquotient of a
product of $S_{18}$'s.

The image has exactly the required line bundle. Write
$\Theta=(\{O\}\times E_0)+(E_0\times\{O\})$, so $L=\mathcal O(3\Theta)$.
The Rosati identity gives $\phi^*[\Theta]=d[\Theta]$.
For the norm line bundle, pulling back its class gives
\[
 \phi^*[\operatorname{Nm}_\phi L]
   =\sum_{g\in G}[t_g^*L]=d^2[L]=\phi^*[L^d].
\]
Pullback by an isogeny is injective on the torsion-free Neron--Severi
group after tensoring with $\mathbf Q$, so the two line bundles
have the same numerical class and differ by a member of $\Pic^0(X)$.
Both are symmetric, and $\Pic^0(X)[2](k)=0$; hence the difference
is trivial. Thus $Z\in|L^d|$ exactly.

Its arithmetic genus is $9d^2+1$, while the normalization has
genus ten and stable rank eight. Therefore it has $9(d^2-1)$
nodes and stable rank $8+9(d^2-1)=9d^2-1$, the full generic rank.
On this boundary family its Lang representation has an exact sequence
\[
                  0\longrightarrow T\longrightarrow V_d
                    \longrightarrow Q\longrightarrow0,
 \qquad Q\simeq V_1,\quad \dim Q=8,
\]
where $T$ is the graph permutation contribution. There is no ambient
stable contribution because $X$ is supersingular.
The quotient image is full $\GL_8(\mathbf F_2)$, whereas the whole
image on $T$ has prime factors at most seventeen.

\par\smallskip\noindent\textbf{Even degree: a fixed odd member and a symmetric node cover.}
Let $d\ge2$ be even. Choose a fixed smooth maximal-rank member
$D\in|L^{d-1}|$, which exists by the odd-degree rank theorem, and
vary $C\in|L|$ through its smooth maximal-rank open, transverse to $D$.
Then $Y=C\cup D$ is nodal and
\[
 f_Y=8+\bigl(9(d-1)^2-1\bigr)+18(d-1)-1=9d^2-3.
\]
This again equals the full generic rank. The boundary Lang system has
a filtration
\[
 0\subset T\subset T_0\subset V_d,
 \qquad T_0/T\simeq H^1_{\et}(D,\mathbf F_2)\text{ constant},
 \qquad V_d/T_0=Q\simeq V_1.
\]
The graph space is $T=\mathbf F_2^N/\mathbf F_2\mathbf 1$,
where $N=18(d-1)$ is the number of intersection points.
Its permutation image is the full $S_N$.

To justify that last assertion, the one-point incidence over $D$ and
the ordered distinct-pair incidence over $D\times D\setminus\Delta$
are irreducible projective bundles: the very ample embedding separates
one and two values. Thus the node action is doubly transitive.
There exists a section with exactly one double zero on $D$ and all
other zeros simple. Choose a general point $P\in D$ whose two
elliptic coordinates are not flex points. At $P$, $L$ separates
second jets, so the quadratic coefficient along the tangent of $D$
can be made nonzero after the value and first tangent derivative
are killed. For $Q\in D\setminus\{P\}$, the two curve jets at
$P,Q$ have combined rank four except for finitely many $Q$, and
rank at least three at those exceptions, by the preceding tensor
first-jet bounds. The exceptional points are finite because $D$ is
integral and ample, hence is not a coordinate fibre. Within the
codimension-two system imposed at $P$, another double zero imposes
two conditions with $Q$ varying in dimension one, or at least one
condition at finitely many exceptions. Its incidence image is proper.
The exact double zero is thus attainable without any second collision.
The point-incidence source is smooth and the base is smooth.
These are precisely the two conditions of
\cite[Proposition~3.1, arXiv version~3, pp.~7--8]{Entin2021}, already
used in Lemma~\ref{lem:all-X-independent-nodes}; they yield a
transposition also in characteristic two and therefore the full $S_N$.
Shrinking to the smooth maximal-rank transverse $C$ open does not
change this finite permutation image.

The action of $S_N$ on $\mathbf F_2^N/\mathbf F_2\mathbf1$ is
faithful for $N\ge3$: the difference of two distinct coordinate
vectors has support two and cannot be the all-ones vector.
The whole image on $T_0$ is an extension of $S_N$ by a two-group,
since $T_0/T$ is constant. Its only nonabelian composition factor
is $A_N$. This is not $\GL_8(\mathbf F_2)$: $11$ divides $|A_N|$
for $N\ge18$, whereas it does not divide
$|\GL_8(\mathbf F_2)|=2^{28}\prod_{i=1}^8(2^i-1)$.

\par\smallskip\noindent\textbf{Extraction of an actual block.}
In either parity let $T_0$ denote the kernel of the quotient to $Q$
(in odd degree take $T_0=T$). The simple group $\GL_8(\mathbf F_2)$
has no quotient in common with the entire $T_0$ image: in odd degree
use the prime $31$, and in even degree use the just determined
composition factors. Goursat's lemma shows that the subgroup acting
trivially on all of $T_0$ still surjects onto $\GL(Q)$.
Its translation kernel in $\operatorname{Hom}(Q,T_0)$ is
$Q^*\otimes T'$ for a subspace $T'\subset T_0$, since this module
is a direct sum of copies of the absolutely simple natural dual.
Modulo that kernel, the remaining block extension is a cocycle with
values in $\operatorname{Hom}(Q,T_0/T')$.
The proved elementary vanishing
$H^1(\GL_8(\mathbf F_2),Q^*)=0$
(Lemma~\ref{lem:char2-group-inputs}) makes it a coboundary.
After conjugation it is zero. The entire translation kernel is
present, so its elements then remove the remaining $T'$-valued
columns. Hence the subgroup contains every block diagonal matrix
$\operatorname{diag}(I_{T_0},A)$ with $A\in\GL(Q)$.
This proves the actual natural-block assertion.

Finally the natural eight-dimensional $\GL_8(\mathbf F_2)$ module
has no invariant covector and no nonzero invariant bilinear form.
For the latter, the stabilizer of any nonzero vector has no fixed
covector: its arbitrary upper row first kills the covector's first
coordinate, and its $\GL_7$ block kills the rest. Thus there is
no nonzero equivariant map $Q\to Q^*$.
Any invariant bilinear form on $V_d=Q\oplus T_0$ therefore vanishes
on $Q\times Q$ and on $Q\times T_0$, so $Q$ lies in its radical.
In even degree every invariant covector vanishes on $Q$; its common
kernel is $Q\oplus(T_0\cap K)$, and the same argument applies there.
This proves the stated form exclusions.
\end{proof}


\subsection{A mixed product with a uniform exact rank formula}
\label{sec:DE-A3-mixed-generic}

The following statement uses one fixed perfect field and one fixed embedded
surface in every degree.  Specialization is used to prove the formula;
the ordinary elliptic factor is not chosen anew as the degree changes.

\begin{theorem}[The geometric generic ordinary elliptic factor]
\label{thm:DE-A3-mixed-generic}
\status{Proved; complete independent root and monodromy proof audits passed.}
Let $k=\mathbf F_2(t)^{\mathrm{perf}}$ be the perfect closure of the
rational function field, and define
\[
 E_t:\ y^2+txy+y=x^3,\qquad
 E_0:\ y^2+y=x^3,\qquad
 X=E_t\times_k E_0,
 \qquad L=\OO_{E_t}(3O)\boxtimes\OO_{E_0}(3O).
\]
Use the cubic embeddings and the Segre map to embed $X$ in
$\mathbf P^8_k$.  The first elliptic factor is ordinary and the second
is supersingular.  For every $d\ge1$, the geometric generic relative
stable rank is
\[
 \boxed{\quad
 \dim E(Y,\OO_Y)_{\ss,\mathrm{gen}}
  =\begin{cases}
       9d^2-1,& d\ \text{odd},\\
       9d^2-2,& d\ \text{even}.
    \end{cases}\quad}                                      \tag{\ddagger}
\]
The relative coherent dimension is always $9d^2-1$.
Thus its generic Frobenius is bijective exactly in odd degree.
The full generic curve defect $g(Y)-f_Y$ is one in odd degree and
two in even degree.

The same formulas hold after replacing $E_t$ by any elliptic curve
over a perfect characteristic-two field whose $j$-invariant is
transcendental over $\mathbf F_2$, and taking geometric generic ranks.
This last assertion does not include arbitrary algebraic
$j$-invariants.
\end{theorem}
\begin{proof}
Consider the smooth elliptic scheme
\[
 \mathcal E\longrightarrow
 S=\Spec\mathbf F_2[t,(t^3+1)^{-1}]
\]
given by the same Weierstrass equation.  The integral formulas with
$a_1=t,a_3=1$ and the other coefficients zero give
$c_4=t^4-24t$ and $\Delta=t^3-27$; in characteristic two these are
$t^4$ and $t^3+1$.  The Hasse coefficient is $a_1=t$.
Thus the generic fibre is ordinary, while the fibre at $t=0$ is
exactly $E_0$.  In particular this is a specialization within smooth
elliptic schemes.

Fix an odd $d\ge1$ and put $m=3d$.  On
$\mathcal X=\mathcal E\times_S(E_0\times S)$ use
$\mathcal N_m=\OO_{\mathcal E}(mO)\boxtimes\OO_{E_0}(mO)$.
Its direct image has rank $m^2$ and commutes with base change,
because both elliptic factor bundles have positive degree and no
higher cohomology.  Let $\mathcal P_m\to S$ be its projective space
of nonzero sections up to scale, and let $\mathcal D_m$ be the
universal divisor.  The relative coherent quotient is a vector
bundle $\mathcal Q_m$ of rank $m^2-1$ on all of $\mathcal P_m$.
Indeed fibrewise K\"unneth gives
\[
 H^1(\mathcal X_s,\mathcal N_{m,s}^{-1})=0,\qquad
 h^2(\mathcal X_s,\mathcal N_{m,s}^{-1})=m^2,
 \qquad h^2(\mathcal X_s,\OO)=1.
\]
Multiplication by a nonzero section maps the middle cohomology
surjectively to the last one, since its Serre-dual map sends $1$
to that nonzero section.  Its kernel is exactly the relative
quotient, including for singular divisors.  This also proves
the required base-change and local-freeness assertions.

Absolute Frobenius gives an $\OO_{\mathcal P_m}$-linear map
$F_{\mathcal P_m}^*\mathcal Q_m\to\mathcal Q_m$.
Its invertibility locus is open.  On the fibre $t=0$, the
supersingular-product theorem
\ref{thm:DE-A3-ss-product-rank} supplies a smooth section where it
is invertible, since $d$ is odd.  The determinant is therefore
a nonzero section of a line bundle on the integral scheme
$\mathcal P_m$.  It cannot vanish on its entire generic fibre:
otherwise it would vanish at the generic point of
$\mathcal P_m$, hence everywhere.  The geometric generic curve
over $\mathbf F_2(t)$ consequently has relative rank $m^2-1$.
Passing to its perfect closure preserves this rank, as does any
further algebraically closed extension.

This argument applies to every odd $d$ at the same generic point
of $S$.  It makes no assertion that an infinite intersection of
nonempty open subsets of $S$ contains a closed point.  The
elliptic curve over the fixed field $k$ is therefore the same
$E_t$ for all of these conclusions.  The smooth-member open is
nonempty in each degree because both factor line bundles have
degree at least three.

Now work over $\bar k$ and let $Y\in|L^d|$ be smooth, with $m=3d$.
The Picard embedding and connected-Prym proof of
Lemma~\ref{lem:DE-A3-ss-prym} applies word for word to
$E_t\times E_0$: it only used negative-degree K\"unneth vanishing
for the embedding, and the diagonal product polarization for
$qj=[m]$.  Thus
\[
 0\longrightarrow P\longrightarrow J_Y\xrightarrow{q}X
 \longrightarrow0
\]
is exact with $P$ an abelian variety.  Here $g(Y)=m^2+1$,
$f_X=1$, $f_Y=f_P+1$, and $\dim P=g(Y)-2$.
If $m$ is even, $j$ embeds $X[2]$ into $P[2]$.  The identity
component of $\ker V_X$ has order $2$, since $\dim X=2$ and
$f_X=1$.  Applying Frobenius twist and functorial finite-group
Verschiebung to that immersion embeds this connected group
into $\ker V_P$.  Therefore
\[
 \dim P-f_P\ge1,\qquad f_Y\le g(Y)-2.
\]
Exactness of stable parts, now subtracting the one-dimensional
ambient stable part, gives the upper bound
$\dim E(Y,\OO_Y)_{\ss}\le m^2-2$.

For even $d\ge2$, choose generic smooth members
$C\in|L|$ and $D\in|L^{d-1}|$ in the odd-degree maximal-rank
opens just proved.  Each has full genus-minus-$2$-rank defect
one, since its ambient stable rank is one.  They can be chosen
transverse, by very ampleness and Bertini.  The nodal union
$C\cup D\in|L^d|$ has full defect two: the normalization
sequence adds the graph rank bijectively for Frobenius and
adds the two component defects.  Constant coherent dimensions
and the open-rank argument in
Theorem~\ref{thm:DE-A3-ss-product-rank} deform it to a smooth
member of relative rank at least $m^2-2$.  The preceding upper
bound makes this equality.  This proves (\ddagger).
Surjectivity from ambient degree-$d$ hypersurfaces to
$H^0(X,L^d)$ follows from the two plane-cubic coordinate rings
and the Segre map exactly as in that theorem.

Finally the rational function
\[
                         j(E_t)=\frac{t^{12}}{t^3+1}
\]
is nonconstant.  If an elliptic curve $E$ has transcendental
$j(E)$, choose a root $u$ of
$u^{12}-j(E)(u^3+1)=0$ in an algebraically closed extension.
Then $u$ is transcendental over $\mathbf F_2$ and $u^3+1\ne0$.
The resulting inclusion $\mathbf F_2(t)\hookrightarrow\bar k$,
$t\mapsto u$, identifies $E$ geometrically with the base change
of $E_t$, since elliptic curves with the same $j$-invariant
are isomorphic over an algebraically closed field.  The origins
and their cubic line bundles are respected.  Geometric generic
ranks are invariant under this field extension, proving the
last assertion.
\end{proof}

\begin{remark}[Algebraic ordinary factors remain a distinct refinement]
\label{rem:DE-A3-mixed-fixed-refinement}
\status{Proved under stated hypotheses}
The theorem does not infer the formula for every ordinary elliptic
curve over $\overline{\mathbf F}_2$.  For the explicit algebraic
factor $y^2+xy=x^3+1$, exact retained Cartier calculations attain
the predicted bounds in degrees $1,\ldots,6$, and in every
rectangular system $(a,3)$ with $3\le a\le31$; the retained run
took $0.7093895580001117$ seconds.  The outstanding uniform
construction and failed sparse-section route are recorded in
\texttt{notes/A3\_mixed\_rank\_research.md}.  These finite probes
are not substituted for the specialization proof or for the
remaining algebraic-factor question.
The separate auxiliary certificate
\path{agents/obstructions_crystals/code/A3_mixed_generic_checks.py}
checks the integral Weierstrass invariants, their characteristic-two
reductions, the nonvanishing denominator in the $j$-invariant argument,
and the exact genus/rank identities for $10{,}000$ degrees.  It passed
in $0.04470579700137023$ seconds; its complete output is retained
beside the program.  These numerical identities support the displayed
construction, whereas the proof above supplies its uniform geometric
existence assertion.
\end{remark}


% Root proof, 2026-11-19. Independent audit requested before promotion.
\subsection{Two-jet geometry of a genus-two theta surface in characteristic two}
\label{subsec:ppav-two-jets}
\status{Proved. The obstruction/crystal agent independently audited the complete
jet and discriminant proofs, and separately audited the supplementary dichotomy
and ordinary-member arguments.}

\begin{lemma}\label{lem:ppav-surface-dichotomy}
\status{Proved under stated hypotheses}
Over an algebraically closed field, a principally polarized abelian surface
with an effective symmetric theta divisor is either a product of two elliptic
curves with its product theta divisor, or the Jacobian of a smooth genus-two
curve with its smooth theta divisor. This assertion includes characteristic two.
\end{lemma}
\begin{proof}
An abelian variety contains no rational curve. Indeed a nonconstant map from
$\mathbf P^1$ factors through a finite number of Frobenius maps until its
differential is nonzero, whereas invariant differentials generate the
cotangent bundle of the abelian variety and $H^0(\mathbf P^1,\Omega^1)=0$.
For any integral curve $B$ on the surface, adjunction therefore gives
$B^2=2p_a(B)-2\ge0$. Also $\Theta^2=2$ and $\Theta\cdot B>0$.
Writing $\Theta=\sum m_iB_i$ shows that a reducible theta divisor must
have exactly two components, each with multiplicity one: their positive
integer intersections with $\Theta$ sum to two. A single component of
multiplicity at least two would give $m_i^2B_i^2=2$, which is impossible.
In the reducible case, evenness of $B_i^2$ and
$B_i^2+B_1\cdot B_2=1$ force $B_i^2=0$ and $B_1\cdot B_2=1$.
Each component is consequently a smooth genus-one curve, hence a translate
of an elliptic subgroup. Addition of the two elliptic subgroups is an
isogeny of degree their intersection number, namely one. It is the
required product isomorphism with its theta divisor.

If $\Theta$ is integral, $p_a(\Theta)=2$. Its normalization cannot have
genus one: a map from an elliptic curve to an abelian variety has a smooth
elliptic translate as image, which would have self-intersection zero.
Thus its normalization has genus two and arithmetic genus two, so
$\Theta$ is smooth. Its inclusion induces an isogeny
$u:\operatorname{Jac}(\Theta)\to A$, since this ample curve generates
$A$. Let $C_J$ denote the genus-two Abel--Jacobi curve in its Jacobian.
Then $u^*\Theta\cdot C_J=2$ and $C_J^2=2$. Hodge index gives
$(u^*\Theta)^2\le2$, whereas this square is $2\deg u$.
Thus $\deg u=1$ and $u$ is an isomorphism. Its restriction maps
$C_J$ onto $\Theta$, so it identifies the principal polarizations.
\end{proof}

We work in this subsection over an algebraically closed field $k$ of
characteristic two. Let $C$ be a smooth genus-two curve, choose a theta
characteristic $\kappa$ with $\kappa^{\otimes2}=\omega_C$, and put
\[
 A=\operatorname{Pic}^0(C),\qquad
 i:C\hookrightarrow A,\quad c\longmapsto\mathcal O_C(c)\otimes\kappa^{-1},
 \qquad \Theta=i(C),\qquad L=\mathcal O_A(3\Theta).
\]
Thus $\Theta$ is smooth and symmetric and $\Theta^2=2$. These Jacobian facts, including the characteristic-two
hyperelliptic description, are recorded in
\cite[Theorem~14.23 and Corollary~14.24, pp.~230--231; Theorem~14.29, p.~233;
Example~14.34, pp.~236--237]{EdixhovenVanDerGeerMoonenAV}.
Write $\Theta_a=\Theta+a$ and
$P_a=t_a^*\mathcal O_A(\Theta)\otimes\mathcal O_A(-\Theta)$.
The principal polarization identifies $a\mapsto P_a$ with an isomorphism
$A\to\operatorname{Pic}^0(A)$. Additions and integer multiples of points
below are in this group law. Put $W=H^0(A,L)$, so $\dim W=9$.
For a point $P\in A(k)$, a first jet means the value in
$L_P\otimes\mathcal O_{A,P}/\mathfrak m_P^2$, of dimension three;
a second jet means the analogous quotient by $\mathfrak m_P^3$, of
dimension six. All jet maps in this subsection are $k$-linear evaluation maps.

\begin{lemma}\label{lem:ppav-theta-preliminaries}
\status{Proved under stated hypotheses}
Every line bundle $\mathcal O_A(2\Theta)\otimes P_a$ is globally generated
and has no positive-degree cohomology. The system $|2\Theta|$ separates
two distinct points $x,y$ whenever $y\ne -x$.
Moreover, under $\operatorname{Pic}^0(C)=A$,
\[
                     i^*P_a=-a.
\]
The Gauss map of $\Theta$ to the projective line of cotangent directions
is the canonical degree-two map of $C$.
\end{lemma}

\begin{proof}
The theorem of the square gives
$\Theta_c+\Theta_{-c}\in|2\Theta|$.
For a specified point $x$, choose $c$ outside the two curves
$\Theta+x$ and $\Theta-x$; the resulting section is nonzero at $x$.
If $y\ne x,-x$, choose $c\in\Theta+x$ outside
$(\Theta+y)\cup(\Theta-y)$. These are distinct irreducible translates:
the translation stabilizer of $\Theta$ is contained in the kernel of
its principal polarization and is trivial. Their intersections with
$\Theta+x$ are finite, so this choice exists. The resulting section
vanishes at $x$ and is nonzero at $y$. Together with a section nonzero
at both points this proves separation. Every twist by $P_a$ is a
translate of $\mathcal O(2\Theta)$, since multiplication by two on
$A(k)$ is surjective, so it is globally generated as well.

For completeness the required ample vanishing does not use
characteristic-zero Kodaira vanishing. If $n$ is odd, the finite etale
map $[n]:A\to A$ has invertible degree $n^4$, and its trace splits the
unit map of structure sheaves. Thus
$H^j(A,T)\hookrightarrow H^j(A,[n]^*T)$ for a line bundle $T$.
For $T=\mathcal O(2\Theta)$, symmetry gives
$[n]^*T\simeq\mathcal O(2n^2\Theta)$; Serre vanishing for arbitrarily
large odd $n$ proves the assertion for $j>0$. Translation proves it
for the twists. The same argument and Riemann--Roch give
$h^0(3\Theta)=(3\Theta)^2/2=9$.

Let $\iota$ be the hyperelliptic involution of $C$; then
$i(\iota c)=-i(c)$ and $c+\iota c$ is canonical. For a generic
$a=i(r)-i(s)$, the intersection $\Theta\cap\Theta_a$ consists of
$i(r)$ and $i(\iota s)$. Indeed these are intersection points, they
are distinct generically, and the intersection has length two.
The difference map $\Theta\times\Theta\to A$ is surjective, since
two ample translates intersect. Restricting $\mathcal O(\Theta_a)$
to $C$ therefore gives generically
$\mathcal O_C(r+\iota s)$. Since
$\mathcal O(\Theta_a)=\mathcal O(\Theta)\otimes P_{-a}$ and
$i^*\mathcal O(\Theta)=\omega_C=\mathcal O_C(s+\iota s)$,
we obtain $i^*P_{-a}=a$. Equality extends from this dense set to the
morphisms $A\to\operatorname{Pic}^0(C)$, proving the formula for all $a$.
Finally the restriction of invariant one-forms on the Jacobian to
$C$ is the defining isomorphism
$H^0(A,\Omega_A^1)\simeq H^0(C,\omega_C)$.
The conormal direction of $\Theta$ is consequently its canonical
map. This map has degree two and is finite and surjective onto
$\mathbf P^1$. No assertion about a characteristic-two Kummer
quotient is needed.
\end{proof}

\begin{lemma}\label{lem:ppav-three-distinct-values}
\status{Proved under stated hypotheses}
In the smooth-theta case, any three distinct points of $A$ impose three
independent value conditions on $H^0(A,3\Theta)$.
\end{lemma}
\begin{proof}
It suffices to separate the value at $P_3$ from zero values at
$P_1,P_2$. Choose $a\in\Theta+P_1$ and $b\in\Theta+P_2$.
Require $a,b\notin\Theta+P_3$ and $P_3+a+b\notin\Theta$.
The first two conditions define dense open subsets of their respective
irreducible curves, because the three points are distinct and a theta
divisor has trivial translation stabilizer. The third defines a
nonempty open subset of their product: the sum of these two translates
of $\Theta$ is all of $A$. These three opens have nonempty intersection.
The product of theta sections for $\Theta_a,\Theta_b,\Theta_{-a-b}$
then vanishes at $P_1,P_2$ and not at $P_3$. Repeating with each
point in the third position proves the assertion.
\end{proof}

\begin{theorem}\label{thm:ppav-three-theta-jets}
\status{Proved under stated hypotheses}
With the data just specified, the following assertions hold.
\begin{enumerate}
\item $L$ separates points and first jets, and it separates second jets
at every point of $A$.
\item For $P\ne Q$, the map
\[
 W\longrightarrow
 L_P\otimes\mathcal O_{A,P}/\mathfrak m_P^2\ \oplus\
 L_Q\otimes\mathcal O_{A,Q}/\mathfrak m_Q^2
\]
has rank at least four. It has rank six outside the dimension-two locus
\[
 R=\{(P,Q):2P+Q\in\Theta,\ P+2Q\in\Theta\}.
\]
Its rank is at least five outside the dimension-at-most-one locus
\[
 R_4=\bigcup_{s\in A[3](k)}
       \{(P,s-P):P+s\in\Theta\}.
\]
The rank-four locus is contained in $R_4$, whether or not the rank
actually drops to four at each listed pair.
\item If $A$ is ordinary, then for every fixed $P$ the fibre $R_P$
is finite. In particular two first jets have full rank six for all
but finitely many $Q$ on any specified integral ample curve.
\end{enumerate}
\end{theorem}

\begin{proof}
For first jets at $P$, choose a theta translate through $P$ and
multiply its equation by a section of the remaining twist of
$2\Theta$ nonzero at $P$. The possible linear terms span both
cotangent directions because the Gauss map of $\Theta$ is
surjective. A product nonzero at $P$ supplies the value. Point
separation follows by choosing a theta translate through one point
and not the other, and a nonzero section of the remaining twist
at the latter point. Thus $L$ is very ample as well.

For second jets, fix $P$ and take $p,q\in\Theta$. Put
$a=P-p$, $b=P-q$, and $c=-a-b$. The product of the three theta
sections for $\Theta_a,\Theta_b,\Theta_c$ belongs to $|3\Theta|$.
Its first two factors vanish at $P$; the third is nonzero precisely
when $3P-p-q\notin\Theta$. This condition defines a nonempty open
subset of $\Theta\times\Theta$, because the addition map to $A$
is surjective. The product of the two Gauss maps is finite and
surjective onto $\mathbf P^1\times\mathbf P^1$. The image of that
open is dense, and the associated products of two linear forms
span all of $\operatorname{Sym}^2T_P^*A$: multiplication of binary
linear forms has dense image in the projective space of binary
quadratic forms, also in characteristic two. The resulting
quadratic terms therefore span
$\mathfrak m_P^2/\mathfrak m_P^3$. Combined with first-jet
surjectivity this proves second-jet surjectivity.

Now write $s=P+Q$. The theta translates through both $P,Q$ are
parametrized by the length-two scheme
\[
                      Z=(\Theta+P)\cap(\Theta+Q).
\]
It is preserved by $a\mapsto s-a$. A fixed point of this involution
has $2a=s$; the corresponding points $P-a,Q-a$ on $\Theta$ are
hyperelliptic conjugates and have equal Gauss direction, so its
intersection multiplicity is at least two. Consequently two
distinct support points of $Z$ are exchanged, and a single support
point is fixed.

Fix one such parameter $a$, and write $p=P-a$, $q=Q-a$ as points
of $C$. The restriction exact sequence is
\[
 0\longrightarrow H^0(2\Theta\otimes P_a)
 \longrightarrow W\longrightarrow H^0(\Theta_a,L|_{\Theta_a})
 \longrightarrow0;
\]
the surjectivity follows from Lemma~\ref{lem:ppav-theta-preliminaries}.
After translation to $C$, its last line bundle is
$M_a=\omega_C^{\otimes3}\otimes(-3a)$, of degree six.
The four along-curve jet conditions have rank at least three.
Indeed their cokernel is $H^1(C,M_a(-2p-2q))$, which has dimension
at most one for this degree-two line bundle. It is one exactly
when $M_a(-2p-2q)=\omega_C$. Since
$\mathcal O_C(p+q)\otimes\omega_C^{-1}$ has class $s-2a$,
this happens exactly when
\[
                              a=2s.
\]
The sections vanishing on $\Theta_a$ supply normal derivatives of
rank at least one. They supply rank two if $a\ne -s$: choose
$b\in A(k)$ with $2b=a$. Their evaluations are those of
$2\Theta$ at $P+b,Q+b$, which are distinct and are not opposite
unless $a=-s$. Lemma~\ref{lem:ppav-theta-preliminaries} applies.
Projection of the surface jet space to the along-curve jet space
has the two normal directions as its kernel. Hence the full rank
is at least the sum of the along-curve rank and the supplied normal
rank. It is therefore at least four always, six if a parameter
$a$ avoids both $-s,2s$, and five unless those two bad parameters
coincide.

If every support point of $Z$ is among $-s,2s$, then $-s$ is a
support point. For two support points this is immediate. For one
support point equal to $2s$, its being fixed gives $4s=s$, so
$3s=0$ and $2s=-s$. Thus rank failure from six implies that
$\Theta_{-s}$ contains both $P,Q$, which is exactly the displayed
condition defining $R$. Rank below five requires $3s=0$, and then
these two containment conditions reduce to $P+s\in\Theta$.
The map $(P,Q)\mapsto(2P+Q,P+2Q)$ is an etale isogeny of
$A\times A$ (its determinant is three). Its inverse image of
$\Theta\times\Theta$ has dimension two. Since $A[3](k)$ has
$3^4=81$ elements, $R_4$ has dimension at most one.

For the final assertion suppose a fibre $R_P$ were positive
dimensional. Its first condition places it on the irreducible curve
$\Theta-2P$; the second would then force
\[
                           [2]\Theta=\Theta+3P.
\]
When $A$ is ordinary, write $[2]=V\circ F_{A/k}$, where $V$ is
etale. The restriction of $F_{A/k}$ to $\Theta$ has degree two.
The restriction of $V$ to $\Theta^{(2)}$ has generic degree one:
any larger degree would make that curve invariant under a nonzero
element of the constant kernel of $V$, whereas a principal theta
divisor has trivial translation stabilizer. Thus $[2]|_\Theta$
has degree two. The asserted image equality contradicts the
projection formula: its left interpretation gives
$([2]_*\Theta)\cdot\Theta=2\Theta^2=4$, while the projection
formula gives $\Theta\cdot[2]^*\Theta=4\Theta^2=8$.
This proves finiteness and the last assertion.
\end{proof}

\begin{corollary}\label{cor:ppav-discriminant-one-node}
\status{Proved under stated hypotheses}
The discriminant of singular divisors in $|3\Theta|$ is irreducible
of dimension seven. Its general member is integral and has exactly
one singularity, an ordinary node. These assertions hold for every
$C$ in the stated characteristic; ordinary $A$ is not needed here.
\end{corollary}

\begin{proof}
The incidence of a section singular at a specified point is a
$\mathbf P^5$-bundle over $A$, hence is irreducible of dimension
seven. Second-jet surjectivity makes an ordinary node a nonempty
open condition on this incidence: the mixed coefficient in a
binary quadratic form is nonzero. The incidence of two distinct
singular points has dimension at most six. Indeed its rank-six
stratum has dimension at most $4+(8-6)=6$, its rank-five stratum
at most $2+(8-5)=5$, and its rank-four stratum at most
$1+(8-4)=5$. A general point of the single-point incidence thus
has exactly one singular point and it is an ordinary node.
Projection to $|3\Theta|$ is generically one-to-one, proving the
dimension and irreducibility assertions.

Such a divisor is reduced, since a multiple component would give
a positive-dimensional singular locus. It cannot be reducible:
if it splits as nonzero effective divisors $B+D\sim3\Theta$,
put $t=\Theta\cdot B$, so $1\le t\le5$. Hodge index gives
$B^2\le t^2/2$, and hence
\[
 B\cdot D=3t-B^2\ge 3t-t^2/2\ge5/2.
\]
Its intersection multiplicity is therefore at least three, whereas
two components of a curve with a unique ordinary node can have
total intersection multiplicity at most one. This contradiction
proves integrality.
\end{proof}

\begin{lemma}\label{lem:ppav-ordinary-members}
\status{Proved under stated hypotheses}
If the genus-two Jacobian $A$ above is ordinary, $|3\Theta|$ contains
an ordinary connected nodal divisor with three smooth theta components.
The general discriminant member is ordinary. For every $d\ge1$, a general
smooth member of $|3d\Theta|$ is ordinary, of genus $9d^2+1$.
\end{lemma}
\begin{proof}
Take the three translates $\Theta_a,\Theta_b,\Theta_{-a-b}$.
Their sum belongs to $|3\Theta|$ by the theorem of the square.
They are distinct and pairwise transverse for general $(a,b)$.
To see transversality without a generic-smoothness assumption in positive
characteristic, the difference map $\Theta^2\to A$ is generically etale:
its differential is invertible when the two Gauss directions differ.
Its critical values thus form a proper closed subset of $A$.
Each of the differences $b-a,-2a-b,-a-2b$ varies surjectively in $A$,
so a general pair avoids all three exceptional subsets.
Triple intersections are absent for a general pair as well. A triple
intersection corresponds to $x,p,q,r$ with $p,q,r\in\Theta$ and
$3x=p+q+r$, followed by $a=x-p,b=x-q$. Since $[3]$ is finite,
this incidence has dimension three, below the four-dimensional
parameter space of $(a,b)$.

The three components are ordinary because their Jacobian is the ordinary
surface $A$. The normalization sequence for a nodal curve expresses its
$H^1(\mathcal O)$ as an extension of the component cohomologies by its
graph cohomology. Frobenius is bijective on both terms, so this union
is ordinary. Its ordinary point belongs to the irreducible discriminant
of Corollary~\ref{cor:ppav-discriminant-one-node}, whose ordinary open
is therefore dense. Ordinarity is also a nonempty open condition in
the entire linear system. Its intersection with the nonempty smooth
open supplies smooth ordinary members of $|3\Theta|$.

Finally take $d$ general such members, meeting transversely without
triple points. Their nodal union lies in $|3d\Theta|$ and is ordinary
by the same normalization sequence. The positive-degree ample vanishing
used above makes $h^1(\mathcal O_Y)=9d^2+1$ constant for this family
of effective divisors. Invertibility of its coherent Frobenius is
therefore open, and intersecting again with the smooth open gives
the asserted general ordinary curve and genus.
\end{proof}

\begin{remark}[Source scope]\label{auto:A_ppav_two_jets:0557c5a87cfe}
\status{Proved under stated hypotheses}

The source search found Bauer--Szemberg's higher-jet statements
\cite[Theorems~1 and~2, conventions p.~2]{BauerSzemberg1997}
and Alvarado--Pareschi's recent three-theta second-jet theorem
\cite[Theorem~D; Section~2 conventions]{AlvaradoPareschi2026},
but their stated conventions are in characteristic zero. We do
not import either statement in characteristic two. The proof
above uses the theta theorem of the square, adjunction on a
genus-two curve, the elementary product sections of $2\Theta$,
and the explicit incidence calculation. The fixed-point finiteness
assertion explicitly requires ordinary $A$.
\end{remark}


% /root/monodromy, 2026-11-19. Proved; independent root and reduction_descent audits passed.
\subsection{A level-three theta symmetrizer for the Cartier pencil}
\label{subsec:ppav-theta-Cartier-symmetrizer}
\status{Proved and checked. The parent and reduction/descent agent
independently audited the complete proof, opened the cited primary sources,
and reran the exact finite Heisenberg calculation. The relative Cartier
and theta-group identifications include the Frobenius twists explicitly.}

\begin{theorem}\label{thm:ppav-level3-Cartier-symmetric}
\status{Proved under stated hypotheses}
Let $(A,\Theta)$ be a principally polarized abelian variety of dimension
$g$ over an algebraically closed field $k$ of characteristic two, and let
$L=\mathcal O_A(3\Theta)$ and $W=H^0(A,L)$.
Choose a nonzero invariant top form $\omega$ and put
$C_{\rm abs}(\omega)=c^{1/2}\omega$, with $c\in k$.
There is a theta basis of $W$, indexed by $\mathbf F_3^g$, in which
the root-removed Cartier multiplication matrix $M(s)$ is symmetric
and has diagonal
\[
                         \operatorname{diag}M(s)=c\,s.
\]
If $A$ is ordinary, normalize $\omega$ by $C_{\rm abs}(\omega)=\omega$.
Then the diagonal coordinates are independent and $\det M(s)$ is a
nonzero geometrically squarefree polynomial of degree $3^g$.
In particular, for $g=2$ the degree-one full coherent Hasse divisor of
$|3\Theta|$ is geometrically reduced and each of its components has
generic linearized Cartier corank one.
\end{theorem}

\begin{proof}
\par\smallskip\noindent\textbf{Relative Cartier and its target.}
Let $F=F_{A/k}:A\to A^{(2)}$ be relative Frobenius. Put
$L^{(2)}=L\otimes_{k,F_k}k$ on the Frobenius twist and
$W^{(2)}=H^0(A^{(2)},L^{(2)})$. The canonical map
$\iota:W\to W^{(2)}$, $w\mapsto w^{(2)}$, is square-semilinear.
The relative top Cartier trace is $k$-linear and satisfies
\[
 C_{\rm rel}(\omega)=c\,\omega^{(2)},\qquad
 C_{\rm rel}(F^*a\,\eta)=a\,C_{\rm rel}(\eta).
\]
Since $L^2=F^*L^{(2)}$, multiplication followed by this trace defines
the $k$-bilinear symmetric map
\begin{equation}\label{eq:ppav-relative-beta}
 \beta:W\otimes W\longrightarrow W^{(2)},\qquad
 C_{\rm rel}(s t\omega)=\beta(s,t)\omega^{(2)}.
\end{equation}
In particular $\beta(s,s)=c\,s^{(2)}$.
These are the relative finite-duality and coefficient-trace formulas
of \cite[Lemma~1.3.6, Proposition~1.3.7, Theorem~1.3.8, pp.~26--28]
{BrionKumar2005}. The target is the Frobenius twist; no identification
of $A$ with $A^{(2)}$ is being made.

\par\smallskip\noindent\textbf{A finite theta subgroup and its two representations.}
The theta group of $L$ is a central extension of
$K(L)=A[3]\simeq\mathbf F_3^{2g}$ by $\mathbf G_m$ with nondegenerate
commutator pairing. Choose a Lagrangian decomposition and lifts of
its generators of order three. They, together with the central
$\mu_3$, generate a finite Heisenberg subgroup $H$ of order
$3^{2g+1}$. A primitive cube root $\zeta\in k$ identifies its
Schrodinger representation on a basis $e_a$, $a\in\mathbf F_3^g$, by
\[
 T_x e_a=e_{a+x},\qquad D_y e_a=\zeta^{y\cdot a}e_a,
 \qquad Z e_a=\zeta e_a.
\]
Here $T_x,D_y,Z$ generate $H$; their commutation convention merely
chooses the sign of the standard symplectic pairing.
The natural theta action on $W$ is this representation, denoted
$\rho$. The precise primary statements are
\cite[Corollary~8.20, Lemma~8.24, Theorem~8.28, description~8.29,
pp.~119, 121--123; Corollary~9.17, pp.~133--134]
{EdixhovenVanDerGeerMoonenAV}. They give the nondegenerate theta
group, its Heisenberg form and the unique weight-one representation
on $H^0(L)$; its dimension is $3^g$.

Relative Cartier is natural under translations. A theta lift on $L$
induces its tensor square on $L^2$; under $L^2=F^*L^{(2)}$, this
tensor square is the pullback of its Frobenius-twisted lift on
$L^{(2)}$. The invariant form is unchanged by translation.
Consequently, in the bases $e_a$ and $e_a^{(2)}$, one has the
fully typed equivariance identity
\begin{equation}\label{eq:ppav-Cartier-equivariance}
 \beta(\rho(h)s,\rho(h)t)=\rho(h)^{[2]}\beta(s,t)
                         \quad(h\in H).
\end{equation}
One can also verify this directly in a local frame: if a lift acts
by a line-bundle multiplier $u$, its action on a product has multiplier
$u^2$, and the Cartier projection formula moves its Frobenius
descendant to the target. On the base the translation descends from
$x$ to $F(x)$. Thus the target action is the coefficient-square
representation, not the original weight-one action.

On the displayed generators of the finite subgroup,
\[
                \rho(h)^{[2]}=\rho(h)^{-t}.
\]
For $T_x$ both sides are the same permutation matrix; for $D_y$ and
$Z$ it is the identity $\zeta^2=\zeta^{-1}$. Since both sides are
representations, it holds for every $h\in H$.
Therefore
\[
 R:W^{(2)}\longrightarrow W^*,\qquad e_a^{(2)}\longmapsto e_a^*,
\]
is an invertible $H$-equivariant linear map. This step uses the finite
central subgroup $\mu_3$. The corresponding statement for the whole
central $\mathbf G_m$ would be false and is not used.

\par\smallskip\noindent\textbf{Cyclic symmetry is a translation in level three.}
Define the trilinear tensor
\[
                  T(s,t,u)=\langle R\beta(s,t),u\rangle.
\]
It is $H$-invariant by \eqref{eq:ppav-Cartier-equivariance} and is
symmetric in its first two arguments by multiplication.
Write $T_{abc}=T(e_a,e_b,e_c)$. Invariance under every modulation
$D_y$ forces
\[
            T_{abc}=0\quad\text{unless }a+b+c=0.
\]
Invariance under $T_x$ gives
$T_{a+x,b+x,c+x}=T_{abc}$.
If $a+b+c=0$, translation by $x=b-a$ sends the ordered triple
$(a,b,c)$ to $(b,c,a)$, because $2=-1$ in the index group
$\mathbf F_3^g$. Thus $T$ is invariant under cyclic permutation.
Together with its first-two symmetry this proves full symmetry
under all permutations of its three arguments.

Let $M(s)$ represent $\beta(s,-)$ from the $e_a$ basis to the
$e_a^{(2)}$ basis. The matrix of $R$ in these chosen bases is the
identity. Full symmetry of $T$ gives $M(s)^t=M(s)$.
Furthermore $\beta(e_a,e_a)=c e_a^{(2)}$, so
$T_{aab}=c\delta_{ab}$. By full symmetry, the $a$th diagonal
entry of $M(s)$ is $c s_a$, proving the asserted formula.

For ordinary $A$, Cartier on its invariant top form is nonzero.
If $C_{\rm abs}(\omega)=u\omega$ initially, replacing $\omega$
by $u^2\omega$ normalizes it; hence $c=1$.
Lemma~\ref{lem:symmetric-independent-diagonal-squarefree} now
applies: the determinant has the full squarefree coordinate
monomial with coefficient one, while every nontrivial matching
term has an even positive exponent. It is nonzero and has no
repeated geometric irreducible factor.
At the generic point of a reduced determinant component the
matrix has corank one: corank at least two would make every
cofactor, and hence every first derivative of the determinant,
vanish. This finishes the proof.
\end{proof}

\begin{remark}[Coordinate changes and the finite check]\label{auto:A_monodromy:bcd64c5e340f}
\status{Proved under stated hypotheses}

In an arbitrary original basis the conclusion is the existence of
a constant invertible matrix $R_0$ for which $R_0M(s)$ is symmetric
with independent diagonal coordinates. The matrix $R_0$ itself
need not be symmetric, because the change of basis on the target
is Frobenius-twisted. This is consistent with the explicit elliptic
product symmetrizers proved earlier.
The retained finite calculation
\path{agents/monodromy/code/level3_theta_tensor_checks.py}
checks the $g=2$ index action, all $243$ finite Heisenberg matrices,
and the full permutation symmetry of the five-dimensional space of
invariant first-two-symmetric tensors. It is a check of the finite
algebra, not a substitute for the geometric equivariance proof.
\end{remark}


% /root/monodromy, 2026-11-19. Proved; independent root, obstructions_crystals and reduction_descent full audits passed.
\subsection{Eventual exact groups on every ordinary principally polarized surface}
\label{subsec:ordinary-ppav-surface-full-groups}
\status{Proved and checked. The parent and both other research agents
independently audited the complete assembly, including every exceptional jet
stratum, the exact invariant quotient, the norm degeneration, actual natural
block extraction and the arithmetic image. The cited source scopes were
checked independently; the level-three finite calculation was also rerun.}

Throughout this subsection $(A,\Theta)$ is an ordinary principally polarized
abelian surface in characteristic two, $L=\mathcal O_A(3\Theta)$, and the
linear systems are complete unless ambient hypersurfaces are explicitly
mentioned. A symmetric representative of the polarization is chosen over the
algebraic closure. Translation of that representative gives an isomorphic
geometric family and does not change its geometric group.

\begin{proposition}[All invariant covectors on an ordinary surface]
\label{prop:ordinary-ppav-all-covectors}
\status{Proved under stated hypotheses}
For every $d\ge1$ the geometric generic relative rank is $r_d=9d^2-1$.
The entire geometric invariant-covector space has dimension zero for odd $d$
and three for even $d$. In even degree it is canonically the quotient
\[
 \mathbf F_2^{H}/\mathbf F_2\mathbf 1,
 \qquad H=\operatorname{Hom}(\ker F_A,\mathbf G_m),\quad |H|=4,
\]
with its natural Galois action. Equivalently its basis can be indexed by the
three nonzero characters, which Galois permutes.
\end{proposition}
\begin{proof}
The ordinary nodal degeneration of Lemma~\ref{lem:ppav-ordinary-members}
proves generic ordinarity and the rank formula in the indecomposable case.
The ordinary product degeneration proves it in the decomposable case.
The coherent relative bundle on $P=\mathbf P H^0(A,L^d)$ is $\Omega_P^1$:
indeed $H^1(A,L^{-d})=0$, $h^2(A,L^{-d})=9d^2$, and
$h^2(A,\mathcal O_A)=1$, and the multiplication map to the last group is
surjective, dual to the inclusion of the defining section.

We give the complete endomorphism reduction, to specify its scope beyond
products. Normalize an invariant two-form by $C(\omega)=\omega$ and put
$W=H^0(A,L^d)$. An invariant covector on the ordinary open extends across
all prime divisors of $P$: its equation in regular frames is
$M^t\lambda=\lambda^{[2]}$, and a largest pole of order $a>0$ would have
order $2a$ on the right and at most $a$ on the left. Normality and local
freeness extend it across codimension two. The dual Euler sequence then
represents it by a linear operator $D\in\End(W)$ modulo scalar identity.
Its Frobenius defect is a constant multiple of $s^{[2]}$; adding a scalar
identity and solving an Artin--Schreier equation removes that defect.
Thus exact lifts are the operators satisfying
\[
                 \beta(s,D(s))=D(s)^{[2]},\qquad s\in W,
\]
modulo $\mathbf F_2\operatorname{id}$, where $\beta$ is relative Cartier
multiplication and $\beta(s,s)=s^{[2]}$.

The eigenvalue and Jordan argument of
Theorem~\ref{thm:product-char2-all-parities-fixed-quotient} applies without
any product assumption: eigenvalues are zero or one, and a Jordan chain
$Dv=w,Dw=0$ would give the nonzero coefficient $\beta(w,w)=w^{[2]}$ on
substituting $v+tw$. Every exact lift is consequently an idempotent, and
its two eigenspaces are orthogonal for $\beta$.

The bundle $\mathcal F=F_{A*}L^d$ is globally generated. For an
indecomposable surface, second-jet surjectivity for $3\Theta$ from
Theorem~\ref{thm:ppav-three-theta-jets} implies surjectivity on a
Frobenius fibre: in regular parameters its algebra is
$k[x,y]/(x^2,y^2)$, a quotient of the second-jet algebra since
$(x,y)^3\subset(x^2,y^2)$. For higher $d$, multiply by a section of
$L^{d-1}$ nonzero at the point. Such a section exists since $L$ is globally
generated. On a decomposable surface the two degree-three elliptic factors
interpolate the two length-two Frobenius fibres, giving the same assertion.
Finite Frobenius duality gives the perfect symmetric Cartier pairing
$\mathcal F\otimes\mathcal F\to L^{d,(2)}$.
Orthogonality therefore makes each of the two eigenspaces preserve the
kernel of every evaluation map. Global generation and perfectness descend
$D$ to a bundle endomorphism of $\mathcal F$, exactly as in the cited
product proof; this reasoning uses only these two properties.

If $d$ is odd, apply \cite[Theorem~1.2(i), printed p.~45; proof
pp.~43--45]{Oda1971} to $F_A$ and $L^d$.
This theorem is stated for abelian varieties of arbitrary dimension and
allows inseparable isogenies. Its hypothesis is that the polarization map
of $L^d$ restricts to an isomorphism of the scheme-theoretic kernel of the
isogeny onto its image. Here this restriction is the restriction of
$3d\lambda_\Theta$ to $\ker F_A$; it is an immersion since $3d$ is odd
and $\lambda_\Theta$ is an isomorphism. Thus
$\End(F_{A*}L^d)=k$. The only exact lifts are zero and identity, and their
quotient gives no invariant covector.

If $d$ is even, write $L^d=F_A^*M$.
The ordinary splitting of $F_{A*}\mathcal O_A$, equivalently the character
decomposition for $\ker F_A\simeq\mu_2^2$, gives
\[
          F_{A*}L^d=\bigoplus_{\chi\in H}M\otimes T_\chi.
\]
The four summands are distinct ample line bundles differing by nontrivial
degree-zero two-torsion line bundles. There are no homomorphisms between
different summands, and every summand has nonzero sections by ample
vanishing and Riemann--Roch. Hence $\End(\mathcal F)=k^4$.
Their Cartier cross pairings vanish, since trace kills nonzero character
weights. Its four projectors are precisely the independent exact lifts:
each projector satisfies the displayed equation, and any idempotent in
$k^4$ is their sum. Modding out the identity gives the stated three-space.
This splitting is \cite[Lemma~4.5, p.~11]{SannaiTanaka2016}, or directly
the character algebra of the ordinary Frobenius torsor.

The construction by projectors is canonical under Galois. The zero
character is fixed, and the three nonzero characters are permuted.
Modulo the sum of all four projectors, those three projectors are a basis.
This proves the final assertion, including its arithmetic description.
\end{proof}

\begin{theorem}[Hasse geometry and the irreducible common kernel]
\label{thm:ordinary-ppav-Hasse-kernel}
\status{Proved under stated hypotheses}
For every ordinary principally polarized abelian surface, the full Hasse
divisor in $|L|$ is geometrically reduced with generic Cartier corank one
on each component. For every $d\ge2$ the Hasse divisor in $|L^d|$ is
geometrically irreducible and reduced, with generic stable loss one.
Its geometric relative group is generated by transvections.
If $K_d$ is the common kernel of its entire invariant-covector space,
then $K_d$ is irreducible, and $V_d$ has no invariant vector.
\end{theorem}
\begin{proof}
The decomposable case was proved in
Theorem~\ref{thm:ordinary-products-Hasse-irreducible} and
Corollary~\ref{cor:ordinary-products-irreducible-kernel}.
In the indecomposable case use the genus-two theta divisor and jet
statements of Theorem~\ref{thm:ppav-three-theta-jets}.
Theorem~\ref{thm:ppav-level3-Cartier-symmetric} makes the degree-one Hasse
polynomial reduced with generic corank one. Its components have smooth
generic curves, since the singular-curve discriminant is irreducible and
has ordinary generic member by Lemma~\ref{lem:ppav-ordinary-members}.

We verify all incidence bounds needed to apply the mixed-tacnode argument,
including the rank-four exceptions. Put $e=d-1\ge1$. Let $H_i$ be a
component of the degree-one Hasse divisor, and choose a general smooth
$C\in H_i$. Its Cartier kernel is one-dimensional. The nonzero tangent
functional $dh_1$ belongs to the dual first-jet space at only finitely many
points of $C$: otherwise choose one such $P$ and another outside the
finite exceptional set $R_P$ of the jet theorem. The two first-jet dual
spaces are then disjoint, a contradiction. Choose $P$ also away from the
zeros of the killed differential and from every theta translate
$\Theta-s$, $s\in A[3]$. An integral curve of class $3\Theta$ is not one
of those theta translates, so these exclusions are finite on $C$.

The affine tangent space $\ker dh_1$ to $H_i$ at $C$ surjects onto the
first jets at $P$. The incidence of $D\in|L^e|$ simply tangent to $C$ at
$P$ therefore dominates $|L^e|$: for an infinitesimal change of $D$, move
$P$ to cancel the value change and use these three freely variable first
jets of $C$ to preserve its value and the common tangent line. This is
the three-equation calculation in
Theorem~\ref{thm:product-Hasse-irred}, on the affine cone; projective
scaling imposes no additional condition.

The good contact locus is nonempty. Every $L^e$ inherits the degree-one
second jets at $P$ and the joint first-jet rank bounds at $P,Q$: multiply
by a section of $L^{e-1}$ nonzero at both points, which exists over the
infinite algebraically closed field. Our choice of $P$ eliminates $R_4$,
so all two-point ranks are at least five and are six except at finitely
many $Q$. Within the two imposed contact conditions at $P$, a singularity
at a different $Q$ imposes total rank five generically with $Q$ varying
in dimension two, and total rank at least four at the finitely many
exceptional $Q$. Both are proper loci. Another contact with $C$ imposes
rank four generically with $Q$ varying in dimension one, and rank at least
three at finitely many exceptions, again proper. Second jets exclude
higher contact and a zero normal derivative at $P$. A general member in
the incidence is thus smooth with just this simple tangency. Dominance
then allows it to be ordinary. A single general ordinary $D$ works for
all the finitely many $H_i$.

For this fixed $D$, the contact divisor $\Delta_D\subset|L|$ is
irreducible: its incidence over $P\in D$ is a projective bundle given by
two independent value/tangent conditions. We justify that its generic
point is one simple contact with smooth $C$. For pairs on $D^2$, the
rank-six locus imposes four contact conditions; its point dimension is
two. The exceptional locus $R\cap D^2$ has dimension at most one, since
$R$ is finite over each first point, and imposes at least three conditions
away from $R_4$. The locus $R_4\cap D^2$ is finite: its first point lies
on some $D\cap(\Theta-s)$. It still imposes at least two independent
value conditions. All three counts give codimension at least two in
$|L|$ for two contacts. A singularity of $C$ at a second point imposes
five conditions generically, with a three-dimensional pair space; on
$R$ it imposes at least four with pair dimension at most one, and on
$R_4$ at least three with pair dimension zero. A singularity at the
contact point and order-three contact each impose three conditions with
that point varying on $D$. Thus they too have codimension at least two.
The generic contact is simple and $C$ smooth. The joint tangency incidence
in $(P,C,D)$ is irreducible and dominates both systems by first-jet
surjectivity, so its two ordinary inverse images meet in a dense open.
For general ordinary $D$, the generic $C\in\Delta_D$ is ordinary.

On the multiplication subspace $s_D|L|\subset|L^d|$, the divisorial
nonordinary support is exactly $\bigcup_i H_i\cup\Delta_D$.
This is the union stable-rank formula
\eqref{eq:product-union-stable-rank}. Nonreduced base members have
codimension at least two by the one-node discriminant theorem. A member
containing $D$ is impossible for $e>1$ and a single projective point for
$e=1$. Each $H_i$ meets $\Delta_D$ at the good mixed contact just
constructed. The pointed Cartier lemma and the direct dualizing-form
calculation of Lemma~\ref{lem:mixed-tacnode-Hasse} apply to any ordinary
abelian surface: its only ambient input is bijective Cartier on
$H^1(\mathcal O_A)$ and $H^2(\mathcal O_A)$, and its deformation input is
first-jet surjectivity. Thus the full Hasse divisor is smooth at every
chosen mixed point.

Every global Hasse component restricts to a nonzero positive-degree
polynomial on this multiplication subspace, since the latter contains
an ordinary point. Its restriction contains one of the displayed prime
divisors. Any component containing $H_i$ or $\Delta_D$ contains their
chosen intersection point. Smoothness there forces all these assignments
to be the same component, even if restriction multiplicities occur.
This proves global irreducibility; that same smooth point proves
multiplicity one. An ordinary pair with a single simple contact has
stable loss exactly one by the union formula. Closedness of the iterated
Frobenius rank loci therefore gives generic stable loss one.
Lemma~\ref{lem:char2-defect-one-inertia} now proves transvection generation.

Finally the coherent bundle is $\Omega^1_{\mathbf P^{9d^2-1}}$.
Lemma~\ref{lem:irreducible-hasse-quotients-trivial} implies that every
nonzero proper invariant subspace has trivial quotient. By
Proposition~\ref{prop:ordinary-ppav-all-covectors}, every such subspace
contains $K_d$, making $K_d$ irreducible. Its dimension is at least $32$,
so an invariant line is impossible. This proves $V_d^G=0$ as well.
\end{proof}

\begin{lemma}[A full degree-one open in surface moduli]
\label{lem:ordinary-ppav-GL8-open}
\status{Proved under stated hypotheses}
There is a nonempty open in the ordinary moduli stack of principally
polarized abelian surfaces on which the degree-one relative geometric
image for $|3\Theta|$ is $\GL_8(\mathbf F_2)$.
\end{lemma}
\begin{proof}
Work over an algebraically closed field. Use full level five and a fixed
primitive fifth root of unity. The corresponding fine moduli variety is
smooth and geometrically irreducible and carries its universal polarized
abelian scheme, by the exact statement
\cite[Theorem~7.3.11, printed p.~76]{ConradAbelianVarietiesModuli}.
This is an opened author exposition of the moduli theorem; no claim that
its original GIT or compactification proof has been reread is required.
Near the level-five point of the certified product
$E^2$, $E:y^2+xy=x^3+1$, choose a line bundle representing the principal
polarization after an etale neighbourhood. Such representatives form a
smooth torsor under the dual abelian scheme; rigidification at zero
supplies an actual line bundle after this localization. Do not require
this family of representatives to be symmetric: that additional torsor
need not be etale in characteristic two. Each geometric fibre is a
translate of a symmetric representative, so its geometric group is the
same. Take an irreducible smooth neighbourhood $B$ containing the chosen
point.

The cube of this line bundle has a rank-nine direct image, and its
projective bundle carries a nonempty smooth ordinary open $U$ with
geometrically integral fibres. Its rank-eight full Lang frame cover
$Z\to U$ is finite etale. The fibre at the product point is connected by
the certified degree-one $\GL_8$ theorem.
The normalization argument of Lemma~\ref{lem:ordinary-pair-GL8-open}
proves that $Z$ has geometrically connected generic fibre over $B$:
a finite extension defining distinct generic components gives a finite
normalization of the excellent base; on the normal integral projective
bundle open each open-and-closed component of the finite etale cover
surjects, contradicting the connected product fibre. The geometric
generic fibre is smooth, hence normal, so it is also irreducible.
By \cite[Lemma~37.27.5]{Stacks0559}, after shrinking $B$ the frame fibres
are geometrically irreducible. Their connectedness means full $\GL_8$
monodromy. The etale map from $B$ to fine moduli is open, and forgetting
level is open; its image is the asserted nonempty open in the moduli
stack. Translating the representing theta line bundle leaves each
geometric family isomorphic, so this property is intrinsic to the
principally polarized surface.
\end{proof}

\begin{lemma}[Nodal images and an actual natural block]
\label{lem:ordinary-ppav-isogeny-block}
\status{Proved under stated hypotheses}
Let $(A',\Theta')$ be an ordinary principally polarized surface with
full degree-one relative image $\GL_8(\mathbf F_2)$, and suppose
$\phi:A'\to A$ is an odd polarized isogeny with
$\phi^*\lambda_\Theta=n\lambda_{\Theta'}$, $n$ odd.
Then $\deg\phi=n^2$, and for every $d\ge n+1$ the degree-$d$
geometric group on $A$ contains a natural $\GL_8(\mathbf F_2)$ on an
actual direct summand, fixing a complementary summand pointwise.
\end{lemma}
\begin{proof}
All work is geometric. Put $L'=\mathcal O(3\Theta')$ and choose symmetric
representatives on these two fixed surfaces. The kernel of $\phi$ is a
constant odd finite group, acting by translations.
A general smooth ordinary $C'\in|L'|$ is preserved by no nonzero kernel
translation: one preserving the line bundle has a nontrivial projective
action in its very ample embedding and hence a proper fixed locus; one
not preserving the line bundle cannot preserve a member. Thus
$C'\to\phi(C')$ is generically one-to-one.

We prove that its image has only ordinary nodes. On an indecomposable
source, for a fixed nonzero kernel element $g$, the two first jets at
$P,P+g$ have rank six except at finitely many $P$. Indeed failure is
contained in
\[
               3P+g\in\Theta',\qquad 3P+2g\in\Theta'.
\]
The two theta translates are distinct because $g\ne0$ and the principal
theta divisor has trivial translation stabilizer. Their intersection is
finite, and $[3]$ is finite. On the rank-six locus, two value conditions
and proportionality of two nonzero gradients impose three independent
conditions. Varying $P$ adds two dimensions, so tangency is generically
absent. The exceptional finite set already imposes two independent value
conditions and is avoided by a general member.

Any three distinct points on an indecomposable $A'$ impose three
independent value conditions on $|3\Theta'|$. To separate $P_3$ from
$P_1,P_2$, choose theta translates $\Theta'_a,\Theta'_b$ through $P_1,P_2$
respectively and not through $P_3$. These choices are dense opens in two
irreducible theta translates. Their addition map dominates $A'$: sums of
two theta curves cover the surface, because any two ample translates
intersect. Thus one can also require the third translate
$\Theta'_{-a-b}$ to avoid $P_3$. The product of these three equations is
a section of $3\Theta'$ vanishing at $P_1,P_2$ but not $P_3$. Cyclically
repeat the construction to prove the three-value rank assertion.
For $P,P+g,P+h$, three independent conditions with $P$ varying in dimension
two exclude triple identifications. There are only finitely many kernel
pairs. On a decomposable source the already proved cubic-product
first-jet and Segre three-value argument gives the identical nodal
conclusion, including its one-dimensional collinear exception. Therefore
$Z=\phi(C')$ is integral nodal with normalization $C'$. For each $g$,
$C'\cap(C'-g)$ has degree $L'^2=18$; its node cover is a quotient of a
union of such degree-eighteen covers. Its permutation group is consequently
a subquotient of a product of copies of $S_{18}$.

The image line bundle is $\operatorname{Nm}_\phi L'$.
Its Neron--Severi class is $3n\Theta$: push forward
$\phi^*\Theta\equiv n\Theta'$ and cancel $n$ in the torsion-free
Neron--Severi group. Both this norm and $L^n$ are symmetric, so their
quotient is a geometric two-torsion line bundle. Since $3n$ is odd,
$\lambda_{L^n}$ identifies the geometric two-torsion of $A$ with that of
its dual. Translating the whole image family by one constant two-torsion
point therefore corrects this quotient and places $Z$ in the exact
linear system $|L^n|$. Denote the corrected finite etale map again by
$\phi$; the translation does not change its kernel identifications.

Fix a general smooth ordinary $D\in|L^{d-n}|$.
Its inverse image $D'=\phi^{-1}D$ is smooth, connected and integral:
it is a smooth ample divisor on the abelian surface, whose connectedness
also follows from ample vanishing in its structure-sheaf sequence.
For varying $C'$, the intersection with $D'$ has degree
\[
                         N=18n(d-n)\ge18.
\]
Its permutation group is the full $S_N$. Here are the required incidence
checks, so that this conclusion does not assume a generic linear-series
monodromy theorem. One-point and ordered two-point incidences are
irreducible projective bundles over $D'$ and $D'^2\setminus\Delta$,
respectively, because $L'$ separates one and two values.
A simple double zero and otherwise simple zeros is attainable. On an
indecomposable source choose $P\in D'$ outside the finite intersections
of $D'$ with the theta translates occurring in the projection of $R_4$.
These intersections are finite since $D'$ is not a theta translate.
The joint surface jets at $P,Q$ have rank at least five for every
$Q\ne P$, and rank six except at finitely many $Q$; their restrictions
to the two tangent directions of $D'$ have rank at least three and
four respectively. After imposing the two conditions for a double zero
at $P$, a second double zero is therefore a proper locus: the generic
second point varies in dimension one and adds two conditions, and each
exceptional point adds at least one. Second jets at $P$ exclude order
three. On a product source choose $P$ also outside the flex fibres;
the same bounds follow from the cubic-product jet lemma, since the
exceptional fibres meet the integral ample $D'$ finitely.
The resulting length-two collision is smooth in the point incidence
and its parameter, since first jets are generated.
The single-collision criterion
\cite[Proposition~3.1, pp.~7--8 of the opened version]{Entin2021}
then gives $S_N$, including characteristic two. Equivalently its proof
uses the irreducible one- and two-point incidences and the nontrivial
transposition inertia of this one ordinary double zero.

After shrinking the coefficient space, $Z$ and $D$ meet transversely,
the image nodes of $Z$ avoid $D$, and $Y=Z\cup D$ is ordinary nodal in
$|L^d|$. The node-avoidance condition is proper: two values at $P,P+g$
with the image point restricted to $D$ impose two conditions with a
one-dimensional point parameter. The normalization sequence and the
isomorphism $H^1_{\et}(A,\mathbf F_2)\xrightarrow{\sim}
H^1_{\et}(A',\mathbf F_2)$ under an odd isogeny give
\[
 0\subset T\subset T_0\subset V_d,\qquad
 T_0/T\text{ constant},\qquad V_d/T_0=Q,
 \qquad\dim Q=8,
\]
where $Q$ is the degree-one relative system of $C'$ and $T$ is graph
cohomology. Its image on $Q$ is full $\GL_8$ even after this shrinking.

The joint node permutation image projects onto $S_N$, with kernel a
subgroup of a product of $S_{18}$'s. The latter has no prime divisor
$31$, whereas $\GL_8(\mathbf F_2)$ does. The factors of $S_N$ are
$A_N$ and cyclic groups of order two; $A_N$ is not $\GL_8(\mathbf F_2)$,
since $11$ divides $|A_N|$ for $N\ge18$ and does not divide
$|\GL_8(\mathbf F_2)|$. Thus the joint image has no composition factor
isomorphic to that simple group. The same is true of its graph image
and of the entire $T_0$ image: the latter only adds a subgroup of the
additive two-group $\Hom(T_0/T,T)$. This step does not assume that the
constant graded pieces split as local systems.

Goursat now shows that the kernel of the entire $T_0$ action still
surjects onto $\GL(Q)$. Inside its block parabolic the translation
kernel is $Q^*\otimes T'$ for a subspace $T'\subset T_0$.
The elementary $H^1(\GL_8(\mathbf F_2),Q^*)=0$ of
Lemma~\ref{lem:char2-group-inputs} splits the remaining cocycle modulo
$T'$ after conjugation. The full translation kernel then removes the
remaining $T'$ columns, producing an actual block-diagonal $\GL_8$.
The whole boundary lies in the coherent ordinary open, so its image is
a subgroup of the same group obtained on the dense smooth ordinary
open. This proves the lemma.
\end{proof}

\begin{theorem}[Eventual exact images for all ordinary principally polarized surfaces]
\label{thm:ordinary-ppav-eventual-full}
\status{Proved under stated hypotheses}
Let $(A,\Theta)$ be an ordinary principally polarized abelian surface over
any perfect field $k$ of characteristic two, with a line bundle representing
$\Theta$ over $k$, and embed it by $L=\mathcal O_A(3\Theta)$.
After a finite extension of $k$, there exists an odd polarized isogeny
$\phi:A'\to A$ with multiplier $n$ whose source has degree-one relative
geometric group $\GL_8(\mathbf F_2)$.
For the complete systems $|L^d|$, put $d_1=n+1$.
For the ambient hypersurface systems of this fixed embedding, replace it by
\[
           d_1=\max\{n+1,\operatorname{reg}(\mathcal I_A)-1\}.
\]
For every $d\ge d_1$, put $r=9d^2-1$. If $d$ is odd, the entire geometric
and arithmetic images both equal $\GL_r(\mathbf F_2)$.
If $d$ is even, put $N=r-3$ and let $\Gamma\subset S_3$ be the actual
Galois image permuting the three nonzero characters of $\ker F_A$.
After compatible choices of bases the entire images are
\[
 G_{{\rm geom},d}=
 \left\{\begin{pmatrix}B&T\\0&I_3\end{pmatrix}:
 B\in\GL_N(\mathbf F_2),\ T\in M_{N\times3}(\mathbf F_2)\right\},
\]
\[
 G_{{\rm arith},d}=
 \left\{\begin{pmatrix}B&T\\0&\gamma\end{pmatrix}:
 B\in\GL_N(\mathbf F_2),\ T\in M_{N\times3}(\mathbf F_2),\
 \gamma\in\Gamma\right\},
\]
where $\Gamma$ acts by its three-dimensional permutation matrices.
The representation is irreducible exactly in odd degree.
The integer $n$, and hence the threshold, can be found by a terminating
procedure over finite fields or effectively presented finitely generated
fields of definition; over arbitrary perfect fields the statement gives
existence without presupposing computable coefficients.
\end{theorem}
\begin{proof}
Chai's theorem \cite[Theorem~2, printed p.~477]{Chai1995} states over
any algebraically closed characteristic-$p$ field that the fixed
$\ell$-power polarized Hecke orbit of every ordinary principally
polarized abelian variety is dense in its moduli stack.
Use $g=2$, $p=2$, $\ell=3$. Its orbit meets the nonempty open of
Lemma~\ref{lem:ordinary-ppav-GL8-open}. By the defining form of this orbit
there is an odd polarized isogeny from a source in that open to $A$;
its multiplier is an odd $n$ and its degree is $n^2$. Its finitely many
defining algebraic coefficients lie in a finite extension of $k$.
This auxiliary extension is only used for the geometric subgroup
construction; the arithmetic image below is computed over the original $k$.

For $d\ge n+1\ge2$, Theorem~\ref{thm:ordinary-ppav-Hasse-kernel} gives
an irreducible common kernel $K_d$ and transvection generation.
Lemma~\ref{lem:ordinary-ppav-isogeny-block} places a natural eight-space
$Q$ inside $K_d$, since a natural $\GL_8$ has no nonzero map to a trivial
module. On $K_d$ the same subgroup fixes a complementary summand.
No invariant nondegenerate alternating form can exist: the natural
$\GL_8$ has no invariant bilinear form and no invariant covector, so
both its internal pairing and its pairing with that trivial complement
are zero. Thus $Q$ would lie in the radical. The opened transvection
classification \cite[Theorem~3.6, printed p.~22]{KantorSeress2001} now
forces the image on $K_d$ to be the full $\GL(K_d)$.
For odd $d$ this is already the entire claimed image.

For even $d$, the translation kernel in
$\Hom(V_d/K_d,K_d)$ is $K_d\otimes U$ for a subspace
$U\subset(V_d/K_d)^*$. The natural-module $H^1$ vanishing of
Lemma~\ref{lem:char2-group-inputs} conjugates the remaining extension
to its block form. Its invariant-vector space has dimension
$3-\dim U$. Theorem~\ref{thm:ordinary-ppav-Hasse-kernel} proved that
there is no invariant vector, so $U$ is the whole three-space and
all translations occur. This proves the geometric image.

Proposition~\ref{prop:ordinary-ppav-all-covectors} identifies the entire
geometric covector space together with its Galois action. Consequently
the arithmetic image is contained in the displayed preimage of $\Gamma$.
It contains the full geometric kernel just proved and projects onto
$\Gamma$, because the parameter open is geometrically connected and
its fundamental group surjects onto the absolute Galois group.
It is therefore exactly that preimage. The action on the quotient is
the dual of the covector permutation representation; permutation matrices
identify these two representations. In odd degree the ambient full
linear group leaves no room for a larger arithmetic image.

For the fixed projective embedding, regularity gives
$H^1(\mathbf P^8,\mathcal I_A(d))=0$ when
$d\ge\operatorname{reg}(\mathcal I_A)-1$.
Ambient sections then surject onto $H^0(A,L^d)$; over the nonzero-section
projective open this is an affine bundle, and the pullback Lang cover
has the same geometric and arithmetic image. This avoids a projective
normality assertion in characteristic two.

Finally enumerate all polarized three-power isogenies from the given
surface, by its finite etale three-power torsion and maximal isotropic
subgroups, and reverse them by polarized duals. Compute the degree-one
geometric Lang frame image of each source by the terminating exact
algorithms proved earlier, dovetailing the tests. Hecke density and the
nonempty full-image open ensure that some source passes. Its retained
multiplier is the required $n$. Regularity of the fixed homogeneous ideal
is computable by a finite graded resolution. All field extensions,
finite subgroup schemes and quotients here are described by finite
algebraic data over an effective field of definition. No undecided
Galois-group oracle or unproved generic-group guess enters this procedure.
The odd full module is irreducible, and in even degree its nonzero proper
common kernel of codimension three is invariant, proving the last
representation-theoretic assertion.
\end{proof}

\begin{remark}[Scope of this structural extension]\label{auto:A_monodromy:20bf13b82135}
\status{Proved under stated hypotheses}

The theorem treats every ordinary principally polarized abelian surface in
characteristic two, including indecomposable surfaces, after a computable
threshold. It does not identify degree-one groups on every such surface,
and it does not assert a structural full-image formula for arbitrary
ambient smooth projective varieties. Those are separate refinements from
the exact finite algorithms and the general residue theorem already proved.
\end{remark}

\begin{remark}[Finite arithmetic and subgroup checks]\label{auto:A_monodromy:03a3e68c4a1c}
\status{Proved under stated hypotheses}

The retained exact program
\path{agents/monodromy/code/ordinary_ppav_group_checks.py}
checks all six $\GL_2(\mathbf F_2)$ character actions on
$\mathbf F_2^4/\mathbf F_2\mathbf1$, their identification with the
three-dimensional permutation representation, and the prime-$31$ and
prime-$11$ separations used for the actual natural block. Its output is
\path{agents/monodromy/outputs/ordinary_ppav_group_checks.json}.
These finite checks support the arithmetic and group calculations; the
geometric incidence arguments and Hecke-density input are proved or cited
separately above.
\end{remark}

